% Options for packages loaded elsewhere
\PassOptionsToPackage{unicode}{hyperref}
\PassOptionsToPackage{hyphens}{url}
\PassOptionsToPackage{dvipsnames,svgnames,x11names}{xcolor}
%
\documentclass[
  10pt,
]{article}
\usepackage{amsmath,amssymb}
\usepackage{iftex}
\ifPDFTeX
  \usepackage[T1]{fontenc}
  \usepackage[utf8]{inputenc}
  \usepackage{textcomp} % provide euro and other symbols
\else % if luatex or xetex
  \usepackage{unicode-math} % this also loads fontspec
  \defaultfontfeatures{Scale=MatchLowercase}
  \defaultfontfeatures[\rmfamily]{Ligatures=TeX,Scale=1}
\fi
\usepackage{lmodern}
\ifPDFTeX\else
  % xetex/luatex font selection
  \setmainfont[]{DejaVu Serif}
  \setsansfont[]{DejaVu Sans}
  \setmonofont[]{DejaVu Sans Mono}
  \setmathfont[]{Latin Modern Math}
\fi
% Use upquote if available, for straight quotes in verbatim environments
\IfFileExists{upquote.sty}{\usepackage{upquote}}{}
\IfFileExists{microtype.sty}{% use microtype if available
  \usepackage[]{microtype}
  \UseMicrotypeSet[protrusion]{basicmath} % disable protrusion for tt fonts
}{}
\makeatletter
\@ifundefined{KOMAClassName}{% if non-KOMA class
  \IfFileExists{parskip.sty}{%
    \usepackage{parskip}
  }{% else
    \setlength{\parindent}{0pt}
    \setlength{\parskip}{6pt plus 2pt minus 1pt}}
}{% if KOMA class
  \KOMAoptions{parskip=half}}
\makeatother
\usepackage{xcolor}
\usepackage[margin=2.0cm]{geometry}
\usepackage{color}
\usepackage{fancyvrb}
\newcommand{\VerbBar}{|}
\newcommand{\VERB}{\Verb[commandchars=\\\{\}]}
\DefineVerbatimEnvironment{Highlighting}{Verbatim}{commandchars=\\\{\}}
% Add ',fontsize=\small' for more characters per line
\newenvironment{Shaded}{}{}
\newcommand{\AlertTok}[1]{\textcolor[rgb]{1.00,0.00,0.00}{\textbf{#1}}}
\newcommand{\AnnotationTok}[1]{\textcolor[rgb]{0.38,0.63,0.69}{\textbf{\textit{#1}}}}
\newcommand{\AttributeTok}[1]{\textcolor[rgb]{0.49,0.56,0.16}{#1}}
\newcommand{\BaseNTok}[1]{\textcolor[rgb]{0.25,0.63,0.44}{#1}}
\newcommand{\BuiltInTok}[1]{\textcolor[rgb]{0.00,0.50,0.00}{#1}}
\newcommand{\CharTok}[1]{\textcolor[rgb]{0.25,0.44,0.63}{#1}}
\newcommand{\CommentTok}[1]{\textcolor[rgb]{0.38,0.63,0.69}{\textit{#1}}}
\newcommand{\CommentVarTok}[1]{\textcolor[rgb]{0.38,0.63,0.69}{\textbf{\textit{#1}}}}
\newcommand{\ConstantTok}[1]{\textcolor[rgb]{0.53,0.00,0.00}{#1}}
\newcommand{\ControlFlowTok}[1]{\textcolor[rgb]{0.00,0.44,0.13}{\textbf{#1}}}
\newcommand{\DataTypeTok}[1]{\textcolor[rgb]{0.56,0.13,0.00}{#1}}
\newcommand{\DecValTok}[1]{\textcolor[rgb]{0.25,0.63,0.44}{#1}}
\newcommand{\DocumentationTok}[1]{\textcolor[rgb]{0.73,0.13,0.13}{\textit{#1}}}
\newcommand{\ErrorTok}[1]{\textcolor[rgb]{1.00,0.00,0.00}{\textbf{#1}}}
\newcommand{\ExtensionTok}[1]{#1}
\newcommand{\FloatTok}[1]{\textcolor[rgb]{0.25,0.63,0.44}{#1}}
\newcommand{\FunctionTok}[1]{\textcolor[rgb]{0.02,0.16,0.49}{#1}}
\newcommand{\ImportTok}[1]{\textcolor[rgb]{0.00,0.50,0.00}{\textbf{#1}}}
\newcommand{\InformationTok}[1]{\textcolor[rgb]{0.38,0.63,0.69}{\textbf{\textit{#1}}}}
\newcommand{\KeywordTok}[1]{\textcolor[rgb]{0.00,0.44,0.13}{\textbf{#1}}}
\newcommand{\NormalTok}[1]{#1}
\newcommand{\OperatorTok}[1]{\textcolor[rgb]{0.40,0.40,0.40}{#1}}
\newcommand{\OtherTok}[1]{\textcolor[rgb]{0.00,0.44,0.13}{#1}}
\newcommand{\PreprocessorTok}[1]{\textcolor[rgb]{0.74,0.48,0.00}{#1}}
\newcommand{\RegionMarkerTok}[1]{#1}
\newcommand{\SpecialCharTok}[1]{\textcolor[rgb]{0.25,0.44,0.63}{#1}}
\newcommand{\SpecialStringTok}[1]{\textcolor[rgb]{0.73,0.40,0.53}{#1}}
\newcommand{\StringTok}[1]{\textcolor[rgb]{0.25,0.44,0.63}{#1}}
\newcommand{\VariableTok}[1]{\textcolor[rgb]{0.10,0.09,0.49}{#1}}
\newcommand{\VerbatimStringTok}[1]{\textcolor[rgb]{0.25,0.44,0.63}{#1}}
\newcommand{\WarningTok}[1]{\textcolor[rgb]{0.38,0.63,0.69}{\textbf{\textit{#1}}}}
\usepackage{longtable,booktabs,array}
\usepackage{calc} % for calculating minipage widths
% Correct order of tables after \paragraph or \subparagraph
\usepackage{etoolbox}
\makeatletter
\patchcmd\longtable{\par}{\if@noskipsec\mbox{}\fi\par}{}{}
\makeatother
% Allow footnotes in longtable head/foot
\IfFileExists{footnotehyper.sty}{\usepackage{footnotehyper}}{\usepackage{footnote}}
\makesavenoteenv{longtable}
\setlength{\emergencystretch}{3em} % prevent overfull lines
\providecommand{\tightlist}{%
  \setlength{\itemsep}{0pt}\setlength{\parskip}{0pt}}
\setcounter{secnumdepth}{-\maxdimen} % remove section numbering
% doc_header.tex -- LaTeX preamble additions used by build_docs.sh for every
% Markdown -> LaTeX -> PDF conversion (supplements, README, submission files).
\usepackage{etoolbox}
% longtable may not be loaded by pandoc when a document has no table
\usepackage{longtable}
% bold vectors: \bm -> \symbf under unicode-math (XeLaTeX), bm package otherwise
\makeatletter
\@ifpackageloaded{unicode-math}{\providecommand{\bm}[1]{\symbf{#1}}}{\usepackage{bm}}
\makeatother
% smaller type in long tables so that wide audit tables fit the text width;
% inside a table a line break is allowed after an underscore of plain text
% (pandoc writes it as \_), so that long identifiers wrap inside their column
% rows are set a little apart, so that symbols with indices do not touch
\AtBeginEnvironment{longtable}{\footnotesize\renewcommand{\arraystretch}{1.45}%
  \let\sptunderscore\_\def\_{\sptunderscore\allowbreak}}
\setlength{\LTcapwidth}{\textwidth}
% allow line breaks after slashes and hyphens inside table cells
\usepackage{xurl}
% code blocks: break over-long lines instead of running into the margin.
% Plain ``` blocks are typeset by pandoc as {verbatim}, highlighted ones
% (```bash) as {Highlighting}; both are re-declared as breakable fancyvrb
% environments.  The continuation line is marked with a small hooked arrow.
\usepackage{fvextra}
\RecustomVerbatimEnvironment{verbatim}{Verbatim}{breaklines,breakanywhere,fontsize=\small}
\makeatletter
\@ifundefined{Highlighting}{}{%
  \RecustomVerbatimEnvironment{Highlighting}{Verbatim}{breaklines,breakanywhere,fontsize=\small,commandchars=\\\{\}}}
\makeatother
% inline code (pandoc writes it as \texttt{...}): never hyphenate -- an inserted
% hyphen could be read as part of a file name -- but allow a line break after
% an underscore, slash, comma, colon or equals sign, so that long paths and
% column lists wrap inside the text width and inside table columns.
% Inside a formula (e.g. the name of a notebook symbol in an equation) the same
% command must first leave math mode, which \text of amsmath does.
\ExplSyntaxOn
\tl_new:N \l__spt_code_tl
\cs_new_protected:Npn \__spt_texttt:n #1
  {
    \tl_set:Nn \l__spt_code_tl {#1}
    \regex_replace_all:nnN { ([\/\,\:\=]) } { \1 \c{allowbreak} } \l__spt_code_tl
    \regex_replace_all:nnN { \c{_} } { \c{_} \c{allowbreak} } \l__spt_code_tl
    \group_begin: \ttfamily \hyphenchar \font = -1 \scan_stop: \l__spt_code_tl \group_end:
  }
\RenewDocumentCommand \texttt { m }
  {
    \mode_if_math:TF { \text { \__spt_texttt:n {#1} } } { \__spt_texttt:n {#1} }
  }
\ExplSyntaxOff
% A figure is written in the Markdown source as an image paragraph followed by
% a caption paragraph ("**Figure S1.** ..."), so that the caption is also
% visible on GitHub.  An image that starts a paragraph is centered.
\usepackage{graphicx}% pandoc loads it only for documents that contain an image
\NewCommandCopy\sptincludegraphics\includegraphics
\RenewDocumentCommand\includegraphics{O{} m}{%
  \ifvmode\begin{center}\sptincludegraphics[#1]{#2}\end{center}%
  \else\sptincludegraphics[#1]{#2}\fi}
% keep figures where they are referenced
\usepackage{float}
\floatplacement{figure}{H}
% running header with the document title
\usepackage{fancyhdr}
\pagestyle{fancy}
\fancyhf{}
\fancyhead[L]{\small PROJECT spacetimes --- documentation and provenance}
\fancyhead[R]{\small\thepage}
\renewcommand{\headrulewidth}{0.4pt}
% 16x16 matrices: tighter columns so that a full matrix fits the text width
\setlength{\arraycolsep}{3.2pt}
% The documents of this project contain long tagged symbols, file paths and
% identifiers that cannot be hyphenated.  A paragraph is therefore allowed to
% have looser lines before TeX gives up and lets a line stick out of the margin.
\AtBeginDocument{\tolerance=2000 \emergencystretch=5em \hbadness=2000}
\ifLuaTeX
  \usepackage{selnolig}  % disable illegal ligatures
\fi
\IfFileExists{bookmark.sty}{\usepackage{bookmark}}{\usepackage{hyperref}}
\IfFileExists{xurl.sty}{\usepackage{xurl}}{} % add URL line breaks if available
\urlstyle{same}
\hypersetup{
  pdftitle={dirac matrices --- the two sets of eight real 16×16 Dirac matrices of PROJECT spacetimes, and the Dirac equations of the fields dirac1644 and dirac1671},
  pdfauthor={Prepared for Patrick Lee Nash, Ph.D.},
  colorlinks=true,
  linkcolor={blue},
  filecolor={Maroon},
  citecolor={blue},
  urlcolor={blue},
  pdfcreator={LaTeX via pandoc}}

\title{dirac matrices --- the two sets of eight real 16×16 Dirac
matrices of PROJECT spacetimes, and the Dirac equations of the fields
dirac1644 and dirac1671}
\usepackage{etoolbox}
\makeatletter
\providecommand{\subtitle}[1]{% add subtitle to \maketitle
  \apptocmd{\@title}{\par {\large #1 \par}}{}{}
}
\makeatother
\subtitle{Provenance: definitions, displays, proofs, machine checks, and
how to reproduce them}
\author{Prepared for Patrick Lee Nash, Ph.D.}
\date{November 12, 2026}

\begin{document}
\maketitle

\hypertarget{what-this-document-proves-and-for-whom}{%
\section{0. What this document proves, and for
whom}\label{what-this-document-proves-and-for-whom}}

This document is written for a reader who knows the Dirac equation in
four dimensions but has never seen the matrices of this project. It has
three parts.

\begin{itemize}
\tightlist
\item
  \textbf{Part 1} treats the eight real 16×16 Dirac matrices of the
  spacetime of signature \textbf{4+4} (four space-like and four
  time-like dimensions), which are defined in the author's notebook
  \texttt{4+4.nb}, and the Dirac equation of the field
  \textbf{dirac1644}.
\item
  \textbf{Part 2} treats the eight real 16×16 Dirac matrices of the
  spacetime of signature \textbf{7+1} (seven space-like dimensions and
  one time-like dimension), which are defined in the notebook
  \texttt{7+1.nb}, and the Dirac equation of the field
  \textbf{dirac1671}.
\item
  \textbf{Part 3} derives from the two sets of matrices a result that
  the rest of the project needs: which pairs of two identical fermions
  can exist.
\end{itemize}

For each set the document (i) displays the eight matrices and the other
16×16 matrices that matter; (ii) proves that the eight matrices are
real; (iii) proves their anticommutation relations; (iv) determines the
group that their scaled commutators generate; (v) re-derives the chiral
projection operators; (vi) writes down a Dirac equation for a
16-component field whose components are complex-valued Grassmann
operators, and checks it; and (vii) uses the chiral projections to write
that equation as two coupled 8-component equations.

\textbf{How the proofs are made.} Every matrix in this document has
entries \(0\), \(+1\), \(-1\) (or, for two projectors, \(0\),
\(\pm\tfrac12\), \(\pm\tfrac i2\)). An identity between such matrices
can be verified with exact integer arithmetic, and a verification of
that kind is a proof. The Rust program \texttt{spacetimes} (directory
\texttt{code/spacetimes-rs}) performs 107 such verifications. Each has
an identifier, such as \texttt{44-C1} (the prefix \texttt{44-} refers to
the 4+4 set, \texttt{71-} to the 7+1 set), which is quoted below at the
place where the statement is used; the complete output of the program is
reproduced verbatim in Appendix A. Where a human-readable argument is
short, it is given as well, so that the reader can follow the logic
without the computer.

\textbf{Three independent constructions.} The matrices are obtained in
three ways, which agree entry by entry (checks \texttt{44-N1},
\texttt{44-N2}, \texttt{44-P3}, \texttt{71-N1} of the Rust program and
\texttt{I44-F}, \texttt{I71-F}, \texttt{I44-N}, \texttt{I71-N} of the
Python script):

\begin{enumerate}
\def\labelenumi{\arabic{enumi}.}
\tightlist
\item
  \emph{as printed by Mathematica}: the script
  \texttt{tools/extract\_nb\_matrices.py} reads the output cells of the
  notebooks and writes \texttt{data/notebook\_matrices\_4+4.txt} and
  \texttt{data/notebook\_matrices\_7+1.txt};
\item
  \emph{rebuilt in Rust} from the definitions in the notebooks' input
  cells (\texttt{code/spacetimes-rs/src/gamma.rs});
\item
  \emph{rebuilt in Python} by a separate implementation
  (\texttt{code/scripts/verify\_matrices\_independent.py}), which also
  repeats the essential proofs by different methods (Appendix A.2).
\end{enumerate}

\begin{longtable}[]{@{}
  >{\raggedright\arraybackslash}p{(\columnwidth - 2\tabcolsep) * \real{0.4000}}
  >{\raggedright\arraybackslash}p{(\columnwidth - 2\tabcolsep) * \real{0.6000}}@{}}
\toprule\noalign{}
\begin{minipage}[b]{\linewidth}\raggedright
Input file
\end{minipage} & \begin{minipage}[b]{\linewidth}\raggedright
SHA-256 checksum
\end{minipage} \\
\midrule\noalign{}
\endhead
\bottomrule\noalign{}
\endlastfoot
\texttt{4+4.nb} &
\texttt{740023233eaa0cd24ca0e2e7010094671b1d314b698114c08db94050d3e0aed5} \\
\texttt{7+1.nb} &
\texttt{41d252f9852f210c228f45e1ac956d2d7ef26faeb0062a5880280e1aae0effc7} \\
\texttt{Dirac\_Einstein-Lovelock\_4+4\_Nash.nb} &
\texttt{cdbb6a2c0b31ffd6a4fd951746a1e1c3ea5748afd371edc5095cea7e34cf79c0} \\
\end{longtable}

\hypertarget{coordinates-confirmation}{%
\subsection{0.1 Coordinates:
confirmation}\label{coordinates-confirmation}}

We use the coordinates \[X=\{x0,\;x1,\;x2,\;x3,\;x4,\;x5,\;x6,\;x7\}\]
of the notebook \texttt{Dirac\_Einstein-Lovelock\_4+4\_Nash.nb}, section
\texttt{DEFCHART}, where they are defined by
\texttt{X\ =\ \{x0{[}{]},\ x1{[}{]},\ x2{[}{]},\ x3{[}{]},\ x4{[}{]},\ x5{[}{]},\ x6{[}{]},\ x7{[}{]}\}}
(with the xAct libraries a coordinate is written \texttt{x0{[}{]}}).
They are never renamed, and their meaning is the same in both signatures
as far as the first five are concerned:

\begin{longtable}[]{@{}
  >{\raggedright\arraybackslash}p{(\columnwidth - 4\tabcolsep) * \real{0.2857}}
  >{\raggedright\arraybackslash}p{(\columnwidth - 4\tabcolsep) * \real{0.3571}}
  >{\raggedright\arraybackslash}p{(\columnwidth - 4\tabcolsep) * \real{0.3571}}@{}}
\toprule\noalign{}
\begin{minipage}[b]{\linewidth}\raggedright
Coordinate
\end{minipage} & \begin{minipage}[b]{\linewidth}\raggedright
4+4 spacetime
\end{minipage} & \begin{minipage}[b]{\linewidth}\raggedright
7+1 spacetime
\end{minipage} \\
\midrule\noalign{}
\endhead
\bottomrule\noalign{}
\endlastfoot
\(x0\) & hidden space-like coordinate & hidden space-like coordinate \\
\(x1,x2,x3\) & the Cartesian space coordinates \(x,y,z\) & the Cartesian
space coordinates \(x,y,z\) \\
\(x4\) & ordinary time & ordinary time \\
\(x5,x6,x7\) & three additional time-like coordinates & three additional
space-like coordinates \\
\end{longtable}

An index \(A,B,C,D\in\{0,\dots,7\}\) always labels the coordinate
\(xA\); we write \(\partial_A\equiv\partial/\partial x^A\) and sum over
repeated indices.

\emph{A difference from \texttt{7+1.nb} that the reader must know.} The
notebook \texttt{7+1.nb} (its section 1.1) uses the same eight symbols
but lets \(x0\) be the time and \(x4\) a space coordinate. On November
12, 2026 the author decided that in this project \(x4\) is the time in
both signatures. The matrices of \texttt{7+1.nb} are therefore used here
with the indices 0 and 4 exchanged. The exchange is denoted
\(\mathrm{swap}_{04}\): \(\mathrm{swap}_{04}(0)=4\),
\(\mathrm{swap}_{04}(4)=0\), and \(\mathrm{swap}_{04}(A)=A\) otherwise.
Section 2.1 states the relabelling explicitly.

\hypertarget{the-unique-symbols-of-this-document}{%
\subsection{0.2 The unique symbols of this
document}\label{the-unique-symbols-of-this-document}}

The author's notebooks use the same names for different objects. Three
examples that concern this document directly: \texttt{γ} is a list of
sixteen 4×4 matrices in \texttt{4+4.nb} but the four-dimensional Dirac
matrices in \texttt{Dirac\_Nash\_Identity.nb}; \texttt{g44} and
\texttt{g71} are the metric tensors of the gravity notebooks; and
\texttt{7+1.nb} itself contains two different sets of matrices and two
different chirality matrices, \texttt{Γ17}, \texttt{χ17} (real set) and
\texttt{Γ71}, \texttt{χ71} (complex set), as well as two metrics,
\texttt{η17} and \texttt{η71} \(=-\)\texttt{η17}. The MASTER provenance
(\texttt{provenance/P1\_master\_symbols}) lists every collision and its
remedy. The remedy adopted throughout the project is a \textbf{tag}:

\begin{itemize}
\tightlist
\item
  the tag \textbf{1644} marks an object of the 4+4 spacetime (the
  spacetime of the field dirac1644: 16 components, signature 4+4);
\item
  the tag \textbf{1671} marks an object of the 7+1 spacetime (the field
  dirac1671).
\end{itemize}

No name of any original file contains either tag, so a tagged symbol
cannot be mistaken for a symbol of a notebook. The symbols of this
document are the following; the second column gives the name used in
program output and in data files.

\begin{longtable}[]{@{}
  >{\raggedright\arraybackslash}p{(\columnwidth - 6\tabcolsep) * \real{0.1765}}
  >{\raggedright\arraybackslash}p{(\columnwidth - 6\tabcolsep) * \real{0.1647}}
  >{\raggedright\arraybackslash}p{(\columnwidth - 6\tabcolsep) * \real{0.4000}}
  >{\raggedright\arraybackslash}p{(\columnwidth - 6\tabcolsep) * \real{0.2588}}@{}}
\toprule\noalign{}
\begin{minipage}[b]{\linewidth}\raggedright
Symbol
\end{minipage} & \begin{minipage}[b]{\linewidth}\raggedright
Name in files
\end{minipage} & \begin{minipage}[b]{\linewidth}\raggedright
Meaning
\end{minipage} & \begin{minipage}[b]{\linewidth}\raggedright
Name in the notebooks
\end{minipage} \\
\midrule\noalign{}
\endhead
\bottomrule\noalign{}
\endlastfoot
\(\eta_{1644}\) & \texttt{eta1644} & metric of 4+4,
\(\mathrm{diag}(+1,+1,+1,+1,-1,-1,-1,-1)\) (space-like \(=+1\)) &
\texttt{η4488} (\texttt{4+4.nb}) \\
\(\gamma_{1644}^A\) & \texttt{gamma1644{[}A{]}} & the eight Dirac
matrices of 4+4 & \texttt{T16\^{}A{[}A{]}} (\texttt{4+4.nb}),
\texttt{T16{[}A{]}} (\texttt{7+1.nb}) \\
\(\omega_{1644}\) & \texttt{omega1644} & volume element of 4+4: the
product of the eight \(\gamma_{1644}^A\) in the order \(A=0,\dots,7\) &
\texttt{T16\^{}A{[}8{]}} \\
\(P_{1644}^{L}\), \(P_{1644}^{R}\) & \texttt{PL1644}, \texttt{PR1644} &
chiral projectors of 4+4 (real) & \texttt{P\_L}, \texttt{P\_R}
(\texttt{4+4.nb}); \texttt{PL44}, \texttt{PR44} (\texttt{7+1.nb}) \\
\(\tau_{1644}^A\), \(\bar\tau_{1644}^A\) & \texttt{tau1644{[}A{]}},
\texttt{taubar1644{[}A{]}} & real 8×8 blocks of \(\gamma_{1644}^A\) &
\texttt{τ{[}A{]}} and τ with an overbar (\texttt{4+4.nb});
\texttt{τ{[}A{]}}, \texttt{τbar{[}A{]}} (\texttt{7+1.nb}) \\
\(\sigma_{88}\) & \texttt{sigma88} & the 8×8 matrix with the 4×4 blocks
\((0,1_4;\,1_4,0)\) & \texttt{σ}, \texttt{σ88} (\texttt{4+4.nb}) \\
\(C_{1644}^{-}\), \(C_{1644}^{+}\) & \texttt{Cm1644}, \texttt{Cp1644} &
the products \(\gamma^0\gamma^1\gamma^2\gamma^3\) and
\(\gamma^4\gamma^5\gamma^6\gamma^7\) of the 4+4 matrices & \texttt{σ16}
(\texttt{4+4.nb}); \texttt{S44} (\texttt{7+1.nb}) \\
\(W_{1644}\) & \texttt{W1644} & \(i\,C_{1644}^{-}\gamma_{1644}^4\), the
grading of Sec. 1.8.4 & --- \\
\(K_{1644}\), \(\mathsf{h}_{1644}\) & \texttt{K1644}, \texttt{h1644} &
generator of the evolution in \(x4\), and Hamiltonian matrix, of 4+4 &
--- \\
\(\psi_{1644}\) & \texttt{dirac1644} & the field \textbf{dirac1644} &
--- \\
\(\psi_{1644}^{L}\), \(\psi_{1644}^{R}\) & \texttt{psiL1644},
\texttt{psiR1644} & the two 8-component chiral parts of dirac1644 &
--- \\
\(m_{1644}\) & \texttt{m1644} & mass of the field dirac1644 & --- \\
\(\eta_{1671}\) & \texttt{eta1671} & metric of 7+1,
\(\mathrm{diag}(-1,-1,-1,-1,+1,-1,-1,-1)\) (time-like \(=+1\), time at
index 4) & \texttt{η17} with the indices 0 and 4 exchanged; \textbf{not}
\texttt{η71} \\
\(\gamma_{1671}^A\) & \texttt{gamma1671{[}A{]}} & the eight Dirac
matrices of 7+1 & \texttt{Γ17up{[}swap04(A){]}}; \textbf{not}
\texttt{Γ71{[}A{]}} \\
\(\omega_{1671}\) & \texttt{omega1671} & volume element of 7+1: the
product of the eight \(\gamma_{1671}^A\) in the order \(A=0,\dots,7\) &
\texttt{Γ17{[}8{]}}, \texttt{ω17} \\
\(\chi_{1671}\) & \texttt{chi1671} & \(i\,\omega_{1671}\), chirality
matrix of 7+1 & \texttt{χ17}; \textbf{not} \texttt{χ71} \\
\(P_{1671}^{L}\), \(P_{1671}^{R}\) & \texttt{PL1671}, \texttt{PR1671} &
chiral projectors of 7+1 (complex) & \texttt{PL17}, \texttt{PR17};
\textbf{not} \texttt{PL71}, \texttt{PR71} \\
\(\tau_{1671}^A\) & \texttt{tau1671{[}A{]}} & real 8×8 blocks of the 7+1
matrices & \texttt{τO{[}swap04(A){]}} \\
\(C_{1671}^{+}\), \(C_{1671}^{-}\) & \texttt{Cp1671}, \texttt{Cm1671} &
\(\gamma_{1671}^4\) and \(\gamma_{1671}^4\omega_{1671}\) & --- \\
\(\beta_{1671}\), \(\alpha_{1671}^j\) & \texttt{beta1671},
\texttt{alpha1671\^{}j} & Hamiltonian matrices of 7+1 & \texttt{βD},
\texttt{αD{[}swap04(j){]}} \\
\(\mathsf{h}_{1671}\) & \texttt{h1671} & Hamiltonian of 7+1 &
\texttt{h} \\
\(E_{1671}^{L}\), \(E_{1671}^{R}\) & \texttt{EL1671}, \texttt{ER1671} &
chiral frames (16×8 matrices) of 7+1 & \texttt{EL}, \texttt{ER} \\
\(\phi_{1671}^{L}\), \(\phi_{1671}^{R}\) & \texttt{phiL1671},
\texttt{phiR1671} & the two 8-component chiral parts of dirac1671 &
\texttt{φL}, \texttt{φR} \\
\(U_{1671}^{c}\) & \texttt{Uc1671} & the unitary matrix of Eq. (2.5) &
--- \\
\(\psi_{1671}\) & \texttt{dirac1671} & the field \textbf{dirac1671} &
\texttt{dirac16complex} \\
\(m_{1671}\) & \texttt{m1671} & mass of the field dirac1671 &
\texttt{m}, \texttt{mD} \\
\(\Sigma_{1644}^{AB}\), \(\Sigma_{1671}^{AB}\) &
\texttt{Sigma1644\^{}(AB)}, \texttt{Sigma1671\^{}(AB)} &
\(\tfrac14[\gamma^A,\gamma^B]\), the scaled commutators of each set &
\texttt{Σ17{[}A,B{]}} (for the \texttt{Γ17}) \\
\(\Delta_{1644}^{A}\), \(\Delta_{1671}^{A}\) & \texttt{DeltaV1644\^{}A},
\texttt{DeltaV1671\^{}A} & the vector pair operators of Part 3 & --- \\
\end{longtable}

\textbf{Scope rule.} Within Part 1 every formula refers to the 4+4
spacetime, and the tag 1644 is dropped from the symbols inside
derivations; within Part 2 the tag 1671 is dropped in the same way. The
boxed equations, which state the results, carry the full tags. In Part 3
and in the appendices, where both signatures appear together, every tag
is written.

In the displays of matrices a dot stands for the entry 0, and thin rules
separate the 8×8 blocks (4×4 blocks in 8×8 matrices). The same matrices
are available as plain integers in \texttt{data/gamma1644.txt} and
\texttt{data/gamma1671.txt}.

\hypertarget{how-to-reproduce-everything-in-this-document}{%
\subsection{0.3 How to reproduce everything in this
document}\label{how-to-reproduce-everything-in-this-document}}

\begin{Shaded}
\begin{Highlighting}[]
\BuiltInTok{cd} \StringTok{"/Users/nsh/Developer/github/spacetimes"}
\CommentTok{\# 1. matrices printed by Mathematica {-}\textgreater{} text files (needs only python3)}
\ExtensionTok{python3}\NormalTok{ tools/extract\_nb\_matrices.py 4+4.nb 314 136 163 165 13 }\OperatorTok{\textgreater{}} \StringTok{"data/notebook\_matrices\_4+4.txt"}
\ExtensionTok{python3}\NormalTok{ tools/extract\_nb\_matrices.py 7+1.nb 204 }\OperatorTok{\textgreater{}} \StringTok{"data/notebook\_matrices\_7+1.txt"}
\CommentTok{\# 2. build the Rust program, run its unit tests, run all checks (a few seconds)}
\ExtensionTok{cargo}\NormalTok{ build }\AttributeTok{{-}{-}release} \AttributeTok{{-}{-}manifest{-}path}\NormalTok{ code/spacetimes{-}rs/Cargo.toml}
\ExtensionTok{cargo}\NormalTok{ test  }\AttributeTok{{-}{-}release} \AttributeTok{{-}{-}manifest{-}path}\NormalTok{ code/spacetimes{-}rs/Cargo.toml}
\ExtensionTok{./code/spacetimes{-}rs/target/release/spacetimes}\NormalTok{ matrices}
\CommentTok{\# 3. independent check (pure Python, exact arithmetic, about 2 seconds)}
\ExtensionTok{python3}\NormalTok{ code/scripts/verify\_matrices\_independent.py}
\CommentTok{\# 4. rebuild this document (Markdown, LaTeX, PDF)}
\ExtensionTok{python3}\NormalTok{ code/scripts/assemble\_docs.py}
\FunctionTok{bash}\NormalTok{ code/scripts/build\_docs.sh }\StringTok{"dirac matrices.md"}
\end{Highlighting}
\end{Shaded}

Step 2 writes \texttt{data/gamma1644.txt}, \texttt{data/gamma1671.txt},
\texttt{data/matrix\_checks.txt}, \texttt{data/matrix\_checks.csv},
\texttt{data/pair\_classification.csv},
\texttt{data/pair\_classification.txt} and the matrix displays of this
document (\texttt{provenance/generated/matrices.md}). The program ends
with the line ``107 of 107 checks passed'' and returns a non-zero exit
status if any check fails. Step 3 writes
\texttt{data/matrix\_checks\_independent.txt}.

\begin{center}\rule{0.5\linewidth}{0.5pt}\end{center}

\hypertarget{part-1.-signature-44-the-matrices-gamma_1644a-of-44.nb}{%
\section{\texorpdfstring{Part 1. Signature 4+4: the matrices
\(\gamma_{1644}^A\) of
\texttt{4+4.nb}}{Part 1. Signature 4+4: the matrices \textbackslash gamma\_\{1644\}\^{}A of 4+4.nb}}\label{part-1.-signature-44-the-matrices-gamma_1644a-of-44.nb}}

\emph{Scope rule: in this Part \(\gamma^A\), \(\eta\), \(\omega\),
\(P^{L}\), \(P^{R}\), \(C^{\mp}\), \(\Sigma^{AB}\), \(\tau^A\),
\(\bar\tau^A\), \(\psi\), \(m\), \(K\), \(W\), \(\mathsf{h}\) stand for
the objects with the tag 1644.}

\hypertarget{definition}{%
\subsection{1.1 Definition}\label{definition}}

The construction of \texttt{4+4.nb} proceeds in four steps, from 4×4
matrices to 16×16 matrices. We repeat it with the notebook's cell
numbers (as counted by \texttt{tools/nb2txt.py}) so that each step can
be compared with the source.

\textbf{Step 1: six antisymmetric 4×4 matrices} (\texttt{4+4.nb} cells
25--36). Let \(\epsilon_{hpq4}\) be the sign of the permutation
\((h,p,q,4)\) of \((1,2,3,4)\), and zero if two indices coincide. For
\(h=1,2,3\) and \(p,q=1,\dots,4\),
\[(s_h)_{pq}=\epsilon_{hpq4}-\big(\delta_{p4}\delta_{qh}-\delta_{ph}\delta_{q4}\big),\qquad
(t_h)_{pq}=\epsilon_{hpq4}+\big(\delta_{p4}\delta_{qh}-\delta_{ph}\delta_{q4}\big).\]
The \(s_h\) are self-dual and the \(t_h\) anti-self-dual
(\texttt{s4by4{[}h{]}}, \texttt{t4by4{[}h{]}} in the notebook).

\textbf{Step 2: seven 8×8 matrices} (cells 74 and 88). With 4×4 blocks,
\[\tau^h=\begin{pmatrix}0&s_h\\s_h&0\end{pmatrix}\;(h=1,2,3),\qquad
\tau^{7-h}=\begin{pmatrix}0&t_h\\-t_h&0\end{pmatrix}\;(h=1,2,3),\qquad
\tau^7=\tau^1\tau^2\tau^3\tau^4\tau^5\tau^6,\] and \(\tau^0=1_8\). (The
notebook first calls these matrices \texttt{ζ{[}A{]}} and then sets
\texttt{τ{[}A{]}\ =\ ζ{[}A{]}}.) One finds
\(\tau^7=\mathrm{diag}(-1_4,1_4)\) and
\(\tau^1\tau^2\tau^3=\tau^4\tau^5\tau^6\tau^7=\sigma_{88}\) (check
\texttt{44-B1}).

\textbf{Step 3: the barred matrices} (cell 92). \(\bar\tau^0=1_8\) and
\(\bar\tau^A=\sigma_{88}\,(\tau^A)^{T}\sigma_{88}\) for \(A=1,\dots,7\),
where \({}^T\) denotes the transpose. They satisfy (check
\texttt{44-B2})
\[\tau^A\bar\tau^B+\tau^B\bar\tau^A=2\,\eta^{AB}\,1_8=\bar\tau^A\tau^B+\bar\tau^B\tau^A .\tag{1.1}\]

\textbf{Step 4: the eight 16×16 matrices} (cell 134). In 8×8 blocks,
\[\boxed{\;\gamma_{1644}^A=\begin{pmatrix}0&\bar\tau_{1644}^A\\ \tau_{1644}^A&0\end{pmatrix},\qquad A=0,1,\dots,7\;}\tag{1.2}\]
(check \texttt{44-B4}). These are the matrices \texttt{T16\^{}A{[}A{]}}
of \texttt{4+4.nb}. Because the notebook itself writes the index as a
superscript, and because these are the matrices that multiply
\(\partial_A\) in the Dirac operator, we regard \(A\) as a contravariant
index.

\hypertarget{the-eight-matrices}{%
\subsection{1.2 The eight matrices}\label{the-eight-matrices}}

\[
\gamma_{1644}^{0} = \left(\begin{array}{rrrrrrrr|rrrrrrrr}
\phantom{-}\cdot & \phantom{-}\cdot & \phantom{-}\cdot & \phantom{-}\cdot & \phantom{-}\cdot & \phantom{-}\cdot & \phantom{-}\cdot & \phantom{-}\cdot & \phantom{-}1 & \phantom{-}\cdot & \phantom{-}\cdot & \phantom{-}\cdot & \phantom{-}\cdot & \phantom{-}\cdot & \phantom{-}\cdot & \phantom{-}\cdot \\[-1.8pt]
\phantom{-}\cdot & \phantom{-}\cdot & \phantom{-}\cdot & \phantom{-}\cdot & \phantom{-}\cdot & \phantom{-}\cdot & \phantom{-}\cdot & \phantom{-}\cdot & \phantom{-}\cdot & \phantom{-}1 & \phantom{-}\cdot & \phantom{-}\cdot & \phantom{-}\cdot & \phantom{-}\cdot & \phantom{-}\cdot & \phantom{-}\cdot \\[-1.8pt]
\phantom{-}\cdot & \phantom{-}\cdot & \phantom{-}\cdot & \phantom{-}\cdot & \phantom{-}\cdot & \phantom{-}\cdot & \phantom{-}\cdot & \phantom{-}\cdot & \phantom{-}\cdot & \phantom{-}\cdot & \phantom{-}1 & \phantom{-}\cdot & \phantom{-}\cdot & \phantom{-}\cdot & \phantom{-}\cdot & \phantom{-}\cdot \\[-1.8pt]
\phantom{-}\cdot & \phantom{-}\cdot & \phantom{-}\cdot & \phantom{-}\cdot & \phantom{-}\cdot & \phantom{-}\cdot & \phantom{-}\cdot & \phantom{-}\cdot & \phantom{-}\cdot & \phantom{-}\cdot & \phantom{-}\cdot & \phantom{-}1 & \phantom{-}\cdot & \phantom{-}\cdot & \phantom{-}\cdot & \phantom{-}\cdot \\[-1.8pt]
\phantom{-}\cdot & \phantom{-}\cdot & \phantom{-}\cdot & \phantom{-}\cdot & \phantom{-}\cdot & \phantom{-}\cdot & \phantom{-}\cdot & \phantom{-}\cdot & \phantom{-}\cdot & \phantom{-}\cdot & \phantom{-}\cdot & \phantom{-}\cdot & \phantom{-}1 & \phantom{-}\cdot & \phantom{-}\cdot & \phantom{-}\cdot \\[-1.8pt]
\phantom{-}\cdot & \phantom{-}\cdot & \phantom{-}\cdot & \phantom{-}\cdot & \phantom{-}\cdot & \phantom{-}\cdot & \phantom{-}\cdot & \phantom{-}\cdot & \phantom{-}\cdot & \phantom{-}\cdot & \phantom{-}\cdot & \phantom{-}\cdot & \phantom{-}\cdot & \phantom{-}1 & \phantom{-}\cdot & \phantom{-}\cdot \\[-1.8pt]
\phantom{-}\cdot & \phantom{-}\cdot & \phantom{-}\cdot & \phantom{-}\cdot & \phantom{-}\cdot & \phantom{-}\cdot & \phantom{-}\cdot & \phantom{-}\cdot & \phantom{-}\cdot & \phantom{-}\cdot & \phantom{-}\cdot & \phantom{-}\cdot & \phantom{-}\cdot & \phantom{-}\cdot & \phantom{-}1 & \phantom{-}\cdot \\[-1.8pt]
\phantom{-}\cdot & \phantom{-}\cdot & \phantom{-}\cdot & \phantom{-}\cdot & \phantom{-}\cdot & \phantom{-}\cdot & \phantom{-}\cdot & \phantom{-}\cdot & \phantom{-}\cdot & \phantom{-}\cdot & \phantom{-}\cdot & \phantom{-}\cdot & \phantom{-}\cdot & \phantom{-}\cdot & \phantom{-}\cdot & \phantom{-}1 \\[-1.8pt]
\hline
\phantom{-}1 & \phantom{-}\cdot & \phantom{-}\cdot & \phantom{-}\cdot & \phantom{-}\cdot & \phantom{-}\cdot & \phantom{-}\cdot & \phantom{-}\cdot & \phantom{-}\cdot & \phantom{-}\cdot & \phantom{-}\cdot & \phantom{-}\cdot & \phantom{-}\cdot & \phantom{-}\cdot & \phantom{-}\cdot & \phantom{-}\cdot \\[-1.8pt]
\phantom{-}\cdot & \phantom{-}1 & \phantom{-}\cdot & \phantom{-}\cdot & \phantom{-}\cdot & \phantom{-}\cdot & \phantom{-}\cdot & \phantom{-}\cdot & \phantom{-}\cdot & \phantom{-}\cdot & \phantom{-}\cdot & \phantom{-}\cdot & \phantom{-}\cdot & \phantom{-}\cdot & \phantom{-}\cdot & \phantom{-}\cdot \\[-1.8pt]
\phantom{-}\cdot & \phantom{-}\cdot & \phantom{-}1 & \phantom{-}\cdot & \phantom{-}\cdot & \phantom{-}\cdot & \phantom{-}\cdot & \phantom{-}\cdot & \phantom{-}\cdot & \phantom{-}\cdot & \phantom{-}\cdot & \phantom{-}\cdot & \phantom{-}\cdot & \phantom{-}\cdot & \phantom{-}\cdot & \phantom{-}\cdot \\[-1.8pt]
\phantom{-}\cdot & \phantom{-}\cdot & \phantom{-}\cdot & \phantom{-}1 & \phantom{-}\cdot & \phantom{-}\cdot & \phantom{-}\cdot & \phantom{-}\cdot & \phantom{-}\cdot & \phantom{-}\cdot & \phantom{-}\cdot & \phantom{-}\cdot & \phantom{-}\cdot & \phantom{-}\cdot & \phantom{-}\cdot & \phantom{-}\cdot \\[-1.8pt]
\phantom{-}\cdot & \phantom{-}\cdot & \phantom{-}\cdot & \phantom{-}\cdot & \phantom{-}1 & \phantom{-}\cdot & \phantom{-}\cdot & \phantom{-}\cdot & \phantom{-}\cdot & \phantom{-}\cdot & \phantom{-}\cdot & \phantom{-}\cdot & \phantom{-}\cdot & \phantom{-}\cdot & \phantom{-}\cdot & \phantom{-}\cdot \\[-1.8pt]
\phantom{-}\cdot & \phantom{-}\cdot & \phantom{-}\cdot & \phantom{-}\cdot & \phantom{-}\cdot & \phantom{-}1 & \phantom{-}\cdot & \phantom{-}\cdot & \phantom{-}\cdot & \phantom{-}\cdot & \phantom{-}\cdot & \phantom{-}\cdot & \phantom{-}\cdot & \phantom{-}\cdot & \phantom{-}\cdot & \phantom{-}\cdot \\[-1.8pt]
\phantom{-}\cdot & \phantom{-}\cdot & \phantom{-}\cdot & \phantom{-}\cdot & \phantom{-}\cdot & \phantom{-}\cdot & \phantom{-}1 & \phantom{-}\cdot & \phantom{-}\cdot & \phantom{-}\cdot & \phantom{-}\cdot & \phantom{-}\cdot & \phantom{-}\cdot & \phantom{-}\cdot & \phantom{-}\cdot & \phantom{-}\cdot \\[-1.8pt]
\phantom{-}\cdot & \phantom{-}\cdot & \phantom{-}\cdot & \phantom{-}\cdot & \phantom{-}\cdot & \phantom{-}\cdot & \phantom{-}\cdot & \phantom{-}1 & \phantom{-}\cdot & \phantom{-}\cdot & \phantom{-}\cdot & \phantom{-}\cdot & \phantom{-}\cdot & \phantom{-}\cdot & \phantom{-}\cdot & \phantom{-}\cdot
\end{array}\right)
\]

\[
\gamma_{1644}^{1} = \left(\begin{array}{rrrrrrrr|rrrrrrrr}
\phantom{-}\cdot & \phantom{-}\cdot & \phantom{-}\cdot & \phantom{-}\cdot & \phantom{-}\cdot & \phantom{-}\cdot & \phantom{-}\cdot & \phantom{-}\cdot & \phantom{-}\cdot & \phantom{-}\cdot & \phantom{-}\cdot & \phantom{-}\cdot & \phantom{-}\cdot & \phantom{-}\cdot & \phantom{-}\cdot & -1 \\[-1.8pt]
\phantom{-}\cdot & \phantom{-}\cdot & \phantom{-}\cdot & \phantom{-}\cdot & \phantom{-}\cdot & \phantom{-}\cdot & \phantom{-}\cdot & \phantom{-}\cdot & \phantom{-}\cdot & \phantom{-}\cdot & \phantom{-}\cdot & \phantom{-}\cdot & \phantom{-}\cdot & \phantom{-}\cdot & -1 & \phantom{-}\cdot \\[-1.8pt]
\phantom{-}\cdot & \phantom{-}\cdot & \phantom{-}\cdot & \phantom{-}\cdot & \phantom{-}\cdot & \phantom{-}\cdot & \phantom{-}\cdot & \phantom{-}\cdot & \phantom{-}\cdot & \phantom{-}\cdot & \phantom{-}\cdot & \phantom{-}\cdot & \phantom{-}\cdot & \phantom{-}1 & \phantom{-}\cdot & \phantom{-}\cdot \\[-1.8pt]
\phantom{-}\cdot & \phantom{-}\cdot & \phantom{-}\cdot & \phantom{-}\cdot & \phantom{-}\cdot & \phantom{-}\cdot & \phantom{-}\cdot & \phantom{-}\cdot & \phantom{-}\cdot & \phantom{-}\cdot & \phantom{-}\cdot & \phantom{-}\cdot & \phantom{-}1 & \phantom{-}\cdot & \phantom{-}\cdot & \phantom{-}\cdot \\[-1.8pt]
\phantom{-}\cdot & \phantom{-}\cdot & \phantom{-}\cdot & \phantom{-}\cdot & \phantom{-}\cdot & \phantom{-}\cdot & \phantom{-}\cdot & \phantom{-}\cdot & \phantom{-}\cdot & \phantom{-}\cdot & \phantom{-}\cdot & -1 & \phantom{-}\cdot & \phantom{-}\cdot & \phantom{-}\cdot & \phantom{-}\cdot \\[-1.8pt]
\phantom{-}\cdot & \phantom{-}\cdot & \phantom{-}\cdot & \phantom{-}\cdot & \phantom{-}\cdot & \phantom{-}\cdot & \phantom{-}\cdot & \phantom{-}\cdot & \phantom{-}\cdot & \phantom{-}\cdot & -1 & \phantom{-}\cdot & \phantom{-}\cdot & \phantom{-}\cdot & \phantom{-}\cdot & \phantom{-}\cdot \\[-1.8pt]
\phantom{-}\cdot & \phantom{-}\cdot & \phantom{-}\cdot & \phantom{-}\cdot & \phantom{-}\cdot & \phantom{-}\cdot & \phantom{-}\cdot & \phantom{-}\cdot & \phantom{-}\cdot & \phantom{-}1 & \phantom{-}\cdot & \phantom{-}\cdot & \phantom{-}\cdot & \phantom{-}\cdot & \phantom{-}\cdot & \phantom{-}\cdot \\[-1.8pt]
\phantom{-}\cdot & \phantom{-}\cdot & \phantom{-}\cdot & \phantom{-}\cdot & \phantom{-}\cdot & \phantom{-}\cdot & \phantom{-}\cdot & \phantom{-}\cdot & \phantom{-}1 & \phantom{-}\cdot & \phantom{-}\cdot & \phantom{-}\cdot & \phantom{-}\cdot & \phantom{-}\cdot & \phantom{-}\cdot & \phantom{-}\cdot \\[-1.8pt]
\hline
\phantom{-}\cdot & \phantom{-}\cdot & \phantom{-}\cdot & \phantom{-}\cdot & \phantom{-}\cdot & \phantom{-}\cdot & \phantom{-}\cdot & \phantom{-}1 & \phantom{-}\cdot & \phantom{-}\cdot & \phantom{-}\cdot & \phantom{-}\cdot & \phantom{-}\cdot & \phantom{-}\cdot & \phantom{-}\cdot & \phantom{-}\cdot \\[-1.8pt]
\phantom{-}\cdot & \phantom{-}\cdot & \phantom{-}\cdot & \phantom{-}\cdot & \phantom{-}\cdot & \phantom{-}\cdot & \phantom{-}1 & \phantom{-}\cdot & \phantom{-}\cdot & \phantom{-}\cdot & \phantom{-}\cdot & \phantom{-}\cdot & \phantom{-}\cdot & \phantom{-}\cdot & \phantom{-}\cdot & \phantom{-}\cdot \\[-1.8pt]
\phantom{-}\cdot & \phantom{-}\cdot & \phantom{-}\cdot & \phantom{-}\cdot & \phantom{-}\cdot & -1 & \phantom{-}\cdot & \phantom{-}\cdot & \phantom{-}\cdot & \phantom{-}\cdot & \phantom{-}\cdot & \phantom{-}\cdot & \phantom{-}\cdot & \phantom{-}\cdot & \phantom{-}\cdot & \phantom{-}\cdot \\[-1.8pt]
\phantom{-}\cdot & \phantom{-}\cdot & \phantom{-}\cdot & \phantom{-}\cdot & -1 & \phantom{-}\cdot & \phantom{-}\cdot & \phantom{-}\cdot & \phantom{-}\cdot & \phantom{-}\cdot & \phantom{-}\cdot & \phantom{-}\cdot & \phantom{-}\cdot & \phantom{-}\cdot & \phantom{-}\cdot & \phantom{-}\cdot \\[-1.8pt]
\phantom{-}\cdot & \phantom{-}\cdot & \phantom{-}\cdot & \phantom{-}1 & \phantom{-}\cdot & \phantom{-}\cdot & \phantom{-}\cdot & \phantom{-}\cdot & \phantom{-}\cdot & \phantom{-}\cdot & \phantom{-}\cdot & \phantom{-}\cdot & \phantom{-}\cdot & \phantom{-}\cdot & \phantom{-}\cdot & \phantom{-}\cdot \\[-1.8pt]
\phantom{-}\cdot & \phantom{-}\cdot & \phantom{-}1 & \phantom{-}\cdot & \phantom{-}\cdot & \phantom{-}\cdot & \phantom{-}\cdot & \phantom{-}\cdot & \phantom{-}\cdot & \phantom{-}\cdot & \phantom{-}\cdot & \phantom{-}\cdot & \phantom{-}\cdot & \phantom{-}\cdot & \phantom{-}\cdot & \phantom{-}\cdot \\[-1.8pt]
\phantom{-}\cdot & -1 & \phantom{-}\cdot & \phantom{-}\cdot & \phantom{-}\cdot & \phantom{-}\cdot & \phantom{-}\cdot & \phantom{-}\cdot & \phantom{-}\cdot & \phantom{-}\cdot & \phantom{-}\cdot & \phantom{-}\cdot & \phantom{-}\cdot & \phantom{-}\cdot & \phantom{-}\cdot & \phantom{-}\cdot \\[-1.8pt]
-1 & \phantom{-}\cdot & \phantom{-}\cdot & \phantom{-}\cdot & \phantom{-}\cdot & \phantom{-}\cdot & \phantom{-}\cdot & \phantom{-}\cdot & \phantom{-}\cdot & \phantom{-}\cdot & \phantom{-}\cdot & \phantom{-}\cdot & \phantom{-}\cdot & \phantom{-}\cdot & \phantom{-}\cdot & \phantom{-}\cdot
\end{array}\right)
\]

\[
\gamma_{1644}^{2} = \left(\begin{array}{rrrrrrrr|rrrrrrrr}
\phantom{-}\cdot & \phantom{-}\cdot & \phantom{-}\cdot & \phantom{-}\cdot & \phantom{-}\cdot & \phantom{-}\cdot & \phantom{-}\cdot & \phantom{-}\cdot & \phantom{-}\cdot & \phantom{-}\cdot & \phantom{-}\cdot & \phantom{-}\cdot & \phantom{-}\cdot & \phantom{-}\cdot & \phantom{-}1 & \phantom{-}\cdot \\[-1.8pt]
\phantom{-}\cdot & \phantom{-}\cdot & \phantom{-}\cdot & \phantom{-}\cdot & \phantom{-}\cdot & \phantom{-}\cdot & \phantom{-}\cdot & \phantom{-}\cdot & \phantom{-}\cdot & \phantom{-}\cdot & \phantom{-}\cdot & \phantom{-}\cdot & \phantom{-}\cdot & \phantom{-}\cdot & \phantom{-}\cdot & -1 \\[-1.8pt]
\phantom{-}\cdot & \phantom{-}\cdot & \phantom{-}\cdot & \phantom{-}\cdot & \phantom{-}\cdot & \phantom{-}\cdot & \phantom{-}\cdot & \phantom{-}\cdot & \phantom{-}\cdot & \phantom{-}\cdot & \phantom{-}\cdot & \phantom{-}\cdot & -1 & \phantom{-}\cdot & \phantom{-}\cdot & \phantom{-}\cdot \\[-1.8pt]
\phantom{-}\cdot & \phantom{-}\cdot & \phantom{-}\cdot & \phantom{-}\cdot & \phantom{-}\cdot & \phantom{-}\cdot & \phantom{-}\cdot & \phantom{-}\cdot & \phantom{-}\cdot & \phantom{-}\cdot & \phantom{-}\cdot & \phantom{-}\cdot & \phantom{-}\cdot & \phantom{-}1 & \phantom{-}\cdot & \phantom{-}\cdot \\[-1.8pt]
\phantom{-}\cdot & \phantom{-}\cdot & \phantom{-}\cdot & \phantom{-}\cdot & \phantom{-}\cdot & \phantom{-}\cdot & \phantom{-}\cdot & \phantom{-}\cdot & \phantom{-}\cdot & \phantom{-}\cdot & \phantom{-}1 & \phantom{-}\cdot & \phantom{-}\cdot & \phantom{-}\cdot & \phantom{-}\cdot & \phantom{-}\cdot \\[-1.8pt]
\phantom{-}\cdot & \phantom{-}\cdot & \phantom{-}\cdot & \phantom{-}\cdot & \phantom{-}\cdot & \phantom{-}\cdot & \phantom{-}\cdot & \phantom{-}\cdot & \phantom{-}\cdot & \phantom{-}\cdot & \phantom{-}\cdot & -1 & \phantom{-}\cdot & \phantom{-}\cdot & \phantom{-}\cdot & \phantom{-}\cdot \\[-1.8pt]
\phantom{-}\cdot & \phantom{-}\cdot & \phantom{-}\cdot & \phantom{-}\cdot & \phantom{-}\cdot & \phantom{-}\cdot & \phantom{-}\cdot & \phantom{-}\cdot & -1 & \phantom{-}\cdot & \phantom{-}\cdot & \phantom{-}\cdot & \phantom{-}\cdot & \phantom{-}\cdot & \phantom{-}\cdot & \phantom{-}\cdot \\[-1.8pt]
\phantom{-}\cdot & \phantom{-}\cdot & \phantom{-}\cdot & \phantom{-}\cdot & \phantom{-}\cdot & \phantom{-}\cdot & \phantom{-}\cdot & \phantom{-}\cdot & \phantom{-}\cdot & \phantom{-}1 & \phantom{-}\cdot & \phantom{-}\cdot & \phantom{-}\cdot & \phantom{-}\cdot & \phantom{-}\cdot & \phantom{-}\cdot \\[-1.8pt]
\hline
\phantom{-}\cdot & \phantom{-}\cdot & \phantom{-}\cdot & \phantom{-}\cdot & \phantom{-}\cdot & \phantom{-}\cdot & -1 & \phantom{-}\cdot & \phantom{-}\cdot & \phantom{-}\cdot & \phantom{-}\cdot & \phantom{-}\cdot & \phantom{-}\cdot & \phantom{-}\cdot & \phantom{-}\cdot & \phantom{-}\cdot \\[-1.8pt]
\phantom{-}\cdot & \phantom{-}\cdot & \phantom{-}\cdot & \phantom{-}\cdot & \phantom{-}\cdot & \phantom{-}\cdot & \phantom{-}\cdot & \phantom{-}1 & \phantom{-}\cdot & \phantom{-}\cdot & \phantom{-}\cdot & \phantom{-}\cdot & \phantom{-}\cdot & \phantom{-}\cdot & \phantom{-}\cdot & \phantom{-}\cdot \\[-1.8pt]
\phantom{-}\cdot & \phantom{-}\cdot & \phantom{-}\cdot & \phantom{-}\cdot & \phantom{-}1 & \phantom{-}\cdot & \phantom{-}\cdot & \phantom{-}\cdot & \phantom{-}\cdot & \phantom{-}\cdot & \phantom{-}\cdot & \phantom{-}\cdot & \phantom{-}\cdot & \phantom{-}\cdot & \phantom{-}\cdot & \phantom{-}\cdot \\[-1.8pt]
\phantom{-}\cdot & \phantom{-}\cdot & \phantom{-}\cdot & \phantom{-}\cdot & \phantom{-}\cdot & -1 & \phantom{-}\cdot & \phantom{-}\cdot & \phantom{-}\cdot & \phantom{-}\cdot & \phantom{-}\cdot & \phantom{-}\cdot & \phantom{-}\cdot & \phantom{-}\cdot & \phantom{-}\cdot & \phantom{-}\cdot \\[-1.8pt]
\phantom{-}\cdot & \phantom{-}\cdot & -1 & \phantom{-}\cdot & \phantom{-}\cdot & \phantom{-}\cdot & \phantom{-}\cdot & \phantom{-}\cdot & \phantom{-}\cdot & \phantom{-}\cdot & \phantom{-}\cdot & \phantom{-}\cdot & \phantom{-}\cdot & \phantom{-}\cdot & \phantom{-}\cdot & \phantom{-}\cdot \\[-1.8pt]
\phantom{-}\cdot & \phantom{-}\cdot & \phantom{-}\cdot & \phantom{-}1 & \phantom{-}\cdot & \phantom{-}\cdot & \phantom{-}\cdot & \phantom{-}\cdot & \phantom{-}\cdot & \phantom{-}\cdot & \phantom{-}\cdot & \phantom{-}\cdot & \phantom{-}\cdot & \phantom{-}\cdot & \phantom{-}\cdot & \phantom{-}\cdot \\[-1.8pt]
\phantom{-}1 & \phantom{-}\cdot & \phantom{-}\cdot & \phantom{-}\cdot & \phantom{-}\cdot & \phantom{-}\cdot & \phantom{-}\cdot & \phantom{-}\cdot & \phantom{-}\cdot & \phantom{-}\cdot & \phantom{-}\cdot & \phantom{-}\cdot & \phantom{-}\cdot & \phantom{-}\cdot & \phantom{-}\cdot & \phantom{-}\cdot \\[-1.8pt]
\phantom{-}\cdot & -1 & \phantom{-}\cdot & \phantom{-}\cdot & \phantom{-}\cdot & \phantom{-}\cdot & \phantom{-}\cdot & \phantom{-}\cdot & \phantom{-}\cdot & \phantom{-}\cdot & \phantom{-}\cdot & \phantom{-}\cdot & \phantom{-}\cdot & \phantom{-}\cdot & \phantom{-}\cdot & \phantom{-}\cdot
\end{array}\right)
\]

\[
\gamma_{1644}^{3} = \left(\begin{array}{rrrrrrrr|rrrrrrrr}
\phantom{-}\cdot & \phantom{-}\cdot & \phantom{-}\cdot & \phantom{-}\cdot & \phantom{-}\cdot & \phantom{-}\cdot & \phantom{-}\cdot & \phantom{-}\cdot & \phantom{-}\cdot & \phantom{-}\cdot & \phantom{-}\cdot & \phantom{-}\cdot & \phantom{-}\cdot & -1 & \phantom{-}\cdot & \phantom{-}\cdot \\[-1.8pt]
\phantom{-}\cdot & \phantom{-}\cdot & \phantom{-}\cdot & \phantom{-}\cdot & \phantom{-}\cdot & \phantom{-}\cdot & \phantom{-}\cdot & \phantom{-}\cdot & \phantom{-}\cdot & \phantom{-}\cdot & \phantom{-}\cdot & \phantom{-}\cdot & \phantom{-}1 & \phantom{-}\cdot & \phantom{-}\cdot & \phantom{-}\cdot \\[-1.8pt]
\phantom{-}\cdot & \phantom{-}\cdot & \phantom{-}\cdot & \phantom{-}\cdot & \phantom{-}\cdot & \phantom{-}\cdot & \phantom{-}\cdot & \phantom{-}\cdot & \phantom{-}\cdot & \phantom{-}\cdot & \phantom{-}\cdot & \phantom{-}\cdot & \phantom{-}\cdot & \phantom{-}\cdot & \phantom{-}\cdot & -1 \\[-1.8pt]
\phantom{-}\cdot & \phantom{-}\cdot & \phantom{-}\cdot & \phantom{-}\cdot & \phantom{-}\cdot & \phantom{-}\cdot & \phantom{-}\cdot & \phantom{-}\cdot & \phantom{-}\cdot & \phantom{-}\cdot & \phantom{-}\cdot & \phantom{-}\cdot & \phantom{-}\cdot & \phantom{-}\cdot & \phantom{-}1 & \phantom{-}\cdot \\[-1.8pt]
\phantom{-}\cdot & \phantom{-}\cdot & \phantom{-}\cdot & \phantom{-}\cdot & \phantom{-}\cdot & \phantom{-}\cdot & \phantom{-}\cdot & \phantom{-}\cdot & \phantom{-}\cdot & -1 & \phantom{-}\cdot & \phantom{-}\cdot & \phantom{-}\cdot & \phantom{-}\cdot & \phantom{-}\cdot & \phantom{-}\cdot \\[-1.8pt]
\phantom{-}\cdot & \phantom{-}\cdot & \phantom{-}\cdot & \phantom{-}\cdot & \phantom{-}\cdot & \phantom{-}\cdot & \phantom{-}\cdot & \phantom{-}\cdot & \phantom{-}1 & \phantom{-}\cdot & \phantom{-}\cdot & \phantom{-}\cdot & \phantom{-}\cdot & \phantom{-}\cdot & \phantom{-}\cdot & \phantom{-}\cdot \\[-1.8pt]
\phantom{-}\cdot & \phantom{-}\cdot & \phantom{-}\cdot & \phantom{-}\cdot & \phantom{-}\cdot & \phantom{-}\cdot & \phantom{-}\cdot & \phantom{-}\cdot & \phantom{-}\cdot & \phantom{-}\cdot & \phantom{-}\cdot & -1 & \phantom{-}\cdot & \phantom{-}\cdot & \phantom{-}\cdot & \phantom{-}\cdot \\[-1.8pt]
\phantom{-}\cdot & \phantom{-}\cdot & \phantom{-}\cdot & \phantom{-}\cdot & \phantom{-}\cdot & \phantom{-}\cdot & \phantom{-}\cdot & \phantom{-}\cdot & \phantom{-}\cdot & \phantom{-}\cdot & \phantom{-}1 & \phantom{-}\cdot & \phantom{-}\cdot & \phantom{-}\cdot & \phantom{-}\cdot & \phantom{-}\cdot \\[-1.8pt]
\hline
\phantom{-}\cdot & \phantom{-}\cdot & \phantom{-}\cdot & \phantom{-}\cdot & \phantom{-}\cdot & \phantom{-}1 & \phantom{-}\cdot & \phantom{-}\cdot & \phantom{-}\cdot & \phantom{-}\cdot & \phantom{-}\cdot & \phantom{-}\cdot & \phantom{-}\cdot & \phantom{-}\cdot & \phantom{-}\cdot & \phantom{-}\cdot \\[-1.8pt]
\phantom{-}\cdot & \phantom{-}\cdot & \phantom{-}\cdot & \phantom{-}\cdot & -1 & \phantom{-}\cdot & \phantom{-}\cdot & \phantom{-}\cdot & \phantom{-}\cdot & \phantom{-}\cdot & \phantom{-}\cdot & \phantom{-}\cdot & \phantom{-}\cdot & \phantom{-}\cdot & \phantom{-}\cdot & \phantom{-}\cdot \\[-1.8pt]
\phantom{-}\cdot & \phantom{-}\cdot & \phantom{-}\cdot & \phantom{-}\cdot & \phantom{-}\cdot & \phantom{-}\cdot & \phantom{-}\cdot & \phantom{-}1 & \phantom{-}\cdot & \phantom{-}\cdot & \phantom{-}\cdot & \phantom{-}\cdot & \phantom{-}\cdot & \phantom{-}\cdot & \phantom{-}\cdot & \phantom{-}\cdot \\[-1.8pt]
\phantom{-}\cdot & \phantom{-}\cdot & \phantom{-}\cdot & \phantom{-}\cdot & \phantom{-}\cdot & \phantom{-}\cdot & -1 & \phantom{-}\cdot & \phantom{-}\cdot & \phantom{-}\cdot & \phantom{-}\cdot & \phantom{-}\cdot & \phantom{-}\cdot & \phantom{-}\cdot & \phantom{-}\cdot & \phantom{-}\cdot \\[-1.8pt]
\phantom{-}\cdot & \phantom{-}1 & \phantom{-}\cdot & \phantom{-}\cdot & \phantom{-}\cdot & \phantom{-}\cdot & \phantom{-}\cdot & \phantom{-}\cdot & \phantom{-}\cdot & \phantom{-}\cdot & \phantom{-}\cdot & \phantom{-}\cdot & \phantom{-}\cdot & \phantom{-}\cdot & \phantom{-}\cdot & \phantom{-}\cdot \\[-1.8pt]
-1 & \phantom{-}\cdot & \phantom{-}\cdot & \phantom{-}\cdot & \phantom{-}\cdot & \phantom{-}\cdot & \phantom{-}\cdot & \phantom{-}\cdot & \phantom{-}\cdot & \phantom{-}\cdot & \phantom{-}\cdot & \phantom{-}\cdot & \phantom{-}\cdot & \phantom{-}\cdot & \phantom{-}\cdot & \phantom{-}\cdot \\[-1.8pt]
\phantom{-}\cdot & \phantom{-}\cdot & \phantom{-}\cdot & \phantom{-}1 & \phantom{-}\cdot & \phantom{-}\cdot & \phantom{-}\cdot & \phantom{-}\cdot & \phantom{-}\cdot & \phantom{-}\cdot & \phantom{-}\cdot & \phantom{-}\cdot & \phantom{-}\cdot & \phantom{-}\cdot & \phantom{-}\cdot & \phantom{-}\cdot \\[-1.8pt]
\phantom{-}\cdot & \phantom{-}\cdot & -1 & \phantom{-}\cdot & \phantom{-}\cdot & \phantom{-}\cdot & \phantom{-}\cdot & \phantom{-}\cdot & \phantom{-}\cdot & \phantom{-}\cdot & \phantom{-}\cdot & \phantom{-}\cdot & \phantom{-}\cdot & \phantom{-}\cdot & \phantom{-}\cdot & \phantom{-}\cdot
\end{array}\right)
\]

\[
\gamma_{1644}^{4} = \left(\begin{array}{rrrrrrrr|rrrrrrrr}
\phantom{-}\cdot & \phantom{-}\cdot & \phantom{-}\cdot & \phantom{-}\cdot & \phantom{-}\cdot & \phantom{-}\cdot & \phantom{-}\cdot & \phantom{-}\cdot & \phantom{-}\cdot & \phantom{-}\cdot & \phantom{-}\cdot & \phantom{-}\cdot & \phantom{-}\cdot & -1 & \phantom{-}\cdot & \phantom{-}\cdot \\[-1.8pt]
\phantom{-}\cdot & \phantom{-}\cdot & \phantom{-}\cdot & \phantom{-}\cdot & \phantom{-}\cdot & \phantom{-}\cdot & \phantom{-}\cdot & \phantom{-}\cdot & \phantom{-}\cdot & \phantom{-}\cdot & \phantom{-}\cdot & \phantom{-}\cdot & \phantom{-}1 & \phantom{-}\cdot & \phantom{-}\cdot & \phantom{-}\cdot \\[-1.8pt]
\phantom{-}\cdot & \phantom{-}\cdot & \phantom{-}\cdot & \phantom{-}\cdot & \phantom{-}\cdot & \phantom{-}\cdot & \phantom{-}\cdot & \phantom{-}\cdot & \phantom{-}\cdot & \phantom{-}\cdot & \phantom{-}\cdot & \phantom{-}\cdot & \phantom{-}\cdot & \phantom{-}\cdot & \phantom{-}\cdot & \phantom{-}1 \\[-1.8pt]
\phantom{-}\cdot & \phantom{-}\cdot & \phantom{-}\cdot & \phantom{-}\cdot & \phantom{-}\cdot & \phantom{-}\cdot & \phantom{-}\cdot & \phantom{-}\cdot & \phantom{-}\cdot & \phantom{-}\cdot & \phantom{-}\cdot & \phantom{-}\cdot & \phantom{-}\cdot & \phantom{-}\cdot & -1 & \phantom{-}\cdot \\[-1.8pt]
\phantom{-}\cdot & \phantom{-}\cdot & \phantom{-}\cdot & \phantom{-}\cdot & \phantom{-}\cdot & \phantom{-}\cdot & \phantom{-}\cdot & \phantom{-}\cdot & \phantom{-}\cdot & \phantom{-}1 & \phantom{-}\cdot & \phantom{-}\cdot & \phantom{-}\cdot & \phantom{-}\cdot & \phantom{-}\cdot & \phantom{-}\cdot \\[-1.8pt]
\phantom{-}\cdot & \phantom{-}\cdot & \phantom{-}\cdot & \phantom{-}\cdot & \phantom{-}\cdot & \phantom{-}\cdot & \phantom{-}\cdot & \phantom{-}\cdot & -1 & \phantom{-}\cdot & \phantom{-}\cdot & \phantom{-}\cdot & \phantom{-}\cdot & \phantom{-}\cdot & \phantom{-}\cdot & \phantom{-}\cdot \\[-1.8pt]
\phantom{-}\cdot & \phantom{-}\cdot & \phantom{-}\cdot & \phantom{-}\cdot & \phantom{-}\cdot & \phantom{-}\cdot & \phantom{-}\cdot & \phantom{-}\cdot & \phantom{-}\cdot & \phantom{-}\cdot & \phantom{-}\cdot & -1 & \phantom{-}\cdot & \phantom{-}\cdot & \phantom{-}\cdot & \phantom{-}\cdot \\[-1.8pt]
\phantom{-}\cdot & \phantom{-}\cdot & \phantom{-}\cdot & \phantom{-}\cdot & \phantom{-}\cdot & \phantom{-}\cdot & \phantom{-}\cdot & \phantom{-}\cdot & \phantom{-}\cdot & \phantom{-}\cdot & \phantom{-}1 & \phantom{-}\cdot & \phantom{-}\cdot & \phantom{-}\cdot & \phantom{-}\cdot & \phantom{-}\cdot \\[-1.8pt]
\hline
\phantom{-}\cdot & \phantom{-}\cdot & \phantom{-}\cdot & \phantom{-}\cdot & \phantom{-}\cdot & \phantom{-}1 & \phantom{-}\cdot & \phantom{-}\cdot & \phantom{-}\cdot & \phantom{-}\cdot & \phantom{-}\cdot & \phantom{-}\cdot & \phantom{-}\cdot & \phantom{-}\cdot & \phantom{-}\cdot & \phantom{-}\cdot \\[-1.8pt]
\phantom{-}\cdot & \phantom{-}\cdot & \phantom{-}\cdot & \phantom{-}\cdot & -1 & \phantom{-}\cdot & \phantom{-}\cdot & \phantom{-}\cdot & \phantom{-}\cdot & \phantom{-}\cdot & \phantom{-}\cdot & \phantom{-}\cdot & \phantom{-}\cdot & \phantom{-}\cdot & \phantom{-}\cdot & \phantom{-}\cdot \\[-1.8pt]
\phantom{-}\cdot & \phantom{-}\cdot & \phantom{-}\cdot & \phantom{-}\cdot & \phantom{-}\cdot & \phantom{-}\cdot & \phantom{-}\cdot & -1 & \phantom{-}\cdot & \phantom{-}\cdot & \phantom{-}\cdot & \phantom{-}\cdot & \phantom{-}\cdot & \phantom{-}\cdot & \phantom{-}\cdot & \phantom{-}\cdot \\[-1.8pt]
\phantom{-}\cdot & \phantom{-}\cdot & \phantom{-}\cdot & \phantom{-}\cdot & \phantom{-}\cdot & \phantom{-}\cdot & \phantom{-}1 & \phantom{-}\cdot & \phantom{-}\cdot & \phantom{-}\cdot & \phantom{-}\cdot & \phantom{-}\cdot & \phantom{-}\cdot & \phantom{-}\cdot & \phantom{-}\cdot & \phantom{-}\cdot \\[-1.8pt]
\phantom{-}\cdot & -1 & \phantom{-}\cdot & \phantom{-}\cdot & \phantom{-}\cdot & \phantom{-}\cdot & \phantom{-}\cdot & \phantom{-}\cdot & \phantom{-}\cdot & \phantom{-}\cdot & \phantom{-}\cdot & \phantom{-}\cdot & \phantom{-}\cdot & \phantom{-}\cdot & \phantom{-}\cdot & \phantom{-}\cdot \\[-1.8pt]
\phantom{-}1 & \phantom{-}\cdot & \phantom{-}\cdot & \phantom{-}\cdot & \phantom{-}\cdot & \phantom{-}\cdot & \phantom{-}\cdot & \phantom{-}\cdot & \phantom{-}\cdot & \phantom{-}\cdot & \phantom{-}\cdot & \phantom{-}\cdot & \phantom{-}\cdot & \phantom{-}\cdot & \phantom{-}\cdot & \phantom{-}\cdot \\[-1.8pt]
\phantom{-}\cdot & \phantom{-}\cdot & \phantom{-}\cdot & \phantom{-}1 & \phantom{-}\cdot & \phantom{-}\cdot & \phantom{-}\cdot & \phantom{-}\cdot & \phantom{-}\cdot & \phantom{-}\cdot & \phantom{-}\cdot & \phantom{-}\cdot & \phantom{-}\cdot & \phantom{-}\cdot & \phantom{-}\cdot & \phantom{-}\cdot \\[-1.8pt]
\phantom{-}\cdot & \phantom{-}\cdot & -1 & \phantom{-}\cdot & \phantom{-}\cdot & \phantom{-}\cdot & \phantom{-}\cdot & \phantom{-}\cdot & \phantom{-}\cdot & \phantom{-}\cdot & \phantom{-}\cdot & \phantom{-}\cdot & \phantom{-}\cdot & \phantom{-}\cdot & \phantom{-}\cdot & \phantom{-}\cdot
\end{array}\right)
\]

\[
\gamma_{1644}^{5} = \left(\begin{array}{rrrrrrrr|rrrrrrrr}
\phantom{-}\cdot & \phantom{-}\cdot & \phantom{-}\cdot & \phantom{-}\cdot & \phantom{-}\cdot & \phantom{-}\cdot & \phantom{-}\cdot & \phantom{-}\cdot & \phantom{-}\cdot & \phantom{-}\cdot & \phantom{-}\cdot & \phantom{-}\cdot & \phantom{-}\cdot & \phantom{-}\cdot & \phantom{-}1 & \phantom{-}\cdot \\[-1.8pt]
\phantom{-}\cdot & \phantom{-}\cdot & \phantom{-}\cdot & \phantom{-}\cdot & \phantom{-}\cdot & \phantom{-}\cdot & \phantom{-}\cdot & \phantom{-}\cdot & \phantom{-}\cdot & \phantom{-}\cdot & \phantom{-}\cdot & \phantom{-}\cdot & \phantom{-}\cdot & \phantom{-}\cdot & \phantom{-}\cdot & \phantom{-}1 \\[-1.8pt]
\phantom{-}\cdot & \phantom{-}\cdot & \phantom{-}\cdot & \phantom{-}\cdot & \phantom{-}\cdot & \phantom{-}\cdot & \phantom{-}\cdot & \phantom{-}\cdot & \phantom{-}\cdot & \phantom{-}\cdot & \phantom{-}\cdot & \phantom{-}\cdot & -1 & \phantom{-}\cdot & \phantom{-}\cdot & \phantom{-}\cdot \\[-1.8pt]
\phantom{-}\cdot & \phantom{-}\cdot & \phantom{-}\cdot & \phantom{-}\cdot & \phantom{-}\cdot & \phantom{-}\cdot & \phantom{-}\cdot & \phantom{-}\cdot & \phantom{-}\cdot & \phantom{-}\cdot & \phantom{-}\cdot & \phantom{-}\cdot & \phantom{-}\cdot & -1 & \phantom{-}\cdot & \phantom{-}\cdot \\[-1.8pt]
\phantom{-}\cdot & \phantom{-}\cdot & \phantom{-}\cdot & \phantom{-}\cdot & \phantom{-}\cdot & \phantom{-}\cdot & \phantom{-}\cdot & \phantom{-}\cdot & \phantom{-}\cdot & \phantom{-}\cdot & -1 & \phantom{-}\cdot & \phantom{-}\cdot & \phantom{-}\cdot & \phantom{-}\cdot & \phantom{-}\cdot \\[-1.8pt]
\phantom{-}\cdot & \phantom{-}\cdot & \phantom{-}\cdot & \phantom{-}\cdot & \phantom{-}\cdot & \phantom{-}\cdot & \phantom{-}\cdot & \phantom{-}\cdot & \phantom{-}\cdot & \phantom{-}\cdot & \phantom{-}\cdot & -1 & \phantom{-}\cdot & \phantom{-}\cdot & \phantom{-}\cdot & \phantom{-}\cdot \\[-1.8pt]
\phantom{-}\cdot & \phantom{-}\cdot & \phantom{-}\cdot & \phantom{-}\cdot & \phantom{-}\cdot & \phantom{-}\cdot & \phantom{-}\cdot & \phantom{-}\cdot & \phantom{-}1 & \phantom{-}\cdot & \phantom{-}\cdot & \phantom{-}\cdot & \phantom{-}\cdot & \phantom{-}\cdot & \phantom{-}\cdot & \phantom{-}\cdot \\[-1.8pt]
\phantom{-}\cdot & \phantom{-}\cdot & \phantom{-}\cdot & \phantom{-}\cdot & \phantom{-}\cdot & \phantom{-}\cdot & \phantom{-}\cdot & \phantom{-}\cdot & \phantom{-}\cdot & \phantom{-}1 & \phantom{-}\cdot & \phantom{-}\cdot & \phantom{-}\cdot & \phantom{-}\cdot & \phantom{-}\cdot & \phantom{-}\cdot \\[-1.8pt]
\hline
\phantom{-}\cdot & \phantom{-}\cdot & \phantom{-}\cdot & \phantom{-}\cdot & \phantom{-}\cdot & \phantom{-}\cdot & -1 & \phantom{-}\cdot & \phantom{-}\cdot & \phantom{-}\cdot & \phantom{-}\cdot & \phantom{-}\cdot & \phantom{-}\cdot & \phantom{-}\cdot & \phantom{-}\cdot & \phantom{-}\cdot \\[-1.8pt]
\phantom{-}\cdot & \phantom{-}\cdot & \phantom{-}\cdot & \phantom{-}\cdot & \phantom{-}\cdot & \phantom{-}\cdot & \phantom{-}\cdot & -1 & \phantom{-}\cdot & \phantom{-}\cdot & \phantom{-}\cdot & \phantom{-}\cdot & \phantom{-}\cdot & \phantom{-}\cdot & \phantom{-}\cdot & \phantom{-}\cdot \\[-1.8pt]
\phantom{-}\cdot & \phantom{-}\cdot & \phantom{-}\cdot & \phantom{-}\cdot & \phantom{-}1 & \phantom{-}\cdot & \phantom{-}\cdot & \phantom{-}\cdot & \phantom{-}\cdot & \phantom{-}\cdot & \phantom{-}\cdot & \phantom{-}\cdot & \phantom{-}\cdot & \phantom{-}\cdot & \phantom{-}\cdot & \phantom{-}\cdot \\[-1.8pt]
\phantom{-}\cdot & \phantom{-}\cdot & \phantom{-}\cdot & \phantom{-}\cdot & \phantom{-}\cdot & \phantom{-}1 & \phantom{-}\cdot & \phantom{-}\cdot & \phantom{-}\cdot & \phantom{-}\cdot & \phantom{-}\cdot & \phantom{-}\cdot & \phantom{-}\cdot & \phantom{-}\cdot & \phantom{-}\cdot & \phantom{-}\cdot \\[-1.8pt]
\phantom{-}\cdot & \phantom{-}\cdot & \phantom{-}1 & \phantom{-}\cdot & \phantom{-}\cdot & \phantom{-}\cdot & \phantom{-}\cdot & \phantom{-}\cdot & \phantom{-}\cdot & \phantom{-}\cdot & \phantom{-}\cdot & \phantom{-}\cdot & \phantom{-}\cdot & \phantom{-}\cdot & \phantom{-}\cdot & \phantom{-}\cdot \\[-1.8pt]
\phantom{-}\cdot & \phantom{-}\cdot & \phantom{-}\cdot & \phantom{-}1 & \phantom{-}\cdot & \phantom{-}\cdot & \phantom{-}\cdot & \phantom{-}\cdot & \phantom{-}\cdot & \phantom{-}\cdot & \phantom{-}\cdot & \phantom{-}\cdot & \phantom{-}\cdot & \phantom{-}\cdot & \phantom{-}\cdot & \phantom{-}\cdot \\[-1.8pt]
-1 & \phantom{-}\cdot & \phantom{-}\cdot & \phantom{-}\cdot & \phantom{-}\cdot & \phantom{-}\cdot & \phantom{-}\cdot & \phantom{-}\cdot & \phantom{-}\cdot & \phantom{-}\cdot & \phantom{-}\cdot & \phantom{-}\cdot & \phantom{-}\cdot & \phantom{-}\cdot & \phantom{-}\cdot & \phantom{-}\cdot \\[-1.8pt]
\phantom{-}\cdot & -1 & \phantom{-}\cdot & \phantom{-}\cdot & \phantom{-}\cdot & \phantom{-}\cdot & \phantom{-}\cdot & \phantom{-}\cdot & \phantom{-}\cdot & \phantom{-}\cdot & \phantom{-}\cdot & \phantom{-}\cdot & \phantom{-}\cdot & \phantom{-}\cdot & \phantom{-}\cdot & \phantom{-}\cdot
\end{array}\right)
\]

\[
\gamma_{1644}^{6} = \left(\begin{array}{rrrrrrrr|rrrrrrrr}
\phantom{-}\cdot & \phantom{-}\cdot & \phantom{-}\cdot & \phantom{-}\cdot & \phantom{-}\cdot & \phantom{-}\cdot & \phantom{-}\cdot & \phantom{-}\cdot & \phantom{-}\cdot & \phantom{-}\cdot & \phantom{-}\cdot & \phantom{-}\cdot & \phantom{-}\cdot & \phantom{-}\cdot & \phantom{-}\cdot & \phantom{-}1 \\[-1.8pt]
\phantom{-}\cdot & \phantom{-}\cdot & \phantom{-}\cdot & \phantom{-}\cdot & \phantom{-}\cdot & \phantom{-}\cdot & \phantom{-}\cdot & \phantom{-}\cdot & \phantom{-}\cdot & \phantom{-}\cdot & \phantom{-}\cdot & \phantom{-}\cdot & \phantom{-}\cdot & \phantom{-}\cdot & -1 & \phantom{-}\cdot \\[-1.8pt]
\phantom{-}\cdot & \phantom{-}\cdot & \phantom{-}\cdot & \phantom{-}\cdot & \phantom{-}\cdot & \phantom{-}\cdot & \phantom{-}\cdot & \phantom{-}\cdot & \phantom{-}\cdot & \phantom{-}\cdot & \phantom{-}\cdot & \phantom{-}\cdot & \phantom{-}\cdot & \phantom{-}1 & \phantom{-}\cdot & \phantom{-}\cdot \\[-1.8pt]
\phantom{-}\cdot & \phantom{-}\cdot & \phantom{-}\cdot & \phantom{-}\cdot & \phantom{-}\cdot & \phantom{-}\cdot & \phantom{-}\cdot & \phantom{-}\cdot & \phantom{-}\cdot & \phantom{-}\cdot & \phantom{-}\cdot & \phantom{-}\cdot & -1 & \phantom{-}\cdot & \phantom{-}\cdot & \phantom{-}\cdot \\[-1.8pt]
\phantom{-}\cdot & \phantom{-}\cdot & \phantom{-}\cdot & \phantom{-}\cdot & \phantom{-}\cdot & \phantom{-}\cdot & \phantom{-}\cdot & \phantom{-}\cdot & \phantom{-}\cdot & \phantom{-}\cdot & \phantom{-}\cdot & -1 & \phantom{-}\cdot & \phantom{-}\cdot & \phantom{-}\cdot & \phantom{-}\cdot \\[-1.8pt]
\phantom{-}\cdot & \phantom{-}\cdot & \phantom{-}\cdot & \phantom{-}\cdot & \phantom{-}\cdot & \phantom{-}\cdot & \phantom{-}\cdot & \phantom{-}\cdot & \phantom{-}\cdot & \phantom{-}\cdot & \phantom{-}1 & \phantom{-}\cdot & \phantom{-}\cdot & \phantom{-}\cdot & \phantom{-}\cdot & \phantom{-}\cdot \\[-1.8pt]
\phantom{-}\cdot & \phantom{-}\cdot & \phantom{-}\cdot & \phantom{-}\cdot & \phantom{-}\cdot & \phantom{-}\cdot & \phantom{-}\cdot & \phantom{-}\cdot & \phantom{-}\cdot & -1 & \phantom{-}\cdot & \phantom{-}\cdot & \phantom{-}\cdot & \phantom{-}\cdot & \phantom{-}\cdot & \phantom{-}\cdot \\[-1.8pt]
\phantom{-}\cdot & \phantom{-}\cdot & \phantom{-}\cdot & \phantom{-}\cdot & \phantom{-}\cdot & \phantom{-}\cdot & \phantom{-}\cdot & \phantom{-}\cdot & \phantom{-}1 & \phantom{-}\cdot & \phantom{-}\cdot & \phantom{-}\cdot & \phantom{-}\cdot & \phantom{-}\cdot & \phantom{-}\cdot & \phantom{-}\cdot \\[-1.8pt]
\hline
\phantom{-}\cdot & \phantom{-}\cdot & \phantom{-}\cdot & \phantom{-}\cdot & \phantom{-}\cdot & \phantom{-}\cdot & \phantom{-}\cdot & -1 & \phantom{-}\cdot & \phantom{-}\cdot & \phantom{-}\cdot & \phantom{-}\cdot & \phantom{-}\cdot & \phantom{-}\cdot & \phantom{-}\cdot & \phantom{-}\cdot \\[-1.8pt]
\phantom{-}\cdot & \phantom{-}\cdot & \phantom{-}\cdot & \phantom{-}\cdot & \phantom{-}\cdot & \phantom{-}\cdot & \phantom{-}1 & \phantom{-}\cdot & \phantom{-}\cdot & \phantom{-}\cdot & \phantom{-}\cdot & \phantom{-}\cdot & \phantom{-}\cdot & \phantom{-}\cdot & \phantom{-}\cdot & \phantom{-}\cdot \\[-1.8pt]
\phantom{-}\cdot & \phantom{-}\cdot & \phantom{-}\cdot & \phantom{-}\cdot & \phantom{-}\cdot & -1 & \phantom{-}\cdot & \phantom{-}\cdot & \phantom{-}\cdot & \phantom{-}\cdot & \phantom{-}\cdot & \phantom{-}\cdot & \phantom{-}\cdot & \phantom{-}\cdot & \phantom{-}\cdot & \phantom{-}\cdot \\[-1.8pt]
\phantom{-}\cdot & \phantom{-}\cdot & \phantom{-}\cdot & \phantom{-}\cdot & \phantom{-}1 & \phantom{-}\cdot & \phantom{-}\cdot & \phantom{-}\cdot & \phantom{-}\cdot & \phantom{-}\cdot & \phantom{-}\cdot & \phantom{-}\cdot & \phantom{-}\cdot & \phantom{-}\cdot & \phantom{-}\cdot & \phantom{-}\cdot \\[-1.8pt]
\phantom{-}\cdot & \phantom{-}\cdot & \phantom{-}\cdot & \phantom{-}1 & \phantom{-}\cdot & \phantom{-}\cdot & \phantom{-}\cdot & \phantom{-}\cdot & \phantom{-}\cdot & \phantom{-}\cdot & \phantom{-}\cdot & \phantom{-}\cdot & \phantom{-}\cdot & \phantom{-}\cdot & \phantom{-}\cdot & \phantom{-}\cdot \\[-1.8pt]
\phantom{-}\cdot & \phantom{-}\cdot & -1 & \phantom{-}\cdot & \phantom{-}\cdot & \phantom{-}\cdot & \phantom{-}\cdot & \phantom{-}\cdot & \phantom{-}\cdot & \phantom{-}\cdot & \phantom{-}\cdot & \phantom{-}\cdot & \phantom{-}\cdot & \phantom{-}\cdot & \phantom{-}\cdot & \phantom{-}\cdot \\[-1.8pt]
\phantom{-}\cdot & \phantom{-}1 & \phantom{-}\cdot & \phantom{-}\cdot & \phantom{-}\cdot & \phantom{-}\cdot & \phantom{-}\cdot & \phantom{-}\cdot & \phantom{-}\cdot & \phantom{-}\cdot & \phantom{-}\cdot & \phantom{-}\cdot & \phantom{-}\cdot & \phantom{-}\cdot & \phantom{-}\cdot & \phantom{-}\cdot \\[-1.8pt]
-1 & \phantom{-}\cdot & \phantom{-}\cdot & \phantom{-}\cdot & \phantom{-}\cdot & \phantom{-}\cdot & \phantom{-}\cdot & \phantom{-}\cdot & \phantom{-}\cdot & \phantom{-}\cdot & \phantom{-}\cdot & \phantom{-}\cdot & \phantom{-}\cdot & \phantom{-}\cdot & \phantom{-}\cdot & \phantom{-}\cdot
\end{array}\right)
\]

\[
\gamma_{1644}^{7} = \left(\begin{array}{rrrrrrrr|rrrrrrrr}
\phantom{-}\cdot & \phantom{-}\cdot & \phantom{-}\cdot & \phantom{-}\cdot & \phantom{-}\cdot & \phantom{-}\cdot & \phantom{-}\cdot & \phantom{-}\cdot & \phantom{-}1 & \phantom{-}\cdot & \phantom{-}\cdot & \phantom{-}\cdot & \phantom{-}\cdot & \phantom{-}\cdot & \phantom{-}\cdot & \phantom{-}\cdot \\[-1.8pt]
\phantom{-}\cdot & \phantom{-}\cdot & \phantom{-}\cdot & \phantom{-}\cdot & \phantom{-}\cdot & \phantom{-}\cdot & \phantom{-}\cdot & \phantom{-}\cdot & \phantom{-}\cdot & \phantom{-}1 & \phantom{-}\cdot & \phantom{-}\cdot & \phantom{-}\cdot & \phantom{-}\cdot & \phantom{-}\cdot & \phantom{-}\cdot \\[-1.8pt]
\phantom{-}\cdot & \phantom{-}\cdot & \phantom{-}\cdot & \phantom{-}\cdot & \phantom{-}\cdot & \phantom{-}\cdot & \phantom{-}\cdot & \phantom{-}\cdot & \phantom{-}\cdot & \phantom{-}\cdot & \phantom{-}1 & \phantom{-}\cdot & \phantom{-}\cdot & \phantom{-}\cdot & \phantom{-}\cdot & \phantom{-}\cdot \\[-1.8pt]
\phantom{-}\cdot & \phantom{-}\cdot & \phantom{-}\cdot & \phantom{-}\cdot & \phantom{-}\cdot & \phantom{-}\cdot & \phantom{-}\cdot & \phantom{-}\cdot & \phantom{-}\cdot & \phantom{-}\cdot & \phantom{-}\cdot & \phantom{-}1 & \phantom{-}\cdot & \phantom{-}\cdot & \phantom{-}\cdot & \phantom{-}\cdot \\[-1.8pt]
\phantom{-}\cdot & \phantom{-}\cdot & \phantom{-}\cdot & \phantom{-}\cdot & \phantom{-}\cdot & \phantom{-}\cdot & \phantom{-}\cdot & \phantom{-}\cdot & \phantom{-}\cdot & \phantom{-}\cdot & \phantom{-}\cdot & \phantom{-}\cdot & -1 & \phantom{-}\cdot & \phantom{-}\cdot & \phantom{-}\cdot \\[-1.8pt]
\phantom{-}\cdot & \phantom{-}\cdot & \phantom{-}\cdot & \phantom{-}\cdot & \phantom{-}\cdot & \phantom{-}\cdot & \phantom{-}\cdot & \phantom{-}\cdot & \phantom{-}\cdot & \phantom{-}\cdot & \phantom{-}\cdot & \phantom{-}\cdot & \phantom{-}\cdot & -1 & \phantom{-}\cdot & \phantom{-}\cdot \\[-1.8pt]
\phantom{-}\cdot & \phantom{-}\cdot & \phantom{-}\cdot & \phantom{-}\cdot & \phantom{-}\cdot & \phantom{-}\cdot & \phantom{-}\cdot & \phantom{-}\cdot & \phantom{-}\cdot & \phantom{-}\cdot & \phantom{-}\cdot & \phantom{-}\cdot & \phantom{-}\cdot & \phantom{-}\cdot & -1 & \phantom{-}\cdot \\[-1.8pt]
\phantom{-}\cdot & \phantom{-}\cdot & \phantom{-}\cdot & \phantom{-}\cdot & \phantom{-}\cdot & \phantom{-}\cdot & \phantom{-}\cdot & \phantom{-}\cdot & \phantom{-}\cdot & \phantom{-}\cdot & \phantom{-}\cdot & \phantom{-}\cdot & \phantom{-}\cdot & \phantom{-}\cdot & \phantom{-}\cdot & -1 \\[-1.8pt]
\hline
-1 & \phantom{-}\cdot & \phantom{-}\cdot & \phantom{-}\cdot & \phantom{-}\cdot & \phantom{-}\cdot & \phantom{-}\cdot & \phantom{-}\cdot & \phantom{-}\cdot & \phantom{-}\cdot & \phantom{-}\cdot & \phantom{-}\cdot & \phantom{-}\cdot & \phantom{-}\cdot & \phantom{-}\cdot & \phantom{-}\cdot \\[-1.8pt]
\phantom{-}\cdot & -1 & \phantom{-}\cdot & \phantom{-}\cdot & \phantom{-}\cdot & \phantom{-}\cdot & \phantom{-}\cdot & \phantom{-}\cdot & \phantom{-}\cdot & \phantom{-}\cdot & \phantom{-}\cdot & \phantom{-}\cdot & \phantom{-}\cdot & \phantom{-}\cdot & \phantom{-}\cdot & \phantom{-}\cdot \\[-1.8pt]
\phantom{-}\cdot & \phantom{-}\cdot & -1 & \phantom{-}\cdot & \phantom{-}\cdot & \phantom{-}\cdot & \phantom{-}\cdot & \phantom{-}\cdot & \phantom{-}\cdot & \phantom{-}\cdot & \phantom{-}\cdot & \phantom{-}\cdot & \phantom{-}\cdot & \phantom{-}\cdot & \phantom{-}\cdot & \phantom{-}\cdot \\[-1.8pt]
\phantom{-}\cdot & \phantom{-}\cdot & \phantom{-}\cdot & -1 & \phantom{-}\cdot & \phantom{-}\cdot & \phantom{-}\cdot & \phantom{-}\cdot & \phantom{-}\cdot & \phantom{-}\cdot & \phantom{-}\cdot & \phantom{-}\cdot & \phantom{-}\cdot & \phantom{-}\cdot & \phantom{-}\cdot & \phantom{-}\cdot \\[-1.8pt]
\phantom{-}\cdot & \phantom{-}\cdot & \phantom{-}\cdot & \phantom{-}\cdot & \phantom{-}1 & \phantom{-}\cdot & \phantom{-}\cdot & \phantom{-}\cdot & \phantom{-}\cdot & \phantom{-}\cdot & \phantom{-}\cdot & \phantom{-}\cdot & \phantom{-}\cdot & \phantom{-}\cdot & \phantom{-}\cdot & \phantom{-}\cdot \\[-1.8pt]
\phantom{-}\cdot & \phantom{-}\cdot & \phantom{-}\cdot & \phantom{-}\cdot & \phantom{-}\cdot & \phantom{-}1 & \phantom{-}\cdot & \phantom{-}\cdot & \phantom{-}\cdot & \phantom{-}\cdot & \phantom{-}\cdot & \phantom{-}\cdot & \phantom{-}\cdot & \phantom{-}\cdot & \phantom{-}\cdot & \phantom{-}\cdot \\[-1.8pt]
\phantom{-}\cdot & \phantom{-}\cdot & \phantom{-}\cdot & \phantom{-}\cdot & \phantom{-}\cdot & \phantom{-}\cdot & \phantom{-}1 & \phantom{-}\cdot & \phantom{-}\cdot & \phantom{-}\cdot & \phantom{-}\cdot & \phantom{-}\cdot & \phantom{-}\cdot & \phantom{-}\cdot & \phantom{-}\cdot & \phantom{-}\cdot \\[-1.8pt]
\phantom{-}\cdot & \phantom{-}\cdot & \phantom{-}\cdot & \phantom{-}\cdot & \phantom{-}\cdot & \phantom{-}\cdot & \phantom{-}\cdot & \phantom{-}1 & \phantom{-}\cdot & \phantom{-}\cdot & \phantom{-}\cdot & \phantom{-}\cdot & \phantom{-}\cdot & \phantom{-}\cdot & \phantom{-}\cdot & \phantom{-}\cdot
\end{array}\right)
\]

These displays were compared by the program, entry by entry, with the
eight matrices that Mathematica printed in \texttt{4+4.nb} (cell 314,
\texttt{Out{[}211{]}}): check \texttt{44-N1}.

\hypertarget{proof-that-the-matrices-are-real}{%
\subsection{1.3 Proof that the matrices are
real}\label{proof-that-the-matrices-are-real}}

The matrices \(s_h\), \(t_h\) are built from the integers
\(\epsilon_{hpq4}\) and Kronecker deltas; \(\tau^A\), \(\bar\tau^A\) and
\(\gamma^A\) are built from them by forming blocks, transposes and
products. No complex number enters at any step. Check \texttt{44-R1}
confirms that every entry of every \(\gamma_{1644}^A\) is one of the
integers \(-1,0,+1\); check \texttt{44-R2} adds that each matrix has
exactly one non-zero entry in every row and in every column (it is a
\emph{signed permutation matrix}). \(\blacksquare\)

\hypertarget{proof-of-the-anticommutation-relations}{%
\subsection{1.4 Proof of the anticommutation
relations}\label{proof-of-the-anticommutation-relations}}

\textbf{Statement.}
\[\boxed{\;\gamma_{1644}^A\gamma_{1644}^B+\gamma_{1644}^B\gamma_{1644}^A=2\,\eta_{1644}^{AB}\,1_{16},\qquad
\eta_{1644}=\mathrm{diag}(+1,+1,+1,+1,-1,-1,-1,-1)\;}\tag{1.3}\] Thus
\(\gamma^0,\gamma^1,\gamma^2,\gamma^3\) (space-like directions) square
to \(+1\), \(\gamma^4,\gamma^5,\gamma^6,\gamma^7\) (time-like
directions) square to \(-1\), and different matrices anticommute.

\textbf{Proof.} From the block form (1.2),
\[\gamma^A\gamma^B+\gamma^B\gamma^A=\begin{pmatrix}\bar\tau^A\tau^B+\bar\tau^B\tau^A&0\\0&\tau^A\bar\tau^B+\tau^B\bar\tau^A\end{pmatrix},\]
and both blocks equal \(2\eta^{AB}1_8\) by Eq. (1.1). Equation (1.1) is
itself an identity between integer matrices; it is verified for all 64
pairs \((A,B)\) by check \texttt{44-B2}, and Eq. (1.3) is verified
directly by check \texttt{44-C1}. \(\blacksquare\)

\textbf{Consequences} (checks \texttt{44-C2}, \texttt{44-C3},
\texttt{44-C4}).

\begin{itemize}
\tightlist
\item
  \((\gamma^A)^T=\eta^{AA}\gamma^A\) (no sum): the space-like matrices
  are symmetric, the time-like ones antisymmetric, and every
  \(\gamma^A\) is an orthogonal matrix.
\item
  The 256 products \(\gamma^{A_1}\cdots\gamma^{A_r}\) with
  \(A_1<\dots<A_r\), \(r=0,\dots,8\), are linearly independent. They
  therefore span all real 16×16 matrices: the matrices generate the full
  matrix algebra, \(\mathrm{Cl}(4,4)\cong\mathbb{R}(16)\), and the
  16-dimensional representation is irreducible over the real numbers.
  (The Python script proves the independence differently: the 256
  products are mutually orthogonal for the scalar product
  \(\langle X,Y\rangle=\mathrm{tr}(X^TY)\), check \texttt{I44-G}.)
\end{itemize}

\hypertarget{the-scaled-commutators-and-the-group-they-generate}{%
\subsection{1.5 The scaled commutators and the group they
generate}\label{the-scaled-commutators-and-the-group-they-generate}}

\textbf{Definition.} The \emph{scaled commutators} are
\[\Sigma^{AB}\equiv\tfrac14\big[\gamma^A,\gamma^B\big]=\tfrac12\gamma^A\gamma^B\quad(A\neq B),\qquad\Sigma^{AB}=-\Sigma^{BA},\tag{1.4}\]
where the second form follows from Eq. (1.3) (check \texttt{44-L1}).
There are \(8\cdot7/2=28\) of them with \(A<B\), and they are linearly
independent (check \texttt{44-L2}).

\textbf{Theorem 1.} Let \(\gamma^A=\gamma_{1644}^A\),
\(\Sigma^{AB}=\Sigma_{1644}^{AB}\) and \(\eta=\eta_{1644}\).

\begin{enumerate}
\def\labelenumi{\arabic{enumi}.}
\tightlist
\item
  \emph{(Lie algebra.)} The scaled commutators satisfy
  \[[\Sigma^{AB},\Sigma^{CD}]=\eta^{BC}\Sigma^{AD}-\eta^{AC}\Sigma^{BD}-\eta^{BD}\Sigma^{AC}+\eta^{AD}\Sigma^{BC},\tag{1.5}\]
  \[[\Sigma^{AB},\gamma^{C}]=\eta^{BC}\gamma^{A}-\eta^{AC}\gamma^{B}.\tag{1.6}\]
  Their real span is a Lie algebra isomorphic to \(\mathfrak{so}(4,4)\).
\item
  \emph{(Identity component.)} The matrices
  \(S=\exp\!\big(\sum_{A<B}\theta_{AB}\Sigma^{AB}\big)\) with real
  parameters \(\theta_{AB}\) generate a connected group,
  \(\mathrm{Spin}_0(4,4)\). Every such \(S\) satisfies
  \(S\gamma^AS^{-1}=\Lambda^A{}_B\gamma^B\) with a real 8×8 matrix
  \(\Lambda\) that obeys \(\Lambda\,\eta\,\Lambda^T=\eta\) and
  \(\det\Lambda=+1\). The map \(S\mapsto\Lambda\) sends
  \(\mathrm{Spin}_0(4,4)\) onto the identity component \(SO_0(4,4)\) of
  the pseudo-orthogonal group, and \(S\) and \(-S\) are sent to the same
  \(\Lambda\).
\item
  \emph{(The whole group.)} Each \(\gamma^A\) is invertible, and
  \(v\mapsto-\gamma^Av(\gamma^A)^{-1}\) maps \(\gamma^B\mapsto\gamma^B\)
  \((B\neq A)\) and \(\gamma^A\mapsto-\gamma^A\): it is the reflection
  that reverses the axis \(xA\). The group generated by
  \(\mathrm{Spin}_0(4,4)\) together with the eight matrices \(\gamma^A\)
  is \(\mathrm{Pin}(4,4)\), which consists of four connected pieces,
  \[\mathrm{Pin}(4,4)=\mathrm{Spin}_0(4,4)\;\cup\;\gamma^0\,\mathrm{Spin}_0(4,4)\;\cup\;\gamma^4\,\mathrm{Spin}_0(4,4)\;\cup\;\gamma^0\gamma^4\,\mathrm{Spin}_0(4,4) .\]
\end{enumerate}

\emph{In words:} the scaled commutators span the Lie algebra of
\(\mathrm{Pin}(4,4)\) and generate, by exponentiation, its identity
component; the three other components are reached by multiplying with
one space-like matrix, one time-like matrix, or both. Exponentials alone
can never leave the identity component, because \(\exp(sY)\),
\(0\le s\le1\), is a continuous path from the unit matrix to
\(\exp(Y)\). This is why item 3 is needed, and it is the precise sense
in which ``the scaled commutators generate \(\mathrm{Pin}(4,4)\)''.

\textbf{Proof.} \emph{Item 1.} Equation (1.6) follows from (1.3):
\([\gamma^A\gamma^B,\gamma^C]=\gamma^A\{\gamma^B,\gamma^C\}-\{\gamma^A,\gamma^C\}\gamma^B=2\eta^{BC}\gamma^A-2\eta^{AC}\gamma^B\);
applying it twice to \(\Sigma^{CD}=\tfrac12\gamma^C\gamma^D\) gives
(1.5). Both are also verified for all index combinations (checks
\texttt{44-L3}, \texttt{44-L4}). Relations (1.5) are the defining
relations of the Lie algebra of the pseudo-orthogonal group of the
metric \(\eta\), and by (1.6) the \(\Sigma^{AB}\) act on the eight
\(\gamma^C\) exactly as the generators of that group act on a vector. To
identify the real form without appeal to the literature, the program
also computes the Killing form
\(B_{\rm Kil}(X,Y)=\mathrm{tr}(\mathrm{ad}_X\mathrm{ad}_Y)\) from the
structure constants: it is diagonal in the basis \(\Sigma^{AB}\), with
\(B_{\rm Kil}(\Sigma^{AB},\Sigma^{AB})=-12\,\eta^{AA}\eta^{BB}\), hence
it has 12 negative and 16 positive eigenvalues (checks \texttt{44-L5},
\texttt{I44-K}). Equivalently, 12 generators are antisymmetric matrices,
which exponentiate to rotations, and 16 are symmetric matrices, which
exponentiate to boosts (check \texttt{44-L6}). The numbers (compact,
non-compact) of the real forms \(\mathfrak{so}(p,q)\), \(p+q=8\), are
\((28,0)\) for \(\mathfrak{so}(8)\), \((21,7)\) for
\(\mathfrak{so}(7,1)\), \((16,12)\) for \(\mathfrak{so}(6,2)\),
\((13,15)\) for \(\mathfrak{so}(5,3)\) and \((12,16)\) for
\(\mathfrak{so}(4,4)\): the algebra is \(\mathfrak{so}(4,4)\).

\emph{Item 2.} For a single plane the exponential is elementary, because
\((\gamma^A\gamma^B)^2=-\eta^{AA}\eta^{BB}\):
\[\exp(\theta\,\Sigma^{AB})=\begin{cases}\cos\frac\theta2+\gamma^A\gamma^B\sin\frac\theta2,&\eta^{AA}\eta^{BB}=+1\ \text{(rotation)},\\[2pt]
\cosh\frac\theta2+\gamma^A\gamma^B\sinh\frac\theta2,&\eta^{AA}\eta^{BB}=-1\ \text{(boost)}\end{cases}\tag{1.7}\]
(check \texttt{44-G2}). By Eq. (1.6),
\(\frac{d}{ds}\big(e^{sY}\gamma^Ce^{-sY}\big)\) is again a linear
combination of the \(\gamma\)'s when \(Y\) is a combination of the
\(\Sigma^{AB}\), so \(S\gamma^CS^{-1}=\Lambda^C{}_D\gamma^D\); applying
\(S(\cdot)S^{-1}\) to Eq. (1.3) shows \(\Lambda\eta\Lambda^T=\eta\).
Check \texttt{44-G3} confirms both statements, and \(\det\Lambda=1\),
numerically for 20 random elements. A rotation by \(\theta=2\pi\) in Eq.
(1.7) gives \(S=-1\) while \(\Lambda=1\): the map is two-to-one, and
\(-1\) is connected to \(+1\) inside \(\mathrm{Spin}_0(4,4)\).

\emph{Item 3.} \((\gamma^A)^{-1}=\eta^{AA}\gamma^A\) and the reflection
property are check \texttt{44-G1}. By definition \(\mathrm{Pin}(4,4)\)
is the group generated, inside the Clifford algebra, by the vectors
\(v=v_A\gamma^A\) with \(v^2=\pm1\); every such \(v\) is obtained from
\(\gamma^0\) (if \(v^2=+1\)) or from \(\gamma^4\) (if \(v^2=-1\)) by
conjugation with an element of \(\mathrm{Spin}_0(4,4)\), so the group
generated by \(\mathrm{Spin}_0(4,4)\) and the \(\gamma^A\) is
\(\mathrm{Pin}(4,4)\). A product of two different \(\gamma\)'s of the
same kind is \(\exp(\pi\Sigma^{AB})\) by Eq. (1.7) and lies in
\(\mathrm{Spin}_0(4,4)\); a product \(\gamma^a\gamma^t\) of a space-like
and a time-like matrix reverses one space axis and one time axis and
does not. Hence the four components listed. \(\blacksquare\)

\hypertarget{the-chiral-projection-operators-re-derived}{%
\subsection{1.6 The chiral projection operators,
re-derived}\label{the-chiral-projection-operators-re-derived}}

\emph{Reference:} \texttt{4+4.nb}, subsection ``Left - Handed vs . Right
- Handed Fermion Fields (Chirality) ; projection operators'' (cells
161--167).

\textbf{Volume element.} Multiply the eight matrices in order:
\[\omega\equiv\gamma^0\gamma^1\gamma^2\gamma^3\gamma^4\gamma^5\gamma^6\gamma^7 .\]
Two properties follow from Eq. (1.3) alone. (a) \(\omega\) anticommutes
with every \(\gamma^A\): moving \(\gamma^A\) through the product costs
seven sign changes (one for each of the other matrices). (b)
\(\omega^2=(-1)^{8\cdot7/2}\prod_A(\gamma^A)^2=(+1)(+1)^4(-1)^4=+1\),
where \((-1)^{28}\) is the sign of reversing the order of eight
anticommuting factors. Because \(\omega\) is real with \(\omega^2=+1\),
the matrices \(\frac12(1\mp\omega)\) are \emph{real} projectors.
Explicitly (checks \texttt{44-V1}, \texttt{44-V2}; identical to
\texttt{T16\^{}A{[}8{]}} printed in cell 136 of the notebook, check
\texttt{44-N2}):

\[
\omega_{1644} = \left(\begin{array}{rrrrrrrr|rrrrrrrr}
-1 & \phantom{-}\cdot & \phantom{-}\cdot & \phantom{-}\cdot & \phantom{-}\cdot & \phantom{-}\cdot & \phantom{-}\cdot & \phantom{-}\cdot & \phantom{-}\cdot & \phantom{-}\cdot & \phantom{-}\cdot & \phantom{-}\cdot & \phantom{-}\cdot & \phantom{-}\cdot & \phantom{-}\cdot & \phantom{-}\cdot \\[-1.8pt]
\phantom{-}\cdot & -1 & \phantom{-}\cdot & \phantom{-}\cdot & \phantom{-}\cdot & \phantom{-}\cdot & \phantom{-}\cdot & \phantom{-}\cdot & \phantom{-}\cdot & \phantom{-}\cdot & \phantom{-}\cdot & \phantom{-}\cdot & \phantom{-}\cdot & \phantom{-}\cdot & \phantom{-}\cdot & \phantom{-}\cdot \\[-1.8pt]
\phantom{-}\cdot & \phantom{-}\cdot & -1 & \phantom{-}\cdot & \phantom{-}\cdot & \phantom{-}\cdot & \phantom{-}\cdot & \phantom{-}\cdot & \phantom{-}\cdot & \phantom{-}\cdot & \phantom{-}\cdot & \phantom{-}\cdot & \phantom{-}\cdot & \phantom{-}\cdot & \phantom{-}\cdot & \phantom{-}\cdot \\[-1.8pt]
\phantom{-}\cdot & \phantom{-}\cdot & \phantom{-}\cdot & -1 & \phantom{-}\cdot & \phantom{-}\cdot & \phantom{-}\cdot & \phantom{-}\cdot & \phantom{-}\cdot & \phantom{-}\cdot & \phantom{-}\cdot & \phantom{-}\cdot & \phantom{-}\cdot & \phantom{-}\cdot & \phantom{-}\cdot & \phantom{-}\cdot \\[-1.8pt]
\phantom{-}\cdot & \phantom{-}\cdot & \phantom{-}\cdot & \phantom{-}\cdot & -1 & \phantom{-}\cdot & \phantom{-}\cdot & \phantom{-}\cdot & \phantom{-}\cdot & \phantom{-}\cdot & \phantom{-}\cdot & \phantom{-}\cdot & \phantom{-}\cdot & \phantom{-}\cdot & \phantom{-}\cdot & \phantom{-}\cdot \\[-1.8pt]
\phantom{-}\cdot & \phantom{-}\cdot & \phantom{-}\cdot & \phantom{-}\cdot & \phantom{-}\cdot & -1 & \phantom{-}\cdot & \phantom{-}\cdot & \phantom{-}\cdot & \phantom{-}\cdot & \phantom{-}\cdot & \phantom{-}\cdot & \phantom{-}\cdot & \phantom{-}\cdot & \phantom{-}\cdot & \phantom{-}\cdot \\[-1.8pt]
\phantom{-}\cdot & \phantom{-}\cdot & \phantom{-}\cdot & \phantom{-}\cdot & \phantom{-}\cdot & \phantom{-}\cdot & -1 & \phantom{-}\cdot & \phantom{-}\cdot & \phantom{-}\cdot & \phantom{-}\cdot & \phantom{-}\cdot & \phantom{-}\cdot & \phantom{-}\cdot & \phantom{-}\cdot & \phantom{-}\cdot \\[-1.8pt]
\phantom{-}\cdot & \phantom{-}\cdot & \phantom{-}\cdot & \phantom{-}\cdot & \phantom{-}\cdot & \phantom{-}\cdot & \phantom{-}\cdot & -1 & \phantom{-}\cdot & \phantom{-}\cdot & \phantom{-}\cdot & \phantom{-}\cdot & \phantom{-}\cdot & \phantom{-}\cdot & \phantom{-}\cdot & \phantom{-}\cdot \\[-1.8pt]
\hline
\phantom{-}\cdot & \phantom{-}\cdot & \phantom{-}\cdot & \phantom{-}\cdot & \phantom{-}\cdot & \phantom{-}\cdot & \phantom{-}\cdot & \phantom{-}\cdot & \phantom{-}1 & \phantom{-}\cdot & \phantom{-}\cdot & \phantom{-}\cdot & \phantom{-}\cdot & \phantom{-}\cdot & \phantom{-}\cdot & \phantom{-}\cdot \\[-1.8pt]
\phantom{-}\cdot & \phantom{-}\cdot & \phantom{-}\cdot & \phantom{-}\cdot & \phantom{-}\cdot & \phantom{-}\cdot & \phantom{-}\cdot & \phantom{-}\cdot & \phantom{-}\cdot & \phantom{-}1 & \phantom{-}\cdot & \phantom{-}\cdot & \phantom{-}\cdot & \phantom{-}\cdot & \phantom{-}\cdot & \phantom{-}\cdot \\[-1.8pt]
\phantom{-}\cdot & \phantom{-}\cdot & \phantom{-}\cdot & \phantom{-}\cdot & \phantom{-}\cdot & \phantom{-}\cdot & \phantom{-}\cdot & \phantom{-}\cdot & \phantom{-}\cdot & \phantom{-}\cdot & \phantom{-}1 & \phantom{-}\cdot & \phantom{-}\cdot & \phantom{-}\cdot & \phantom{-}\cdot & \phantom{-}\cdot \\[-1.8pt]
\phantom{-}\cdot & \phantom{-}\cdot & \phantom{-}\cdot & \phantom{-}\cdot & \phantom{-}\cdot & \phantom{-}\cdot & \phantom{-}\cdot & \phantom{-}\cdot & \phantom{-}\cdot & \phantom{-}\cdot & \phantom{-}\cdot & \phantom{-}1 & \phantom{-}\cdot & \phantom{-}\cdot & \phantom{-}\cdot & \phantom{-}\cdot \\[-1.8pt]
\phantom{-}\cdot & \phantom{-}\cdot & \phantom{-}\cdot & \phantom{-}\cdot & \phantom{-}\cdot & \phantom{-}\cdot & \phantom{-}\cdot & \phantom{-}\cdot & \phantom{-}\cdot & \phantom{-}\cdot & \phantom{-}\cdot & \phantom{-}\cdot & \phantom{-}1 & \phantom{-}\cdot & \phantom{-}\cdot & \phantom{-}\cdot \\[-1.8pt]
\phantom{-}\cdot & \phantom{-}\cdot & \phantom{-}\cdot & \phantom{-}\cdot & \phantom{-}\cdot & \phantom{-}\cdot & \phantom{-}\cdot & \phantom{-}\cdot & \phantom{-}\cdot & \phantom{-}\cdot & \phantom{-}\cdot & \phantom{-}\cdot & \phantom{-}\cdot & \phantom{-}1 & \phantom{-}\cdot & \phantom{-}\cdot \\[-1.8pt]
\phantom{-}\cdot & \phantom{-}\cdot & \phantom{-}\cdot & \phantom{-}\cdot & \phantom{-}\cdot & \phantom{-}\cdot & \phantom{-}\cdot & \phantom{-}\cdot & \phantom{-}\cdot & \phantom{-}\cdot & \phantom{-}\cdot & \phantom{-}\cdot & \phantom{-}\cdot & \phantom{-}\cdot & \phantom{-}1 & \phantom{-}\cdot \\[-1.8pt]
\phantom{-}\cdot & \phantom{-}\cdot & \phantom{-}\cdot & \phantom{-}\cdot & \phantom{-}\cdot & \phantom{-}\cdot & \phantom{-}\cdot & \phantom{-}\cdot & \phantom{-}\cdot & \phantom{-}\cdot & \phantom{-}\cdot & \phantom{-}\cdot & \phantom{-}\cdot & \phantom{-}\cdot & \phantom{-}\cdot & \phantom{-}1
\end{array}\right)
\]

\textbf{Projectors.}
\[\boxed{\;P_{1644}^{L}=\tfrac12\big(1-\omega_{1644}\big)=\begin{pmatrix}1_8&0\\0&0\end{pmatrix},\qquad
P_{1644}^{R}=\tfrac12\big(1+\omega_{1644}\big)=\begin{pmatrix}0&0\\0&1_8\end{pmatrix}\;}\tag{1.8}\]

\[
P_{1644}^{L} = \left(\begin{array}{rrrrrrrr|rrrrrrrr}
\phantom{-}1 & \phantom{-}\cdot & \phantom{-}\cdot & \phantom{-}\cdot & \phantom{-}\cdot & \phantom{-}\cdot & \phantom{-}\cdot & \phantom{-}\cdot & \phantom{-}\cdot & \phantom{-}\cdot & \phantom{-}\cdot & \phantom{-}\cdot & \phantom{-}\cdot & \phantom{-}\cdot & \phantom{-}\cdot & \phantom{-}\cdot \\[-1.8pt]
\phantom{-}\cdot & \phantom{-}1 & \phantom{-}\cdot & \phantom{-}\cdot & \phantom{-}\cdot & \phantom{-}\cdot & \phantom{-}\cdot & \phantom{-}\cdot & \phantom{-}\cdot & \phantom{-}\cdot & \phantom{-}\cdot & \phantom{-}\cdot & \phantom{-}\cdot & \phantom{-}\cdot & \phantom{-}\cdot & \phantom{-}\cdot \\[-1.8pt]
\phantom{-}\cdot & \phantom{-}\cdot & \phantom{-}1 & \phantom{-}\cdot & \phantom{-}\cdot & \phantom{-}\cdot & \phantom{-}\cdot & \phantom{-}\cdot & \phantom{-}\cdot & \phantom{-}\cdot & \phantom{-}\cdot & \phantom{-}\cdot & \phantom{-}\cdot & \phantom{-}\cdot & \phantom{-}\cdot & \phantom{-}\cdot \\[-1.8pt]
\phantom{-}\cdot & \phantom{-}\cdot & \phantom{-}\cdot & \phantom{-}1 & \phantom{-}\cdot & \phantom{-}\cdot & \phantom{-}\cdot & \phantom{-}\cdot & \phantom{-}\cdot & \phantom{-}\cdot & \phantom{-}\cdot & \phantom{-}\cdot & \phantom{-}\cdot & \phantom{-}\cdot & \phantom{-}\cdot & \phantom{-}\cdot \\[-1.8pt]
\phantom{-}\cdot & \phantom{-}\cdot & \phantom{-}\cdot & \phantom{-}\cdot & \phantom{-}1 & \phantom{-}\cdot & \phantom{-}\cdot & \phantom{-}\cdot & \phantom{-}\cdot & \phantom{-}\cdot & \phantom{-}\cdot & \phantom{-}\cdot & \phantom{-}\cdot & \phantom{-}\cdot & \phantom{-}\cdot & \phantom{-}\cdot \\[-1.8pt]
\phantom{-}\cdot & \phantom{-}\cdot & \phantom{-}\cdot & \phantom{-}\cdot & \phantom{-}\cdot & \phantom{-}1 & \phantom{-}\cdot & \phantom{-}\cdot & \phantom{-}\cdot & \phantom{-}\cdot & \phantom{-}\cdot & \phantom{-}\cdot & \phantom{-}\cdot & \phantom{-}\cdot & \phantom{-}\cdot & \phantom{-}\cdot \\[-1.8pt]
\phantom{-}\cdot & \phantom{-}\cdot & \phantom{-}\cdot & \phantom{-}\cdot & \phantom{-}\cdot & \phantom{-}\cdot & \phantom{-}1 & \phantom{-}\cdot & \phantom{-}\cdot & \phantom{-}\cdot & \phantom{-}\cdot & \phantom{-}\cdot & \phantom{-}\cdot & \phantom{-}\cdot & \phantom{-}\cdot & \phantom{-}\cdot \\[-1.8pt]
\phantom{-}\cdot & \phantom{-}\cdot & \phantom{-}\cdot & \phantom{-}\cdot & \phantom{-}\cdot & \phantom{-}\cdot & \phantom{-}\cdot & \phantom{-}1 & \phantom{-}\cdot & \phantom{-}\cdot & \phantom{-}\cdot & \phantom{-}\cdot & \phantom{-}\cdot & \phantom{-}\cdot & \phantom{-}\cdot & \phantom{-}\cdot \\[-1.8pt]
\hline
\phantom{-}\cdot & \phantom{-}\cdot & \phantom{-}\cdot & \phantom{-}\cdot & \phantom{-}\cdot & \phantom{-}\cdot & \phantom{-}\cdot & \phantom{-}\cdot & \phantom{-}\cdot & \phantom{-}\cdot & \phantom{-}\cdot & \phantom{-}\cdot & \phantom{-}\cdot & \phantom{-}\cdot & \phantom{-}\cdot & \phantom{-}\cdot \\[-1.8pt]
\phantom{-}\cdot & \phantom{-}\cdot & \phantom{-}\cdot & \phantom{-}\cdot & \phantom{-}\cdot & \phantom{-}\cdot & \phantom{-}\cdot & \phantom{-}\cdot & \phantom{-}\cdot & \phantom{-}\cdot & \phantom{-}\cdot & \phantom{-}\cdot & \phantom{-}\cdot & \phantom{-}\cdot & \phantom{-}\cdot & \phantom{-}\cdot \\[-1.8pt]
\phantom{-}\cdot & \phantom{-}\cdot & \phantom{-}\cdot & \phantom{-}\cdot & \phantom{-}\cdot & \phantom{-}\cdot & \phantom{-}\cdot & \phantom{-}\cdot & \phantom{-}\cdot & \phantom{-}\cdot & \phantom{-}\cdot & \phantom{-}\cdot & \phantom{-}\cdot & \phantom{-}\cdot & \phantom{-}\cdot & \phantom{-}\cdot \\[-1.8pt]
\phantom{-}\cdot & \phantom{-}\cdot & \phantom{-}\cdot & \phantom{-}\cdot & \phantom{-}\cdot & \phantom{-}\cdot & \phantom{-}\cdot & \phantom{-}\cdot & \phantom{-}\cdot & \phantom{-}\cdot & \phantom{-}\cdot & \phantom{-}\cdot & \phantom{-}\cdot & \phantom{-}\cdot & \phantom{-}\cdot & \phantom{-}\cdot \\[-1.8pt]
\phantom{-}\cdot & \phantom{-}\cdot & \phantom{-}\cdot & \phantom{-}\cdot & \phantom{-}\cdot & \phantom{-}\cdot & \phantom{-}\cdot & \phantom{-}\cdot & \phantom{-}\cdot & \phantom{-}\cdot & \phantom{-}\cdot & \phantom{-}\cdot & \phantom{-}\cdot & \phantom{-}\cdot & \phantom{-}\cdot & \phantom{-}\cdot \\[-1.8pt]
\phantom{-}\cdot & \phantom{-}\cdot & \phantom{-}\cdot & \phantom{-}\cdot & \phantom{-}\cdot & \phantom{-}\cdot & \phantom{-}\cdot & \phantom{-}\cdot & \phantom{-}\cdot & \phantom{-}\cdot & \phantom{-}\cdot & \phantom{-}\cdot & \phantom{-}\cdot & \phantom{-}\cdot & \phantom{-}\cdot & \phantom{-}\cdot \\[-1.8pt]
\phantom{-}\cdot & \phantom{-}\cdot & \phantom{-}\cdot & \phantom{-}\cdot & \phantom{-}\cdot & \phantom{-}\cdot & \phantom{-}\cdot & \phantom{-}\cdot & \phantom{-}\cdot & \phantom{-}\cdot & \phantom{-}\cdot & \phantom{-}\cdot & \phantom{-}\cdot & \phantom{-}\cdot & \phantom{-}\cdot & \phantom{-}\cdot \\[-1.8pt]
\phantom{-}\cdot & \phantom{-}\cdot & \phantom{-}\cdot & \phantom{-}\cdot & \phantom{-}\cdot & \phantom{-}\cdot & \phantom{-}\cdot & \phantom{-}\cdot & \phantom{-}\cdot & \phantom{-}\cdot & \phantom{-}\cdot & \phantom{-}\cdot & \phantom{-}\cdot & \phantom{-}\cdot & \phantom{-}\cdot & \phantom{-}\cdot
\end{array}\right)
\]

\[
P_{1644}^{R} = \left(\begin{array}{rrrrrrrr|rrrrrrrr}
\phantom{-}\cdot & \phantom{-}\cdot & \phantom{-}\cdot & \phantom{-}\cdot & \phantom{-}\cdot & \phantom{-}\cdot & \phantom{-}\cdot & \phantom{-}\cdot & \phantom{-}\cdot & \phantom{-}\cdot & \phantom{-}\cdot & \phantom{-}\cdot & \phantom{-}\cdot & \phantom{-}\cdot & \phantom{-}\cdot & \phantom{-}\cdot \\[-1.8pt]
\phantom{-}\cdot & \phantom{-}\cdot & \phantom{-}\cdot & \phantom{-}\cdot & \phantom{-}\cdot & \phantom{-}\cdot & \phantom{-}\cdot & \phantom{-}\cdot & \phantom{-}\cdot & \phantom{-}\cdot & \phantom{-}\cdot & \phantom{-}\cdot & \phantom{-}\cdot & \phantom{-}\cdot & \phantom{-}\cdot & \phantom{-}\cdot \\[-1.8pt]
\phantom{-}\cdot & \phantom{-}\cdot & \phantom{-}\cdot & \phantom{-}\cdot & \phantom{-}\cdot & \phantom{-}\cdot & \phantom{-}\cdot & \phantom{-}\cdot & \phantom{-}\cdot & \phantom{-}\cdot & \phantom{-}\cdot & \phantom{-}\cdot & \phantom{-}\cdot & \phantom{-}\cdot & \phantom{-}\cdot & \phantom{-}\cdot \\[-1.8pt]
\phantom{-}\cdot & \phantom{-}\cdot & \phantom{-}\cdot & \phantom{-}\cdot & \phantom{-}\cdot & \phantom{-}\cdot & \phantom{-}\cdot & \phantom{-}\cdot & \phantom{-}\cdot & \phantom{-}\cdot & \phantom{-}\cdot & \phantom{-}\cdot & \phantom{-}\cdot & \phantom{-}\cdot & \phantom{-}\cdot & \phantom{-}\cdot \\[-1.8pt]
\phantom{-}\cdot & \phantom{-}\cdot & \phantom{-}\cdot & \phantom{-}\cdot & \phantom{-}\cdot & \phantom{-}\cdot & \phantom{-}\cdot & \phantom{-}\cdot & \phantom{-}\cdot & \phantom{-}\cdot & \phantom{-}\cdot & \phantom{-}\cdot & \phantom{-}\cdot & \phantom{-}\cdot & \phantom{-}\cdot & \phantom{-}\cdot \\[-1.8pt]
\phantom{-}\cdot & \phantom{-}\cdot & \phantom{-}\cdot & \phantom{-}\cdot & \phantom{-}\cdot & \phantom{-}\cdot & \phantom{-}\cdot & \phantom{-}\cdot & \phantom{-}\cdot & \phantom{-}\cdot & \phantom{-}\cdot & \phantom{-}\cdot & \phantom{-}\cdot & \phantom{-}\cdot & \phantom{-}\cdot & \phantom{-}\cdot \\[-1.8pt]
\phantom{-}\cdot & \phantom{-}\cdot & \phantom{-}\cdot & \phantom{-}\cdot & \phantom{-}\cdot & \phantom{-}\cdot & \phantom{-}\cdot & \phantom{-}\cdot & \phantom{-}\cdot & \phantom{-}\cdot & \phantom{-}\cdot & \phantom{-}\cdot & \phantom{-}\cdot & \phantom{-}\cdot & \phantom{-}\cdot & \phantom{-}\cdot \\[-1.8pt]
\phantom{-}\cdot & \phantom{-}\cdot & \phantom{-}\cdot & \phantom{-}\cdot & \phantom{-}\cdot & \phantom{-}\cdot & \phantom{-}\cdot & \phantom{-}\cdot & \phantom{-}\cdot & \phantom{-}\cdot & \phantom{-}\cdot & \phantom{-}\cdot & \phantom{-}\cdot & \phantom{-}\cdot & \phantom{-}\cdot & \phantom{-}\cdot \\[-1.8pt]
\hline
\phantom{-}\cdot & \phantom{-}\cdot & \phantom{-}\cdot & \phantom{-}\cdot & \phantom{-}\cdot & \phantom{-}\cdot & \phantom{-}\cdot & \phantom{-}\cdot & \phantom{-}1 & \phantom{-}\cdot & \phantom{-}\cdot & \phantom{-}\cdot & \phantom{-}\cdot & \phantom{-}\cdot & \phantom{-}\cdot & \phantom{-}\cdot \\[-1.8pt]
\phantom{-}\cdot & \phantom{-}\cdot & \phantom{-}\cdot & \phantom{-}\cdot & \phantom{-}\cdot & \phantom{-}\cdot & \phantom{-}\cdot & \phantom{-}\cdot & \phantom{-}\cdot & \phantom{-}1 & \phantom{-}\cdot & \phantom{-}\cdot & \phantom{-}\cdot & \phantom{-}\cdot & \phantom{-}\cdot & \phantom{-}\cdot \\[-1.8pt]
\phantom{-}\cdot & \phantom{-}\cdot & \phantom{-}\cdot & \phantom{-}\cdot & \phantom{-}\cdot & \phantom{-}\cdot & \phantom{-}\cdot & \phantom{-}\cdot & \phantom{-}\cdot & \phantom{-}\cdot & \phantom{-}1 & \phantom{-}\cdot & \phantom{-}\cdot & \phantom{-}\cdot & \phantom{-}\cdot & \phantom{-}\cdot \\[-1.8pt]
\phantom{-}\cdot & \phantom{-}\cdot & \phantom{-}\cdot & \phantom{-}\cdot & \phantom{-}\cdot & \phantom{-}\cdot & \phantom{-}\cdot & \phantom{-}\cdot & \phantom{-}\cdot & \phantom{-}\cdot & \phantom{-}\cdot & \phantom{-}1 & \phantom{-}\cdot & \phantom{-}\cdot & \phantom{-}\cdot & \phantom{-}\cdot \\[-1.8pt]
\phantom{-}\cdot & \phantom{-}\cdot & \phantom{-}\cdot & \phantom{-}\cdot & \phantom{-}\cdot & \phantom{-}\cdot & \phantom{-}\cdot & \phantom{-}\cdot & \phantom{-}\cdot & \phantom{-}\cdot & \phantom{-}\cdot & \phantom{-}\cdot & \phantom{-}1 & \phantom{-}\cdot & \phantom{-}\cdot & \phantom{-}\cdot \\[-1.8pt]
\phantom{-}\cdot & \phantom{-}\cdot & \phantom{-}\cdot & \phantom{-}\cdot & \phantom{-}\cdot & \phantom{-}\cdot & \phantom{-}\cdot & \phantom{-}\cdot & \phantom{-}\cdot & \phantom{-}\cdot & \phantom{-}\cdot & \phantom{-}\cdot & \phantom{-}\cdot & \phantom{-}1 & \phantom{-}\cdot & \phantom{-}\cdot \\[-1.8pt]
\phantom{-}\cdot & \phantom{-}\cdot & \phantom{-}\cdot & \phantom{-}\cdot & \phantom{-}\cdot & \phantom{-}\cdot & \phantom{-}\cdot & \phantom{-}\cdot & \phantom{-}\cdot & \phantom{-}\cdot & \phantom{-}\cdot & \phantom{-}\cdot & \phantom{-}\cdot & \phantom{-}\cdot & \phantom{-}1 & \phantom{-}\cdot \\[-1.8pt]
\phantom{-}\cdot & \phantom{-}\cdot & \phantom{-}\cdot & \phantom{-}\cdot & \phantom{-}\cdot & \phantom{-}\cdot & \phantom{-}\cdot & \phantom{-}\cdot & \phantom{-}\cdot & \phantom{-}\cdot & \phantom{-}\cdot & \phantom{-}\cdot & \phantom{-}\cdot & \phantom{-}\cdot & \phantom{-}\cdot & \phantom{-}1
\end{array}\right)
\]

These are the definitions of the notebook (cell 162:
\texttt{P\_L\ =\ (1/2)(ID16\ -\ T16\^{}A{[}8{]})}), and the matrices are
identical to \texttt{P\_L}, \texttt{P\_R} printed in its cells 163 and
165 (check \texttt{44-P3}). They satisfy
\[P^{L}+P^{R}=1,\qquad P^{L}P^{L}=P^{L},\qquad P^{R}P^{R}=P^{R},\qquad P^{L}P^{R}=P^{R}P^{L}=0,\qquad\mathrm{tr}\,P^{L}=\mathrm{tr}\,P^{R}=8\]
(checks \texttt{44-P1}, \texttt{44-P2}); the notebook verifies the
first, the second and the fourth of these relations in its cell 166,
with the output \texttt{\{True,\ True,\ True\}}. Two further properties
make them \emph{chiral} projectors:

\begin{itemize}
\tightlist
\item
  \(\gamma^AP^{L}=P^{R}\gamma^A\) and \(\gamma^AP^{R}=P^{L}\gamma^A\)
  (check \texttt{44-P4}): every Dirac matrix turns a left-handed spinor
  into a right-handed one, and conversely;
\item
  \([P^{L},\Sigma^{AB}]=[P^{R},\Sigma^{AB}]=0\) (check \texttt{44-P5}):
  the splitting \(\psi=P^{L}\psi+P^{R}\psi\) is the same in every
  Lorentz frame.
\end{itemize}

In this basis the splitting is simply ``upper eight components, lower
eight components''.

\hypertarget{other-important-1616-matrices-of-44}{%
\subsection{1.7 Other important 16×16 matrices of
4+4}\label{other-important-1616-matrices-of-44}}

\textbf{The two transposition matrices.} Let
\[C^{-}\equiv\gamma^0\gamma^1\gamma^2\gamma^3=\begin{pmatrix}-\sigma_{88}&0\\0&\sigma_{88}\end{pmatrix},\qquad
C^{+}\equiv\gamma^4\gamma^5\gamma^6\gamma^7=\begin{pmatrix}\sigma_{88}&0\\0&\sigma_{88}\end{pmatrix},\qquad
\omega=C^{-}C^{+}\] (checks \texttt{44-S1}, \texttt{44-S2}).
\(C_{1644}^{-}\) is the matrix \texttt{σ16} of \texttt{4+4.nb};
\(C_{1644}^{+}\) is the matrix \texttt{S44} of \texttt{7+1.nb}. Both are
real, symmetric and square to \(1\), and
\[C^{-}\gamma^A(C^{-})^{-1}=-(\gamma^A)^T,\qquad C^{+}\gamma^A(C^{+})^{-1}=+(\gamma^A)^T\tag{1.9}\]
(check \texttt{44-S3}). They play two roles: \(C^{-}\) defines the Dirac
adjoint of Sec. 1.8, and both define the bilinear forms
\(\psi^TC^{\mp}\,\cdots\,\psi\) from which fermion pairs are built (Part
3).

\[
C_{1644}^{-} = \left(\begin{array}{rrrrrrrr|rrrrrrrr}
\phantom{-}\cdot & \phantom{-}\cdot & \phantom{-}\cdot & \phantom{-}\cdot & -1 & \phantom{-}\cdot & \phantom{-}\cdot & \phantom{-}\cdot & \phantom{-}\cdot & \phantom{-}\cdot & \phantom{-}\cdot & \phantom{-}\cdot & \phantom{-}\cdot & \phantom{-}\cdot & \phantom{-}\cdot & \phantom{-}\cdot \\[-1.8pt]
\phantom{-}\cdot & \phantom{-}\cdot & \phantom{-}\cdot & \phantom{-}\cdot & \phantom{-}\cdot & -1 & \phantom{-}\cdot & \phantom{-}\cdot & \phantom{-}\cdot & \phantom{-}\cdot & \phantom{-}\cdot & \phantom{-}\cdot & \phantom{-}\cdot & \phantom{-}\cdot & \phantom{-}\cdot & \phantom{-}\cdot \\[-1.8pt]
\phantom{-}\cdot & \phantom{-}\cdot & \phantom{-}\cdot & \phantom{-}\cdot & \phantom{-}\cdot & \phantom{-}\cdot & -1 & \phantom{-}\cdot & \phantom{-}\cdot & \phantom{-}\cdot & \phantom{-}\cdot & \phantom{-}\cdot & \phantom{-}\cdot & \phantom{-}\cdot & \phantom{-}\cdot & \phantom{-}\cdot \\[-1.8pt]
\phantom{-}\cdot & \phantom{-}\cdot & \phantom{-}\cdot & \phantom{-}\cdot & \phantom{-}\cdot & \phantom{-}\cdot & \phantom{-}\cdot & -1 & \phantom{-}\cdot & \phantom{-}\cdot & \phantom{-}\cdot & \phantom{-}\cdot & \phantom{-}\cdot & \phantom{-}\cdot & \phantom{-}\cdot & \phantom{-}\cdot \\[-1.8pt]
-1 & \phantom{-}\cdot & \phantom{-}\cdot & \phantom{-}\cdot & \phantom{-}\cdot & \phantom{-}\cdot & \phantom{-}\cdot & \phantom{-}\cdot & \phantom{-}\cdot & \phantom{-}\cdot & \phantom{-}\cdot & \phantom{-}\cdot & \phantom{-}\cdot & \phantom{-}\cdot & \phantom{-}\cdot & \phantom{-}\cdot \\[-1.8pt]
\phantom{-}\cdot & -1 & \phantom{-}\cdot & \phantom{-}\cdot & \phantom{-}\cdot & \phantom{-}\cdot & \phantom{-}\cdot & \phantom{-}\cdot & \phantom{-}\cdot & \phantom{-}\cdot & \phantom{-}\cdot & \phantom{-}\cdot & \phantom{-}\cdot & \phantom{-}\cdot & \phantom{-}\cdot & \phantom{-}\cdot \\[-1.8pt]
\phantom{-}\cdot & \phantom{-}\cdot & -1 & \phantom{-}\cdot & \phantom{-}\cdot & \phantom{-}\cdot & \phantom{-}\cdot & \phantom{-}\cdot & \phantom{-}\cdot & \phantom{-}\cdot & \phantom{-}\cdot & \phantom{-}\cdot & \phantom{-}\cdot & \phantom{-}\cdot & \phantom{-}\cdot & \phantom{-}\cdot \\[-1.8pt]
\phantom{-}\cdot & \phantom{-}\cdot & \phantom{-}\cdot & -1 & \phantom{-}\cdot & \phantom{-}\cdot & \phantom{-}\cdot & \phantom{-}\cdot & \phantom{-}\cdot & \phantom{-}\cdot & \phantom{-}\cdot & \phantom{-}\cdot & \phantom{-}\cdot & \phantom{-}\cdot & \phantom{-}\cdot & \phantom{-}\cdot \\[-1.8pt]
\hline
\phantom{-}\cdot & \phantom{-}\cdot & \phantom{-}\cdot & \phantom{-}\cdot & \phantom{-}\cdot & \phantom{-}\cdot & \phantom{-}\cdot & \phantom{-}\cdot & \phantom{-}\cdot & \phantom{-}\cdot & \phantom{-}\cdot & \phantom{-}\cdot & \phantom{-}1 & \phantom{-}\cdot & \phantom{-}\cdot & \phantom{-}\cdot \\[-1.8pt]
\phantom{-}\cdot & \phantom{-}\cdot & \phantom{-}\cdot & \phantom{-}\cdot & \phantom{-}\cdot & \phantom{-}\cdot & \phantom{-}\cdot & \phantom{-}\cdot & \phantom{-}\cdot & \phantom{-}\cdot & \phantom{-}\cdot & \phantom{-}\cdot & \phantom{-}\cdot & \phantom{-}1 & \phantom{-}\cdot & \phantom{-}\cdot \\[-1.8pt]
\phantom{-}\cdot & \phantom{-}\cdot & \phantom{-}\cdot & \phantom{-}\cdot & \phantom{-}\cdot & \phantom{-}\cdot & \phantom{-}\cdot & \phantom{-}\cdot & \phantom{-}\cdot & \phantom{-}\cdot & \phantom{-}\cdot & \phantom{-}\cdot & \phantom{-}\cdot & \phantom{-}\cdot & \phantom{-}1 & \phantom{-}\cdot \\[-1.8pt]
\phantom{-}\cdot & \phantom{-}\cdot & \phantom{-}\cdot & \phantom{-}\cdot & \phantom{-}\cdot & \phantom{-}\cdot & \phantom{-}\cdot & \phantom{-}\cdot & \phantom{-}\cdot & \phantom{-}\cdot & \phantom{-}\cdot & \phantom{-}\cdot & \phantom{-}\cdot & \phantom{-}\cdot & \phantom{-}\cdot & \phantom{-}1 \\[-1.8pt]
\phantom{-}\cdot & \phantom{-}\cdot & \phantom{-}\cdot & \phantom{-}\cdot & \phantom{-}\cdot & \phantom{-}\cdot & \phantom{-}\cdot & \phantom{-}\cdot & \phantom{-}1 & \phantom{-}\cdot & \phantom{-}\cdot & \phantom{-}\cdot & \phantom{-}\cdot & \phantom{-}\cdot & \phantom{-}\cdot & \phantom{-}\cdot \\[-1.8pt]
\phantom{-}\cdot & \phantom{-}\cdot & \phantom{-}\cdot & \phantom{-}\cdot & \phantom{-}\cdot & \phantom{-}\cdot & \phantom{-}\cdot & \phantom{-}\cdot & \phantom{-}\cdot & \phantom{-}1 & \phantom{-}\cdot & \phantom{-}\cdot & \phantom{-}\cdot & \phantom{-}\cdot & \phantom{-}\cdot & \phantom{-}\cdot \\[-1.8pt]
\phantom{-}\cdot & \phantom{-}\cdot & \phantom{-}\cdot & \phantom{-}\cdot & \phantom{-}\cdot & \phantom{-}\cdot & \phantom{-}\cdot & \phantom{-}\cdot & \phantom{-}\cdot & \phantom{-}\cdot & \phantom{-}1 & \phantom{-}\cdot & \phantom{-}\cdot & \phantom{-}\cdot & \phantom{-}\cdot & \phantom{-}\cdot \\[-1.8pt]
\phantom{-}\cdot & \phantom{-}\cdot & \phantom{-}\cdot & \phantom{-}\cdot & \phantom{-}\cdot & \phantom{-}\cdot & \phantom{-}\cdot & \phantom{-}\cdot & \phantom{-}\cdot & \phantom{-}\cdot & \phantom{-}\cdot & \phantom{-}1 & \phantom{-}\cdot & \phantom{-}\cdot & \phantom{-}\cdot & \phantom{-}\cdot
\end{array}\right)
\]

\[
C_{1644}^{+} = \left(\begin{array}{rrrrrrrr|rrrrrrrr}
\phantom{-}\cdot & \phantom{-}\cdot & \phantom{-}\cdot & \phantom{-}\cdot & \phantom{-}1 & \phantom{-}\cdot & \phantom{-}\cdot & \phantom{-}\cdot & \phantom{-}\cdot & \phantom{-}\cdot & \phantom{-}\cdot & \phantom{-}\cdot & \phantom{-}\cdot & \phantom{-}\cdot & \phantom{-}\cdot & \phantom{-}\cdot \\[-1.8pt]
\phantom{-}\cdot & \phantom{-}\cdot & \phantom{-}\cdot & \phantom{-}\cdot & \phantom{-}\cdot & \phantom{-}1 & \phantom{-}\cdot & \phantom{-}\cdot & \phantom{-}\cdot & \phantom{-}\cdot & \phantom{-}\cdot & \phantom{-}\cdot & \phantom{-}\cdot & \phantom{-}\cdot & \phantom{-}\cdot & \phantom{-}\cdot \\[-1.8pt]
\phantom{-}\cdot & \phantom{-}\cdot & \phantom{-}\cdot & \phantom{-}\cdot & \phantom{-}\cdot & \phantom{-}\cdot & \phantom{-}1 & \phantom{-}\cdot & \phantom{-}\cdot & \phantom{-}\cdot & \phantom{-}\cdot & \phantom{-}\cdot & \phantom{-}\cdot & \phantom{-}\cdot & \phantom{-}\cdot & \phantom{-}\cdot \\[-1.8pt]
\phantom{-}\cdot & \phantom{-}\cdot & \phantom{-}\cdot & \phantom{-}\cdot & \phantom{-}\cdot & \phantom{-}\cdot & \phantom{-}\cdot & \phantom{-}1 & \phantom{-}\cdot & \phantom{-}\cdot & \phantom{-}\cdot & \phantom{-}\cdot & \phantom{-}\cdot & \phantom{-}\cdot & \phantom{-}\cdot & \phantom{-}\cdot \\[-1.8pt]
\phantom{-}1 & \phantom{-}\cdot & \phantom{-}\cdot & \phantom{-}\cdot & \phantom{-}\cdot & \phantom{-}\cdot & \phantom{-}\cdot & \phantom{-}\cdot & \phantom{-}\cdot & \phantom{-}\cdot & \phantom{-}\cdot & \phantom{-}\cdot & \phantom{-}\cdot & \phantom{-}\cdot & \phantom{-}\cdot & \phantom{-}\cdot \\[-1.8pt]
\phantom{-}\cdot & \phantom{-}1 & \phantom{-}\cdot & \phantom{-}\cdot & \phantom{-}\cdot & \phantom{-}\cdot & \phantom{-}\cdot & \phantom{-}\cdot & \phantom{-}\cdot & \phantom{-}\cdot & \phantom{-}\cdot & \phantom{-}\cdot & \phantom{-}\cdot & \phantom{-}\cdot & \phantom{-}\cdot & \phantom{-}\cdot \\[-1.8pt]
\phantom{-}\cdot & \phantom{-}\cdot & \phantom{-}1 & \phantom{-}\cdot & \phantom{-}\cdot & \phantom{-}\cdot & \phantom{-}\cdot & \phantom{-}\cdot & \phantom{-}\cdot & \phantom{-}\cdot & \phantom{-}\cdot & \phantom{-}\cdot & \phantom{-}\cdot & \phantom{-}\cdot & \phantom{-}\cdot & \phantom{-}\cdot \\[-1.8pt]
\phantom{-}\cdot & \phantom{-}\cdot & \phantom{-}\cdot & \phantom{-}1 & \phantom{-}\cdot & \phantom{-}\cdot & \phantom{-}\cdot & \phantom{-}\cdot & \phantom{-}\cdot & \phantom{-}\cdot & \phantom{-}\cdot & \phantom{-}\cdot & \phantom{-}\cdot & \phantom{-}\cdot & \phantom{-}\cdot & \phantom{-}\cdot \\[-1.8pt]
\hline
\phantom{-}\cdot & \phantom{-}\cdot & \phantom{-}\cdot & \phantom{-}\cdot & \phantom{-}\cdot & \phantom{-}\cdot & \phantom{-}\cdot & \phantom{-}\cdot & \phantom{-}\cdot & \phantom{-}\cdot & \phantom{-}\cdot & \phantom{-}\cdot & \phantom{-}1 & \phantom{-}\cdot & \phantom{-}\cdot & \phantom{-}\cdot \\[-1.8pt]
\phantom{-}\cdot & \phantom{-}\cdot & \phantom{-}\cdot & \phantom{-}\cdot & \phantom{-}\cdot & \phantom{-}\cdot & \phantom{-}\cdot & \phantom{-}\cdot & \phantom{-}\cdot & \phantom{-}\cdot & \phantom{-}\cdot & \phantom{-}\cdot & \phantom{-}\cdot & \phantom{-}1 & \phantom{-}\cdot & \phantom{-}\cdot \\[-1.8pt]
\phantom{-}\cdot & \phantom{-}\cdot & \phantom{-}\cdot & \phantom{-}\cdot & \phantom{-}\cdot & \phantom{-}\cdot & \phantom{-}\cdot & \phantom{-}\cdot & \phantom{-}\cdot & \phantom{-}\cdot & \phantom{-}\cdot & \phantom{-}\cdot & \phantom{-}\cdot & \phantom{-}\cdot & \phantom{-}1 & \phantom{-}\cdot \\[-1.8pt]
\phantom{-}\cdot & \phantom{-}\cdot & \phantom{-}\cdot & \phantom{-}\cdot & \phantom{-}\cdot & \phantom{-}\cdot & \phantom{-}\cdot & \phantom{-}\cdot & \phantom{-}\cdot & \phantom{-}\cdot & \phantom{-}\cdot & \phantom{-}\cdot & \phantom{-}\cdot & \phantom{-}\cdot & \phantom{-}\cdot & \phantom{-}1 \\[-1.8pt]
\phantom{-}\cdot & \phantom{-}\cdot & \phantom{-}\cdot & \phantom{-}\cdot & \phantom{-}\cdot & \phantom{-}\cdot & \phantom{-}\cdot & \phantom{-}\cdot & \phantom{-}1 & \phantom{-}\cdot & \phantom{-}\cdot & \phantom{-}\cdot & \phantom{-}\cdot & \phantom{-}\cdot & \phantom{-}\cdot & \phantom{-}\cdot \\[-1.8pt]
\phantom{-}\cdot & \phantom{-}\cdot & \phantom{-}\cdot & \phantom{-}\cdot & \phantom{-}\cdot & \phantom{-}\cdot & \phantom{-}\cdot & \phantom{-}\cdot & \phantom{-}\cdot & \phantom{-}1 & \phantom{-}\cdot & \phantom{-}\cdot & \phantom{-}\cdot & \phantom{-}\cdot & \phantom{-}\cdot & \phantom{-}\cdot \\[-1.8pt]
\phantom{-}\cdot & \phantom{-}\cdot & \phantom{-}\cdot & \phantom{-}\cdot & \phantom{-}\cdot & \phantom{-}\cdot & \phantom{-}\cdot & \phantom{-}\cdot & \phantom{-}\cdot & \phantom{-}\cdot & \phantom{-}1 & \phantom{-}\cdot & \phantom{-}\cdot & \phantom{-}\cdot & \phantom{-}\cdot & \phantom{-}\cdot \\[-1.8pt]
\phantom{-}\cdot & \phantom{-}\cdot & \phantom{-}\cdot & \phantom{-}\cdot & \phantom{-}\cdot & \phantom{-}\cdot & \phantom{-}\cdot & \phantom{-}\cdot & \phantom{-}\cdot & \phantom{-}\cdot & \phantom{-}\cdot & \phantom{-}1 & \phantom{-}\cdot & \phantom{-}\cdot & \phantom{-}\cdot & \phantom{-}\cdot
\end{array}\right)
\]

\textbf{The complex structure.} The matrix \(\gamma_{1644}^4\) (called
\texttt{JT4} in \texttt{4+4.nb}, cell 670) is real, squares to \(-1\),
and satisfies \((\gamma^4)^TC^{-}\gamma^4=C^{-}\) (check
\texttt{44-S5}). It will reappear as the generator of the evolution in
the time \(x4\) of a particle at rest.

\textbf{The 8×8 blocks.} \(\sigma_{88}\), \(\tau_{1644}^A\) and
\(\bar\tau_{1644}^A\) are listed in Appendix B.

\hypertarget{the-dirac-equation-of-the-field-dirac1644}{%
\subsection{1.8 The Dirac equation of the field
dirac1644}\label{the-dirac-equation-of-the-field-dirac1644}}

\hypertarget{the-equation}{%
\subsubsection{1.8.1 The equation}\label{the-equation}}

The field \textbf{dirac1644} is the column
\[\psi_{1644}(x)=\big(\psi_1(x),\dots,\psi_{16}(x)\big)^T,\qquad x=(x0,\dots,x7),\]
whose sixteen components are complex-valued Grassmann operators
(anticommuting, and not Hermitian: \(\psi_a^\dagger\neq\psi_a\)). Its
Dirac equation is
\[\boxed{\;\big(\gamma_{1644}^A\,\partial_A-m_{1644}\big)\,\psi_{1644}=0\;}\tag{1.10}\]
with a real mass \(m_{1644}>0\) and the purely real matrices
\(\gamma_{1644}^A\) of Sec. 1.2.

\emph{Why there is no factor \(i\).} The Dirac operator must be a square
root of the wave operator. By Eq. (1.3),
\[(\gamma^A\partial_A)^2=\eta^{AB}\partial_A\partial_B\equiv\Box_{1644}
=\partial_0^2+\partial_1^2+\partial_2^2+\partial_3^2-\partial_4^2-\partial_5^2-\partial_6^2-\partial_7^2,\]
so that
\[\big(\gamma^A\partial_A+m\big)\big(\gamma^B\partial_B-m\big)=\Box_{1644}-m^2 .\tag{1.11}\]
Every component of a solution of (1.10) therefore satisfies the
Klein--Gordon equation \((\Box_{1644}-m^2)\psi_a=0\). For a plane wave
\(\psi\propto e^{ik_Ax^A}\) this reads
\[k_4^2+k_5^2+k_6^2+k_7^2=k_0^2+k_1^2+k_2^2+k_3^2+m^2,\tag{1.12}\] which
is the mass shell of a particle of mass \(m\) when \(x4,\dots,x7\) are
the time-like coordinates. With the metric convention of \texttt{4+4.nb}
(space-like \(=+1\)) the correct sign of \(m^2\) is obtained
\emph{without} an \(i\); an operator \(i\gamma^A\partial_A-m\) would
give \(\Box_{1644}+m^2\), that is, Eq. (1.12) with space and time
exchanged. The operator in Eq. (1.10) is real; the field is complex
because a Lagrangian with a mass term requires it (Sec. 1.8.3).

\textbf{Checks.} In momentum space (\(\partial_A\to ik_A\)) the
statements above are identities between Gaussian-integer matrices for
integer \(k_A\) and \(m\): checks \texttt{44-D1}, \texttt{44-D2}. Check
\texttt{44-D3} is a worked example: for \(k=(0,2,2,4;\,5,0,0,0)\) and
\(m=1\), which satisfies (1.12) because \(25=24+1\), the matrix
\(\Pi_{1644}=(i\gamma^Ak_A+m)/2m\) is a projector of trace 8 with
\((i\gamma^Ak_A-m)\Pi_{1644}=0\): the plane-wave solutions
\(u\,e^{ik_Ax^A}\), \(u=\Pi_{1644}w\), form an 8-dimensional space.

\hypertarget{lorentz-covariance-and-the-conserved-current}{%
\subsubsection{1.8.2 Lorentz covariance and the conserved
current}\label{lorentz-covariance-and-the-conserved-current}}

\emph{Covariance.} Let \(S\in\mathrm{Spin}_0(4,4)\) with
\(S\gamma^AS^{-1}=\Lambda^A{}_B\gamma^B\) (Theorem 1), so that
\(S^{-1}\gamma^AS=(\Lambda^{-1})^A{}_B\gamma^B\). If \(\psi(x)\) solves
(1.10), then so does
\[\psi'(x)=S\,\psi(\Lambda x),\qquad(\Lambda x)^A=\Lambda^A{}_Bx^B .\]
Indeed
\(\partial_A\psi'(x)=\Lambda^B{}_A\,S\,(\partial_B\psi)(\Lambda x)\),
hence
\[\begin{aligned}\gamma^A\partial_A\psi'(x)&=S\,(S^{-1}\gamma^AS)\,\Lambda^B{}_A\,(\partial_B\psi)(\Lambda x)
=S\,(\Lambda^{-1})^A{}_C\,\Lambda^B{}_A\,\gamma^C(\partial_B\psi)(\Lambda x)\\
&=S\,\gamma^B(\partial_B\psi)(\Lambda x)=m\,S\,\psi(\Lambda x)=m\,\psi'(x).\end{aligned}\]
The relation \(\Lambda\eta\Lambda^T=\eta\) of Theorem 1 states that
\(x\mapsto\Lambda x\) preserves the interval \(\eta_{AB}x^Ax^B\),
i.e.~that it is a symmetry of the 4+4 spacetime.

\emph{Adjoint.} Define \(\tilde\psi\equiv\psi^\dagger C^{-}\). From Eq.
(1.9) and the reality of the matrices,
\((\gamma^A)^\dagger C^{-}=-C^{-}\gamma^A\). The conjugate transpose of
(1.10) is \(\partial_A\psi^\dagger(\gamma^A)^\dagger=m\psi^\dagger\);
multiplying by \(C^{-}\) from the right gives the adjoint equation
\(-\partial_A\tilde\psi\,\gamma^A=m\tilde\psi\).

\emph{Current.} Adding the two equations,
\[\partial_A\big(\tilde\psi\gamma^A\psi\big)=(\partial_A\tilde\psi\gamma^A)\psi+\tilde\psi\gamma^A\partial_A\psi=-m\tilde\psi\psi+m\tilde\psi\psi=0 .\]
Since \(C^{-}\gamma^A\) is a real antisymmetric matrix (check
\texttt{44-S4}), \(\tilde\psi\gamma^A\psi\) is anti-Hermitian, and the
Hermitian conserved current is
\[j_{1644}^A=i\,\tilde\psi\,\gamma^A\psi=\psi^\dagger\big(i\,C^{-}\gamma^A\big)\psi,\qquad\partial_Aj_{1644}^A=0 .\tag{1.13}\]
It is the Noether current of the phase symmetry
\(\psi\to e^{i\vartheta}\psi\) (\emph{fermion number}).

\hypertarget{lagrangian}{%
\subsubsection{1.8.3 Lagrangian}\label{lagrangian}}

\[\mathcal{L}_{1644}=\tilde\psi\big(\gamma^A\partial_A-m\big)\psi,\qquad\tilde\psi=\psi^\dagger C^{-} .\tag{1.14}\]
It is invariant under \(\mathrm{Spin}_0(4,4)\), because
\(S^TC^{-}S=C^{-}\) for the real matrices \(S\) of Theorem 1 (a
consequence of Eq. (1.9): \((\Sigma^{AB})^TC^{-}=-C^{-}\Sigma^{AB}\));
it is Hermitian up to a total derivative, because \(C^{-}\gamma^A\) is
antisymmetric and \(C^{-}\) symmetric; and varying \(\tilde\psi\) gives
Eq. (1.10).

\emph{Why the components must be complex.} For a real (Majorana)
anticommuting field the mass term would have to be \(\psi^TC\psi\) with
an invariant matrix \(C\). The invariant matrices are the combinations
of \(C^{-}\) and \(C^{+}\), and both are symmetric (Sec. 1.7). For
anticommuting components
\(\psi^TC\psi=\sum_{ab}\psi_aC_{ab}\psi_b=-\sum_{ab}\psi_bC_{ab}\psi_a=-\psi^TC^T\psi\),
so \(\psi^TC^{\mp}\psi=0\): a real 16-component Grassmann field in 4+4
is necessarily massless. The complex field avoids this because
\(\psi^\dagger C^{-}\psi\) does not vanish.

\hypertarget{evolution-in-the-ordinary-time-x4-and-what-quantized-can-mean}{%
\subsubsection{\texorpdfstring{1.8.4 Evolution in the ordinary time
\(x4\), and what ``quantized'' can
mean}{1.8.4 Evolution in the ordinary time x4, and what ``quantized'' can mean}}\label{evolution-in-the-ordinary-time-x4-and-what-quantized-can-mean}}

Here the four time-like dimensions show their consequences, and they
must be stated without embellishment.

Solve Eq. (1.10) for \(\partial_4\psi\), using
\((\gamma^4)^{-1}=-\gamma^4\). For a field proportional to
\(\exp\!\big(i\sum_{a=0}^3k_ax^a+i\sum_{t=5}^7\kappa_tx^t\big)\),
\[\partial_4\psi=K\psi,\qquad K=-m\gamma^4+i\sum_{a=0}^{3}k_a\,\gamma^4\gamma^a+i\sum_{t=5}^{7}\kappa_t\,\gamma^4\gamma^t .\tag{1.15}\]
Here \(k_a\) (\(a=0,1,2,3\)) are the four momenta along the space-like
coordinates and \(\kappa_t\) (\(t=5,6,7\)) the three frequencies along
the additional times. The matrices \(\gamma^4\gamma^a\) are real
symmetric, while \(\gamma^4\gamma^t\) and \(\gamma^4\) are real
antisymmetric (check \texttt{44-D6}). Two exact results follow (checks
\texttt{44-D7}, \texttt{44-D8}):
\[K^2=-\big(m^2+|k|^2-|\kappa|^2\big)\,1,\qquad K+K^\dagger=2i\sum_{t=5}^{7}\kappa_t\,\gamma^4\gamma^t .\tag{1.16}\]

\begin{enumerate}
\def\labelenumi{\arabic{enumi}.}
\tightlist
\item
  \textbf{If the field does not depend on \(x5,x6,x7\)} (\(\kappa=0\)),
  \(K\) is anti-Hermitian:
  \(\mathsf{h}\equiv iK=-im\gamma^4-k_a\gamma^4\gamma^a\) is a Hermitian
  matrix with \(\mathsf{h}^2=(m^2+|k|^2)1\). The evolution in \(x4\) is
  unitary, the positive quantity \(\psi^\dagger\psi\) is conserved, and
  the energies are \(\pm E\), \(E=\sqrt{m^2+|k|^2}\), each eight times.
\item
  \textbf{If the field depends on the additional times}
  (\(\kappa\neq0\)), \(K\) is not anti-Hermitian and
  \(\psi^\dagger\psi\) is not conserved. For \(|\kappa|^2>m^2+|k|^2\)
  the frequencies in \(x4\) are imaginary: the modes grow exponentially.
  Check \texttt{44-D9} exhibits one: \(m=3\), \(k=0\),
  \(\kappa=(5,0,0)\) gives \(K^2=+16\), so eight solutions grow as
  \(e^{+4\,x4}\).
\item
  \textbf{The covariant charge is not positive.} The density of the
  conserved current (1.13) is \(j_{1644}^4=\psi^\dagger W\psi\) with
  \(W\equiv i\,C^{-}\gamma^4\), a Hermitian matrix with \(W^2=1\) and
  trace 0, i.e.~with eigenvalues \(+1\) and \(-1\), eight times each
  (checks \texttt{44-D10}, \texttt{44-D11}).
\end{enumerate}

These are not defects of the matrices; they are the known price of more
than one time dimension: the wave operator is ultra-hyperbolic, and no
positive conserved density is invariant under the whole group
\(\mathrm{Spin}(4,4)\).

\textbf{Definition of the quantized field in this project (Stage 1).}
The author's hypothesis is that the three additional times deflate away.
In flat spacetime, without gravity, we implement this hypothesis as a
restriction: \emph{the field dirac1644 does not depend on \(x5\),
\(x6\), \(x7\)}. On such fields \(W\) commutes with \(K\) (check
\texttt{44-D10}), the Lagrangian (1.14) equals
\(-\psi^\dagger W(i\partial_4-\mathsf{h})\psi\), and we quantize with
the positive Lagrangian
\[\mathcal{L}^{\rm red}_{1644}=\psi^\dagger\big(i\partial_4-\mathsf{h}\big)\psi,\qquad\mathsf{h}=-im\gamma^4+i\gamma^4\gamma^a\partial_a,\quad a=0,1,2,3,\tag{1.17}\]
which gives the canonical anticommutators
\[\{\psi_a(x4,\vec x),\psi_b^\dagger(x4,\vec y)\}=\delta_{ab}\,\delta^4(\vec x-\vec y),\qquad\{\psi_a,\psi_b\}=0,\qquad\vec x=(x0,x1,x2,x3).\tag{1.18}\]
\(\mathcal{L}^{\rm red}_{1644}\) agrees with (1.14) on the half of the
components with \(W=-1\) and has the opposite sign on the half with
\(W=+1\). It is invariant under the rotations of the four space
coordinates, under boosts between them and \(x4\), and under rotations
of the three additional times, i.e.~under
\(\mathrm{Spin}(4,1)\times\mathrm{Spin}(3)\), because \(W\) commutes
with the generators of these transformations; it is not invariant under
boosts into the additional times, with whose generators \(W\)
anticommutes. Whether gravity (Stage 2 of the author's programme)
justifies this restriction is outside the present document; it is
recorded as an assumption wherever it is used.

\hypertarget{the-two-coupled-8-component-equations}{%
\subsubsection{1.8.5 The two coupled 8-component
equations}\label{the-two-coupled-8-component-equations}}

Write \(\psi_{1644}=(\psi^{L},\psi^{R})\), where \(\psi^{L}\) consists
of the upper and \(\psi^{R}\) of the lower eight components; by Eq.
(1.8), \(P^{L}\psi=(\psi^{L},0)\) and \(P^{R}\psi=(0,\psi^{R})\).
Inserting the block form (1.2) into (1.10) gives (check \texttt{44-D4})
\[\boxed{\;\bar\tau_{1644}^A\,\partial_A\,\psi_{1644}^{R}=m_{1644}\,\psi_{1644}^{L},\qquad\tau_{1644}^A\,\partial_A\,\psi_{1644}^{L}=m_{1644}\,\psi_{1644}^{R}\;}\tag{1.19}\]
with the real 8×8 matrices of Sec. 1.1 (\(\tau^0=\bar\tau^0=1_8\), so
the first terms are \(\partial_0\psi^{R}\) and \(\partial_0\psi^{L}\)).
The mass couples the two 8-component fields; for \(m=0\) they decouple
into \(\bar\tau^A\partial_A\psi^{R}=0\) and
\(\tau^A\partial_A\psi^{L}=0\). Each field obeys the Klein--Gordon
equation: applying \(\tau^B\partial_B\) to the first equation and using
the second and Eq. (1.1),
\(\Box_{1644}\psi^{R}=m\,\tau^B\partial_B\psi^{L}=m^2\psi^{R}\) (check
\texttt{44-D5}). Because the projectors and all matrices are real, Eq.
(1.19) is a pair of real equations: the Hermitian and the anti-Hermitian
parts of \(\psi^{L}\), \(\psi^{R}\) do not mix.

\begin{center}\rule{0.5\linewidth}{0.5pt}\end{center}

\hypertarget{part-2.-signature-71-the-matrices-gamma_1671a-of-71.nb}{%
\section{\texorpdfstring{Part 2. Signature 7+1: the matrices
\(\gamma_{1671}^A\) of
\texttt{7+1.nb}}{Part 2. Signature 7+1: the matrices \textbackslash gamma\_\{1671\}\^{}A of 7+1.nb}}\label{part-2.-signature-71-the-matrices-gamma_1671a-of-71.nb}}

\emph{Scope rule: in this Part \(\gamma^A\), \(\eta\), \(\omega\),
\(\chi\), \(P^{L}\), \(P^{R}\), \(C^{\pm}\), \(\Sigma^{AB}\),
\(\tau^A\), \(\alpha^j\), \(\beta\), \(\mathsf{h}\), \(E^{L}\),
\(E^{R}\), \(\phi^{L}\), \(\phi^{R}\), \(\psi\), \(m\) stand for the
objects with the tag 1671. Symbols of Part 1 that are needed here carry
their tag 1644.}

\hypertarget{definition-and-the-relation-to-the-notebook}{%
\subsection{2.1 Definition, and the relation to the
notebook}\label{definition-and-the-relation-to-the-notebook}}

\texttt{7+1.nb} builds its matrices from those of \texttt{4+4.nb} in
three steps (its sections 3.1, 4.1 and 4.2; cells 90, 187 and 197).

\textbf{Step 1: octonionic 8×8 matrices.} \texttt{τO{[}0{]}} \(=1_8\);
\texttt{τO{[}A{]}} \(=\tau_{1644}^A\) for \(A=1,2,3\);
\texttt{τO{[}A{]}} \(=\tau_{1644}^A\sigma_{88}\) for \(A=4,\dots,7\).
The seven matrices \texttt{τO{[}1..7{]}} are real, antisymmetric,
mutually anticommuting, and square to \(-1\) (check \texttt{71-B1}).

\textbf{Step 2: sixteen-dimensional Euclidean generators.}
\texttt{T16O{[}A{]}}
\(=\begin{pmatrix}0&\texttt{τO[A]}^T\\\texttt{τO[A]}&0\end{pmatrix}\).

\textbf{Step 3: the 7+1 matrices of the notebook.} \texttt{Γ17{[}0{]}}
\(=\) \texttt{T16O{[}0{]}} and \texttt{Γ17{[}i{]}} \(=\)
\texttt{T16O{[}0{]}.T16O{[}i{]}} for \(i=1,\dots,7\). They obey
\texttt{Γ17{[}A{]}.Γ17{[}B{]}\ +\ Γ17{[}B{]}.Γ17{[}A{]}}
\(=2\,\)\texttt{η17{[}{[}A+1,B+1{]}{]}} with \texttt{η17}
\(=\mathrm{diag}(+1,-1,\dots,-1)\) (check \texttt{71-N3}), and the
notebook raises the index with this metric: \texttt{Γ17up{[}A{]}} \(=\)
\texttt{η17{[}{[}A+1,A+1{]}{]}\ Γ17{[}A{]}} (its sections 5.3 and 7.1),
so \texttt{Γ17up{[}0{]}\ =\ Γ17{[}0{]}} and
\texttt{Γ17up{[}i{]}\ =\ -Γ17{[}i{]}}. The notebook's Dirac equation is
written with \texttt{Γ17up}.

\textbf{The matrices of this project.} \[\boxed{\;\begin{aligned}
&\gamma_{1671}^A\equiv\texttt{Γ17up}[\mathrm{swap}_{04}(A)],\qquad A=0,1,\dots,7,\qquad\text{that is,}\\
&\gamma_{1671}^4=\texttt{Γ17[0]},\qquad\gamma_{1671}^0=-\texttt{Γ17[4]},\qquad\gamma_{1671}^j=-\texttt{Γ17[}j\texttt{]}\quad(j=1,2,3,5,6,7)
\end{aligned}\;}\tag{2.1}\] (check \texttt{71-N2}). Thus
\(\gamma_{1671}^4\) belongs to the time \(x4\), and the matrix that the
notebook attaches to its space coordinate ``x4'' now belongs to the
hidden coordinate \(x0\). The eight matrices \texttt{Γ17{[}0..7{]}}
rebuilt by the program are identical to those printed by Mathematica in
cell 204 of \texttt{7+1.nb} (check \texttt{71-N1}).

\textbf{In 8×8 blocks.} Let \(\tau_{1671}^A\equiv\)
\texttt{τO{[}swap04(A){]}}, so that \(\tau^4=1_8\) and the seven
\(\tau^j\) with \(j\in\{0,1,2,3,5,6,7\}\) are real antisymmetric with
\(\tau^i\tau^j+\tau^j\tau^i=-2\delta^{ij}\) (check \texttt{71-D6}). Then
\[\gamma_{1671}^4=\begin{pmatrix}0&1_8\\1_8&0\end{pmatrix},\qquad
\gamma_{1671}^j=\begin{pmatrix}-\tau_{1671}^j&0\\0&\tau_{1671}^j\end{pmatrix}\quad(j=0,1,2,3,5,6,7)\tag{2.2}\]
(check \texttt{71-B2}, with the sign of \texttt{Γ17up}; verified again
by \texttt{I71-D2}).

\textbf{Built from the 4+4 matrices} (check \texttt{71-B3}):
\texttt{Γ17{[}0{]}} \(=\gamma_{1644}^0\), \texttt{Γ17{[}i{]}}
\(=\gamma_{1644}^0\gamma_{1644}^i\) for \(i=1,2,3\), and
\texttt{Γ17{[}k{]}} \(=\gamma_{1644}^0\gamma_{1644}^kC_{1644}^{+}\) for
\(k=4,\dots,7\). In the labels of this project, and with every product
reduced to distinct factors (check \texttt{71-B4}): \[\begin{aligned}
\gamma_{1671}^4&=\gamma_{1644}^0,\qquad\gamma_{1671}^j=-\gamma_{1644}^0\gamma_{1644}^j\quad(j=1,2,3),\\
\gamma_{1671}^0&=+\gamma_{1644}^0\gamma_{1644}^5\gamma_{1644}^6\gamma_{1644}^7,\qquad
\gamma_{1671}^5=-\gamma_{1644}^0\gamma_{1644}^4\gamma_{1644}^6\gamma_{1644}^7,\\
\gamma_{1671}^6&=+\gamma_{1644}^0\gamma_{1644}^4\gamma_{1644}^5\gamma_{1644}^7,\qquad
\gamma_{1671}^7=-\gamma_{1644}^0\gamma_{1644}^4\gamma_{1644}^5\gamma_{1644}^6 .
\end{aligned}\tag{2.2a}\] Both sets therefore live in the same algebra
of real 16×16 matrices. \textbf{The time-like matrix of 7+1 is the
matrix of the hidden axis of 4+4}; the three matrices of ordinary space
are products of two, and the four matrices of the extra axes are
products of four, of the 4+4 matrices, and these seven are \emph{not}
linear combinations of the 4+4 matrices. (An earlier version of this
paragraph said that all eight are products of two and of six of the 4+4
matrices and that none is a linear combination of them; that was
imprecise, and wrong for \(\gamma_{1671}^4\). Error log, E36.)

\hypertarget{the-eight-matrices-1}{%
\subsection{2.2 The eight matrices}\label{the-eight-matrices-1}}

\[
\gamma_{1671}^{0} = \left(\begin{array}{rrrrrrrr|rrrrrrrr}
\phantom{-}\cdot & -1 & \phantom{-}\cdot & \phantom{-}\cdot & \phantom{-}\cdot & \phantom{-}\cdot & \phantom{-}\cdot & \phantom{-}\cdot & \phantom{-}\cdot & \phantom{-}\cdot & \phantom{-}\cdot & \phantom{-}\cdot & \phantom{-}\cdot & \phantom{-}\cdot & \phantom{-}\cdot & \phantom{-}\cdot \\[-1.8pt]
\phantom{-}1 & \phantom{-}\cdot & \phantom{-}\cdot & \phantom{-}\cdot & \phantom{-}\cdot & \phantom{-}\cdot & \phantom{-}\cdot & \phantom{-}\cdot & \phantom{-}\cdot & \phantom{-}\cdot & \phantom{-}\cdot & \phantom{-}\cdot & \phantom{-}\cdot & \phantom{-}\cdot & \phantom{-}\cdot & \phantom{-}\cdot \\[-1.8pt]
\phantom{-}\cdot & \phantom{-}\cdot & \phantom{-}\cdot & \phantom{-}1 & \phantom{-}\cdot & \phantom{-}\cdot & \phantom{-}\cdot & \phantom{-}\cdot & \phantom{-}\cdot & \phantom{-}\cdot & \phantom{-}\cdot & \phantom{-}\cdot & \phantom{-}\cdot & \phantom{-}\cdot & \phantom{-}\cdot & \phantom{-}\cdot \\[-1.8pt]
\phantom{-}\cdot & \phantom{-}\cdot & -1 & \phantom{-}\cdot & \phantom{-}\cdot & \phantom{-}\cdot & \phantom{-}\cdot & \phantom{-}\cdot & \phantom{-}\cdot & \phantom{-}\cdot & \phantom{-}\cdot & \phantom{-}\cdot & \phantom{-}\cdot & \phantom{-}\cdot & \phantom{-}\cdot & \phantom{-}\cdot \\[-1.8pt]
\phantom{-}\cdot & \phantom{-}\cdot & \phantom{-}\cdot & \phantom{-}\cdot & \phantom{-}\cdot & \phantom{-}1 & \phantom{-}\cdot & \phantom{-}\cdot & \phantom{-}\cdot & \phantom{-}\cdot & \phantom{-}\cdot & \phantom{-}\cdot & \phantom{-}\cdot & \phantom{-}\cdot & \phantom{-}\cdot & \phantom{-}\cdot \\[-1.8pt]
\phantom{-}\cdot & \phantom{-}\cdot & \phantom{-}\cdot & \phantom{-}\cdot & -1 & \phantom{-}\cdot & \phantom{-}\cdot & \phantom{-}\cdot & \phantom{-}\cdot & \phantom{-}\cdot & \phantom{-}\cdot & \phantom{-}\cdot & \phantom{-}\cdot & \phantom{-}\cdot & \phantom{-}\cdot & \phantom{-}\cdot \\[-1.8pt]
\phantom{-}\cdot & \phantom{-}\cdot & \phantom{-}\cdot & \phantom{-}\cdot & \phantom{-}\cdot & \phantom{-}\cdot & \phantom{-}\cdot & -1 & \phantom{-}\cdot & \phantom{-}\cdot & \phantom{-}\cdot & \phantom{-}\cdot & \phantom{-}\cdot & \phantom{-}\cdot & \phantom{-}\cdot & \phantom{-}\cdot \\[-1.8pt]
\phantom{-}\cdot & \phantom{-}\cdot & \phantom{-}\cdot & \phantom{-}\cdot & \phantom{-}\cdot & \phantom{-}\cdot & \phantom{-}1 & \phantom{-}\cdot & \phantom{-}\cdot & \phantom{-}\cdot & \phantom{-}\cdot & \phantom{-}\cdot & \phantom{-}\cdot & \phantom{-}\cdot & \phantom{-}\cdot & \phantom{-}\cdot \\[-1.8pt]
\hline
\phantom{-}\cdot & \phantom{-}\cdot & \phantom{-}\cdot & \phantom{-}\cdot & \phantom{-}\cdot & \phantom{-}\cdot & \phantom{-}\cdot & \phantom{-}\cdot & \phantom{-}\cdot & \phantom{-}1 & \phantom{-}\cdot & \phantom{-}\cdot & \phantom{-}\cdot & \phantom{-}\cdot & \phantom{-}\cdot & \phantom{-}\cdot \\[-1.8pt]
\phantom{-}\cdot & \phantom{-}\cdot & \phantom{-}\cdot & \phantom{-}\cdot & \phantom{-}\cdot & \phantom{-}\cdot & \phantom{-}\cdot & \phantom{-}\cdot & -1 & \phantom{-}\cdot & \phantom{-}\cdot & \phantom{-}\cdot & \phantom{-}\cdot & \phantom{-}\cdot & \phantom{-}\cdot & \phantom{-}\cdot \\[-1.8pt]
\phantom{-}\cdot & \phantom{-}\cdot & \phantom{-}\cdot & \phantom{-}\cdot & \phantom{-}\cdot & \phantom{-}\cdot & \phantom{-}\cdot & \phantom{-}\cdot & \phantom{-}\cdot & \phantom{-}\cdot & \phantom{-}\cdot & -1 & \phantom{-}\cdot & \phantom{-}\cdot & \phantom{-}\cdot & \phantom{-}\cdot \\[-1.8pt]
\phantom{-}\cdot & \phantom{-}\cdot & \phantom{-}\cdot & \phantom{-}\cdot & \phantom{-}\cdot & \phantom{-}\cdot & \phantom{-}\cdot & \phantom{-}\cdot & \phantom{-}\cdot & \phantom{-}\cdot & \phantom{-}1 & \phantom{-}\cdot & \phantom{-}\cdot & \phantom{-}\cdot & \phantom{-}\cdot & \phantom{-}\cdot \\[-1.8pt]
\phantom{-}\cdot & \phantom{-}\cdot & \phantom{-}\cdot & \phantom{-}\cdot & \phantom{-}\cdot & \phantom{-}\cdot & \phantom{-}\cdot & \phantom{-}\cdot & \phantom{-}\cdot & \phantom{-}\cdot & \phantom{-}\cdot & \phantom{-}\cdot & \phantom{-}\cdot & -1 & \phantom{-}\cdot & \phantom{-}\cdot \\[-1.8pt]
\phantom{-}\cdot & \phantom{-}\cdot & \phantom{-}\cdot & \phantom{-}\cdot & \phantom{-}\cdot & \phantom{-}\cdot & \phantom{-}\cdot & \phantom{-}\cdot & \phantom{-}\cdot & \phantom{-}\cdot & \phantom{-}\cdot & \phantom{-}\cdot & \phantom{-}1 & \phantom{-}\cdot & \phantom{-}\cdot & \phantom{-}\cdot \\[-1.8pt]
\phantom{-}\cdot & \phantom{-}\cdot & \phantom{-}\cdot & \phantom{-}\cdot & \phantom{-}\cdot & \phantom{-}\cdot & \phantom{-}\cdot & \phantom{-}\cdot & \phantom{-}\cdot & \phantom{-}\cdot & \phantom{-}\cdot & \phantom{-}\cdot & \phantom{-}\cdot & \phantom{-}\cdot & \phantom{-}\cdot & \phantom{-}1 \\[-1.8pt]
\phantom{-}\cdot & \phantom{-}\cdot & \phantom{-}\cdot & \phantom{-}\cdot & \phantom{-}\cdot & \phantom{-}\cdot & \phantom{-}\cdot & \phantom{-}\cdot & \phantom{-}\cdot & \phantom{-}\cdot & \phantom{-}\cdot & \phantom{-}\cdot & \phantom{-}\cdot & \phantom{-}\cdot & -1 & \phantom{-}\cdot
\end{array}\right)
\]

\[
\gamma_{1671}^{1} = \left(\begin{array}{rrrrrrrr|rrrrrrrr}
\phantom{-}\cdot & \phantom{-}\cdot & \phantom{-}\cdot & \phantom{-}\cdot & \phantom{-}\cdot & \phantom{-}\cdot & \phantom{-}\cdot & -1 & \phantom{-}\cdot & \phantom{-}\cdot & \phantom{-}\cdot & \phantom{-}\cdot & \phantom{-}\cdot & \phantom{-}\cdot & \phantom{-}\cdot & \phantom{-}\cdot \\[-1.8pt]
\phantom{-}\cdot & \phantom{-}\cdot & \phantom{-}\cdot & \phantom{-}\cdot & \phantom{-}\cdot & \phantom{-}\cdot & -1 & \phantom{-}\cdot & \phantom{-}\cdot & \phantom{-}\cdot & \phantom{-}\cdot & \phantom{-}\cdot & \phantom{-}\cdot & \phantom{-}\cdot & \phantom{-}\cdot & \phantom{-}\cdot \\[-1.8pt]
\phantom{-}\cdot & \phantom{-}\cdot & \phantom{-}\cdot & \phantom{-}\cdot & \phantom{-}\cdot & \phantom{-}1 & \phantom{-}\cdot & \phantom{-}\cdot & \phantom{-}\cdot & \phantom{-}\cdot & \phantom{-}\cdot & \phantom{-}\cdot & \phantom{-}\cdot & \phantom{-}\cdot & \phantom{-}\cdot & \phantom{-}\cdot \\[-1.8pt]
\phantom{-}\cdot & \phantom{-}\cdot & \phantom{-}\cdot & \phantom{-}\cdot & \phantom{-}1 & \phantom{-}\cdot & \phantom{-}\cdot & \phantom{-}\cdot & \phantom{-}\cdot & \phantom{-}\cdot & \phantom{-}\cdot & \phantom{-}\cdot & \phantom{-}\cdot & \phantom{-}\cdot & \phantom{-}\cdot & \phantom{-}\cdot \\[-1.8pt]
\phantom{-}\cdot & \phantom{-}\cdot & \phantom{-}\cdot & -1 & \phantom{-}\cdot & \phantom{-}\cdot & \phantom{-}\cdot & \phantom{-}\cdot & \phantom{-}\cdot & \phantom{-}\cdot & \phantom{-}\cdot & \phantom{-}\cdot & \phantom{-}\cdot & \phantom{-}\cdot & \phantom{-}\cdot & \phantom{-}\cdot \\[-1.8pt]
\phantom{-}\cdot & \phantom{-}\cdot & -1 & \phantom{-}\cdot & \phantom{-}\cdot & \phantom{-}\cdot & \phantom{-}\cdot & \phantom{-}\cdot & \phantom{-}\cdot & \phantom{-}\cdot & \phantom{-}\cdot & \phantom{-}\cdot & \phantom{-}\cdot & \phantom{-}\cdot & \phantom{-}\cdot & \phantom{-}\cdot \\[-1.8pt]
\phantom{-}\cdot & \phantom{-}1 & \phantom{-}\cdot & \phantom{-}\cdot & \phantom{-}\cdot & \phantom{-}\cdot & \phantom{-}\cdot & \phantom{-}\cdot & \phantom{-}\cdot & \phantom{-}\cdot & \phantom{-}\cdot & \phantom{-}\cdot & \phantom{-}\cdot & \phantom{-}\cdot & \phantom{-}\cdot & \phantom{-}\cdot \\[-1.8pt]
\phantom{-}1 & \phantom{-}\cdot & \phantom{-}\cdot & \phantom{-}\cdot & \phantom{-}\cdot & \phantom{-}\cdot & \phantom{-}\cdot & \phantom{-}\cdot & \phantom{-}\cdot & \phantom{-}\cdot & \phantom{-}\cdot & \phantom{-}\cdot & \phantom{-}\cdot & \phantom{-}\cdot & \phantom{-}\cdot & \phantom{-}\cdot \\[-1.8pt]
\hline
\phantom{-}\cdot & \phantom{-}\cdot & \phantom{-}\cdot & \phantom{-}\cdot & \phantom{-}\cdot & \phantom{-}\cdot & \phantom{-}\cdot & \phantom{-}\cdot & \phantom{-}\cdot & \phantom{-}\cdot & \phantom{-}\cdot & \phantom{-}\cdot & \phantom{-}\cdot & \phantom{-}\cdot & \phantom{-}\cdot & \phantom{-}1 \\[-1.8pt]
\phantom{-}\cdot & \phantom{-}\cdot & \phantom{-}\cdot & \phantom{-}\cdot & \phantom{-}\cdot & \phantom{-}\cdot & \phantom{-}\cdot & \phantom{-}\cdot & \phantom{-}\cdot & \phantom{-}\cdot & \phantom{-}\cdot & \phantom{-}\cdot & \phantom{-}\cdot & \phantom{-}\cdot & \phantom{-}1 & \phantom{-}\cdot \\[-1.8pt]
\phantom{-}\cdot & \phantom{-}\cdot & \phantom{-}\cdot & \phantom{-}\cdot & \phantom{-}\cdot & \phantom{-}\cdot & \phantom{-}\cdot & \phantom{-}\cdot & \phantom{-}\cdot & \phantom{-}\cdot & \phantom{-}\cdot & \phantom{-}\cdot & \phantom{-}\cdot & -1 & \phantom{-}\cdot & \phantom{-}\cdot \\[-1.8pt]
\phantom{-}\cdot & \phantom{-}\cdot & \phantom{-}\cdot & \phantom{-}\cdot & \phantom{-}\cdot & \phantom{-}\cdot & \phantom{-}\cdot & \phantom{-}\cdot & \phantom{-}\cdot & \phantom{-}\cdot & \phantom{-}\cdot & \phantom{-}\cdot & -1 & \phantom{-}\cdot & \phantom{-}\cdot & \phantom{-}\cdot \\[-1.8pt]
\phantom{-}\cdot & \phantom{-}\cdot & \phantom{-}\cdot & \phantom{-}\cdot & \phantom{-}\cdot & \phantom{-}\cdot & \phantom{-}\cdot & \phantom{-}\cdot & \phantom{-}\cdot & \phantom{-}\cdot & \phantom{-}\cdot & \phantom{-}1 & \phantom{-}\cdot & \phantom{-}\cdot & \phantom{-}\cdot & \phantom{-}\cdot \\[-1.8pt]
\phantom{-}\cdot & \phantom{-}\cdot & \phantom{-}\cdot & \phantom{-}\cdot & \phantom{-}\cdot & \phantom{-}\cdot & \phantom{-}\cdot & \phantom{-}\cdot & \phantom{-}\cdot & \phantom{-}\cdot & \phantom{-}1 & \phantom{-}\cdot & \phantom{-}\cdot & \phantom{-}\cdot & \phantom{-}\cdot & \phantom{-}\cdot \\[-1.8pt]
\phantom{-}\cdot & \phantom{-}\cdot & \phantom{-}\cdot & \phantom{-}\cdot & \phantom{-}\cdot & \phantom{-}\cdot & \phantom{-}\cdot & \phantom{-}\cdot & \phantom{-}\cdot & -1 & \phantom{-}\cdot & \phantom{-}\cdot & \phantom{-}\cdot & \phantom{-}\cdot & \phantom{-}\cdot & \phantom{-}\cdot \\[-1.8pt]
\phantom{-}\cdot & \phantom{-}\cdot & \phantom{-}\cdot & \phantom{-}\cdot & \phantom{-}\cdot & \phantom{-}\cdot & \phantom{-}\cdot & \phantom{-}\cdot & -1 & \phantom{-}\cdot & \phantom{-}\cdot & \phantom{-}\cdot & \phantom{-}\cdot & \phantom{-}\cdot & \phantom{-}\cdot & \phantom{-}\cdot
\end{array}\right)
\]

\[
\gamma_{1671}^{2} = \left(\begin{array}{rrrrrrrr|rrrrrrrr}
\phantom{-}\cdot & \phantom{-}\cdot & \phantom{-}\cdot & \phantom{-}\cdot & \phantom{-}\cdot & \phantom{-}\cdot & \phantom{-}1 & \phantom{-}\cdot & \phantom{-}\cdot & \phantom{-}\cdot & \phantom{-}\cdot & \phantom{-}\cdot & \phantom{-}\cdot & \phantom{-}\cdot & \phantom{-}\cdot & \phantom{-}\cdot \\[-1.8pt]
\phantom{-}\cdot & \phantom{-}\cdot & \phantom{-}\cdot & \phantom{-}\cdot & \phantom{-}\cdot & \phantom{-}\cdot & \phantom{-}\cdot & -1 & \phantom{-}\cdot & \phantom{-}\cdot & \phantom{-}\cdot & \phantom{-}\cdot & \phantom{-}\cdot & \phantom{-}\cdot & \phantom{-}\cdot & \phantom{-}\cdot \\[-1.8pt]
\phantom{-}\cdot & \phantom{-}\cdot & \phantom{-}\cdot & \phantom{-}\cdot & -1 & \phantom{-}\cdot & \phantom{-}\cdot & \phantom{-}\cdot & \phantom{-}\cdot & \phantom{-}\cdot & \phantom{-}\cdot & \phantom{-}\cdot & \phantom{-}\cdot & \phantom{-}\cdot & \phantom{-}\cdot & \phantom{-}\cdot \\[-1.8pt]
\phantom{-}\cdot & \phantom{-}\cdot & \phantom{-}\cdot & \phantom{-}\cdot & \phantom{-}\cdot & \phantom{-}1 & \phantom{-}\cdot & \phantom{-}\cdot & \phantom{-}\cdot & \phantom{-}\cdot & \phantom{-}\cdot & \phantom{-}\cdot & \phantom{-}\cdot & \phantom{-}\cdot & \phantom{-}\cdot & \phantom{-}\cdot \\[-1.8pt]
\phantom{-}\cdot & \phantom{-}\cdot & \phantom{-}1 & \phantom{-}\cdot & \phantom{-}\cdot & \phantom{-}\cdot & \phantom{-}\cdot & \phantom{-}\cdot & \phantom{-}\cdot & \phantom{-}\cdot & \phantom{-}\cdot & \phantom{-}\cdot & \phantom{-}\cdot & \phantom{-}\cdot & \phantom{-}\cdot & \phantom{-}\cdot \\[-1.8pt]
\phantom{-}\cdot & \phantom{-}\cdot & \phantom{-}\cdot & -1 & \phantom{-}\cdot & \phantom{-}\cdot & \phantom{-}\cdot & \phantom{-}\cdot & \phantom{-}\cdot & \phantom{-}\cdot & \phantom{-}\cdot & \phantom{-}\cdot & \phantom{-}\cdot & \phantom{-}\cdot & \phantom{-}\cdot & \phantom{-}\cdot \\[-1.8pt]
-1 & \phantom{-}\cdot & \phantom{-}\cdot & \phantom{-}\cdot & \phantom{-}\cdot & \phantom{-}\cdot & \phantom{-}\cdot & \phantom{-}\cdot & \phantom{-}\cdot & \phantom{-}\cdot & \phantom{-}\cdot & \phantom{-}\cdot & \phantom{-}\cdot & \phantom{-}\cdot & \phantom{-}\cdot & \phantom{-}\cdot \\[-1.8pt]
\phantom{-}\cdot & \phantom{-}1 & \phantom{-}\cdot & \phantom{-}\cdot & \phantom{-}\cdot & \phantom{-}\cdot & \phantom{-}\cdot & \phantom{-}\cdot & \phantom{-}\cdot & \phantom{-}\cdot & \phantom{-}\cdot & \phantom{-}\cdot & \phantom{-}\cdot & \phantom{-}\cdot & \phantom{-}\cdot & \phantom{-}\cdot \\[-1.8pt]
\hline
\phantom{-}\cdot & \phantom{-}\cdot & \phantom{-}\cdot & \phantom{-}\cdot & \phantom{-}\cdot & \phantom{-}\cdot & \phantom{-}\cdot & \phantom{-}\cdot & \phantom{-}\cdot & \phantom{-}\cdot & \phantom{-}\cdot & \phantom{-}\cdot & \phantom{-}\cdot & \phantom{-}\cdot & -1 & \phantom{-}\cdot \\[-1.8pt]
\phantom{-}\cdot & \phantom{-}\cdot & \phantom{-}\cdot & \phantom{-}\cdot & \phantom{-}\cdot & \phantom{-}\cdot & \phantom{-}\cdot & \phantom{-}\cdot & \phantom{-}\cdot & \phantom{-}\cdot & \phantom{-}\cdot & \phantom{-}\cdot & \phantom{-}\cdot & \phantom{-}\cdot & \phantom{-}\cdot & \phantom{-}1 \\[-1.8pt]
\phantom{-}\cdot & \phantom{-}\cdot & \phantom{-}\cdot & \phantom{-}\cdot & \phantom{-}\cdot & \phantom{-}\cdot & \phantom{-}\cdot & \phantom{-}\cdot & \phantom{-}\cdot & \phantom{-}\cdot & \phantom{-}\cdot & \phantom{-}\cdot & \phantom{-}1 & \phantom{-}\cdot & \phantom{-}\cdot & \phantom{-}\cdot \\[-1.8pt]
\phantom{-}\cdot & \phantom{-}\cdot & \phantom{-}\cdot & \phantom{-}\cdot & \phantom{-}\cdot & \phantom{-}\cdot & \phantom{-}\cdot & \phantom{-}\cdot & \phantom{-}\cdot & \phantom{-}\cdot & \phantom{-}\cdot & \phantom{-}\cdot & \phantom{-}\cdot & -1 & \phantom{-}\cdot & \phantom{-}\cdot \\[-1.8pt]
\phantom{-}\cdot & \phantom{-}\cdot & \phantom{-}\cdot & \phantom{-}\cdot & \phantom{-}\cdot & \phantom{-}\cdot & \phantom{-}\cdot & \phantom{-}\cdot & \phantom{-}\cdot & \phantom{-}\cdot & -1 & \phantom{-}\cdot & \phantom{-}\cdot & \phantom{-}\cdot & \phantom{-}\cdot & \phantom{-}\cdot \\[-1.8pt]
\phantom{-}\cdot & \phantom{-}\cdot & \phantom{-}\cdot & \phantom{-}\cdot & \phantom{-}\cdot & \phantom{-}\cdot & \phantom{-}\cdot & \phantom{-}\cdot & \phantom{-}\cdot & \phantom{-}\cdot & \phantom{-}\cdot & \phantom{-}1 & \phantom{-}\cdot & \phantom{-}\cdot & \phantom{-}\cdot & \phantom{-}\cdot \\[-1.8pt]
\phantom{-}\cdot & \phantom{-}\cdot & \phantom{-}\cdot & \phantom{-}\cdot & \phantom{-}\cdot & \phantom{-}\cdot & \phantom{-}\cdot & \phantom{-}\cdot & \phantom{-}1 & \phantom{-}\cdot & \phantom{-}\cdot & \phantom{-}\cdot & \phantom{-}\cdot & \phantom{-}\cdot & \phantom{-}\cdot & \phantom{-}\cdot \\[-1.8pt]
\phantom{-}\cdot & \phantom{-}\cdot & \phantom{-}\cdot & \phantom{-}\cdot & \phantom{-}\cdot & \phantom{-}\cdot & \phantom{-}\cdot & \phantom{-}\cdot & \phantom{-}\cdot & -1 & \phantom{-}\cdot & \phantom{-}\cdot & \phantom{-}\cdot & \phantom{-}\cdot & \phantom{-}\cdot & \phantom{-}\cdot
\end{array}\right)
\]

\[
\gamma_{1671}^{3} = \left(\begin{array}{rrrrrrrr|rrrrrrrr}
\phantom{-}\cdot & \phantom{-}\cdot & \phantom{-}\cdot & \phantom{-}\cdot & \phantom{-}\cdot & -1 & \phantom{-}\cdot & \phantom{-}\cdot & \phantom{-}\cdot & \phantom{-}\cdot & \phantom{-}\cdot & \phantom{-}\cdot & \phantom{-}\cdot & \phantom{-}\cdot & \phantom{-}\cdot & \phantom{-}\cdot \\[-1.8pt]
\phantom{-}\cdot & \phantom{-}\cdot & \phantom{-}\cdot & \phantom{-}\cdot & \phantom{-}1 & \phantom{-}\cdot & \phantom{-}\cdot & \phantom{-}\cdot & \phantom{-}\cdot & \phantom{-}\cdot & \phantom{-}\cdot & \phantom{-}\cdot & \phantom{-}\cdot & \phantom{-}\cdot & \phantom{-}\cdot & \phantom{-}\cdot \\[-1.8pt]
\phantom{-}\cdot & \phantom{-}\cdot & \phantom{-}\cdot & \phantom{-}\cdot & \phantom{-}\cdot & \phantom{-}\cdot & \phantom{-}\cdot & -1 & \phantom{-}\cdot & \phantom{-}\cdot & \phantom{-}\cdot & \phantom{-}\cdot & \phantom{-}\cdot & \phantom{-}\cdot & \phantom{-}\cdot & \phantom{-}\cdot \\[-1.8pt]
\phantom{-}\cdot & \phantom{-}\cdot & \phantom{-}\cdot & \phantom{-}\cdot & \phantom{-}\cdot & \phantom{-}\cdot & \phantom{-}1 & \phantom{-}\cdot & \phantom{-}\cdot & \phantom{-}\cdot & \phantom{-}\cdot & \phantom{-}\cdot & \phantom{-}\cdot & \phantom{-}\cdot & \phantom{-}\cdot & \phantom{-}\cdot \\[-1.8pt]
\phantom{-}\cdot & -1 & \phantom{-}\cdot & \phantom{-}\cdot & \phantom{-}\cdot & \phantom{-}\cdot & \phantom{-}\cdot & \phantom{-}\cdot & \phantom{-}\cdot & \phantom{-}\cdot & \phantom{-}\cdot & \phantom{-}\cdot & \phantom{-}\cdot & \phantom{-}\cdot & \phantom{-}\cdot & \phantom{-}\cdot \\[-1.8pt]
\phantom{-}1 & \phantom{-}\cdot & \phantom{-}\cdot & \phantom{-}\cdot & \phantom{-}\cdot & \phantom{-}\cdot & \phantom{-}\cdot & \phantom{-}\cdot & \phantom{-}\cdot & \phantom{-}\cdot & \phantom{-}\cdot & \phantom{-}\cdot & \phantom{-}\cdot & \phantom{-}\cdot & \phantom{-}\cdot & \phantom{-}\cdot \\[-1.8pt]
\phantom{-}\cdot & \phantom{-}\cdot & \phantom{-}\cdot & -1 & \phantom{-}\cdot & \phantom{-}\cdot & \phantom{-}\cdot & \phantom{-}\cdot & \phantom{-}\cdot & \phantom{-}\cdot & \phantom{-}\cdot & \phantom{-}\cdot & \phantom{-}\cdot & \phantom{-}\cdot & \phantom{-}\cdot & \phantom{-}\cdot \\[-1.8pt]
\phantom{-}\cdot & \phantom{-}\cdot & \phantom{-}1 & \phantom{-}\cdot & \phantom{-}\cdot & \phantom{-}\cdot & \phantom{-}\cdot & \phantom{-}\cdot & \phantom{-}\cdot & \phantom{-}\cdot & \phantom{-}\cdot & \phantom{-}\cdot & \phantom{-}\cdot & \phantom{-}\cdot & \phantom{-}\cdot & \phantom{-}\cdot \\[-1.8pt]
\hline
\phantom{-}\cdot & \phantom{-}\cdot & \phantom{-}\cdot & \phantom{-}\cdot & \phantom{-}\cdot & \phantom{-}\cdot & \phantom{-}\cdot & \phantom{-}\cdot & \phantom{-}\cdot & \phantom{-}\cdot & \phantom{-}\cdot & \phantom{-}\cdot & \phantom{-}\cdot & \phantom{-}1 & \phantom{-}\cdot & \phantom{-}\cdot \\[-1.8pt]
\phantom{-}\cdot & \phantom{-}\cdot & \phantom{-}\cdot & \phantom{-}\cdot & \phantom{-}\cdot & \phantom{-}\cdot & \phantom{-}\cdot & \phantom{-}\cdot & \phantom{-}\cdot & \phantom{-}\cdot & \phantom{-}\cdot & \phantom{-}\cdot & -1 & \phantom{-}\cdot & \phantom{-}\cdot & \phantom{-}\cdot \\[-1.8pt]
\phantom{-}\cdot & \phantom{-}\cdot & \phantom{-}\cdot & \phantom{-}\cdot & \phantom{-}\cdot & \phantom{-}\cdot & \phantom{-}\cdot & \phantom{-}\cdot & \phantom{-}\cdot & \phantom{-}\cdot & \phantom{-}\cdot & \phantom{-}\cdot & \phantom{-}\cdot & \phantom{-}\cdot & \phantom{-}\cdot & \phantom{-}1 \\[-1.8pt]
\phantom{-}\cdot & \phantom{-}\cdot & \phantom{-}\cdot & \phantom{-}\cdot & \phantom{-}\cdot & \phantom{-}\cdot & \phantom{-}\cdot & \phantom{-}\cdot & \phantom{-}\cdot & \phantom{-}\cdot & \phantom{-}\cdot & \phantom{-}\cdot & \phantom{-}\cdot & \phantom{-}\cdot & -1 & \phantom{-}\cdot \\[-1.8pt]
\phantom{-}\cdot & \phantom{-}\cdot & \phantom{-}\cdot & \phantom{-}\cdot & \phantom{-}\cdot & \phantom{-}\cdot & \phantom{-}\cdot & \phantom{-}\cdot & \phantom{-}\cdot & \phantom{-}1 & \phantom{-}\cdot & \phantom{-}\cdot & \phantom{-}\cdot & \phantom{-}\cdot & \phantom{-}\cdot & \phantom{-}\cdot \\[-1.8pt]
\phantom{-}\cdot & \phantom{-}\cdot & \phantom{-}\cdot & \phantom{-}\cdot & \phantom{-}\cdot & \phantom{-}\cdot & \phantom{-}\cdot & \phantom{-}\cdot & -1 & \phantom{-}\cdot & \phantom{-}\cdot & \phantom{-}\cdot & \phantom{-}\cdot & \phantom{-}\cdot & \phantom{-}\cdot & \phantom{-}\cdot \\[-1.8pt]
\phantom{-}\cdot & \phantom{-}\cdot & \phantom{-}\cdot & \phantom{-}\cdot & \phantom{-}\cdot & \phantom{-}\cdot & \phantom{-}\cdot & \phantom{-}\cdot & \phantom{-}\cdot & \phantom{-}\cdot & \phantom{-}\cdot & \phantom{-}1 & \phantom{-}\cdot & \phantom{-}\cdot & \phantom{-}\cdot & \phantom{-}\cdot \\[-1.8pt]
\phantom{-}\cdot & \phantom{-}\cdot & \phantom{-}\cdot & \phantom{-}\cdot & \phantom{-}\cdot & \phantom{-}\cdot & \phantom{-}\cdot & \phantom{-}\cdot & \phantom{-}\cdot & \phantom{-}\cdot & -1 & \phantom{-}\cdot & \phantom{-}\cdot & \phantom{-}\cdot & \phantom{-}\cdot & \phantom{-}\cdot
\end{array}\right)
\]

\[
\gamma_{1671}^{4} = \left(\begin{array}{rrrrrrrr|rrrrrrrr}
\phantom{-}\cdot & \phantom{-}\cdot & \phantom{-}\cdot & \phantom{-}\cdot & \phantom{-}\cdot & \phantom{-}\cdot & \phantom{-}\cdot & \phantom{-}\cdot & \phantom{-}1 & \phantom{-}\cdot & \phantom{-}\cdot & \phantom{-}\cdot & \phantom{-}\cdot & \phantom{-}\cdot & \phantom{-}\cdot & \phantom{-}\cdot \\[-1.8pt]
\phantom{-}\cdot & \phantom{-}\cdot & \phantom{-}\cdot & \phantom{-}\cdot & \phantom{-}\cdot & \phantom{-}\cdot & \phantom{-}\cdot & \phantom{-}\cdot & \phantom{-}\cdot & \phantom{-}1 & \phantom{-}\cdot & \phantom{-}\cdot & \phantom{-}\cdot & \phantom{-}\cdot & \phantom{-}\cdot & \phantom{-}\cdot \\[-1.8pt]
\phantom{-}\cdot & \phantom{-}\cdot & \phantom{-}\cdot & \phantom{-}\cdot & \phantom{-}\cdot & \phantom{-}\cdot & \phantom{-}\cdot & \phantom{-}\cdot & \phantom{-}\cdot & \phantom{-}\cdot & \phantom{-}1 & \phantom{-}\cdot & \phantom{-}\cdot & \phantom{-}\cdot & \phantom{-}\cdot & \phantom{-}\cdot \\[-1.8pt]
\phantom{-}\cdot & \phantom{-}\cdot & \phantom{-}\cdot & \phantom{-}\cdot & \phantom{-}\cdot & \phantom{-}\cdot & \phantom{-}\cdot & \phantom{-}\cdot & \phantom{-}\cdot & \phantom{-}\cdot & \phantom{-}\cdot & \phantom{-}1 & \phantom{-}\cdot & \phantom{-}\cdot & \phantom{-}\cdot & \phantom{-}\cdot \\[-1.8pt]
\phantom{-}\cdot & \phantom{-}\cdot & \phantom{-}\cdot & \phantom{-}\cdot & \phantom{-}\cdot & \phantom{-}\cdot & \phantom{-}\cdot & \phantom{-}\cdot & \phantom{-}\cdot & \phantom{-}\cdot & \phantom{-}\cdot & \phantom{-}\cdot & \phantom{-}1 & \phantom{-}\cdot & \phantom{-}\cdot & \phantom{-}\cdot \\[-1.8pt]
\phantom{-}\cdot & \phantom{-}\cdot & \phantom{-}\cdot & \phantom{-}\cdot & \phantom{-}\cdot & \phantom{-}\cdot & \phantom{-}\cdot & \phantom{-}\cdot & \phantom{-}\cdot & \phantom{-}\cdot & \phantom{-}\cdot & \phantom{-}\cdot & \phantom{-}\cdot & \phantom{-}1 & \phantom{-}\cdot & \phantom{-}\cdot \\[-1.8pt]
\phantom{-}\cdot & \phantom{-}\cdot & \phantom{-}\cdot & \phantom{-}\cdot & \phantom{-}\cdot & \phantom{-}\cdot & \phantom{-}\cdot & \phantom{-}\cdot & \phantom{-}\cdot & \phantom{-}\cdot & \phantom{-}\cdot & \phantom{-}\cdot & \phantom{-}\cdot & \phantom{-}\cdot & \phantom{-}1 & \phantom{-}\cdot \\[-1.8pt]
\phantom{-}\cdot & \phantom{-}\cdot & \phantom{-}\cdot & \phantom{-}\cdot & \phantom{-}\cdot & \phantom{-}\cdot & \phantom{-}\cdot & \phantom{-}\cdot & \phantom{-}\cdot & \phantom{-}\cdot & \phantom{-}\cdot & \phantom{-}\cdot & \phantom{-}\cdot & \phantom{-}\cdot & \phantom{-}\cdot & \phantom{-}1 \\[-1.8pt]
\hline
\phantom{-}1 & \phantom{-}\cdot & \phantom{-}\cdot & \phantom{-}\cdot & \phantom{-}\cdot & \phantom{-}\cdot & \phantom{-}\cdot & \phantom{-}\cdot & \phantom{-}\cdot & \phantom{-}\cdot & \phantom{-}\cdot & \phantom{-}\cdot & \phantom{-}\cdot & \phantom{-}\cdot & \phantom{-}\cdot & \phantom{-}\cdot \\[-1.8pt]
\phantom{-}\cdot & \phantom{-}1 & \phantom{-}\cdot & \phantom{-}\cdot & \phantom{-}\cdot & \phantom{-}\cdot & \phantom{-}\cdot & \phantom{-}\cdot & \phantom{-}\cdot & \phantom{-}\cdot & \phantom{-}\cdot & \phantom{-}\cdot & \phantom{-}\cdot & \phantom{-}\cdot & \phantom{-}\cdot & \phantom{-}\cdot \\[-1.8pt]
\phantom{-}\cdot & \phantom{-}\cdot & \phantom{-}1 & \phantom{-}\cdot & \phantom{-}\cdot & \phantom{-}\cdot & \phantom{-}\cdot & \phantom{-}\cdot & \phantom{-}\cdot & \phantom{-}\cdot & \phantom{-}\cdot & \phantom{-}\cdot & \phantom{-}\cdot & \phantom{-}\cdot & \phantom{-}\cdot & \phantom{-}\cdot \\[-1.8pt]
\phantom{-}\cdot & \phantom{-}\cdot & \phantom{-}\cdot & \phantom{-}1 & \phantom{-}\cdot & \phantom{-}\cdot & \phantom{-}\cdot & \phantom{-}\cdot & \phantom{-}\cdot & \phantom{-}\cdot & \phantom{-}\cdot & \phantom{-}\cdot & \phantom{-}\cdot & \phantom{-}\cdot & \phantom{-}\cdot & \phantom{-}\cdot \\[-1.8pt]
\phantom{-}\cdot & \phantom{-}\cdot & \phantom{-}\cdot & \phantom{-}\cdot & \phantom{-}1 & \phantom{-}\cdot & \phantom{-}\cdot & \phantom{-}\cdot & \phantom{-}\cdot & \phantom{-}\cdot & \phantom{-}\cdot & \phantom{-}\cdot & \phantom{-}\cdot & \phantom{-}\cdot & \phantom{-}\cdot & \phantom{-}\cdot \\[-1.8pt]
\phantom{-}\cdot & \phantom{-}\cdot & \phantom{-}\cdot & \phantom{-}\cdot & \phantom{-}\cdot & \phantom{-}1 & \phantom{-}\cdot & \phantom{-}\cdot & \phantom{-}\cdot & \phantom{-}\cdot & \phantom{-}\cdot & \phantom{-}\cdot & \phantom{-}\cdot & \phantom{-}\cdot & \phantom{-}\cdot & \phantom{-}\cdot \\[-1.8pt]
\phantom{-}\cdot & \phantom{-}\cdot & \phantom{-}\cdot & \phantom{-}\cdot & \phantom{-}\cdot & \phantom{-}\cdot & \phantom{-}1 & \phantom{-}\cdot & \phantom{-}\cdot & \phantom{-}\cdot & \phantom{-}\cdot & \phantom{-}\cdot & \phantom{-}\cdot & \phantom{-}\cdot & \phantom{-}\cdot & \phantom{-}\cdot \\[-1.8pt]
\phantom{-}\cdot & \phantom{-}\cdot & \phantom{-}\cdot & \phantom{-}\cdot & \phantom{-}\cdot & \phantom{-}\cdot & \phantom{-}\cdot & \phantom{-}1 & \phantom{-}\cdot & \phantom{-}\cdot & \phantom{-}\cdot & \phantom{-}\cdot & \phantom{-}\cdot & \phantom{-}\cdot & \phantom{-}\cdot & \phantom{-}\cdot
\end{array}\right)
\]

\[
\gamma_{1671}^{5} = \left(\begin{array}{rrrrrrrr|rrrrrrrr}
\phantom{-}\cdot & \phantom{-}\cdot & \phantom{-}1 & \phantom{-}\cdot & \phantom{-}\cdot & \phantom{-}\cdot & \phantom{-}\cdot & \phantom{-}\cdot & \phantom{-}\cdot & \phantom{-}\cdot & \phantom{-}\cdot & \phantom{-}\cdot & \phantom{-}\cdot & \phantom{-}\cdot & \phantom{-}\cdot & \phantom{-}\cdot \\[-1.8pt]
\phantom{-}\cdot & \phantom{-}\cdot & \phantom{-}\cdot & \phantom{-}1 & \phantom{-}\cdot & \phantom{-}\cdot & \phantom{-}\cdot & \phantom{-}\cdot & \phantom{-}\cdot & \phantom{-}\cdot & \phantom{-}\cdot & \phantom{-}\cdot & \phantom{-}\cdot & \phantom{-}\cdot & \phantom{-}\cdot & \phantom{-}\cdot \\[-1.8pt]
-1 & \phantom{-}\cdot & \phantom{-}\cdot & \phantom{-}\cdot & \phantom{-}\cdot & \phantom{-}\cdot & \phantom{-}\cdot & \phantom{-}\cdot & \phantom{-}\cdot & \phantom{-}\cdot & \phantom{-}\cdot & \phantom{-}\cdot & \phantom{-}\cdot & \phantom{-}\cdot & \phantom{-}\cdot & \phantom{-}\cdot \\[-1.8pt]
\phantom{-}\cdot & -1 & \phantom{-}\cdot & \phantom{-}\cdot & \phantom{-}\cdot & \phantom{-}\cdot & \phantom{-}\cdot & \phantom{-}\cdot & \phantom{-}\cdot & \phantom{-}\cdot & \phantom{-}\cdot & \phantom{-}\cdot & \phantom{-}\cdot & \phantom{-}\cdot & \phantom{-}\cdot & \phantom{-}\cdot \\[-1.8pt]
\phantom{-}\cdot & \phantom{-}\cdot & \phantom{-}\cdot & \phantom{-}\cdot & \phantom{-}\cdot & \phantom{-}\cdot & -1 & \phantom{-}\cdot & \phantom{-}\cdot & \phantom{-}\cdot & \phantom{-}\cdot & \phantom{-}\cdot & \phantom{-}\cdot & \phantom{-}\cdot & \phantom{-}\cdot & \phantom{-}\cdot \\[-1.8pt]
\phantom{-}\cdot & \phantom{-}\cdot & \phantom{-}\cdot & \phantom{-}\cdot & \phantom{-}\cdot & \phantom{-}\cdot & \phantom{-}\cdot & -1 & \phantom{-}\cdot & \phantom{-}\cdot & \phantom{-}\cdot & \phantom{-}\cdot & \phantom{-}\cdot & \phantom{-}\cdot & \phantom{-}\cdot & \phantom{-}\cdot \\[-1.8pt]
\phantom{-}\cdot & \phantom{-}\cdot & \phantom{-}\cdot & \phantom{-}\cdot & \phantom{-}1 & \phantom{-}\cdot & \phantom{-}\cdot & \phantom{-}\cdot & \phantom{-}\cdot & \phantom{-}\cdot & \phantom{-}\cdot & \phantom{-}\cdot & \phantom{-}\cdot & \phantom{-}\cdot & \phantom{-}\cdot & \phantom{-}\cdot \\[-1.8pt]
\phantom{-}\cdot & \phantom{-}\cdot & \phantom{-}\cdot & \phantom{-}\cdot & \phantom{-}\cdot & \phantom{-}1 & \phantom{-}\cdot & \phantom{-}\cdot & \phantom{-}\cdot & \phantom{-}\cdot & \phantom{-}\cdot & \phantom{-}\cdot & \phantom{-}\cdot & \phantom{-}\cdot & \phantom{-}\cdot & \phantom{-}\cdot \\[-1.8pt]
\hline
\phantom{-}\cdot & \phantom{-}\cdot & \phantom{-}\cdot & \phantom{-}\cdot & \phantom{-}\cdot & \phantom{-}\cdot & \phantom{-}\cdot & \phantom{-}\cdot & \phantom{-}\cdot & \phantom{-}\cdot & -1 & \phantom{-}\cdot & \phantom{-}\cdot & \phantom{-}\cdot & \phantom{-}\cdot & \phantom{-}\cdot \\[-1.8pt]
\phantom{-}\cdot & \phantom{-}\cdot & \phantom{-}\cdot & \phantom{-}\cdot & \phantom{-}\cdot & \phantom{-}\cdot & \phantom{-}\cdot & \phantom{-}\cdot & \phantom{-}\cdot & \phantom{-}\cdot & \phantom{-}\cdot & -1 & \phantom{-}\cdot & \phantom{-}\cdot & \phantom{-}\cdot & \phantom{-}\cdot \\[-1.8pt]
\phantom{-}\cdot & \phantom{-}\cdot & \phantom{-}\cdot & \phantom{-}\cdot & \phantom{-}\cdot & \phantom{-}\cdot & \phantom{-}\cdot & \phantom{-}\cdot & \phantom{-}1 & \phantom{-}\cdot & \phantom{-}\cdot & \phantom{-}\cdot & \phantom{-}\cdot & \phantom{-}\cdot & \phantom{-}\cdot & \phantom{-}\cdot \\[-1.8pt]
\phantom{-}\cdot & \phantom{-}\cdot & \phantom{-}\cdot & \phantom{-}\cdot & \phantom{-}\cdot & \phantom{-}\cdot & \phantom{-}\cdot & \phantom{-}\cdot & \phantom{-}\cdot & \phantom{-}1 & \phantom{-}\cdot & \phantom{-}\cdot & \phantom{-}\cdot & \phantom{-}\cdot & \phantom{-}\cdot & \phantom{-}\cdot \\[-1.8pt]
\phantom{-}\cdot & \phantom{-}\cdot & \phantom{-}\cdot & \phantom{-}\cdot & \phantom{-}\cdot & \phantom{-}\cdot & \phantom{-}\cdot & \phantom{-}\cdot & \phantom{-}\cdot & \phantom{-}\cdot & \phantom{-}\cdot & \phantom{-}\cdot & \phantom{-}\cdot & \phantom{-}\cdot & \phantom{-}1 & \phantom{-}\cdot \\[-1.8pt]
\phantom{-}\cdot & \phantom{-}\cdot & \phantom{-}\cdot & \phantom{-}\cdot & \phantom{-}\cdot & \phantom{-}\cdot & \phantom{-}\cdot & \phantom{-}\cdot & \phantom{-}\cdot & \phantom{-}\cdot & \phantom{-}\cdot & \phantom{-}\cdot & \phantom{-}\cdot & \phantom{-}\cdot & \phantom{-}\cdot & \phantom{-}1 \\[-1.8pt]
\phantom{-}\cdot & \phantom{-}\cdot & \phantom{-}\cdot & \phantom{-}\cdot & \phantom{-}\cdot & \phantom{-}\cdot & \phantom{-}\cdot & \phantom{-}\cdot & \phantom{-}\cdot & \phantom{-}\cdot & \phantom{-}\cdot & \phantom{-}\cdot & -1 & \phantom{-}\cdot & \phantom{-}\cdot & \phantom{-}\cdot \\[-1.8pt]
\phantom{-}\cdot & \phantom{-}\cdot & \phantom{-}\cdot & \phantom{-}\cdot & \phantom{-}\cdot & \phantom{-}\cdot & \phantom{-}\cdot & \phantom{-}\cdot & \phantom{-}\cdot & \phantom{-}\cdot & \phantom{-}\cdot & \phantom{-}\cdot & \phantom{-}\cdot & -1 & \phantom{-}\cdot & \phantom{-}\cdot
\end{array}\right)
\]

\[
\gamma_{1671}^{6} = \left(\begin{array}{rrrrrrrr|rrrrrrrr}
\phantom{-}\cdot & \phantom{-}\cdot & \phantom{-}\cdot & \phantom{-}1 & \phantom{-}\cdot & \phantom{-}\cdot & \phantom{-}\cdot & \phantom{-}\cdot & \phantom{-}\cdot & \phantom{-}\cdot & \phantom{-}\cdot & \phantom{-}\cdot & \phantom{-}\cdot & \phantom{-}\cdot & \phantom{-}\cdot & \phantom{-}\cdot \\[-1.8pt]
\phantom{-}\cdot & \phantom{-}\cdot & -1 & \phantom{-}\cdot & \phantom{-}\cdot & \phantom{-}\cdot & \phantom{-}\cdot & \phantom{-}\cdot & \phantom{-}\cdot & \phantom{-}\cdot & \phantom{-}\cdot & \phantom{-}\cdot & \phantom{-}\cdot & \phantom{-}\cdot & \phantom{-}\cdot & \phantom{-}\cdot \\[-1.8pt]
\phantom{-}\cdot & \phantom{-}1 & \phantom{-}\cdot & \phantom{-}\cdot & \phantom{-}\cdot & \phantom{-}\cdot & \phantom{-}\cdot & \phantom{-}\cdot & \phantom{-}\cdot & \phantom{-}\cdot & \phantom{-}\cdot & \phantom{-}\cdot & \phantom{-}\cdot & \phantom{-}\cdot & \phantom{-}\cdot & \phantom{-}\cdot \\[-1.8pt]
-1 & \phantom{-}\cdot & \phantom{-}\cdot & \phantom{-}\cdot & \phantom{-}\cdot & \phantom{-}\cdot & \phantom{-}\cdot & \phantom{-}\cdot & \phantom{-}\cdot & \phantom{-}\cdot & \phantom{-}\cdot & \phantom{-}\cdot & \phantom{-}\cdot & \phantom{-}\cdot & \phantom{-}\cdot & \phantom{-}\cdot \\[-1.8pt]
\phantom{-}\cdot & \phantom{-}\cdot & \phantom{-}\cdot & \phantom{-}\cdot & \phantom{-}\cdot & \phantom{-}\cdot & \phantom{-}\cdot & -1 & \phantom{-}\cdot & \phantom{-}\cdot & \phantom{-}\cdot & \phantom{-}\cdot & \phantom{-}\cdot & \phantom{-}\cdot & \phantom{-}\cdot & \phantom{-}\cdot \\[-1.8pt]
\phantom{-}\cdot & \phantom{-}\cdot & \phantom{-}\cdot & \phantom{-}\cdot & \phantom{-}\cdot & \phantom{-}\cdot & \phantom{-}1 & \phantom{-}\cdot & \phantom{-}\cdot & \phantom{-}\cdot & \phantom{-}\cdot & \phantom{-}\cdot & \phantom{-}\cdot & \phantom{-}\cdot & \phantom{-}\cdot & \phantom{-}\cdot \\[-1.8pt]
\phantom{-}\cdot & \phantom{-}\cdot & \phantom{-}\cdot & \phantom{-}\cdot & \phantom{-}\cdot & -1 & \phantom{-}\cdot & \phantom{-}\cdot & \phantom{-}\cdot & \phantom{-}\cdot & \phantom{-}\cdot & \phantom{-}\cdot & \phantom{-}\cdot & \phantom{-}\cdot & \phantom{-}\cdot & \phantom{-}\cdot \\[-1.8pt]
\phantom{-}\cdot & \phantom{-}\cdot & \phantom{-}\cdot & \phantom{-}\cdot & \phantom{-}1 & \phantom{-}\cdot & \phantom{-}\cdot & \phantom{-}\cdot & \phantom{-}\cdot & \phantom{-}\cdot & \phantom{-}\cdot & \phantom{-}\cdot & \phantom{-}\cdot & \phantom{-}\cdot & \phantom{-}\cdot & \phantom{-}\cdot \\[-1.8pt]
\hline
\phantom{-}\cdot & \phantom{-}\cdot & \phantom{-}\cdot & \phantom{-}\cdot & \phantom{-}\cdot & \phantom{-}\cdot & \phantom{-}\cdot & \phantom{-}\cdot & \phantom{-}\cdot & \phantom{-}\cdot & \phantom{-}\cdot & -1 & \phantom{-}\cdot & \phantom{-}\cdot & \phantom{-}\cdot & \phantom{-}\cdot \\[-1.8pt]
\phantom{-}\cdot & \phantom{-}\cdot & \phantom{-}\cdot & \phantom{-}\cdot & \phantom{-}\cdot & \phantom{-}\cdot & \phantom{-}\cdot & \phantom{-}\cdot & \phantom{-}\cdot & \phantom{-}\cdot & \phantom{-}1 & \phantom{-}\cdot & \phantom{-}\cdot & \phantom{-}\cdot & \phantom{-}\cdot & \phantom{-}\cdot \\[-1.8pt]
\phantom{-}\cdot & \phantom{-}\cdot & \phantom{-}\cdot & \phantom{-}\cdot & \phantom{-}\cdot & \phantom{-}\cdot & \phantom{-}\cdot & \phantom{-}\cdot & \phantom{-}\cdot & -1 & \phantom{-}\cdot & \phantom{-}\cdot & \phantom{-}\cdot & \phantom{-}\cdot & \phantom{-}\cdot & \phantom{-}\cdot \\[-1.8pt]
\phantom{-}\cdot & \phantom{-}\cdot & \phantom{-}\cdot & \phantom{-}\cdot & \phantom{-}\cdot & \phantom{-}\cdot & \phantom{-}\cdot & \phantom{-}\cdot & \phantom{-}1 & \phantom{-}\cdot & \phantom{-}\cdot & \phantom{-}\cdot & \phantom{-}\cdot & \phantom{-}\cdot & \phantom{-}\cdot & \phantom{-}\cdot \\[-1.8pt]
\phantom{-}\cdot & \phantom{-}\cdot & \phantom{-}\cdot & \phantom{-}\cdot & \phantom{-}\cdot & \phantom{-}\cdot & \phantom{-}\cdot & \phantom{-}\cdot & \phantom{-}\cdot & \phantom{-}\cdot & \phantom{-}\cdot & \phantom{-}\cdot & \phantom{-}\cdot & \phantom{-}\cdot & \phantom{-}\cdot & \phantom{-}1 \\[-1.8pt]
\phantom{-}\cdot & \phantom{-}\cdot & \phantom{-}\cdot & \phantom{-}\cdot & \phantom{-}\cdot & \phantom{-}\cdot & \phantom{-}\cdot & \phantom{-}\cdot & \phantom{-}\cdot & \phantom{-}\cdot & \phantom{-}\cdot & \phantom{-}\cdot & \phantom{-}\cdot & \phantom{-}\cdot & -1 & \phantom{-}\cdot \\[-1.8pt]
\phantom{-}\cdot & \phantom{-}\cdot & \phantom{-}\cdot & \phantom{-}\cdot & \phantom{-}\cdot & \phantom{-}\cdot & \phantom{-}\cdot & \phantom{-}\cdot & \phantom{-}\cdot & \phantom{-}\cdot & \phantom{-}\cdot & \phantom{-}\cdot & \phantom{-}\cdot & \phantom{-}1 & \phantom{-}\cdot & \phantom{-}\cdot \\[-1.8pt]
\phantom{-}\cdot & \phantom{-}\cdot & \phantom{-}\cdot & \phantom{-}\cdot & \phantom{-}\cdot & \phantom{-}\cdot & \phantom{-}\cdot & \phantom{-}\cdot & \phantom{-}\cdot & \phantom{-}\cdot & \phantom{-}\cdot & \phantom{-}\cdot & -1 & \phantom{-}\cdot & \phantom{-}\cdot & \phantom{-}\cdot
\end{array}\right)
\]

\[
\gamma_{1671}^{7} = \left(\begin{array}{rrrrrrrr|rrrrrrrr}
\phantom{-}\cdot & \phantom{-}\cdot & \phantom{-}\cdot & \phantom{-}\cdot & \phantom{-}1 & \phantom{-}\cdot & \phantom{-}\cdot & \phantom{-}\cdot & \phantom{-}\cdot & \phantom{-}\cdot & \phantom{-}\cdot & \phantom{-}\cdot & \phantom{-}\cdot & \phantom{-}\cdot & \phantom{-}\cdot & \phantom{-}\cdot \\[-1.8pt]
\phantom{-}\cdot & \phantom{-}\cdot & \phantom{-}\cdot & \phantom{-}\cdot & \phantom{-}\cdot & \phantom{-}1 & \phantom{-}\cdot & \phantom{-}\cdot & \phantom{-}\cdot & \phantom{-}\cdot & \phantom{-}\cdot & \phantom{-}\cdot & \phantom{-}\cdot & \phantom{-}\cdot & \phantom{-}\cdot & \phantom{-}\cdot \\[-1.8pt]
\phantom{-}\cdot & \phantom{-}\cdot & \phantom{-}\cdot & \phantom{-}\cdot & \phantom{-}\cdot & \phantom{-}\cdot & \phantom{-}1 & \phantom{-}\cdot & \phantom{-}\cdot & \phantom{-}\cdot & \phantom{-}\cdot & \phantom{-}\cdot & \phantom{-}\cdot & \phantom{-}\cdot & \phantom{-}\cdot & \phantom{-}\cdot \\[-1.8pt]
\phantom{-}\cdot & \phantom{-}\cdot & \phantom{-}\cdot & \phantom{-}\cdot & \phantom{-}\cdot & \phantom{-}\cdot & \phantom{-}\cdot & \phantom{-}1 & \phantom{-}\cdot & \phantom{-}\cdot & \phantom{-}\cdot & \phantom{-}\cdot & \phantom{-}\cdot & \phantom{-}\cdot & \phantom{-}\cdot & \phantom{-}\cdot \\[-1.8pt]
-1 & \phantom{-}\cdot & \phantom{-}\cdot & \phantom{-}\cdot & \phantom{-}\cdot & \phantom{-}\cdot & \phantom{-}\cdot & \phantom{-}\cdot & \phantom{-}\cdot & \phantom{-}\cdot & \phantom{-}\cdot & \phantom{-}\cdot & \phantom{-}\cdot & \phantom{-}\cdot & \phantom{-}\cdot & \phantom{-}\cdot \\[-1.8pt]
\phantom{-}\cdot & -1 & \phantom{-}\cdot & \phantom{-}\cdot & \phantom{-}\cdot & \phantom{-}\cdot & \phantom{-}\cdot & \phantom{-}\cdot & \phantom{-}\cdot & \phantom{-}\cdot & \phantom{-}\cdot & \phantom{-}\cdot & \phantom{-}\cdot & \phantom{-}\cdot & \phantom{-}\cdot & \phantom{-}\cdot \\[-1.8pt]
\phantom{-}\cdot & \phantom{-}\cdot & -1 & \phantom{-}\cdot & \phantom{-}\cdot & \phantom{-}\cdot & \phantom{-}\cdot & \phantom{-}\cdot & \phantom{-}\cdot & \phantom{-}\cdot & \phantom{-}\cdot & \phantom{-}\cdot & \phantom{-}\cdot & \phantom{-}\cdot & \phantom{-}\cdot & \phantom{-}\cdot \\[-1.8pt]
\phantom{-}\cdot & \phantom{-}\cdot & \phantom{-}\cdot & -1 & \phantom{-}\cdot & \phantom{-}\cdot & \phantom{-}\cdot & \phantom{-}\cdot & \phantom{-}\cdot & \phantom{-}\cdot & \phantom{-}\cdot & \phantom{-}\cdot & \phantom{-}\cdot & \phantom{-}\cdot & \phantom{-}\cdot & \phantom{-}\cdot \\[-1.8pt]
\hline
\phantom{-}\cdot & \phantom{-}\cdot & \phantom{-}\cdot & \phantom{-}\cdot & \phantom{-}\cdot & \phantom{-}\cdot & \phantom{-}\cdot & \phantom{-}\cdot & \phantom{-}\cdot & \phantom{-}\cdot & \phantom{-}\cdot & \phantom{-}\cdot & -1 & \phantom{-}\cdot & \phantom{-}\cdot & \phantom{-}\cdot \\[-1.8pt]
\phantom{-}\cdot & \phantom{-}\cdot & \phantom{-}\cdot & \phantom{-}\cdot & \phantom{-}\cdot & \phantom{-}\cdot & \phantom{-}\cdot & \phantom{-}\cdot & \phantom{-}\cdot & \phantom{-}\cdot & \phantom{-}\cdot & \phantom{-}\cdot & \phantom{-}\cdot & -1 & \phantom{-}\cdot & \phantom{-}\cdot \\[-1.8pt]
\phantom{-}\cdot & \phantom{-}\cdot & \phantom{-}\cdot & \phantom{-}\cdot & \phantom{-}\cdot & \phantom{-}\cdot & \phantom{-}\cdot & \phantom{-}\cdot & \phantom{-}\cdot & \phantom{-}\cdot & \phantom{-}\cdot & \phantom{-}\cdot & \phantom{-}\cdot & \phantom{-}\cdot & -1 & \phantom{-}\cdot \\[-1.8pt]
\phantom{-}\cdot & \phantom{-}\cdot & \phantom{-}\cdot & \phantom{-}\cdot & \phantom{-}\cdot & \phantom{-}\cdot & \phantom{-}\cdot & \phantom{-}\cdot & \phantom{-}\cdot & \phantom{-}\cdot & \phantom{-}\cdot & \phantom{-}\cdot & \phantom{-}\cdot & \phantom{-}\cdot & \phantom{-}\cdot & -1 \\[-1.8pt]
\phantom{-}\cdot & \phantom{-}\cdot & \phantom{-}\cdot & \phantom{-}\cdot & \phantom{-}\cdot & \phantom{-}\cdot & \phantom{-}\cdot & \phantom{-}\cdot & \phantom{-}1 & \phantom{-}\cdot & \phantom{-}\cdot & \phantom{-}\cdot & \phantom{-}\cdot & \phantom{-}\cdot & \phantom{-}\cdot & \phantom{-}\cdot \\[-1.8pt]
\phantom{-}\cdot & \phantom{-}\cdot & \phantom{-}\cdot & \phantom{-}\cdot & \phantom{-}\cdot & \phantom{-}\cdot & \phantom{-}\cdot & \phantom{-}\cdot & \phantom{-}\cdot & \phantom{-}1 & \phantom{-}\cdot & \phantom{-}\cdot & \phantom{-}\cdot & \phantom{-}\cdot & \phantom{-}\cdot & \phantom{-}\cdot \\[-1.8pt]
\phantom{-}\cdot & \phantom{-}\cdot & \phantom{-}\cdot & \phantom{-}\cdot & \phantom{-}\cdot & \phantom{-}\cdot & \phantom{-}\cdot & \phantom{-}\cdot & \phantom{-}\cdot & \phantom{-}\cdot & \phantom{-}1 & \phantom{-}\cdot & \phantom{-}\cdot & \phantom{-}\cdot & \phantom{-}\cdot & \phantom{-}\cdot \\[-1.8pt]
\phantom{-}\cdot & \phantom{-}\cdot & \phantom{-}\cdot & \phantom{-}\cdot & \phantom{-}\cdot & \phantom{-}\cdot & \phantom{-}\cdot & \phantom{-}\cdot & \phantom{-}\cdot & \phantom{-}\cdot & \phantom{-}\cdot & \phantom{-}1 & \phantom{-}\cdot & \phantom{-}\cdot & \phantom{-}\cdot & \phantom{-}\cdot
\end{array}\right)
\]

\hypertarget{proof-that-the-matrices-are-real-1}{%
\subsection{2.3 Proof that the matrices are
real}\label{proof-that-the-matrices-are-real-1}}

The matrices are products and blocks of the integer matrices of Part 1
and of \(\sigma_{88}\); no complex number enters. Check \texttt{71-R1}:
every entry is \(-1\), \(0\) or \(+1\); check \texttt{71-R2}: every
matrix is a signed permutation matrix. \(\blacksquare\)

\hypertarget{proof-of-the-anticommutation-relations-1}{%
\subsection{2.4 Proof of the anticommutation
relations}\label{proof-of-the-anticommutation-relations-1}}

\textbf{Statement.}
\[\boxed{\;\gamma_{1671}^A\gamma_{1671}^B+\gamma_{1671}^B\gamma_{1671}^A=2\,\eta_{1671}^{AB}\,1_{16},\qquad
\eta_{1671}=\mathrm{diag}(-1,-1,-1,-1,+1,-1,-1,-1)\;}\tag{2.3}\] The
time-like matrix \(\gamma^4\) squares to \(+1\) and the seven space-like
matrices square to \(-1\).

\textbf{Proof.} From the block form (2.2): \((\gamma^4)^2=1\);
\((\gamma^j)^2=\mathrm{diag}((\tau^j)^2,(\tau^j)^2)=-1\); for
\(i\ne j\),
\(\gamma^i\gamma^j+\gamma^j\gamma^i=\mathrm{diag}(\tau^i\tau^j+\tau^j\tau^i,\ \tau^i\tau^j+\tau^j\tau^i)=0\);
and
\[\gamma^4\gamma^j+\gamma^j\gamma^4=\begin{pmatrix}0&\tau^j\\-\tau^j&0\end{pmatrix}+\begin{pmatrix}0&-\tau^j\\ \tau^j&0\end{pmatrix}=0 .\]
The properties of the \(\tau^j\) used here are check \texttt{71-D6}; Eq.
(2.3) itself is check \texttt{71-C1}. \(\blacksquare\)

\textbf{Why the sign convention differs from Part 1.} In Part 1
space-like directions square to \(+1\); here they square to \(-1\). This
is forced. \texttt{7+1.nb} proves (its section 4.4) that eight
\emph{real} 16×16 matrices with seven squares \(+1\) and one square
\(-1\) do not exist, because that Clifford algebra is
\(\mathrm{Cl}(7,1)\cong\mathbb{H}(8)\), the 8×8 matrices over the
quaternions, whereas \(\mathrm{Cl}(1,7)\cong\mathbb{R}(16)\). For 4+4
the two conventions are equivalent
(\(\mathrm{Cl}(4,4)\cong\mathbb{R}(16)\) either way). The consequence
for the Dirac equation is the factor \(i\) in Eq. (2.9), which is absent
in Eq. (1.10).

\textbf{Consequences} (checks \texttt{71-C2}, \texttt{71-C3},
\texttt{71-C4}, \texttt{I71-G}): \((\gamma^A)^T=\eta^{AA}\gamma^A\), so
\(\gamma^4\) is symmetric and the seven \(\gamma^j\) are antisymmetric;
the 256 ordered products are linearly independent, so the matrices
generate all real 16×16 matrices and the representation is irreducible
over the reals.

\hypertarget{the-scaled-commutators-and-the-group-they-generate-1}{%
\subsection{2.5 The scaled commutators and the group they
generate}\label{the-scaled-commutators-and-the-group-they-generate-1}}

Define \(\Sigma^{AB}=\frac14[\gamma^A,\gamma^B]\) as in Eq. (1.4), now
with the matrices \(\gamma_{1671}^A\).

\textbf{Theorem 2.} Statements 1--3 of Theorem 1 hold for
\(\gamma^A=\gamma_{1671}^A\), \(\Sigma^{AB}=\Sigma_{1671}^{AB}\) and
\(\eta=\eta_{1671}\), with the following changes: the Lie algebra is
\(\mathfrak{so}(7,1)\); the identity component generated by the
exponentials is \(\mathrm{Spin}_0(7,1)\), which covers the proper
orthochronous Lorentz group \(SO_0(7,1)\) twice; and the group generated
by it together with the eight matrices is the pin group of the 7+1
spacetime,
\[\boxed{\;\mathrm{Pin}(7,1)=\mathrm{Spin}_0(7,1)\;\cup\;\gamma_{1671}^1\,\mathrm{Spin}_0(7,1)\;\cup\;\gamma_{1671}^4\,\mathrm{Spin}_0(7,1)\;\cup\;\gamma_{1671}^1\gamma_{1671}^4\,\mathrm{Spin}_0(7,1)\;}\]

\textbf{Proof.} Identical to that of Theorem 1, with checks
\texttt{71-L1} to \texttt{71-L6} and \texttt{71-G1} to \texttt{71-G3}.
The Killing form is
\(B_{\rm Kil}(\Sigma^{AB},\Sigma^{AB})=-12\eta^{AA}\eta^{BB}\); now 21
generators are compact (the rotations among the seven space coordinates,
which are antisymmetric matrices) and 7 are non-compact (the boosts
\(\Sigma^{4j}\), symmetric matrices) (checks \texttt{71-L5},
\texttt{71-L6}, \texttt{I71-K}). \(\blacksquare\)

\textbf{What ``Pin(7,1)'' means here.} The name is used as in the task
specification: the pin group of the spacetime with seven space-like
dimensions and one time-like dimension, realized by the real matrices
\(\gamma_{1671}^A\). Two remarks make this precise.

\begin{enumerate}
\def\labelenumi{\arabic{enumi}.}
\tightlist
\item
  \emph{Convention.} The matrices \(\gamma_{1671}^A\) generate the
  Clifford algebra in which the time-like generator squares to \(+1\)
  and the seven space-like ones to \(-1\),
  \(\mathrm{Cl}(1,7)\cong\mathbb{R}(16)\) (Sec. 2.4); in the notation
  that records these signs the group of Theorem 2 is written
  \(\mathrm{Pin}(1,7)\). The full Lorentz group \(O(7,1)\) has a second
  double cover, the pin group of the algebra
  \(\mathrm{Cl}(7,1)\cong\mathbb{H}(8)\) with the opposite signs, in
  which a reflection of a space axis squares to \(+1\) instead of
  \(-1\). Its generators would be eight real 16×16 matrices with seven
  squares \(+1\) and one square \(-1\), and such matrices do not exist
  (Sec. 2.4). The two groups have the same identity component
  \(\mathrm{Spin}_0(7,1)\) and the same Lie algebra; they differ only in
  the components that contain reflections. With purely real 16×16 Dirac
  matrices, as the task demands, the group of Theorem 2 is the only
  possibility.
\item
  \emph{It is not the group of Part 1.} The Lie algebra spanned by the
  scaled commutators of the 7+1 matrices has 21 compact and 7
  non-compact generators, while \(\mathfrak{so}(4,4)\) has 12 and 16.
  The signature of the Killing form (checks \texttt{44-L5},
  \texttt{71-L5}) is an invariant of a real Lie algebra, so the two
  algebras, and with them \(\mathrm{Pin}(7,1)\) and
  \(\mathrm{Pin}(4,4)\), are different, although both sets of matrices
  live in the same algebra of real 16×16 matrices. (The first version of
  the task specification repeated ``Pin(4,4)'' at this place; the
  author's corrected specification of November 12, 2026 reads Pin(7,1),
  which is what Theorem 2 proves.)
\end{enumerate}

\hypertarget{chirality-in-71-no-real-projector-the-complex-projectors}{%
\subsection{2.6 Chirality in 7+1: no real projector; the complex
projectors}\label{chirality-in-71-no-real-projector-the-complex-projectors}}

\textbf{Volume element.}
\(\omega\equiv\gamma^0\gamma^1\cdots\gamma^7=\begin{pmatrix}0&-1_8\\1_8&0\end{pmatrix}\),
which equals \texttt{Γ17{[}8{]}}, the \texttt{ω17} of the notebook
(check \texttt{71-V1}):

\[
\omega_{1671} = \left(\begin{array}{rrrrrrrr|rrrrrrrr}
\phantom{-}\cdot & \phantom{-}\cdot & \phantom{-}\cdot & \phantom{-}\cdot & \phantom{-}\cdot & \phantom{-}\cdot & \phantom{-}\cdot & \phantom{-}\cdot & -1 & \phantom{-}\cdot & \phantom{-}\cdot & \phantom{-}\cdot & \phantom{-}\cdot & \phantom{-}\cdot & \phantom{-}\cdot & \phantom{-}\cdot \\[-1.8pt]
\phantom{-}\cdot & \phantom{-}\cdot & \phantom{-}\cdot & \phantom{-}\cdot & \phantom{-}\cdot & \phantom{-}\cdot & \phantom{-}\cdot & \phantom{-}\cdot & \phantom{-}\cdot & -1 & \phantom{-}\cdot & \phantom{-}\cdot & \phantom{-}\cdot & \phantom{-}\cdot & \phantom{-}\cdot & \phantom{-}\cdot \\[-1.8pt]
\phantom{-}\cdot & \phantom{-}\cdot & \phantom{-}\cdot & \phantom{-}\cdot & \phantom{-}\cdot & \phantom{-}\cdot & \phantom{-}\cdot & \phantom{-}\cdot & \phantom{-}\cdot & \phantom{-}\cdot & -1 & \phantom{-}\cdot & \phantom{-}\cdot & \phantom{-}\cdot & \phantom{-}\cdot & \phantom{-}\cdot \\[-1.8pt]
\phantom{-}\cdot & \phantom{-}\cdot & \phantom{-}\cdot & \phantom{-}\cdot & \phantom{-}\cdot & \phantom{-}\cdot & \phantom{-}\cdot & \phantom{-}\cdot & \phantom{-}\cdot & \phantom{-}\cdot & \phantom{-}\cdot & -1 & \phantom{-}\cdot & \phantom{-}\cdot & \phantom{-}\cdot & \phantom{-}\cdot \\[-1.8pt]
\phantom{-}\cdot & \phantom{-}\cdot & \phantom{-}\cdot & \phantom{-}\cdot & \phantom{-}\cdot & \phantom{-}\cdot & \phantom{-}\cdot & \phantom{-}\cdot & \phantom{-}\cdot & \phantom{-}\cdot & \phantom{-}\cdot & \phantom{-}\cdot & -1 & \phantom{-}\cdot & \phantom{-}\cdot & \phantom{-}\cdot \\[-1.8pt]
\phantom{-}\cdot & \phantom{-}\cdot & \phantom{-}\cdot & \phantom{-}\cdot & \phantom{-}\cdot & \phantom{-}\cdot & \phantom{-}\cdot & \phantom{-}\cdot & \phantom{-}\cdot & \phantom{-}\cdot & \phantom{-}\cdot & \phantom{-}\cdot & \phantom{-}\cdot & -1 & \phantom{-}\cdot & \phantom{-}\cdot \\[-1.8pt]
\phantom{-}\cdot & \phantom{-}\cdot & \phantom{-}\cdot & \phantom{-}\cdot & \phantom{-}\cdot & \phantom{-}\cdot & \phantom{-}\cdot & \phantom{-}\cdot & \phantom{-}\cdot & \phantom{-}\cdot & \phantom{-}\cdot & \phantom{-}\cdot & \phantom{-}\cdot & \phantom{-}\cdot & -1 & \phantom{-}\cdot \\[-1.8pt]
\phantom{-}\cdot & \phantom{-}\cdot & \phantom{-}\cdot & \phantom{-}\cdot & \phantom{-}\cdot & \phantom{-}\cdot & \phantom{-}\cdot & \phantom{-}\cdot & \phantom{-}\cdot & \phantom{-}\cdot & \phantom{-}\cdot & \phantom{-}\cdot & \phantom{-}\cdot & \phantom{-}\cdot & \phantom{-}\cdot & -1 \\[-1.8pt]
\hline
\phantom{-}1 & \phantom{-}\cdot & \phantom{-}\cdot & \phantom{-}\cdot & \phantom{-}\cdot & \phantom{-}\cdot & \phantom{-}\cdot & \phantom{-}\cdot & \phantom{-}\cdot & \phantom{-}\cdot & \phantom{-}\cdot & \phantom{-}\cdot & \phantom{-}\cdot & \phantom{-}\cdot & \phantom{-}\cdot & \phantom{-}\cdot \\[-1.8pt]
\phantom{-}\cdot & \phantom{-}1 & \phantom{-}\cdot & \phantom{-}\cdot & \phantom{-}\cdot & \phantom{-}\cdot & \phantom{-}\cdot & \phantom{-}\cdot & \phantom{-}\cdot & \phantom{-}\cdot & \phantom{-}\cdot & \phantom{-}\cdot & \phantom{-}\cdot & \phantom{-}\cdot & \phantom{-}\cdot & \phantom{-}\cdot \\[-1.8pt]
\phantom{-}\cdot & \phantom{-}\cdot & \phantom{-}1 & \phantom{-}\cdot & \phantom{-}\cdot & \phantom{-}\cdot & \phantom{-}\cdot & \phantom{-}\cdot & \phantom{-}\cdot & \phantom{-}\cdot & \phantom{-}\cdot & \phantom{-}\cdot & \phantom{-}\cdot & \phantom{-}\cdot & \phantom{-}\cdot & \phantom{-}\cdot \\[-1.8pt]
\phantom{-}\cdot & \phantom{-}\cdot & \phantom{-}\cdot & \phantom{-}1 & \phantom{-}\cdot & \phantom{-}\cdot & \phantom{-}\cdot & \phantom{-}\cdot & \phantom{-}\cdot & \phantom{-}\cdot & \phantom{-}\cdot & \phantom{-}\cdot & \phantom{-}\cdot & \phantom{-}\cdot & \phantom{-}\cdot & \phantom{-}\cdot \\[-1.8pt]
\phantom{-}\cdot & \phantom{-}\cdot & \phantom{-}\cdot & \phantom{-}\cdot & \phantom{-}1 & \phantom{-}\cdot & \phantom{-}\cdot & \phantom{-}\cdot & \phantom{-}\cdot & \phantom{-}\cdot & \phantom{-}\cdot & \phantom{-}\cdot & \phantom{-}\cdot & \phantom{-}\cdot & \phantom{-}\cdot & \phantom{-}\cdot \\[-1.8pt]
\phantom{-}\cdot & \phantom{-}\cdot & \phantom{-}\cdot & \phantom{-}\cdot & \phantom{-}\cdot & \phantom{-}1 & \phantom{-}\cdot & \phantom{-}\cdot & \phantom{-}\cdot & \phantom{-}\cdot & \phantom{-}\cdot & \phantom{-}\cdot & \phantom{-}\cdot & \phantom{-}\cdot & \phantom{-}\cdot & \phantom{-}\cdot \\[-1.8pt]
\phantom{-}\cdot & \phantom{-}\cdot & \phantom{-}\cdot & \phantom{-}\cdot & \phantom{-}\cdot & \phantom{-}\cdot & \phantom{-}1 & \phantom{-}\cdot & \phantom{-}\cdot & \phantom{-}\cdot & \phantom{-}\cdot & \phantom{-}\cdot & \phantom{-}\cdot & \phantom{-}\cdot & \phantom{-}\cdot & \phantom{-}\cdot \\[-1.8pt]
\phantom{-}\cdot & \phantom{-}\cdot & \phantom{-}\cdot & \phantom{-}\cdot & \phantom{-}\cdot & \phantom{-}\cdot & \phantom{-}\cdot & \phantom{-}1 & \phantom{-}\cdot & \phantom{-}\cdot & \phantom{-}\cdot & \phantom{-}\cdot & \phantom{-}\cdot & \phantom{-}\cdot & \phantom{-}\cdot & \phantom{-}\cdot
\end{array}\right)
\]

As in Sec. 1.6, \(\omega\) anticommutes with every \(\gamma^A\) (check
\texttt{71-V2}), but now \(\omega^2=(-1)^{28}\,(+1)(-1)^7=-1\).

\textbf{Theorem 3 (no real chirality in 7+1).} A real 16×16 matrix that
anticommutes with all eight \(\gamma_{1671}^A\) is a real multiple of
\(\omega_{1671}\). Since \(\omega_{1671}^2=-1\), no real matrix with
square \(+1\) anticommutes with all of them, and no real chiral
projector exists.

\textbf{Proof.} The condition \(X\gamma^A+\gamma^AX=0\) for
\(A=0,\dots,7\) is a homogeneous linear system for the 256 entries of
\(X\). Its rank is 255, and \(X=\omega\) is a solution (check
\texttt{71-V3}); hence the solutions are the multiples \(c\,\omega\),
whose squares are \(-c^2\le0\). \emph{A second proof, without a linear
system} (check \texttt{I71-A}): expand \(X\) in the 256 products
\(\gamma^{A_1}\cdots\gamma^{A_r}\). A product of rank \(r\) anticommutes
with \(\gamma^A\) if \(A\) is among its indices and \(r\) is even, or if
\(A\) is not among them and \(r\) is odd; it commutes otherwise. A
product anticommutes with all eight matrices only if every index occurs
and \(r\) is even, i.e.~only for \(r=8\). Finally, the map
\(X\mapsto X\gamma^A+\gamma^AX\) sends each product either to zero or to
\(\pm2\) times another product, and different products to different
products, so the terms of a general \(X\) cannot cancel one another:
\(X\) anticommutes with all eight matrices only if it is a multiple of
the product of rank 8, which is \(\omega\). \(\blacksquare\)

\textbf{The complex projectors.} With complex numbers a chirality matrix
exists,
\[\chi\equiv i\,\omega=\begin{pmatrix}0&-i\,1_8\\ i\,1_8&0\end{pmatrix},\qquad\chi^\dagger=\chi,\qquad\chi^2=1,\]
anticommuting with every \(\gamma^A\) (check \texttt{71-P1});
\(\chi_{1671}\) is the matrix \texttt{χ17} of the notebook. The chiral
projectors are
\[\boxed{\;P_{1671}^{L}=\tfrac12\big(1-\chi_{1671}\big)=\tfrac12\begin{pmatrix}1_8&i\,1_8\\-i\,1_8&1_8\end{pmatrix},\qquad
P_{1671}^{R}=\tfrac12\big(1+\chi_{1671}\big)=\tfrac12\begin{pmatrix}1_8&-i\,1_8\\ i\,1_8&1_8\end{pmatrix}\;}\tag{2.4}\]
(checks \texttt{71-P2}, \texttt{71-P4}): Hermitian, \(P^{L}+P^{R}=1\),
\(P^{L}P^{L}=P^{L}\), \(P^{R}P^{R}=P^{R}\), \(P^{L}P^{R}=0\), trace 8
each, \(\gamma^AP^{L}=P^{R}\gamma^A\), and commuting with every
\(\Sigma^{AB}\) (check \texttt{71-P3}). They are \texttt{PL17},
\texttt{PR17} of \texttt{7+1.nb} (its sections 4.6 and 7.10). The Dirac
matrices remain real; only the projectors are complex.

\[
2\,P_{1671}^{L} = \left(\begin{array}{rrrrrrrr|rrrrrrrr}
\phantom{-}1 & \phantom{-}\cdot & \phantom{-}\cdot & \phantom{-}\cdot & \phantom{-}\cdot & \phantom{-}\cdot & \phantom{-}\cdot & \phantom{-}\cdot & \phantom{-}i & \phantom{-}\cdot & \phantom{-}\cdot & \phantom{-}\cdot & \phantom{-}\cdot & \phantom{-}\cdot & \phantom{-}\cdot & \phantom{-}\cdot \\[-1.8pt]
\phantom{-}\cdot & \phantom{-}1 & \phantom{-}\cdot & \phantom{-}\cdot & \phantom{-}\cdot & \phantom{-}\cdot & \phantom{-}\cdot & \phantom{-}\cdot & \phantom{-}\cdot & \phantom{-}i & \phantom{-}\cdot & \phantom{-}\cdot & \phantom{-}\cdot & \phantom{-}\cdot & \phantom{-}\cdot & \phantom{-}\cdot \\[-1.8pt]
\phantom{-}\cdot & \phantom{-}\cdot & \phantom{-}1 & \phantom{-}\cdot & \phantom{-}\cdot & \phantom{-}\cdot & \phantom{-}\cdot & \phantom{-}\cdot & \phantom{-}\cdot & \phantom{-}\cdot & \phantom{-}i & \phantom{-}\cdot & \phantom{-}\cdot & \phantom{-}\cdot & \phantom{-}\cdot & \phantom{-}\cdot \\[-1.8pt]
\phantom{-}\cdot & \phantom{-}\cdot & \phantom{-}\cdot & \phantom{-}1 & \phantom{-}\cdot & \phantom{-}\cdot & \phantom{-}\cdot & \phantom{-}\cdot & \phantom{-}\cdot & \phantom{-}\cdot & \phantom{-}\cdot & \phantom{-}i & \phantom{-}\cdot & \phantom{-}\cdot & \phantom{-}\cdot & \phantom{-}\cdot \\[-1.8pt]
\phantom{-}\cdot & \phantom{-}\cdot & \phantom{-}\cdot & \phantom{-}\cdot & \phantom{-}1 & \phantom{-}\cdot & \phantom{-}\cdot & \phantom{-}\cdot & \phantom{-}\cdot & \phantom{-}\cdot & \phantom{-}\cdot & \phantom{-}\cdot & \phantom{-}i & \phantom{-}\cdot & \phantom{-}\cdot & \phantom{-}\cdot \\[-1.8pt]
\phantom{-}\cdot & \phantom{-}\cdot & \phantom{-}\cdot & \phantom{-}\cdot & \phantom{-}\cdot & \phantom{-}1 & \phantom{-}\cdot & \phantom{-}\cdot & \phantom{-}\cdot & \phantom{-}\cdot & \phantom{-}\cdot & \phantom{-}\cdot & \phantom{-}\cdot & \phantom{-}i & \phantom{-}\cdot & \phantom{-}\cdot \\[-1.8pt]
\phantom{-}\cdot & \phantom{-}\cdot & \phantom{-}\cdot & \phantom{-}\cdot & \phantom{-}\cdot & \phantom{-}\cdot & \phantom{-}1 & \phantom{-}\cdot & \phantom{-}\cdot & \phantom{-}\cdot & \phantom{-}\cdot & \phantom{-}\cdot & \phantom{-}\cdot & \phantom{-}\cdot & \phantom{-}i & \phantom{-}\cdot \\[-1.8pt]
\phantom{-}\cdot & \phantom{-}\cdot & \phantom{-}\cdot & \phantom{-}\cdot & \phantom{-}\cdot & \phantom{-}\cdot & \phantom{-}\cdot & \phantom{-}1 & \phantom{-}\cdot & \phantom{-}\cdot & \phantom{-}\cdot & \phantom{-}\cdot & \phantom{-}\cdot & \phantom{-}\cdot & \phantom{-}\cdot & \phantom{-}i \\[-1.8pt]
\hline
-i & \phantom{-}\cdot & \phantom{-}\cdot & \phantom{-}\cdot & \phantom{-}\cdot & \phantom{-}\cdot & \phantom{-}\cdot & \phantom{-}\cdot & \phantom{-}1 & \phantom{-}\cdot & \phantom{-}\cdot & \phantom{-}\cdot & \phantom{-}\cdot & \phantom{-}\cdot & \phantom{-}\cdot & \phantom{-}\cdot \\[-1.8pt]
\phantom{-}\cdot & -i & \phantom{-}\cdot & \phantom{-}\cdot & \phantom{-}\cdot & \phantom{-}\cdot & \phantom{-}\cdot & \phantom{-}\cdot & \phantom{-}\cdot & \phantom{-}1 & \phantom{-}\cdot & \phantom{-}\cdot & \phantom{-}\cdot & \phantom{-}\cdot & \phantom{-}\cdot & \phantom{-}\cdot \\[-1.8pt]
\phantom{-}\cdot & \phantom{-}\cdot & -i & \phantom{-}\cdot & \phantom{-}\cdot & \phantom{-}\cdot & \phantom{-}\cdot & \phantom{-}\cdot & \phantom{-}\cdot & \phantom{-}\cdot & \phantom{-}1 & \phantom{-}\cdot & \phantom{-}\cdot & \phantom{-}\cdot & \phantom{-}\cdot & \phantom{-}\cdot \\[-1.8pt]
\phantom{-}\cdot & \phantom{-}\cdot & \phantom{-}\cdot & -i & \phantom{-}\cdot & \phantom{-}\cdot & \phantom{-}\cdot & \phantom{-}\cdot & \phantom{-}\cdot & \phantom{-}\cdot & \phantom{-}\cdot & \phantom{-}1 & \phantom{-}\cdot & \phantom{-}\cdot & \phantom{-}\cdot & \phantom{-}\cdot \\[-1.8pt]
\phantom{-}\cdot & \phantom{-}\cdot & \phantom{-}\cdot & \phantom{-}\cdot & -i & \phantom{-}\cdot & \phantom{-}\cdot & \phantom{-}\cdot & \phantom{-}\cdot & \phantom{-}\cdot & \phantom{-}\cdot & \phantom{-}\cdot & \phantom{-}1 & \phantom{-}\cdot & \phantom{-}\cdot & \phantom{-}\cdot \\[-1.8pt]
\phantom{-}\cdot & \phantom{-}\cdot & \phantom{-}\cdot & \phantom{-}\cdot & \phantom{-}\cdot & -i & \phantom{-}\cdot & \phantom{-}\cdot & \phantom{-}\cdot & \phantom{-}\cdot & \phantom{-}\cdot & \phantom{-}\cdot & \phantom{-}\cdot & \phantom{-}1 & \phantom{-}\cdot & \phantom{-}\cdot \\[-1.8pt]
\phantom{-}\cdot & \phantom{-}\cdot & \phantom{-}\cdot & \phantom{-}\cdot & \phantom{-}\cdot & \phantom{-}\cdot & -i & \phantom{-}\cdot & \phantom{-}\cdot & \phantom{-}\cdot & \phantom{-}\cdot & \phantom{-}\cdot & \phantom{-}\cdot & \phantom{-}\cdot & \phantom{-}1 & \phantom{-}\cdot \\[-1.8pt]
\phantom{-}\cdot & \phantom{-}\cdot & \phantom{-}\cdot & \phantom{-}\cdot & \phantom{-}\cdot & \phantom{-}\cdot & \phantom{-}\cdot & -i & \phantom{-}\cdot & \phantom{-}\cdot & \phantom{-}\cdot & \phantom{-}\cdot & \phantom{-}\cdot & \phantom{-}\cdot & \phantom{-}\cdot & \phantom{-}1
\end{array}\right)
\]

\[
2\,P_{1671}^{R} = \left(\begin{array}{rrrrrrrr|rrrrrrrr}
\phantom{-}1 & \phantom{-}\cdot & \phantom{-}\cdot & \phantom{-}\cdot & \phantom{-}\cdot & \phantom{-}\cdot & \phantom{-}\cdot & \phantom{-}\cdot & -i & \phantom{-}\cdot & \phantom{-}\cdot & \phantom{-}\cdot & \phantom{-}\cdot & \phantom{-}\cdot & \phantom{-}\cdot & \phantom{-}\cdot \\[-1.8pt]
\phantom{-}\cdot & \phantom{-}1 & \phantom{-}\cdot & \phantom{-}\cdot & \phantom{-}\cdot & \phantom{-}\cdot & \phantom{-}\cdot & \phantom{-}\cdot & \phantom{-}\cdot & -i & \phantom{-}\cdot & \phantom{-}\cdot & \phantom{-}\cdot & \phantom{-}\cdot & \phantom{-}\cdot & \phantom{-}\cdot \\[-1.8pt]
\phantom{-}\cdot & \phantom{-}\cdot & \phantom{-}1 & \phantom{-}\cdot & \phantom{-}\cdot & \phantom{-}\cdot & \phantom{-}\cdot & \phantom{-}\cdot & \phantom{-}\cdot & \phantom{-}\cdot & -i & \phantom{-}\cdot & \phantom{-}\cdot & \phantom{-}\cdot & \phantom{-}\cdot & \phantom{-}\cdot \\[-1.8pt]
\phantom{-}\cdot & \phantom{-}\cdot & \phantom{-}\cdot & \phantom{-}1 & \phantom{-}\cdot & \phantom{-}\cdot & \phantom{-}\cdot & \phantom{-}\cdot & \phantom{-}\cdot & \phantom{-}\cdot & \phantom{-}\cdot & -i & \phantom{-}\cdot & \phantom{-}\cdot & \phantom{-}\cdot & \phantom{-}\cdot \\[-1.8pt]
\phantom{-}\cdot & \phantom{-}\cdot & \phantom{-}\cdot & \phantom{-}\cdot & \phantom{-}1 & \phantom{-}\cdot & \phantom{-}\cdot & \phantom{-}\cdot & \phantom{-}\cdot & \phantom{-}\cdot & \phantom{-}\cdot & \phantom{-}\cdot & -i & \phantom{-}\cdot & \phantom{-}\cdot & \phantom{-}\cdot \\[-1.8pt]
\phantom{-}\cdot & \phantom{-}\cdot & \phantom{-}\cdot & \phantom{-}\cdot & \phantom{-}\cdot & \phantom{-}1 & \phantom{-}\cdot & \phantom{-}\cdot & \phantom{-}\cdot & \phantom{-}\cdot & \phantom{-}\cdot & \phantom{-}\cdot & \phantom{-}\cdot & -i & \phantom{-}\cdot & \phantom{-}\cdot \\[-1.8pt]
\phantom{-}\cdot & \phantom{-}\cdot & \phantom{-}\cdot & \phantom{-}\cdot & \phantom{-}\cdot & \phantom{-}\cdot & \phantom{-}1 & \phantom{-}\cdot & \phantom{-}\cdot & \phantom{-}\cdot & \phantom{-}\cdot & \phantom{-}\cdot & \phantom{-}\cdot & \phantom{-}\cdot & -i & \phantom{-}\cdot \\[-1.8pt]
\phantom{-}\cdot & \phantom{-}\cdot & \phantom{-}\cdot & \phantom{-}\cdot & \phantom{-}\cdot & \phantom{-}\cdot & \phantom{-}\cdot & \phantom{-}1 & \phantom{-}\cdot & \phantom{-}\cdot & \phantom{-}\cdot & \phantom{-}\cdot & \phantom{-}\cdot & \phantom{-}\cdot & \phantom{-}\cdot & -i \\[-1.8pt]
\hline
\phantom{-}i & \phantom{-}\cdot & \phantom{-}\cdot & \phantom{-}\cdot & \phantom{-}\cdot & \phantom{-}\cdot & \phantom{-}\cdot & \phantom{-}\cdot & \phantom{-}1 & \phantom{-}\cdot & \phantom{-}\cdot & \phantom{-}\cdot & \phantom{-}\cdot & \phantom{-}\cdot & \phantom{-}\cdot & \phantom{-}\cdot \\[-1.8pt]
\phantom{-}\cdot & \phantom{-}i & \phantom{-}\cdot & \phantom{-}\cdot & \phantom{-}\cdot & \phantom{-}\cdot & \phantom{-}\cdot & \phantom{-}\cdot & \phantom{-}\cdot & \phantom{-}1 & \phantom{-}\cdot & \phantom{-}\cdot & \phantom{-}\cdot & \phantom{-}\cdot & \phantom{-}\cdot & \phantom{-}\cdot \\[-1.8pt]
\phantom{-}\cdot & \phantom{-}\cdot & \phantom{-}i & \phantom{-}\cdot & \phantom{-}\cdot & \phantom{-}\cdot & \phantom{-}\cdot & \phantom{-}\cdot & \phantom{-}\cdot & \phantom{-}\cdot & \phantom{-}1 & \phantom{-}\cdot & \phantom{-}\cdot & \phantom{-}\cdot & \phantom{-}\cdot & \phantom{-}\cdot \\[-1.8pt]
\phantom{-}\cdot & \phantom{-}\cdot & \phantom{-}\cdot & \phantom{-}i & \phantom{-}\cdot & \phantom{-}\cdot & \phantom{-}\cdot & \phantom{-}\cdot & \phantom{-}\cdot & \phantom{-}\cdot & \phantom{-}\cdot & \phantom{-}1 & \phantom{-}\cdot & \phantom{-}\cdot & \phantom{-}\cdot & \phantom{-}\cdot \\[-1.8pt]
\phantom{-}\cdot & \phantom{-}\cdot & \phantom{-}\cdot & \phantom{-}\cdot & \phantom{-}i & \phantom{-}\cdot & \phantom{-}\cdot & \phantom{-}\cdot & \phantom{-}\cdot & \phantom{-}\cdot & \phantom{-}\cdot & \phantom{-}\cdot & \phantom{-}1 & \phantom{-}\cdot & \phantom{-}\cdot & \phantom{-}\cdot \\[-1.8pt]
\phantom{-}\cdot & \phantom{-}\cdot & \phantom{-}\cdot & \phantom{-}\cdot & \phantom{-}\cdot & \phantom{-}i & \phantom{-}\cdot & \phantom{-}\cdot & \phantom{-}\cdot & \phantom{-}\cdot & \phantom{-}\cdot & \phantom{-}\cdot & \phantom{-}\cdot & \phantom{-}1 & \phantom{-}\cdot & \phantom{-}\cdot \\[-1.8pt]
\phantom{-}\cdot & \phantom{-}\cdot & \phantom{-}\cdot & \phantom{-}\cdot & \phantom{-}\cdot & \phantom{-}\cdot & \phantom{-}i & \phantom{-}\cdot & \phantom{-}\cdot & \phantom{-}\cdot & \phantom{-}\cdot & \phantom{-}\cdot & \phantom{-}\cdot & \phantom{-}\cdot & \phantom{-}1 & \phantom{-}\cdot \\[-1.8pt]
\phantom{-}\cdot & \phantom{-}\cdot & \phantom{-}\cdot & \phantom{-}\cdot & \phantom{-}\cdot & \phantom{-}\cdot & \phantom{-}\cdot & \phantom{-}i & \phantom{-}\cdot & \phantom{-}\cdot & \phantom{-}\cdot & \phantom{-}\cdot & \phantom{-}\cdot & \phantom{-}\cdot & \phantom{-}\cdot & \phantom{-}1
\end{array}\right)
\]

\textbf{Relation to the projectors of \texttt{4+4.nb}.} The
specification asks, in its part {[}2{]} as in its part {[}1{]}, to
re-derive the projectors of the subsection ``Left - Handed vs . Right -
Handed Fermion Fields (Chirality) ; projection operators'' of
\texttt{4+4.nb}. That was done in Sec. 1.6. For the 7+1 matrices those
real projectors are not chiral projectors (Theorem 3), but they are
unitarily equivalent to the complex ones: with the unitary matrix
\(U_{1671}^{c}=(1-i\,\texttt{Γ17[0]})/\sqrt2\),
\[\begin{aligned}U_{1671}^{c}\,\chi_{1671}\,(U_{1671}^{c})^{-1}&=-\omega_{1644},\\
U_{1671}^{c}\,P_{1671}^{L}\,(U_{1671}^{c})^{-1}&=P_{1644}^{R},\qquad U_{1671}^{c}\,P_{1671}^{R}\,(U_{1671}^{c})^{-1}=P_{1644}^{L}\end{aligned}\tag{2.5}\]
(checks \texttt{71-P5}, \texttt{I71-U}). \(U_{1671}^{c}\) is the matrix
by which \texttt{7+1.nb} (its section 4.3) passes from the real set
\texttt{Γ17} to its complex set \texttt{Γ71}, which has the block form
of \texttt{4+4.nb}. The exchange of the names L and R in Eq. (2.5) is a
matter of convention: the notebook defines \texttt{χ17}
\(=+i\,\)\texttt{Γ17{[}8{]}} for the real set but \texttt{χ71}
\(=-i\,\)\texttt{Γ71{[}8{]}} for the complex set. The project's
\(\chi_{1671}\) is \texttt{χ17}.

\hypertarget{other-important-1616-matrices-of-71}{%
\subsection{2.7 Other important 16×16 matrices of
7+1}\label{other-important-1616-matrices-of-71}}

\textbf{The Dirac adjoint matrix} is the time-like matrix itself,
\(\gamma_{1671}^4\): it is real, symmetric, squares to \(1\), and
\((\gamma^A)^T=\gamma^4\gamma^A\gamma^4\) (check \texttt{71-S1}). We
write \(\bar\psi\equiv\psi^\dagger\gamma^4\).

\textbf{The two transposition matrices} (check \texttt{71-K0}):
\[C_{1671}^{+}\equiv\gamma_{1671}^4,\qquad C_{1671}^{-}\equiv\gamma_{1671}^4\,\omega_{1671}=\begin{pmatrix}1_8&0\\0&-1_8\end{pmatrix},\qquad
C^{\pm}\gamma^A(C^{\pm})^{-1}=\pm(\gamma^A)^T .\tag{2.6}\] Both are real
and symmetric.

\[
C_{1671}^{+} = \left(\begin{array}{rrrrrrrr|rrrrrrrr}
\phantom{-}\cdot & \phantom{-}\cdot & \phantom{-}\cdot & \phantom{-}\cdot & \phantom{-}\cdot & \phantom{-}\cdot & \phantom{-}\cdot & \phantom{-}\cdot & \phantom{-}1 & \phantom{-}\cdot & \phantom{-}\cdot & \phantom{-}\cdot & \phantom{-}\cdot & \phantom{-}\cdot & \phantom{-}\cdot & \phantom{-}\cdot \\[-1.8pt]
\phantom{-}\cdot & \phantom{-}\cdot & \phantom{-}\cdot & \phantom{-}\cdot & \phantom{-}\cdot & \phantom{-}\cdot & \phantom{-}\cdot & \phantom{-}\cdot & \phantom{-}\cdot & \phantom{-}1 & \phantom{-}\cdot & \phantom{-}\cdot & \phantom{-}\cdot & \phantom{-}\cdot & \phantom{-}\cdot & \phantom{-}\cdot \\[-1.8pt]
\phantom{-}\cdot & \phantom{-}\cdot & \phantom{-}\cdot & \phantom{-}\cdot & \phantom{-}\cdot & \phantom{-}\cdot & \phantom{-}\cdot & \phantom{-}\cdot & \phantom{-}\cdot & \phantom{-}\cdot & \phantom{-}1 & \phantom{-}\cdot & \phantom{-}\cdot & \phantom{-}\cdot & \phantom{-}\cdot & \phantom{-}\cdot \\[-1.8pt]
\phantom{-}\cdot & \phantom{-}\cdot & \phantom{-}\cdot & \phantom{-}\cdot & \phantom{-}\cdot & \phantom{-}\cdot & \phantom{-}\cdot & \phantom{-}\cdot & \phantom{-}\cdot & \phantom{-}\cdot & \phantom{-}\cdot & \phantom{-}1 & \phantom{-}\cdot & \phantom{-}\cdot & \phantom{-}\cdot & \phantom{-}\cdot \\[-1.8pt]
\phantom{-}\cdot & \phantom{-}\cdot & \phantom{-}\cdot & \phantom{-}\cdot & \phantom{-}\cdot & \phantom{-}\cdot & \phantom{-}\cdot & \phantom{-}\cdot & \phantom{-}\cdot & \phantom{-}\cdot & \phantom{-}\cdot & \phantom{-}\cdot & \phantom{-}1 & \phantom{-}\cdot & \phantom{-}\cdot & \phantom{-}\cdot \\[-1.8pt]
\phantom{-}\cdot & \phantom{-}\cdot & \phantom{-}\cdot & \phantom{-}\cdot & \phantom{-}\cdot & \phantom{-}\cdot & \phantom{-}\cdot & \phantom{-}\cdot & \phantom{-}\cdot & \phantom{-}\cdot & \phantom{-}\cdot & \phantom{-}\cdot & \phantom{-}\cdot & \phantom{-}1 & \phantom{-}\cdot & \phantom{-}\cdot \\[-1.8pt]
\phantom{-}\cdot & \phantom{-}\cdot & \phantom{-}\cdot & \phantom{-}\cdot & \phantom{-}\cdot & \phantom{-}\cdot & \phantom{-}\cdot & \phantom{-}\cdot & \phantom{-}\cdot & \phantom{-}\cdot & \phantom{-}\cdot & \phantom{-}\cdot & \phantom{-}\cdot & \phantom{-}\cdot & \phantom{-}1 & \phantom{-}\cdot \\[-1.8pt]
\phantom{-}\cdot & \phantom{-}\cdot & \phantom{-}\cdot & \phantom{-}\cdot & \phantom{-}\cdot & \phantom{-}\cdot & \phantom{-}\cdot & \phantom{-}\cdot & \phantom{-}\cdot & \phantom{-}\cdot & \phantom{-}\cdot & \phantom{-}\cdot & \phantom{-}\cdot & \phantom{-}\cdot & \phantom{-}\cdot & \phantom{-}1 \\[-1.8pt]
\hline
\phantom{-}1 & \phantom{-}\cdot & \phantom{-}\cdot & \phantom{-}\cdot & \phantom{-}\cdot & \phantom{-}\cdot & \phantom{-}\cdot & \phantom{-}\cdot & \phantom{-}\cdot & \phantom{-}\cdot & \phantom{-}\cdot & \phantom{-}\cdot & \phantom{-}\cdot & \phantom{-}\cdot & \phantom{-}\cdot & \phantom{-}\cdot \\[-1.8pt]
\phantom{-}\cdot & \phantom{-}1 & \phantom{-}\cdot & \phantom{-}\cdot & \phantom{-}\cdot & \phantom{-}\cdot & \phantom{-}\cdot & \phantom{-}\cdot & \phantom{-}\cdot & \phantom{-}\cdot & \phantom{-}\cdot & \phantom{-}\cdot & \phantom{-}\cdot & \phantom{-}\cdot & \phantom{-}\cdot & \phantom{-}\cdot \\[-1.8pt]
\phantom{-}\cdot & \phantom{-}\cdot & \phantom{-}1 & \phantom{-}\cdot & \phantom{-}\cdot & \phantom{-}\cdot & \phantom{-}\cdot & \phantom{-}\cdot & \phantom{-}\cdot & \phantom{-}\cdot & \phantom{-}\cdot & \phantom{-}\cdot & \phantom{-}\cdot & \phantom{-}\cdot & \phantom{-}\cdot & \phantom{-}\cdot \\[-1.8pt]
\phantom{-}\cdot & \phantom{-}\cdot & \phantom{-}\cdot & \phantom{-}1 & \phantom{-}\cdot & \phantom{-}\cdot & \phantom{-}\cdot & \phantom{-}\cdot & \phantom{-}\cdot & \phantom{-}\cdot & \phantom{-}\cdot & \phantom{-}\cdot & \phantom{-}\cdot & \phantom{-}\cdot & \phantom{-}\cdot & \phantom{-}\cdot \\[-1.8pt]
\phantom{-}\cdot & \phantom{-}\cdot & \phantom{-}\cdot & \phantom{-}\cdot & \phantom{-}1 & \phantom{-}\cdot & \phantom{-}\cdot & \phantom{-}\cdot & \phantom{-}\cdot & \phantom{-}\cdot & \phantom{-}\cdot & \phantom{-}\cdot & \phantom{-}\cdot & \phantom{-}\cdot & \phantom{-}\cdot & \phantom{-}\cdot \\[-1.8pt]
\phantom{-}\cdot & \phantom{-}\cdot & \phantom{-}\cdot & \phantom{-}\cdot & \phantom{-}\cdot & \phantom{-}1 & \phantom{-}\cdot & \phantom{-}\cdot & \phantom{-}\cdot & \phantom{-}\cdot & \phantom{-}\cdot & \phantom{-}\cdot & \phantom{-}\cdot & \phantom{-}\cdot & \phantom{-}\cdot & \phantom{-}\cdot \\[-1.8pt]
\phantom{-}\cdot & \phantom{-}\cdot & \phantom{-}\cdot & \phantom{-}\cdot & \phantom{-}\cdot & \phantom{-}\cdot & \phantom{-}1 & \phantom{-}\cdot & \phantom{-}\cdot & \phantom{-}\cdot & \phantom{-}\cdot & \phantom{-}\cdot & \phantom{-}\cdot & \phantom{-}\cdot & \phantom{-}\cdot & \phantom{-}\cdot \\[-1.8pt]
\phantom{-}\cdot & \phantom{-}\cdot & \phantom{-}\cdot & \phantom{-}\cdot & \phantom{-}\cdot & \phantom{-}\cdot & \phantom{-}\cdot & \phantom{-}1 & \phantom{-}\cdot & \phantom{-}\cdot & \phantom{-}\cdot & \phantom{-}\cdot & \phantom{-}\cdot & \phantom{-}\cdot & \phantom{-}\cdot & \phantom{-}\cdot
\end{array}\right)
\]

\[
C_{1671}^{-} = \left(\begin{array}{rrrrrrrr|rrrrrrrr}
\phantom{-}1 & \phantom{-}\cdot & \phantom{-}\cdot & \phantom{-}\cdot & \phantom{-}\cdot & \phantom{-}\cdot & \phantom{-}\cdot & \phantom{-}\cdot & \phantom{-}\cdot & \phantom{-}\cdot & \phantom{-}\cdot & \phantom{-}\cdot & \phantom{-}\cdot & \phantom{-}\cdot & \phantom{-}\cdot & \phantom{-}\cdot \\[-1.8pt]
\phantom{-}\cdot & \phantom{-}1 & \phantom{-}\cdot & \phantom{-}\cdot & \phantom{-}\cdot & \phantom{-}\cdot & \phantom{-}\cdot & \phantom{-}\cdot & \phantom{-}\cdot & \phantom{-}\cdot & \phantom{-}\cdot & \phantom{-}\cdot & \phantom{-}\cdot & \phantom{-}\cdot & \phantom{-}\cdot & \phantom{-}\cdot \\[-1.8pt]
\phantom{-}\cdot & \phantom{-}\cdot & \phantom{-}1 & \phantom{-}\cdot & \phantom{-}\cdot & \phantom{-}\cdot & \phantom{-}\cdot & \phantom{-}\cdot & \phantom{-}\cdot & \phantom{-}\cdot & \phantom{-}\cdot & \phantom{-}\cdot & \phantom{-}\cdot & \phantom{-}\cdot & \phantom{-}\cdot & \phantom{-}\cdot \\[-1.8pt]
\phantom{-}\cdot & \phantom{-}\cdot & \phantom{-}\cdot & \phantom{-}1 & \phantom{-}\cdot & \phantom{-}\cdot & \phantom{-}\cdot & \phantom{-}\cdot & \phantom{-}\cdot & \phantom{-}\cdot & \phantom{-}\cdot & \phantom{-}\cdot & \phantom{-}\cdot & \phantom{-}\cdot & \phantom{-}\cdot & \phantom{-}\cdot \\[-1.8pt]
\phantom{-}\cdot & \phantom{-}\cdot & \phantom{-}\cdot & \phantom{-}\cdot & \phantom{-}1 & \phantom{-}\cdot & \phantom{-}\cdot & \phantom{-}\cdot & \phantom{-}\cdot & \phantom{-}\cdot & \phantom{-}\cdot & \phantom{-}\cdot & \phantom{-}\cdot & \phantom{-}\cdot & \phantom{-}\cdot & \phantom{-}\cdot \\[-1.8pt]
\phantom{-}\cdot & \phantom{-}\cdot & \phantom{-}\cdot & \phantom{-}\cdot & \phantom{-}\cdot & \phantom{-}1 & \phantom{-}\cdot & \phantom{-}\cdot & \phantom{-}\cdot & \phantom{-}\cdot & \phantom{-}\cdot & \phantom{-}\cdot & \phantom{-}\cdot & \phantom{-}\cdot & \phantom{-}\cdot & \phantom{-}\cdot \\[-1.8pt]
\phantom{-}\cdot & \phantom{-}\cdot & \phantom{-}\cdot & \phantom{-}\cdot & \phantom{-}\cdot & \phantom{-}\cdot & \phantom{-}1 & \phantom{-}\cdot & \phantom{-}\cdot & \phantom{-}\cdot & \phantom{-}\cdot & \phantom{-}\cdot & \phantom{-}\cdot & \phantom{-}\cdot & \phantom{-}\cdot & \phantom{-}\cdot \\[-1.8pt]
\phantom{-}\cdot & \phantom{-}\cdot & \phantom{-}\cdot & \phantom{-}\cdot & \phantom{-}\cdot & \phantom{-}\cdot & \phantom{-}\cdot & \phantom{-}1 & \phantom{-}\cdot & \phantom{-}\cdot & \phantom{-}\cdot & \phantom{-}\cdot & \phantom{-}\cdot & \phantom{-}\cdot & \phantom{-}\cdot & \phantom{-}\cdot \\[-1.8pt]
\hline
\phantom{-}\cdot & \phantom{-}\cdot & \phantom{-}\cdot & \phantom{-}\cdot & \phantom{-}\cdot & \phantom{-}\cdot & \phantom{-}\cdot & \phantom{-}\cdot & -1 & \phantom{-}\cdot & \phantom{-}\cdot & \phantom{-}\cdot & \phantom{-}\cdot & \phantom{-}\cdot & \phantom{-}\cdot & \phantom{-}\cdot \\[-1.8pt]
\phantom{-}\cdot & \phantom{-}\cdot & \phantom{-}\cdot & \phantom{-}\cdot & \phantom{-}\cdot & \phantom{-}\cdot & \phantom{-}\cdot & \phantom{-}\cdot & \phantom{-}\cdot & -1 & \phantom{-}\cdot & \phantom{-}\cdot & \phantom{-}\cdot & \phantom{-}\cdot & \phantom{-}\cdot & \phantom{-}\cdot \\[-1.8pt]
\phantom{-}\cdot & \phantom{-}\cdot & \phantom{-}\cdot & \phantom{-}\cdot & \phantom{-}\cdot & \phantom{-}\cdot & \phantom{-}\cdot & \phantom{-}\cdot & \phantom{-}\cdot & \phantom{-}\cdot & -1 & \phantom{-}\cdot & \phantom{-}\cdot & \phantom{-}\cdot & \phantom{-}\cdot & \phantom{-}\cdot \\[-1.8pt]
\phantom{-}\cdot & \phantom{-}\cdot & \phantom{-}\cdot & \phantom{-}\cdot & \phantom{-}\cdot & \phantom{-}\cdot & \phantom{-}\cdot & \phantom{-}\cdot & \phantom{-}\cdot & \phantom{-}\cdot & \phantom{-}\cdot & -1 & \phantom{-}\cdot & \phantom{-}\cdot & \phantom{-}\cdot & \phantom{-}\cdot \\[-1.8pt]
\phantom{-}\cdot & \phantom{-}\cdot & \phantom{-}\cdot & \phantom{-}\cdot & \phantom{-}\cdot & \phantom{-}\cdot & \phantom{-}\cdot & \phantom{-}\cdot & \phantom{-}\cdot & \phantom{-}\cdot & \phantom{-}\cdot & \phantom{-}\cdot & -1 & \phantom{-}\cdot & \phantom{-}\cdot & \phantom{-}\cdot \\[-1.8pt]
\phantom{-}\cdot & \phantom{-}\cdot & \phantom{-}\cdot & \phantom{-}\cdot & \phantom{-}\cdot & \phantom{-}\cdot & \phantom{-}\cdot & \phantom{-}\cdot & \phantom{-}\cdot & \phantom{-}\cdot & \phantom{-}\cdot & \phantom{-}\cdot & \phantom{-}\cdot & -1 & \phantom{-}\cdot & \phantom{-}\cdot \\[-1.8pt]
\phantom{-}\cdot & \phantom{-}\cdot & \phantom{-}\cdot & \phantom{-}\cdot & \phantom{-}\cdot & \phantom{-}\cdot & \phantom{-}\cdot & \phantom{-}\cdot & \phantom{-}\cdot & \phantom{-}\cdot & \phantom{-}\cdot & \phantom{-}\cdot & \phantom{-}\cdot & \phantom{-}\cdot & -1 & \phantom{-}\cdot \\[-1.8pt]
\phantom{-}\cdot & \phantom{-}\cdot & \phantom{-}\cdot & \phantom{-}\cdot & \phantom{-}\cdot & \phantom{-}\cdot & \phantom{-}\cdot & \phantom{-}\cdot & \phantom{-}\cdot & \phantom{-}\cdot & \phantom{-}\cdot & \phantom{-}\cdot & \phantom{-}\cdot & \phantom{-}\cdot & \phantom{-}\cdot & -1
\end{array}\right)
\]

\textbf{The Hamiltonian matrices.} \(\beta\equiv\gamma^4\) and
\(\alpha^j\equiv\gamma^4\gamma^j\) (\(j=0,1,2,3,5,6,7\)) are real
symmetric, with \(\alpha^i\alpha^j+\alpha^j\alpha^i=2\delta^{ij}\),
\(\alpha^j\beta+\beta\alpha^j=0\), \(\beta^2=1\) (check \texttt{71-D4}).
In blocks,
\[\beta_{1671}=\begin{pmatrix}0&1_8\\1_8&0\end{pmatrix},\qquad\alpha_{1671}^j=\begin{pmatrix}0&\tau_{1671}^j\\-\tau_{1671}^j&0\end{pmatrix}\tag{2.7}\]
(check \texttt{71-D6}); they are \texttt{βD} and
\texttt{αD{[}swap04(j){]}} of \texttt{7+1.nb}, section 7.1.

\textbf{The 8×8 blocks} \(\tau_{1671}^A\) are listed in Appendix B.

\hypertarget{the-dirac-equation-of-the-field-dirac1671}{%
\subsection{2.8 The Dirac equation of the field
dirac1671}\label{the-dirac-equation-of-the-field-dirac1671}}

\hypertarget{the-equation-1}{%
\subsubsection{2.8.1 The equation}\label{the-equation-1}}

The field \textbf{dirac1671} is the column
\(\psi_{1671}(x)=(\psi_1(x),\dots,\psi_{16}(x))^T\) of sixteen
complex-valued Grassmann operators. Its Dirac equation is
\[\boxed{\;\big(i\,\gamma_{1671}^A\,\partial_A-m_{1671}\big)\,\psi_{1671}=0\;}\tag{2.9}\]
i.e.~\(i\gamma^4\partial_4\psi+i\sum_j\gamma^j\partial_j\psi-m\psi=0\)
with the real matrices of Sec. 2.2 and a real mass \(m_{1671}>0\). With
Eq. (2.3),
\((\gamma^A\partial_A)^2=\eta^{AB}\partial_A\partial_B=\partial_4^2-\nabla^2\),
where
\(\nabla^2=\partial_0^2+\partial_1^2+\partial_2^2+\partial_3^2+\partial_5^2+\partial_6^2+\partial_7^2\),
so
\[\big(i\gamma^A\partial_A+m\big)\big(i\gamma^B\partial_B-m\big)=-\big(\partial_4^2-\nabla^2+m^2\big) :\tag{2.10}\]
every component obeys the Klein--Gordon equation of 7+1 spacetime. For a
plane wave \(\psi=u\,e^{-ik_Ax^A}\) the operator becomes
\(\gamma^Ak_A-m\) and the mass shell is
\(k_4^2=k_0^2+k_1^2+k_2^2+k_3^2+k_5^2+k_6^2+k_7^2+m^2\) (checks
\texttt{71-D1}, \texttt{71-D2}). Equation (2.9) is the Dirac equation of
\texttt{7+1.nb} (its section 7.2) for the field that the notebook calls
\texttt{dirac16complex}, written with the coordinates relabelled by
\(\mathrm{swap}_{04}\).

\emph{The factor \(i\) is unavoidable here.} A real first-order operator
\(\gamma^A\partial_A-m\) squares to \(\partial_4^2-\nabla^2-m^2\), the
wrong sign of the mass term. A real equation with the correct mass term
would require eight real 16×16 matrices \(M^A\) with
\(M^AM^B+M^BM^A=-2\eta_{1671}^{AB}\) (seven squares \(+1\), one square
\(-1\)), and by Sec. 2.4 no such matrices exist. The matrices in Eq.
(2.9) are real; only the number \(i\) in front of them is not.

\textbf{Worked example} (check \texttt{71-D3}, the example of the
notebook): energy \(E=5\), momentum \((p^1,p^2,p^3)=(2,2,4)\), \(m=1\);
then \(\Pi_{1671}=(\gamma^Ak_A+m)/2m\) is a projector of trace 8 with
\((\gamma^Ak_A-m)\Pi_{1671}=0\): eight independent positive-energy plane
waves.

\hypertarget{hamiltonian-form-plane-waves-quantization}{%
\subsubsection{2.8.2 Hamiltonian form, plane waves,
quantization}\label{hamiltonian-form-plane-waves-quantization}}

Multiplying Eq. (2.9) by \(\gamma^4\) gives
\[i\,\partial_4\psi=\mathsf{h}\,\psi,\qquad\mathsf{h}_{1671}=-i\,\alpha_{1671}^j\partial_j+\beta_{1671}\,m_{1671},\tag{2.11}\]
with the real symmetric matrices of Eq. (2.7) and summation over the
seven space indices \(j\). In momentum space
\(\mathsf{h}(\vec p)=p^j\alpha^j+m\beta\) is a real symmetric matrix
with \(\mathsf{h}(\vec p)^2=(|\vec p|^2+m^2)\,1\) (check
\texttt{71-D5}): the evolution in \(x4\) is unitary for every momentum,
and the energies are \(\pm E\), \(E=\sqrt{|\vec p|^2+m^2}\), eight times
each. There is one time, and none of the difficulties of Sec. 1.8.4
arises.

\emph{Real spinors.} With the real 8×8 matrix
\(M_{1671}(\vec p)=m+p^j\tau^j\), which satisfies
\(M_{1671}M_{1671}^T=E^2\), the eigenvectors of \(\mathsf{h}\) are real:
\[u_s(\vec p)=\frac1{\sqrt E}\begin{pmatrix}E\,e_s\\M_{1671}^Te_s\end{pmatrix},\quad\mathsf{h}(\vec p)u_s=+Eu_s;\qquad
v_s(\vec p)=\frac1{\sqrt E}\begin{pmatrix}E\,e_s\\-M_{1671}e_s\end{pmatrix},\quad\mathsf{h}(-\vec p)v_s=-Ev_s,\]
where \(e_s\) (\(s=1,\dots,8\)) are the unit vectors of
\(\mathbb{R}^8\); they are normalized to
\(u_s^Tu_r=v_s^Tv_r=2E\delta_{sr}\) (check \texttt{71-D8}, exact for the
worked example).

\emph{Lagrangian, current, quantization.} With
\(\bar\psi=\psi^\dagger\gamma^4\),
\[\mathcal{L}_{1671}=\bar\psi\big(i\gamma^A\partial_A-m\big)\psi,\qquad j_{1671}^A=\bar\psi\gamma^A\psi,\qquad\partial_Aj^A_{1671}=0,\]
and the density \(j_{1671}^4=\psi^\dagger\psi\) is positive. The
quantized field is, in a box of volume \(V\),
\[\psi(x4,\vec x)=\sum_{\vec p}\sum_{s=1}^{8}\frac{1}{\sqrt{2EV}}\Big[b_s(\vec p)\,u_s(\vec p)\,e^{-i(E\,x4-\vec p\cdot\vec x)}+d_s^\dagger(\vec p)\,v_s(\vec p)\,e^{+i(E\,x4-\vec p\cdot\vec x)}\Big],\]
with
\(\{b_s(\vec p),b_r^\dagger(\vec q)\}=\{d_s(\vec p),d_r^\dagger(\vec q)\}=\delta_{sr}\delta_{\vec p\vec q}\)
and all other anticommutators zero, which is equivalent to
\(\{\psi_a(x4,\vec x),\psi_b^\dagger(x4,\vec y)\}=\delta_{ab}\delta^7(\vec x-\vec y)\).
A particle has eight spin states and so has its antiparticle. These
statements are derived and checked symbolically in \texttt{7+1.nb} (its
sections 7.4--7.7); here they are restated in the relabelled
coordinates, and the algebraic identities on which they rest are checks
\texttt{71-D4}, \texttt{71-D5}, \texttt{71-D8}.

\emph{Charge conjugation} (check \texttt{71-D9}):
\(\psi_c=\omega\,\psi^*\), where \(\psi^*\) is the column of the
Hermitian conjugates of the components, solves Eq. (2.9) whenever
\(\psi\) does, and \((\psi_c)_c=-\psi\); a 16-component Majorana
condition \(\psi_c=\psi\) has no non-zero solution.

\hypertarget{the-two-coupled-8-component-equations-1}{%
\subsubsection{2.8.3 The two coupled 8-component
equations}\label{the-two-coupled-8-component-equations-1}}

The eigenvectors of \(\chi\) give two 16×8 matrices with orthonormal
columns,
\[E^{L}=\frac1{\sqrt2}\begin{pmatrix}1_8\\-i\,1_8\end{pmatrix}\;(\chi E^{L}=-E^{L}),\qquad
E^{R}=\frac1{\sqrt2}\begin{pmatrix}-i\,1_8\\1_8\end{pmatrix}\;(\chi E^{R}=+E^{R}),\]
with \(E^{L}(E^{L})^\dagger=P^{L}\) and \(E^{R}(E^{R})^\dagger=P^{R}\).
The chiral parts of the field are the 8-component fields
\[\phi^{L}=(E^{L})^\dagger\psi=\frac{\psi_{\rm up}+i\,\psi_{\rm down}}{\sqrt2},\qquad\phi^{R}=(E^{R})^\dagger\psi=\frac{i\,\psi_{\rm up}+\psi_{\rm down}}{\sqrt2},\]
where \(\psi_{\rm up}\), \(\psi_{\rm down}\) are the upper and lower
eight components, so that \(P^{L}\psi=E^{L}\phi^{L}\) and
\(P^{R}\psi=E^{R}\phi^{R}\). In this basis
\((E^{L})^\dagger\mathsf{h}E^{L}=-ip^j\tau^j\),
\((E^{R})^\dagger\mathsf{h}E^{R}=+ip^j\tau^j\),
\((E^{L})^\dagger\mathsf{h}E^{R}=(E^{R})^\dagger\mathsf{h}E^{L}=m\)
(checks \texttt{71-D7}, \texttt{I71-D3}), and Eq. (2.11) becomes
\[\boxed{\;\big(i\,\partial_4+\tau_{1671}^j\partial_j\big)\phi_{1671}^{L}=m_{1671}\,\phi_{1671}^{R},\qquad\big(i\,\partial_4-\tau_{1671}^j\partial_j\big)\phi_{1671}^{R}=m_{1671}\,\phi_{1671}^{L}\;}\tag{2.12}\]
with summation over the seven space indices \(j\in\{0,1,2,3,5,6,7\}\)
and the real antisymmetric 8×8 matrices \(\tau_{1671}^j\). The mass
couples the two 8-component fields; for \(m=0\) they decouple into two
Weyl equations. Each field obeys the Klein--Gordon equation, because
\((\tau^j\partial_j)^2=-\nabla^2\):
\((i\partial_4-\tau^i\partial_i)(i\partial_4+\tau^j\partial_j)=-(\partial_4^2-\nabla^2)\).
Equations (2.12) are those of section 7.11 of \texttt{7+1.nb} with the
coordinates \(x0\) and \(x4\) exchanged.

\begin{center}\rule{0.5\linewidth}{0.5pt}\end{center}

\hypertarget{part-3.-which-fermion-pairs-exist-a-consequence-of-the-two-sets-of-matrices}{%
\section{Part 3. Which fermion pairs exist: a consequence of the two
sets of
matrices}\label{part-3.-which-fermion-pairs-exist-a-consequence-of-the-two-sets-of-matrices}}

\emph{All tags are written in this Part. The field \(\psi\) is
\(\psi_{1644}\) in statements about 4+4 and \(\psi_{1671}\) in
statements about 7+1.}

The main task of the project is a superfluid built from pairs of quanta
of the 16-component field. Which pairs are possible is decided by the
matrices of Parts 1 and 2, and the result is needed in every later
document, so it is recorded here.

\hypertarget{no-lorentz-scalar-pair}{%
\subsection{3.1 No Lorentz-scalar pair}\label{no-lorentz-scalar-pair}}

A pair operator of two field operators of the same kind has the form
\(\psi^TX\psi=\sum_{a,b}\psi_aX_{ab}\psi_b\). Because the components
anticommute, only the antisymmetric part of the 16×16 matrix \(X\)
contributes. The matrices \(X=C\,\gamma^{A_1}\cdots\gamma^{A_r}\)
(\(A_1<\dots<A_r\)), with \(C=C^{-}\) or \(C^{+}\), transform under
Lorentz transformations as the components of antisymmetric tensors of
rank \(r\) (by Eqs. (1.9) and (2.6):
\(S^TC\,\gamma^{A_1}\cdots\gamma^{A_r}S=C\,S^{-1}\gamma^{A_1}\cdots\gamma^{A_r}S\)).
The program determines the symmetry of all 256 of them, for each of the
four matrices \(C\) (checks \texttt{44-K1}, \texttt{71-K1},
\texttt{I-P1}). A symmetric matrix gives a pair operator that vanishes
identically:

\begin{longtable}[]{@{}
  >{\centering\arraybackslash}p{(\columnwidth - 10\tabcolsep) * \real{0.1667}}
  >{\centering\arraybackslash}p{(\columnwidth - 10\tabcolsep) * \real{0.1667}}
  >{\centering\arraybackslash}p{(\columnwidth - 10\tabcolsep) * \real{0.1667}}
  >{\centering\arraybackslash}p{(\columnwidth - 10\tabcolsep) * \real{0.1667}}
  >{\centering\arraybackslash}p{(\columnwidth - 10\tabcolsep) * \real{0.1667}}
  >{\centering\arraybackslash}p{(\columnwidth - 10\tabcolsep) * \real{0.1667}}@{}}
\toprule\noalign{}
\begin{minipage}[b]{\linewidth}\centering
rank \(r\)
\end{minipage} & \begin{minipage}[b]{\linewidth}\centering
number of matrices
\end{minipage} & \begin{minipage}[b]{\linewidth}\centering
\(C_{1644}^{-}\,\gamma^{A_1}\cdots\gamma^{A_r}\)
\end{minipage} & \begin{minipage}[b]{\linewidth}\centering
\(C_{1644}^{+}\,\gamma^{A_1}\cdots\gamma^{A_r}\)
\end{minipage} & \begin{minipage}[b]{\linewidth}\centering
\(C_{1671}^{-}\,\gamma^{A_1}\cdots\gamma^{A_r}\)
\end{minipage} & \begin{minipage}[b]{\linewidth}\centering
\(C_{1671}^{+}\,\gamma^{A_1}\cdots\gamma^{A_r}\)
\end{minipage} \\
\midrule\noalign{}
\endhead
\bottomrule\noalign{}
\endlastfoot
0 & 1 & symmetric & symmetric & symmetric & symmetric \\
1 & 8 & \textbf{antisymmetric} & symmetric & \textbf{antisymmetric} &
symmetric \\
2 & 28 & \textbf{antisymmetric} & \textbf{antisymmetric} &
\textbf{antisymmetric} & \textbf{antisymmetric} \\
3 & 56 & symmetric & \textbf{antisymmetric} & symmetric &
\textbf{antisymmetric} \\
4 & 70 & symmetric & symmetric & symmetric & symmetric \\
5 & 56 & \textbf{antisymmetric} & symmetric & \textbf{antisymmetric} &
symmetric \\
6 & 28 & \textbf{antisymmetric} & \textbf{antisymmetric} &
\textbf{antisymmetric} & \textbf{antisymmetric} \\
7 & 8 & symmetric & \textbf{antisymmetric} & symmetric &
\textbf{antisymmetric} \\
8 & 1 & symmetric & symmetric & symmetric & symmetric \\
\end{longtable}

The pattern is the same for both signatures; 120 matrices are
antisymmetric, as many as there are independent products
\(\psi_a\psi_b\) with \(a<b\).

\textbf{Theorem 4.} In both spacetimes the Lorentz-scalar pair
operators, \(\psi^TC^{-}\psi\) (rank 0) and
\(\psi^TC^{-}\omega\,\psi=\pm\psi^TC^{+}\psi\) (rank 8), vanish
identically (checks \texttt{44-K2}, \texttt{71-K2}). A single
16-component field has no Lorentz-invariant fermion pair. The pair
operators that exist are, in terms of \(C^{-}\), a vector (\(r=1\), 8
components), a 2-form (28), a 5-form (56) and a 6-form (28).

\hypertarget{the-pairs-of-two-particles-at-rest}{%
\subsection{3.2 The pairs of two particles at
rest}\label{the-pairs-of-two-particles-at-rest}}

A superfluid is formed in a medium, and the medium singles out a rest
frame. The relevant question is then: which pairs of two identical
particles at rest, in a state of zero orbital angular momentum (an
\emph{s wave}), are allowed, and how do they transform under the
rotations that the medium leaves intact? One particle at rest has eight
internal states; the Pauli principle admits the \(8\cdot7/2=28\)
antisymmetric combinations, which are described by the antisymmetric 8×8
\emph{pair matrices}.

\textbf{7+1} (checks \texttt{71-Q1} to \texttt{71-Q10}, \texttt{I71-P2}
to \texttt{I71-P4}). On particles at rest the rotation generators act as
\(\frac12\tau_{1671}^i\tau_{1671}^j\), and the 28 antisymmetric 8×8
matrices are spanned by the seven \(\tau_{1671}^j\) and the twenty-one
\(\tau_{1671}^i\tau_{1671}^j\) (\(i<j\); space indices).

\begin{itemize}
\tightlist
\item
  The only matrix that commutes with all rotations is the unit matrix,
  which is symmetric. \textbf{There is no s-wave pair that is a singlet
  under the rotations \(SO(7)\) of the seven space dimensions} (check
  \texttt{71-Q6}).
\item
  Under \(SO(7)\) the 28 pairs form a vector (7, spanned by the
  \(\tau_{1671}^j\)) and an adjoint (21, spanned by the
  \(\tau_{1671}^i\tau_{1671}^j\)) (check \texttt{71-Q5}).
\item
  Split the seven space coordinates into the ordinary ones, \(x1,x2,x3\)
  (rotations \(SO(3)\), generating ordinary spin), and the four others,
  \(x0,x5,x6,x7\) (rotations \(SO(4)\), locally
  \(SU(2)_A\times SU(2)_B\)). \textbf{Exactly ten pairs have ordinary
  spin 0}: the four matrices
  \(\tau_{1671}^0,\tau_{1671}^5,\tau_{1671}^6,\tau_{1671}^7\), which
  form a 4-vector of \(SO(4)\), i.e.~the representation \((2,2)\) of
  \(SU(2)_A\times SU(2)_B\); and the six matrices
  \(\tau_{1671}^a\tau_{1671}^b\) (\(a<b\) in \(\{0,5,6,7\}\)), which
  form \((3,1)+(1,3)\) (checks \texttt{71-Q2}, \texttt{71-Q7}).
\item
  \textbf{Exactly one pair is a singlet both of the ordinary rotations
  and of the rotations of \(x5,x6,x7\)}: the pair matrix
  \(\tau_{1671}^0\) (checks \texttt{71-Q10}, \texttt{I71-P4}).
\item
  The corresponding field operators are
  \(\psi^TC_{1671}^{-}\gamma_{1671}^j\psi\) for the vector and
  \(\psi^TC_{1671}^{+}\gamma_{1671}^i\gamma_{1671}^j\psi\) for the
  adjoint (check \texttt{71-Q9}).
\end{itemize}

\textbf{4+4, fields independent of \(x5,x6,x7\)} (checks \texttt{44-Q1}
to \texttt{44-Q4}, \texttt{I44-P2}, \texttt{I44-P3}). A particle at rest
satisfies \(\gamma_{1644}^4u=iu\). The symmetry is \(SO(4)\) (rotations
of the space \(x0,\dots,x3\); locally \(SU(2)_X\times SU(2)_Y\)) times
\(SO(3)_T\) (rotations of the additional times) times the sign of
\(W_{1644}\). The group \(SU(2)_X\times SU(2)_Y\) must not be confused
with the group \(SU(2)_A\times SU(2)_B\) of 7+1 above: it rotates the
hidden coordinate into ordinary space, and it is not an internal
symmetry.

\begin{itemize}
\tightlist
\item
  The pairs that are singlets under all rotations of the four space
  dimensions are six in number: two triplets of \(SO(3)_T\). None of
  them is a singlet of \(SO(3)_T\).
\item
  Ten pairs have ordinary spin 0, as in 7+1.
\item
  \textbf{Exactly one pair is a singlet both of the ordinary rotations
  and of the rotations of the additional times.} It is the two-particle
  content of \(\psi^TC_{1644}^{-}\gamma_{1644}^0\psi\) (check
  \texttt{44-Q4}).
\end{itemize}

The complete lists are in \texttt{data/pair\_classification.csv}.

\hypertarget{the-common-result-the-hidden-axis-pair}{%
\subsection{3.3 The common result: the hidden-axis
pair}\label{the-common-result-the-hidden-axis-pair}}

In both signatures there is a \emph{vector pair operator},
\[\Delta_{1644}^{A}\equiv\psi_{1644}^{T}\,C_{1644}^{-}\,\gamma_{1644}^{A}\,\psi_{1644},\qquad
\Delta_{1671}^{A}\equiv\psi_{1671}^{T}\,C_{1671}^{-}\,\gamma_{1671}^{A}\,\psi_{1671},\qquad A=0,\dots,7,\]
and in both signatures its component along the hidden coordinate \(x0\),
\[\boxed{\;\Delta_{1644}^{0}=\psi_{1644}^{T}\,C_{1644}^{-}\,\gamma_{1644}^{0}\,\psi_{1644},\qquad
\Delta_{1671}^{0}=\psi_{1671}^{T}\,C_{1671}^{-}\,\gamma_{1671}^{0}\,\psi_{1671}\;}\tag{3.1}\]
has three properties:

\begin{enumerate}
\def\labelenumi{\arabic{enumi}.}
\tightlist
\item
  its two-particle content has ordinary spin 0;
\item
  it is the only s-wave pair that is invariant under the ordinary
  rotations and under the rotations of \(x5,x6,x7\) (Sec. 3.2);
\item
  it is invariant under all 21 generators \(\Sigma^{AB}\) with
  \(A,B\in\{1,\dots,7\}\), and only the 7 generators \(\Sigma^{0B}\)
  that involve the hidden axis change it (checks \texttt{44-K3},
  \texttt{71-K3}, \texttt{I-P5}). In other words \(\Delta^{0}\) is a
  \emph{scalar} under the Lorentz group of the seven coordinates other
  than \(x0\) --- in particular under the Lorentz group of ordinary
  spacetime \((x1,x2,x3;x4)\) --- although it is the component of a
  vector of the eight-dimensional spacetime.
\end{enumerate}

Property 3 follows from Eq. (1.6): a transformation \(\psi\to S\psi\)
generated by \(\Sigma^{AB}\) changes \(\Delta^{0}\) by
\(-\psi^TC^{-}[\Sigma^{AB},\gamma^0]\psi\), which vanishes unless \(A\)
or \(B\) is 0.

The operator (3.1) carries fermion number 2 and is the natural order
parameter of a spin-0 (singlet) superfluid of the field. Its dynamics is
not the subject of this document; it is developed in the supplements of
the two manuscripts.

\begin{center}\rule{0.5\linewidth}{0.5pt}\end{center}

\hypertarget{appendix-a.-complete-output-of-the-programs}{%
\section{Appendix A. Complete output of the
programs}\label{appendix-a.-complete-output-of-the-programs}}

\hypertarget{a.1-the-rust-program}{%
\subsection{A.1 The Rust program}\label{a.1-the-rust-program}}

The following is the file \texttt{data/matrix\_checks.txt}, written by
\texttt{spacetimes\ matrices}. ``exact'' means that the statement is an
identity between integer or Gaussian-integer matrices, verified in exact
arithmetic; ``numerical'' means that it involves real numbers and is
verified to the accuracy stated in the line. Ranks are computed modulo
the prime \(2^{61}-1\); the documentation of the function \texttt{rank}
in \texttt{code/spacetimes-rs/src/imat.rs} explains why a full rank
found in this way proves linear independence over the rational numbers.

\begin{Shaded}
\begin{Highlighting}[]
\NormalTok{PROJECT spacetimes: checks of the real 16x16 Dirac matrices (4+4: gamma1644[A] = T16[A] of 4+4.nb; 7+1: gamma1671[A] = Gamma17up[swap04(A)] of 7+1.nb, swap04 = (0 4)); identifiers 44{-}.. refer to the 4+4 set, 71{-}.. to the 7+1 set}
\NormalTok{107 of 107 checks passed}

\NormalTok{[PASS] 44{-}R1  (exact) every entry of gamma1644[A], A = 0..7, is one of the integers {-}1, 0, +1; hence the eight matrices are real}
\NormalTok{[PASS] 44{-}R2  (exact) every gamma1644[A] is a signed permutation matrix (one entry +1 or {-}1 in each row and in each column)}
\NormalTok{[PASS] 44{-}C1  (exact) gamma1644[A].gamma1644[B] + gamma1644[B].gamma1644[A] = 2 eta1644\^{}(AB) 1 for all 64 pairs (A, B), with eta1644 = diag[1, 1, 1, 1, {-}1, {-}1, {-}1, {-}1]}
\NormalTok{[PASS] 44{-}C2  (exact) transpose(gamma1644[A]) = eta1644\^{}(AA) gamma1644[A]: matrices with square +1 are symmetric, matrices with square {-}1 are antisymmetric; hence transpose(gamma1644[A]) = inverse(gamma1644[A])}
\NormalTok{[PASS] 44{-}C3  (exact) the 256 ordered products gamma1644[A1]...gamma1644[Ak] (A1 \textless{} ... \textless{} Ak, k = 0..8) are linearly independent: they span all real 16x16 matrices, so the representation is irreducible over the real numbers}
\NormalTok{[PASS] 44{-}C4  (exact) every one of the 255 non{-}empty products has trace 0, and the square of every product is +1 or {-}1}
\NormalTok{[PASS] 44{-}L1  (exact) Sigma1644\^{}(AB) := (1/4)[gamma1644[A], gamma1644[B]] equals (1/2) gamma1644[A].gamma1644[B] for A != B, and Sigma1644\^{}(AB) = {-}Sigma1644\^{}(BA)}
\NormalTok{[PASS] 44{-}L2  (exact) the 28 matrices Sigma1644\^{}(AB), A \textless{} B, are linearly independent}
\NormalTok{[PASS] 44{-}L3  (exact) [Sigma1644\^{}(AB), gamma1644[C]] = eta1644\^{}(BC) gamma1644[A] {-} eta1644\^{}(AC) gamma1644[B] for all A, B, C: the Sigma1644\^{}(AB) rotate the eight matrices as the components of a vector}
\NormalTok{[PASS] 44{-}L4  (exact) [Sigma1644\^{}(AB), Sigma1644\^{}(CD)] = eta1644\^{}(BC) Sigma1644\^{}(AD) {-} eta1644\^{}(AC) Sigma1644\^{}(BD) {-} eta1644\^{}(BD) Sigma1644\^{}(AC) + eta1644\^{}(AD) Sigma1644\^{}(BC) for all 4096 index combinations: the commutation relations of the Lie algebra so(p,q) of the metric eta1644}
\NormalTok{[PASS] 44{-}L5  (exact) the structure constants are integers and the Killing form is diagonal in the basis Sigma1644\^{}(AB): B(Sigma1644\^{}(AB), Sigma1644\^{}(AB)) = {-}12 eta1644\^{}(AA) eta1644\^{}(BB); it has 16 positive and 12 negative eigenvalues (no zero eigenvalue: the algebra is semisimple)}
\NormalTok{[PASS] 44{-}L6  (exact) 12 generators are antisymmetric matrices (compact: they generate rotations) and 16 are symmetric (non{-}compact: they generate boosts); these are the numbers of the real form so(4,4), which has 12 compact and 16 non{-}compact generators}
\NormalTok{[PASS] 44{-}G1  (exact) each gamma1644[A] is invertible with inverse eta1644\^{}(AA) gamma1644[A], and {-}gamma1644[A].gamma1644[B].inverse(gamma1644[A]) = gamma1644[B] for B != A, = {-}gamma1644[A] for B = A: gamma1644[A] acts on the vector index as the reflection that reverses the axis xA}
\NormalTok{[PASS] 44{-}G2  (numerical) exp(theta Sigma1644\^{}(AB)) = cos(theta/2) 1 + sin(theta/2) gamma1644[A].gamma1644[B] when (gamma1644[A].gamma1644[B])\^{}2 = {-}1 (rotation), and = cosh(theta/2) 1 + sinh(theta/2) gamma1644[A].gamma1644[B] when (gamma1644[A].gamma1644[B])\^{}2 = +1 (boost), for the 28 planes and random theta (largest deviation 2.3e{-}14)}
\NormalTok{[PASS] 44{-}G3  (numerical) for 20 random elements S = exp(sum\_(A\textless{}B) theta\_AB Sigma1644\^{}(AB)) (theta\_AB random real numbers):  S.gamma1644[C].inverse(S) = sum\_D Lambda\^{}C\_D gamma1644[D] with a real 8x8 matrix Lambda that preserves eta1644 (Lambda.eta1644.transpose(Lambda) = eta1644) and has determinant +1 (largest deviations 2.0e{-}13, 2.2e{-}15, 2.7e{-}15)}
\NormalTok{[PASS] 44{-}N1  (exact) the eight matrices rebuilt from the definitions equal, entry by entry, the eight matrices T16[0..7] printed by Mathematica in 4+4.nb (cell 314, Out[211])}
\NormalTok{[PASS] 44{-}N2  (exact) the rebuilt volume element equals T16[8] printed in 4+4.nb (cell 136, Out[99])}
\NormalTok{[PASS] 44{-}N3  (exact) the metric printed in 4+4.nb (cell 13, Out[9]) is eta4488 = diag(1,1,1,1,{-}1,{-}1,{-}1,{-}1) (PROJECT name: eta1644)}
\NormalTok{[PASS] 44{-}B1  (exact) the 8x8 matrix sigma (= sigma88) of 4+4.nb, PROJECT name sigma88: sigma88 = tau[1].tau[2].tau[3] = tau[4].tau[5].tau[6].tau[7], sigma88.sigma88 = 1, tau[1].tau[2]...tau[7] = 1, and tau[7] = diag({-}1\_4, +1\_4) (4+4.nb Out[72], Out[89], Out[87], Out[63])}
\NormalTok{[PASS] 44{-}B2  (exact) tau[A].taubar[B] + tau[B].taubar[A] = 2 eta1644\^{}(AB) 1 and taubar[A].tau[B] + taubar[B].tau[A] = 2 eta1644\^{}(AB) 1 (8x8), A, B = 0..7 (4+4.nb Out[77])}
\NormalTok{[PASS] 44{-}B3  (exact) sigma88.taubar[A] = transpose(sigma88.tau[A]) for A = 0..7 (4+4.nb Out[76])}
\NormalTok{[PASS] 44{-}B4  (exact) T16[A] = [[0, taubar[A]], [tau[A], 0]] in 8x8 blocks, with tau[0] = taubar[0] = 1.  PROJECT names of these objects of 4+4.nb: gamma1644[A] := T16[A], tau1644[A] := tau[A], taubar1644[A] := taubar[A]}
\NormalTok{[PASS] 44{-}V1  (exact) omega1644 := gamma1644[0]...gamma1644[7] (= T16[8]) = diag({-}1\_8, +1\_8): real, diagonal, symmetric, with square +1}
\NormalTok{[PASS] 44{-}V2  (exact) omega1644 anticommutes with every gamma1644[A] and commutes with every Sigma1644\^{}(AB)}
\NormalTok{[PASS] 44{-}P1  (exact) PL1644 := (1 {-} omega1644)/2 = diag(1\_8, 0\_8) and PR1644 := (1 + omega1644)/2 = diag(0\_8, 1\_8) are real matrices with integer entries}
\NormalTok{[PASS] 44{-}P2  (exact) PL1644 + PR1644 = 1, PL1644.PL1644 = PL1644, PR1644.PR1644 = PR1644, PL1644.PR1644 = PR1644.PL1644 = 0, trace PL1644 = trace PR1644 = 8 (4+4.nb Out[114])}
\NormalTok{[PASS] 44{-}P3  (exact) the rebuilt projectors PL1644 and PR1644 equal the matrices P\_L and P\_R printed in 4+4.nb (cells 163 and 165, Out[112] and Out[113])}
\NormalTok{[PASS] 44{-}P4  (exact) gamma1644[A].PL1644 = PR1644.gamma1644[A] and gamma1644[A].PR1644 = PL1644.gamma1644[A]: every Dirac matrix maps left{-}handed spinors to right{-}handed ones and conversely}
\NormalTok{[PASS] 44{-}P5  (exact) PL1644 and PR1644 commute with every Sigma1644\^{}(AB): chirality is invariant under Spin(4,4)}
\NormalTok{[PASS] 44{-}S1  (exact) Cm1644 := gamma1644[0].gamma1644[1].gamma1644[2].gamma1644[3] (the matrix sigma16 of 4+4.nb) = diag({-}sigma88, sigma88), symmetric, Cm1644.Cm1644 = 1 (4+4.nb Out[106])}
\NormalTok{[PASS] 44{-}S2  (exact) Cp1644 := gamma1644[4].gamma1644[5].gamma1644[6].gamma1644[7] (the matrix S44 of 7+1.nb) = diag(sigma88, sigma88), symmetric, Cp1644.Cp1644 = 1, and omega1644 = Cm1644.Cp1644}
\NormalTok{[PASS] 44{-}S3  (exact) Cm1644.gamma1644[A].Cm1644 = {-}transpose(gamma1644[A]) and Cp1644.gamma1644[A].Cp1644 = +transpose(gamma1644[A]) for A = 0..7}
\NormalTok{[PASS] 44{-}S4  (exact) Cm1644.gamma1644[A] is antisymmetric for A = 0..7 (4+4.nb Out[105]) and Cp1644.gamma1644[A] is symmetric}
\NormalTok{[PASS] 44{-}S5  (exact) gamma1644[4] (the matrix JT4 = T16[4] of 4+4.nb) is a complex structure: gamma1644[4].gamma1644[4] = {-}1, and transpose(gamma1644[4]).Cm1644.gamma1644[4] = Cm1644}
\NormalTok{[PASS] 71{-}R1  (exact) every entry of gamma1671[A], A = 0..7, is one of the integers {-}1, 0, +1; hence the eight matrices are real}
\NormalTok{[PASS] 71{-}R2  (exact) every gamma1671[A] is a signed permutation matrix (one entry +1 or {-}1 in each row and in each column)}
\NormalTok{[PASS] 71{-}C1  (exact) gamma1671[A].gamma1671[B] + gamma1671[B].gamma1671[A] = 2 eta1671\^{}(AB) 1 for all 64 pairs (A, B), with eta1671 = diag[{-}1, {-}1, {-}1, {-}1, 1, {-}1, {-}1, {-}1]}
\NormalTok{[PASS] 71{-}C2  (exact) transpose(gamma1671[A]) = eta1671\^{}(AA) gamma1671[A]: matrices with square +1 are symmetric, matrices with square {-}1 are antisymmetric; hence transpose(gamma1671[A]) = inverse(gamma1671[A])}
\NormalTok{[PASS] 71{-}C3  (exact) the 256 ordered products gamma1671[A1]...gamma1671[Ak] (A1 \textless{} ... \textless{} Ak, k = 0..8) are linearly independent: they span all real 16x16 matrices, so the representation is irreducible over the real numbers}
\NormalTok{[PASS] 71{-}C4  (exact) every one of the 255 non{-}empty products has trace 0, and the square of every product is +1 or {-}1}
\NormalTok{[PASS] 71{-}L1  (exact) Sigma1671\^{}(AB) := (1/4)[gamma1671[A], gamma1671[B]] equals (1/2) gamma1671[A].gamma1671[B] for A != B, and Sigma1671\^{}(AB) = {-}Sigma1671\^{}(BA)}
\NormalTok{[PASS] 71{-}L2  (exact) the 28 matrices Sigma1671\^{}(AB), A \textless{} B, are linearly independent}
\NormalTok{[PASS] 71{-}L3  (exact) [Sigma1671\^{}(AB), gamma1671[C]] = eta1671\^{}(BC) gamma1671[A] {-} eta1671\^{}(AC) gamma1671[B] for all A, B, C: the Sigma1671\^{}(AB) rotate the eight matrices as the components of a vector}
\NormalTok{[PASS] 71{-}L4  (exact) [Sigma1671\^{}(AB), Sigma1671\^{}(CD)] = eta1671\^{}(BC) Sigma1671\^{}(AD) {-} eta1671\^{}(AC) Sigma1671\^{}(BD) {-} eta1671\^{}(BD) Sigma1671\^{}(AC) + eta1671\^{}(AD) Sigma1671\^{}(BC) for all 4096 index combinations: the commutation relations of the Lie algebra so(p,q) of the metric eta1671}
\NormalTok{[PASS] 71{-}L5  (exact) the structure constants are integers and the Killing form is diagonal in the basis Sigma1671\^{}(AB): B(Sigma1671\^{}(AB), Sigma1671\^{}(AB)) = {-}12 eta1671\^{}(AA) eta1671\^{}(BB); it has 7 positive and 21 negative eigenvalues (no zero eigenvalue: the algebra is semisimple)}
\NormalTok{[PASS] 71{-}L6  (exact) 21 generators are antisymmetric matrices (compact: they generate rotations) and 7 are symmetric (non{-}compact: they generate boosts); these are the numbers of the real form so(7,1), which has 21 compact and 7 non{-}compact generators}
\NormalTok{[PASS] 71{-}G1  (exact) each gamma1671[A] is invertible with inverse eta1671\^{}(AA) gamma1671[A], and {-}gamma1671[A].gamma1671[B].inverse(gamma1671[A]) = gamma1671[B] for B != A, = {-}gamma1671[A] for B = A: gamma1671[A] acts on the vector index as the reflection that reverses the axis xA}
\NormalTok{[PASS] 71{-}G2  (numerical) exp(theta Sigma1671\^{}(AB)) = cos(theta/2) 1 + sin(theta/2) gamma1671[A].gamma1671[B] when (gamma1671[A].gamma1671[B])\^{}2 = {-}1 (rotation), and = cosh(theta/2) 1 + sinh(theta/2) gamma1671[A].gamma1671[B] when (gamma1671[A].gamma1671[B])\^{}2 = +1 (boost), for the 28 planes and random theta (largest deviation 7.1e{-}15)}
\NormalTok{[PASS] 71{-}G3  (numerical) for 20 random elements S = exp(sum\_(A\textless{}B) theta\_AB Sigma1671\^{}(AB)) (theta\_AB random real numbers):  S.gamma1671[C].inverse(S) = sum\_D Lambda\^{}C\_D gamma1671[D] with a real 8x8 matrix Lambda that preserves eta1671 (Lambda.eta1671.transpose(Lambda) = eta1671) and has determinant +1 (largest deviations 5.0e{-}14, 2.0e{-}15, 2.0e{-}15)}
\NormalTok{[PASS] 71{-}N1  (exact) the eight matrices rebuilt from the definitions equal, entry by entry, the eight matrices Gamma17[0..7] printed by Mathematica in 7+1.nb (cell 204, Out[178])}
\NormalTok{[PASS] 71{-}N2  (exact) PROJECT matrices: gamma1671[A] = Gamma17up[swap04(A)] with swap04 = (0 4) and Gamma17up[A] = eta17\^{}(AA) Gamma17[A] (7+1.nb section 5.3); explicitly gamma1671[4] = Gamma17[0] (time{-}like), gamma1671[0] = {-}Gamma17[4], gamma1671[j] = {-}Gamma17[j] for j = 1, 2, 3, 5, 6, 7.  (gamma1671 is NOT the complex set Gamma71 of 7+1.nb section 4.3, and eta1671 is NOT the matrix eta71 of 7+1.nb)}
\NormalTok{[PASS] 71{-}N3  (exact) in the labelling of 7+1.nb: Gamma17[A].Gamma17[B] + Gamma17[B].Gamma17[A] = 2 eta17\^{}(AB) 1 with eta17 = diag(+1,{-}1,{-}1,{-}1,{-}1,{-}1,{-}1,{-}1) (7+1.nb Out[173])}
\NormalTok{[PASS] 71{-}B1  (exact) tauO[A] = tau[A] (A = 1, 2, 3), tauO[A] = tau[A].sigma (A = 4..7): real, antisymmetric, tauO[A].tauO[A] = {-}1, mutually anticommuting, tauO[1]...tauO[7] = 1 (7+1.nb section 3.1)}
\NormalTok{[PASS] 71{-}B2  (exact) block form: Gamma17[0] = [[0, 1], [1, 0]] and Gamma17[i] = [[tauO[i], 0], [0, tauObar[i]]] with tauObar[i] = transpose(tauO[i]) = {-}tauO[i], i = 1..7 (7+1.nb Out[175])}
\NormalTok{[PASS] 71{-}B3  (exact) built from the 4+4 matrices only: Gamma17[0] = T16[0], Gamma17[i] = T16[0].T16[i] (i = 1, 2, 3), Gamma17[k] = T16[0].T16[k].S44 (k = 4..7) (7+1.nb Out[176])}
\NormalTok{[PASS] 71{-}B4  (exact) in the labels of this project: gamma1671[4] = gamma1644[0]; gamma1671[j] = {-}gamma1644[0].gamma1644[j] (j = 1, 2, 3); gamma1671[0] = +gamma1644[0].gamma1644[5].gamma1644[6].gamma1644[7]; gamma1671[5] = {-}gamma1644[0].gamma1644[4].gamma1644[6].gamma1644[7]; gamma1671[6] = +gamma1644[0].gamma1644[4].gamma1644[5].gamma1644[7]; gamma1671[7] = {-}gamma1644[0].gamma1644[4].gamma1644[5].gamma1644[6].  The time{-}like matrix of 7+1 IS the matrix of the hidden axis of 4+4; the other seven are products of two or of four of the 4+4 matrices and therefore not linear combinations of them}
\NormalTok{[PASS] 71{-}V1  (exact) omega1671 := gamma1671[0]...gamma1671[7] = Gamma17[8] = [[0, {-}1], [1, 0]] in 8x8 blocks (the omega17 of 7+1.nb: the seven sign changes of Gamma17up and the odd permutation swap04 compensate): real, antisymmetric, with square {-}1}
\NormalTok{[PASS] 71{-}V2  (exact) omega1671 anticommutes with every gamma1671[A] and commutes with every Sigma1671\^{}(AB)}
\NormalTok{[PASS] 71{-}V3  (exact) the linear system X.gamma1671[A] + gamma1671[A].X = 0 (A = 0..7) for a real 16x16 matrix X has rank 255, and omega1671 solves it (71{-}V2): the real matrices that anticommute with all eight gamma1671[A] are exactly the multiples of omega1671; since omega1671.omega1671 = {-}1, no real matrix with square +1 anticommutes with all of them: there is NO real chirality matrix in 7+1}
\NormalTok{[PASS] 71{-}P1  (exact) chi1671 := i omega1671 = i Gamma17[8] (= chi17 of 7+1.nb) is Hermitian, chi1671.chi1671 = 1, anticommutes with every gamma1671[A]}
\NormalTok{[PASS] 71{-}P2  (exact) PL1671 := (1 {-} chi1671)/2 and PR1671 := (1 + chi1671)/2: PL1671 + PR1671 = 1, PL1671.PL1671 = PL1671, PR1671.PR1671 = PR1671, PL1671.PR1671 = 0, Hermitian, trace 8 each; they are complex (not real) matrices}
\NormalTok{[PASS] 71{-}P3  (exact) gamma1671[A].PL1671 = PR1671.gamma1671[A] for A = 0..7, and PL1671, PR1671 commute with every Sigma1671\^{}(AB)}
\NormalTok{[PASS] 71{-}P4  (exact) in 8x8 blocks chi1671 = [[0, {-}i], [i, 0]], 2 PL1671 = [[1, i], [{-}i, 1]] and 2 PR1671 = [[1, {-}i], [i, 1]]; these are chi17, PL17 and PR17 of 7+1.nb (sections 4.6 and 7.10)}
\NormalTok{[PASS] 71{-}P5  (exact) with s := 1 {-} i Gamma17[0] (s.s\^{}dagger = 2, so Uc1671 := s/sqrt(2) is unitary): s.chi1671 = {-}T16[8].s, i.e. Uc1671.chi1671.inverse(Uc1671) = {-}omega1644.  Hence Uc1671.PL1671.inverse(Uc1671) = PR1644 and Uc1671.PR1671.inverse(Uc1671) = PL1644: the complex chiral projectors of 7+1 are unitarily equivalent to the real projectors of 4+4.nb, with the names L and R exchanged.  (7+1.nb defines chi17 = +i Gamma17[8] for its real matrices and chi71 = {-}i Gamma71[8] for its complex matrices Gamma71; the PROJECT matrix chi1671 equals chi17, NOT chi71)}
\NormalTok{[PASS] 71{-}S1  (exact) Dirac adjoint: transpose(gamma1671[A]) = gamma1671[4].gamma1671[A].gamma1671[4] for A = 0..7 (gamma1671[4] is the time{-}like matrix; it is symmetric and squares to +1)}
\NormalTok{[PASS] 44{-}K1  (exact) 4+4: the matrices Cm1644.gamma1644[A1]...gamma1644[Ar] (A1 \textless{} ... \textless{} Ar) of rank r = 0..8 are all symmetric (+) or all antisymmetric ({-}): signs [1, {-}1, {-}1, 1, 1, {-}1, {-}1, 1, 1]; the matrices Cp1644.gamma1644[A1]...gamma1644[Ar]: signs [1, 1, {-}1, {-}1, 1, 1, {-}1, {-}1, 1]; in both cases exactly 120 of the 256 matrices are antisymmetric}
\NormalTok{[PASS] 44{-}K2  (exact) 4+4: the two Lorentz{-}scalar pair matrices Cm1644 (rank 0) and Cm1644.omega1644 (rank 8; equal to Cp1644 up to sign) are SYMMETRIC: the Lorentz{-}scalar bilinears dirac1644\^{}T Cm1644 dirac1644 and dirac1644\^{}T Cp1644 dirac1644 of one anticommuting 16{-}component field vanish identically}
\NormalTok{[PASS] 71{-}K0  (exact) 7+1: Cp1671 := gamma1671[4] satisfies Cp1671.gamma1671[A] = +transpose(gamma1671[A]).Cp1671 and Cm1671 := gamma1671[4].omega1671 satisfies Cm1671.gamma1671[A] = {-}transpose(gamma1671[A]).Cm1671; both are symmetric matrices}
\NormalTok{[PASS] 71{-}K1  (exact) 7+1: signs (+ symmetric, {-} antisymmetric) of the matrices Cm1671.gamma1671[A1]...gamma1671[Ar] (A1 \textless{} ... \textless{} Ar) by rank r = 0..8: [1, {-}1, {-}1, 1, 1, {-}1, {-}1, 1, 1]; of Cp1671.gamma1671[A1]...gamma1671[Ar]: [1, 1, {-}1, {-}1, 1, 1, {-}1, {-}1, 1]; exactly 120 of the 256 matrices are antisymmetric in both cases}
\NormalTok{[PASS] 71{-}K2  (exact) 7+1: the two Lorentz{-}scalar pair matrices Cm1671 (rank 0) and Cm1671.omega1671 = {-}Cp1671 (rank 8) are SYMMETRIC: the Lorentz{-}scalar bilinears dirac1671\^{}T Cm1671 dirac1671 and dirac1671\^{}T Cp1671 dirac1671 of one anticommuting 16{-}component field vanish identically}
\NormalTok{[PASS] 44{-}K3  (exact) 4+4: for all A != B in 1..7, transpose(Sigma1644\^{}(AB)).Cm1644.gamma1644[0] + Cm1644.gamma1644[0].Sigma1644\^{}(AB) = 0, while the same expression with A = 0 is not zero: the bilinear DeltaV1644\^{}0 = dirac1644\^{}T Cm1644 gamma1644\^{}0 dirac1644 is INVARIANT under the 21 generators of the pseudo{-}rotations of the seven coordinates x1..x7 (the group Spin(3,4), which contains the Lorentz group of x1 x2 x3 x4), and only the 7 generators that involve the hidden axis x0 change it}
\NormalTok{[PASS] 71{-}K3  (exact) 7+1: for all A != B in 1..7, transpose(Sigma1671\^{}(AB)).Cm1671.gamma1671[0] + Cm1671.gamma1671[0].Sigma1671\^{}(AB) = 0, while the same expression with A = 0 is not zero: the bilinear DeltaV1671\^{}0 = dirac1671\^{}T Cm1671 gamma1671\^{}0 dirac1671 is INVARIANT under the 21 generators of the Lorentz group of the seven coordinates x1..x7 (the group Spin(6,1), which contains the Lorentz group of x1 x2 x3 x4), and only the 7 generators that involve the hidden axis x0 change it}
\NormalTok{[PASS] 44{-}D1  (exact) (gamma1644\^{}A k\_A)\^{}2 = (eta1644\^{}(AB) k\_A k\_B) 1 for 40 random integer vectors k: the Dirac operator gamma1644\^{}A d\_A squares to the wave operator eta1644\^{}(AB) d\_A d\_B = d0\^{}2 + d1\^{}2 + d2\^{}2 + d3\^{}2 {-} d4\^{}2 {-} d5\^{}2 {-} d6\^{}2 {-} d7\^{}2}
\NormalTok{[PASS] 44{-}D2  (exact) (i gamma1644.k + m1644)(i gamma1644.k {-} m1644) = {-}(k.k + m1644\^{}2) 1 for 40 random integer (k, m1644): every component of a solution of (gamma1644\^{}A d\_A {-} m1644) dirac1644 = 0 obeys the Klein{-}Gordon equation (eta1644\^{}(AB) d\_A d\_B {-} m1644\^{}2) dirac1644 = 0}
\NormalTok{[PASS] 44{-}D3  (exact) worked example k = (0,2,2,4; 5,0,0,0), m1644 = 1 (on the mass shell: k4\^{}2+k5\^{}2+k6\^{}2+k7\^{}2 = k0\^{}2+k1\^{}2+k2\^{}2+k3\^{}2 + m1644\^{}2): Pi1644 := (i gamma1644.k + m1644)/(2 m1644) satisfies Pi1644.Pi1644 = Pi1644, trace Pi1644 = 8 and (i gamma1644.k {-} m1644).Pi1644 = 0, so the plane{-}wave solutions u exp(i k.x) form an 8{-}dimensional space, u = Pi1644 w}
\NormalTok{[PASS] 44{-}D4  (exact) PL1644.gamma1644[A].PR1644 = [[0, taubar1644[A]], [0, 0]], PR1644.gamma1644[A].PL1644 = [[0, 0], [tau1644[A], 0]], PL1644.gamma1644[A].PL1644 = PR1644.gamma1644[A].PR1644 = 0: with dirac1644 = (psiL1644, psiR1644) the 16{-}component equation is the pair  taubar1644\^{}A d\_A psiR1644 = m1644 psiL1644,  tau1644\^{}A d\_A psiL1644 = m1644 psiR1644  of 8{-}component equations}
\NormalTok{[PASS] 44{-}D5  (exact) (tau1644\^{}A k\_A)(taubar1644\^{}B k\_B) = (taubar1644\^{}A k\_A)(tau1644\^{}B k\_B) = (eta1644\^{}(AB) k\_A k\_B) 1\_8 for 40 random integer k: each 8{-}component field obeys the Klein{-}Gordon equation}
\NormalTok{[PASS] 44{-}D6  (exact) alpha1644\^{}a := gamma1644[4].gamma1644[a] (a = 0..3) are real symmetric with square +1; gamma1644[4].gamma1644[t] (t = 5, 6, 7) are real antisymmetric with square {-}1; gamma1644[4] is real antisymmetric with square {-}1}
\NormalTok{[PASS] 44{-}D7  (exact) x4{-}evolution d\_4 dirac1644 = K1644 dirac1644 with K1644 = {-}m1644 gamma1644\^{}4 + i k\_a gamma1644\^{}4 gamma1644\^{}a + i kappa\_t gamma1644\^{}4 gamma1644\^{}t (a = 0..3 space, t = 5, 6, 7 extra times): K1644.K1644 = {-}(m1644\^{}2 + |k|\^{}2 {-} |kappa|\^{}2) 1 for 40 random integer (k, kappa, m1644): the frequencies in x4 are +{-}sqrt(m1644\^{}2 + |k|\^{}2 {-} |kappa|\^{}2)}
\NormalTok{[PASS] 44{-}D8  (exact) K1644 + K1644\^{}dagger = 2 i kappa\_t gamma1644\^{}4 gamma1644\^{}t: K1644 is anti{-}Hermitian (so that h1644 := i K1644 is Hermitian and the x4{-}evolution preserves dirac1644\^{}dagger dirac1644) if and only if kappa = 0, i.e. if the field does not depend on the extra times x5, x6, x7}
\NormalTok{[PASS] 44{-}D9  (exact) example of the instability: m1644 = 3, k = 0, kappa = (5, 0, 0) gives K1644.K1644 = +16, (K1644 {-} 4)(K1644 + 4) = 0 with trace K1644 = 0: K1644 has the real eigenvalues +4 and {-}4, eight times each, so eight solutions grow like exp(+4 x4)}
\NormalTok{[PASS] 44{-}D10 (exact) W1644 := i Cm1644.gamma1644[4] is Hermitian with W1644.W1644 = 1 and trace 0 (eigenvalues +1 and {-}1, eight times each); W1644 commutes with gamma1644[4] and with gamma1644[4].gamma1644[a] (a = 0..3), hence with K1644 at kappa = 0, and anticommutes with gamma1644[4].gamma1644[t] (t = 5, 6, 7)}
\NormalTok{[PASS] 44{-}D11 (exact) the density of the conserved current of the Spin(4,4){-}invariant Lagrangian, j1644\^{}4 = dirac1644\^{}dagger (i Cm1644.gamma1644[4]) dirac1644 = dirac1644\^{}dagger W1644 dirac1644, is NOT positive: W1644 has eight eigenvalues {-}1; the positive density dirac1644\^{}dagger dirac1644 is conserved only for kappa = 0 (44{-}D8)}
\NormalTok{[PASS] 71{-}D1  (exact) (gamma1671\^{}A k\_A)\^{}2 = (eta1671\^{}(AB) k\_A k\_B) 1 = (k4\^{}2 {-} k0\^{}2 {-} k1\^{}2 {-} k2\^{}2 {-} k3\^{}2 {-} k5\^{}2 {-} k6\^{}2 {-} k7\^{}2) 1 for 40 random integer vectors k}
\NormalTok{[PASS] 71{-}D2  (exact) (gamma1671.k + m1671)(gamma1671.k {-} m1671) = (k.k {-} m1671\^{}2) 1 for 40 random integer (k, m1671): with dirac1671 = u exp({-}i k\_A x\^{}A) the operator i gamma1671\^{}A d\_A {-} m1671 becomes gamma1671.k {-} m1671, and every component of a solution obeys the Klein{-}Gordon equation (d4\^{}2 {-} d0\^{}2 {-} d1\^{}2 {-} d2\^{}2 {-} d3\^{}2 {-} d5\^{}2 {-} d6\^{}2 {-} d7\^{}2 + m1671\^{}2) dirac1671 = 0}
\NormalTok{[PASS] 71{-}D3  (exact) worked example E = 5, momentum p = (p1, p2, p3) = (2, 2, 4), m1671 = 1 (on the mass shell, 25 = 24 + 1): Pi1671 := (gamma1671.k + m1671)/(2 m1671) satisfies Pi1671.Pi1671 = Pi1671, trace Pi1671 = 8, (gamma1671.k {-} m1671).Pi1671 = 0: eight independent positive{-}energy plane waves}
\NormalTok{[PASS] 71{-}D4  (exact) beta1671 := gamma1671[4] and alpha1671\^{}j := gamma1671[4].gamma1671[j] (j = 0, 1, 2, 3, 5, 6, 7) are real symmetric; alpha1671\^{}i.alpha1671\^{}j + alpha1671\^{}j.alpha1671\^{}i = 2 delta\^{}(ij), alpha1671\^{}j.beta1671 + beta1671.alpha1671\^{}j = 0, beta1671.beta1671 = 1: the Hamiltonian h1671 = {-}i alpha1671\^{}j d\_j + beta1671 m1671 is Hermitian}
\NormalTok{[PASS] 71{-}D5  (exact) h1671(p) = p\^{}j alpha1671\^{}j + m1671 beta1671 is a real symmetric matrix with h1671(p).h1671(p) = (|p|\^{}2 + m1671\^{}2) 1 for 40 random integer (p, m1671): the energies are +E and {-}E, E = sqrt(|p|\^{}2 + m1671\^{}2), eight times each}
\NormalTok{[PASS] 71{-}D6  (exact) tau1671\^{}j := tauO[swap04(j)] (swap04 = (0 4)): tau1671\^{}4 = 1\_8; for the seven space indices tau1671\^{}j is real, antisymmetric, tau1671\^{}i.tau1671\^{}j + tau1671\^{}j.tau1671\^{}i = {-}2 delta\^{}(ij); in 8x8 blocks beta1671 = [[0, 1], [1, 0]] and alpha1671\^{}j = [[0, tau1671\^{}j], [{-}tau1671\^{}j, 0]] (beta1671 = betaD and alpha1671\^{}j = alphaD[swap04(j)] of 7+1.nb section 7.1)}
\NormalTok{[PASS] 71{-}D7  (exact) chiral form: with EL1671 = (1; {-}i)/sqrt(2), ER1671 = ({-}i; 1)/sqrt(2) (16x8, in 8x8 blocks; chi1671.EL1671 = {-}EL1671, chi1671.ER1671 = +ER1671, orthonormal): EL1671\^{}dagger.h1671(p).EL1671 = {-}i p\^{}j tau1671\^{}j, ER1671\^{}dagger.h1671(p).ER1671 = +i p\^{}j tau1671\^{}j, EL1671\^{}dagger.h1671(p).ER1671 = ER1671\^{}dagger.h1671(p).EL1671 = m1671 (20 random integer (p, m1671)); hence (i d\_4 + tau1671\^{}j d\_j) phiL1671 = m1671 phiR1671 and (i d\_4 {-} tau1671\^{}j d\_j) phiR1671 = m1671 phiL1671 for phiL1671 = EL1671\^{}dagger dirac1671, phiR1671 = ER1671\^{}dagger dirac1671 (EL1671, ER1671, phiL1671, phiR1671 are the objects EL, ER, phiL, phiR of 7+1.nb sections 7.10 and 7.11, in the PROJECT labelling of the coordinates)}
\NormalTok{[PASS] 71{-}D8  (exact) worked example E = 5, p = (2, 2, 4), m1671 = 1: with M1671 = m1671 + p\^{}j tau1671\^{}j (M1671.transpose(M1671) = E\^{}2 1\_8), the columns of Upos1671 := (E 1; transpose(M1671)) satisfy h1671(p).Upos1671 = +E Upos1671 and those of Vneg1671 := (E 1; {-}M1671) satisfy h1671({-}p).Vneg1671 = {-}E Vneg1671; transpose(Upos1671).Upos1671 = transpose(Vneg1671).Vneg1671 = 2 E\^{}2 1\_8: the spinors u\_s = Upos1671 e\_s/sqrt(E), v\_s = Vneg1671 e\_s/sqrt(E) are REAL and normalized to 2E}
\NormalTok{[PASS] 71{-}D9  (exact) charge conjugation dirac1671\_c := omega1671 conj(dirac1671): omega1671 is real and omega1671.gamma1671[A] = {-}gamma1671[A].omega1671, so (i gamma1671\^{}A d\_A {-} m1671) dirac1671\_c = omega1671 conj((i gamma1671\^{}A d\_A {-} m1671) dirac1671); omega1671.omega1671 = {-}1 gives (dirac1671\_c)\_c = {-}dirac1671: a 16{-}component Majorana condition dirac1671\_c = dirac1671 has only the solution dirac1671 = 0}
\NormalTok{[PASS] 71{-}Q1  (numerical) 7+1, one particle at rest: the 8 spin states form the spinor of SO(7) (Casimir 21/4) and split under SO(3) x SU(2)\_A x SU(2)\_B (ordinary rotations x rotations of the extra space x0 x5 x6 x7) as (2; 2, 1) + (2; 1, 2); [A\_a, B\_b] = 0 to 0e0, residual 6e{-}16}
\NormalTok{[PASS] 71{-}Q2  (numerical) 7+1, two identical particles in an s wave (28 states allowed by the Pauli principle): under SO(7) they form 7 + 21 (Casimir 6 and 10); there is NO SO(7) singlet.  Under SO(3) x SU(2)\_A x SU(2)\_B:  7 = (1; 2, 2) + (3; 1, 1),  21 = (1; 3, 1) + (1; 1, 3) + (3; 1, 1) + (3; 2, 2).  The pairs of ordinary spin 0 are the 10 states (1; 2, 2) + (1; 3, 1) + (1; 1, 3) (residual 3e{-}13)}
\NormalTok{[PASS] 71{-}Q3  (exact) 7+1: for space indices i != j, gamma1671[i].gamma1671[j] = diag(tau1671\^{}i tau1671\^{}j, tau1671\^{}i tau1671\^{}j) in 8x8 blocks, so on the particles at rest, u = (e; e)/sqrt(2), the rotation generator Sigma1671\^{}(ij) acts as the real antisymmetric 8x8 matrix (1/2) tau1671\^{}i tau1671\^{}j}
\NormalTok{[PASS] 71{-}Q4  (exact) 7+1: the 7 matrices tau1671\^{}j and the 21 matrices tau1671\^{}i tau1671\^{}j (i \textless{} j; space indices) are antisymmetric and linearly independent: together they span all 28 antisymmetric 8x8 matrices, i.e. all s{-}wave pair amplitudes}
\NormalTok{[PASS] 71{-}Q5  (exact) 7+1: the quadratic Casimir operator of SO(7), Cas7(Phi) := {-}sum\_(i\textless{}j) [Sigma1671\^{}(ij), [Sigma1671\^{}(ij), Phi]], gives Cas7(Phi) = 6 Phi for every Phi = tau1671\^{}j (the vector, 7) and Cas7(Phi) = 10 Phi for every Phi = tau1671\^{}i tau1671\^{}j (the adjoint, 21)}
\NormalTok{[PASS] 71{-}Q6  (exact) 7+1: the linear system [tau1671\^{}i tau1671\^{}j, X] = 0 (all 21 planes) for an 8x8 matrix X has rank 63 and is solved by X = 1: the only rotation{-}invariant pair matrix is the unit matrix, which is SYMMETRIC and therefore excluded by the Pauli principle.  Two identical particles of the field cannot form an s{-}wave pair that is a singlet under all rotations of the seven space dimensions}
\NormalTok{[PASS] 71{-}Q7  (exact) 7+1: the ten antisymmetric matrices tau1671\^{}0, tau1671\^{}5, tau1671\^{}6, tau1671\^{}7 and tau1671\^{}a tau1671\^{}b (a \textless{} b in \{0, 5, 6, 7\}) are linearly independent and commute with the generators tau1671\^{}2 tau1671\^{}3, tau1671\^{}3 tau1671\^{}1, tau1671\^{}1 tau1671\^{}2 of ordinary rotations; the system [tau1671\^{}i tau1671\^{}j, Phi] = 0 (ordinary rotations) on antisymmetric Phi has rank 18 = 28 {-} 10: these ten matrices are ALL the s{-}wave pairs of ordinary spin 0}
\NormalTok{[PASS] 71{-}Q8  (exact) 7+1: each of tau1671\^{}j (any space index j) and tau1671\^{}i tau1671\^{}j (i != j) is an orthogonal matrix: Phi.transpose(Phi) = 1.  A condensate whose pair matrix is proportional to tau1671\^{}j or to tau1671\^{}i tau1671\^{}j therefore gives the SAME gap to all eight spin states (a \textquotesingle{}unitary\textquotesingle{} state)}
\NormalTok{[PASS] 71{-}Q9  (exact) 7+1: two{-}particle content of the field bilinears, u\_s\^{}T X u\_r for particles at rest: X = Cm1671.gamma1671[j] gives {-}tau1671\^{}j (space index j) and X = Cm1671.gamma1671[4] gives 0; X = Cp1671.gamma1671[i].gamma1671[j] gives tau1671\^{}i tau1671\^{}j.  Hence the vector pair field is DeltaV1671\^{}j := dirac1671\^{}T Cm1671 gamma1671\^{}j dirac1671 (its time component contains no pair of two particles) and the adjoint pair field is DeltaT1671\^{}(ij) := dirac1671\^{}T Cp1671 gamma1671\^{}i gamma1671\^{}j dirac1671}
\NormalTok{[PASS] 71{-}Q10 (exact) 7+1: the system [tau1671\^{}i tau1671\^{}j, Phi] = 0 for the six rotation planes (1,2), (1,3), (2,3) of ordinary space and (5,6), (5,7), (6,7) of the extra space, on antisymmetric Phi, has rank 27 = 28 {-} 1, and Phi = tau1671\^{}0 solves it: EXACTLY ONE s{-}wave pair is a singlet both of the ordinary rotations and of the rotations of x5 x6 x7.  Its pair matrix is tau1671\^{}0, the two{-}particle content of DeltaV1671\^{}0 = dirac1671\^{}T Cm1671 gamma1671\^{}0 dirac1671, the component of the vector pair field along the hidden space dimension x0 (the same result as in 4+4, check 44{-}Q4)}
\NormalTok{[PASS] 44{-}Q1  (numerical) 4+4, fields independent of x5 x6 x7, one particle at rest (gamma1644[4] u = i u, 8 states; deviation 0e0): under SU(2)\_X x SU(2)\_Y (the two factors of the rotation group SO(4) of the four space coordinates x0 x1 x2 x3; they are NOT the factors SU(2)\_A x SU(2)\_B of the internal symmetry of the zero{-}mode sector, which do not act on ordinary space) x SU(2)\_T (rotations of the extra times) the states are (2, 1; 2) with one sign of W1644 and (1, 2; 2) with the other; every state has ordinary spin 1/2 (residual 2e{-}16)}
\NormalTok{[PASS] 44{-}Q2  (numerical) 4+4, fields independent of x5 x6 x7, two identical particles in an s wave (28 states): (1, 1; 3) twice, (3, 1; 1), (1, 3; 1), (2, 2; 1), (2, 2; 3) under SU(2)\_X x SU(2)\_Y x SU(2)\_T (X, Y as in 44{-}Q1).  The pairs that are singlets under ALL rotations of the four space dimensions are the 6 states (1, 1; 3) + (1, 1; 3): two triplets of the extra{-}time rotations, none of them a singlet of SU(2)\_T.  10 states have ordinary spin 0 (residual 1e{-}14)}
\NormalTok{[PASS] 44{-}Q3  (numerical) 4+4: among the 120 Pauli{-}allowed bilinears dirac1644\^{}T Cm1644 gamma1644[A1]...gamma1644[Ar] dirac1644 (A1 \textless{} ... \textless{} Ar), those whose two{-}particle content (particles at rest) lies entirely in the SO(4){-}singlet pairs have the index sets \{A1,...,Ar\} = \{5\} \{0,1,2,3,5\} \{4,5\} \{0,1,2,3,4,5\} \{6\} \{0,1,2,3,6\} \{4,6\} \{0,1,2,3,4,6\} \{7\} \{0,1,2,3,7\} \{4,7\} \{0,1,2,3,4,7\}; they include the components A = 5, 6, 7 (the extra{-}time components) of the vector pair field DeltaV1644\^{}A := dirac1644\^{}T Cm1644 gamma1644\^{}A dirac1644}
\NormalTok{[PASS] 44{-}Q4  (numerical) 4+4: exactly 1 of the 28 pair states is a singlet both of the ordinary rotations (x1 x2 x3) and of the extra{-}time rotations (x5 x6 x7); it belongs to the space vector (2, 2; 1), is made of one particle with W1644 = +1 and one with W1644 = {-}1, and is the two{-}particle content of the bilinears with the index sets \{A1,...,Ar\} = \{0\} \{0,4\}: in particular of DeltaV1644\^{}0 = dirac1644\^{}T Cm1644 gamma1644\^{}0 dirac1644, the component of the vector pair field along the hidden space dimension x0}
\NormalTok{[PASS] ALL    (exact) every check above passed}
\end{Highlighting}
\end{Shaded}

The classification of the pair states
(\texttt{data/pair\_classification.txt}; for an \(SU(2)\) multiplet of
spin \(j\) the program prints \(j(j+1)\), so that 0, 0.75 and 2 denote a
singlet, a doublet and a triplet). The letters name the groups: J the
ordinary rotations; in the lines for 7+1, A and B the two \(SU(2)\)
factors of the rotation group of the four extra space axes
\(x0,x5,x6,x7\); in the lines for 4+4, X and Y the two \(SU(2)\) factors
of the rotation group of the four space coordinates \(x0,x1,x2,x3\) ---
a different group, which contains the ordinary rotations and is
therefore given different letters --- and T the rotations of the three
additional times:

\begin{Shaded}
\begin{Highlighting}[]
\NormalTok{7+1 one particle at rest (8 states): [J(J+1)=0.75, A(A+1)=0, B(B+1)=0.75, C\_SO(7)=5.25 : 4 states] [J(J+1)=0.75, A(A+1)=0.75, B(B+1)=0, C\_SO(7)=5.25 : 4 states] }
\NormalTok{7+1 two particles, s wave (28 states): [J(J+1)=0, A(A+1)=0, B(B+1)=2, C\_SO(7)=10 : 3 states] [J(J+1)=0, A(A+1)=0.75, B(B+1)=0.75, C\_SO(7)=6 : 4 states] [J(J+1)=0, A(A+1)=2, B(B+1)=0, C\_SO(7)=10 : 3 states] [J(J+1)=2, A(A+1)=0, B(B+1)=0, C\_SO(7)=6 : 3 states] [J(J+1)=2, A(A+1)=0, B(B+1)=0, C\_SO(7)=10 : 3 states] [J(J+1)=2, A(A+1)=0.75, B(B+1)=0.75, C\_SO(7)=10 : 12 states] }
\NormalTok{4+4 (fields independent of x5,x6,x7) one particle at rest (8 states): [X(X+1)=0, Y(Y+1)=0.75, T(T+1)=0.75, J(J+1)=0.75, W=1 : 4 states] [X(X+1)=0.75, Y(Y+1)=0, T(T+1)=0.75, J(J+1)=0.75, W={-}1 : 4 states] }
\NormalTok{4+4 two particles, s wave (28 states): [X(X+1)=0, Y(Y+1)=0, T(T+1)=2, J(J+1)=0, WW=1 : 6 states] [X(X+1)=0, Y(Y+1)=2, T(T+1)=0, J(J+1)=2, WW=1 : 3 states] [X(X+1)=0.75, Y(Y+1)=0.75, T(T+1)=0, J(J+1)=0, WW={-}1 : 1 states] [X(X+1)=0.75, Y(Y+1)=0.75, T(T+1)=0, J(J+1)=2, WW={-}1 : 3 states] [X(X+1)=0.75, Y(Y+1)=0.75, T(T+1)=2, J(J+1)=0, WW={-}1 : 3 states] [X(X+1)=0.75, Y(Y+1)=0.75, T(T+1)=2, J(J+1)=2, WW={-}1 : 9 states] [X(X+1)=2, Y(Y+1)=0, T(T+1)=0, J(J+1)=2, WW=1 : 3 states] }
\NormalTok{4+4: bilinears dirac1644\^{}T Cm1644 gamma1644[A1]...gamma1644[Ar] dirac1644 whose two{-}particle content lies entirely in the SO(4){-}singlet pairs: index sets \{A1,...,Ar\} = \{5\} \{0,1,2,3,5\} \{4,5\} \{0,1,2,3,4,5\} \{6\} \{0,1,2,3,6\} \{4,6\} \{0,1,2,3,4,6\} \{7\} \{0,1,2,3,7\} \{4,7\} \{0,1,2,3,4,7\}}
\NormalTok{4+4: bilinears dirac1644\^{}T Cm1644 gamma1644[A1]...gamma1644[Ar] dirac1644 whose two{-}particle content is the unique pair state with ordinary spin 0 and extra{-}time isospin 0: index sets \{A1,...,Ar\} = \{0\} \{0,4\}}
\end{Highlighting}
\end{Shaded}

\hypertarget{a.2-the-independent-python-script}{%
\subsection{A.2 The independent Python
script}\label{a.2-the-independent-python-script}}

The following is the file \texttt{data/matrix\_checks\_independent.txt},
written by \texttt{code/scripts/verify\_matrices\_independent.py}. All
arithmetic is exact (integers, rational numbers, Gaussian integers).

\begin{Shaded}
\begin{Highlighting}[]
\NormalTok{PROJECT spacetimes: INDEPENDENT checks (pure Python, exact arithmetic) of the two sets of real 16x16 Dirac matrices}
\NormalTok{36 of 36 checks passed}

\NormalTok{[PASS] I44{-}F   the Python matrices gamma1644[0..7], omega1644, PL1644, PR1644, Cm1644, Cp1644, sigma88, tau1644[0..7], taubar1644[0..7] equal those in data/gamma1644.txt (written by the Rust program), entry by entry}
\NormalTok{[PASS] I71{-}F   the Python matrices gamma1671[0..7], omega1671, 2 PL1671, 2 PR1671, Cp1671, Cm1671, tau1671[0..7] equal those in data/gamma1671.txt, entry by entry}
\NormalTok{[PASS] I44{-}N   gamma1644[0..7], omega1644, PL1644, PR1644 and eta1644 equal the matrices printed by Mathematica in 4+4.nb (cells 314, 136, 163, 165, 13)}
\NormalTok{[PASS] I71{-}N   Gamma17[0..7], rebuilt from the definitions of 7+1.nb, equal the matrices printed by Mathematica in 7+1.nb (cell 204); gamma1671[4] = Gamma17[0], gamma1671[0] = {-}Gamma17[4], gamma1671[j] = {-}Gamma17[j] (j = 1,2,3,5,6,7)}
\NormalTok{[PASS] I44{-}R   every entry of gamma1644[A] is {-}1, 0 or +1 (the matrices are real), and every matrix has exactly one non{-}zero entry in each row and column}
\NormalTok{[PASS] I44{-}C   gamma1644[A].gamma1644[B] + gamma1644[B].gamma1644[A] = 2 eta1644\^{}(AB) 1 for all 64 pairs, eta1644 = diag[1, 1, 1, 1, {-}1, {-}1, {-}1, {-}1]}
\NormalTok{[PASS] I44{-}G   Gram matrix of the 256 ordered products of the gamma1644[A] under \textless{}X, Y\textgreater{} = trace(transpose(X).Y) equals 16 times the unit matrix: the products are orthogonal, hence linearly independent, hence a basis of all real 16x16 matrices (METHOD differs from the Rust check C3, which uses a rank)}
\NormalTok{[PASS] I44{-}A   of the 256 basis products exactly one anticommutes with all eight gamma1644[A]: the product of all eight (index set [(0, 1, 2, 3, 4, 5, 6, 7)]); exactly one commutes with all eight: the unit matrix.  Because X.gamma1644[A] +{-} gamma1644[A].X maps different basis products to different basis products, a general real matrix X anticommutes (commutes) with all eight only if it is a multiple of that product (of the unit matrix)}
\NormalTok{[PASS] I44{-}L1  [gamma1644[A], gamma1644[B]] = 4 Sigma1644\^{}(AB) with Sigma1644\^{}(AB) = (1/2) gamma1644[A].gamma1644[B] (A != B), and [Sigma1644\^{}(AB), gamma1644[C]] = eta1644\^{}(BC) gamma1644[A] {-} eta1644\^{}(AC) gamma1644[B] for all 512 index triples}
\NormalTok{[PASS] I44{-}L2  [Sigma1644\^{}(AB), Sigma1644\^{}(CD)] = eta1644\^{}(BC) Sigma1644\^{}(AD) {-} eta1644\^{}(AC) Sigma1644\^{}(BD) {-} eta1644\^{}(BD) Sigma1644\^{}(AC) + eta1644\^{}(AD) Sigma1644\^{}(BC) for all A \textless{} B, C \textless{} D (784 cases)}
\NormalTok{[PASS] I44{-}K   4+4: the scaled commutators close under commutation; the Killing form, computed from the adjoint representation, is diagonal with B(Sigma1644\^{}(AB), Sigma1644\^{}(AB)) = {-}12 eta1644\^{}(AA) eta1644\^{}(BB): 16 positive (non{-}compact) and 12 negative (compact) directions, the signature of so(4,4)}
\NormalTok{[PASS] I71{-}R   every entry of gamma1671[A] is {-}1, 0 or +1 (the matrices are real), and every matrix has exactly one non{-}zero entry in each row and column}
\NormalTok{[PASS] I71{-}C   gamma1671[A].gamma1671[B] + gamma1671[B].gamma1671[A] = 2 eta1671\^{}(AB) 1 for all 64 pairs, eta1671 = diag[{-}1, {-}1, {-}1, {-}1, 1, {-}1, {-}1, {-}1]}
\NormalTok{[PASS] I71{-}G   Gram matrix of the 256 ordered products of the gamma1671[A] under \textless{}X, Y\textgreater{} = trace(transpose(X).Y) equals 16 times the unit matrix: the products are orthogonal, hence linearly independent, hence a basis of all real 16x16 matrices (METHOD differs from the Rust check C3, which uses a rank)}
\NormalTok{[PASS] I71{-}A   of the 256 basis products exactly one anticommutes with all eight gamma1671[A]: the product of all eight (index set [(0, 1, 2, 3, 4, 5, 6, 7)]); exactly one commutes with all eight: the unit matrix.  Because X.gamma1671[A] +{-} gamma1671[A].X maps different basis products to different basis products, a general real matrix X anticommutes (commutes) with all eight only if it is a multiple of that product (of the unit matrix)}
\NormalTok{[PASS] I71{-}L1  [gamma1671[A], gamma1671[B]] = 4 Sigma1671\^{}(AB) with Sigma1671\^{}(AB) = (1/2) gamma1671[A].gamma1671[B] (A != B), and [Sigma1671\^{}(AB), gamma1671[C]] = eta1671\^{}(BC) gamma1671[A] {-} eta1671\^{}(AC) gamma1671[B] for all 512 index triples}
\NormalTok{[PASS] I71{-}L2  [Sigma1671\^{}(AB), Sigma1671\^{}(CD)] = eta1671\^{}(BC) Sigma1671\^{}(AD) {-} eta1671\^{}(AC) Sigma1671\^{}(BD) {-} eta1671\^{}(BD) Sigma1671\^{}(AC) + eta1671\^{}(AD) Sigma1671\^{}(BC) for all A \textless{} B, C \textless{} D (784 cases)}
\NormalTok{[PASS] I71{-}K   7+1: the scaled commutators close under commutation; the Killing form is diagonal with B(Sigma1671\^{}(AB), Sigma1671\^{}(AB)) = {-}12 eta1671\^{}(AA) eta1671\^{}(BB): 7 positive (non{-}compact) and 21 negative (compact) directions, the signature of so(7,1), NOT of so(4,4)}
\NormalTok{[PASS] I44{-}V   omega1644 = diag({-}1\_8, +1\_8), omega1644.omega1644 = 1; PL1644 = diag(1\_8, 0), PR1644 = diag(0, 1\_8); PL1644 + PR1644 = 1, PL1644.PL1644 = PL1644, PR1644.PR1644 = PR1644, PL1644.PR1644 = 0; gamma1644[A].PL1644 = PR1644.gamma1644[A]}
\NormalTok{[PASS] I44{-}S   Cm1644 = diag({-}sigma88, sigma88) and Cp1644 = diag(sigma88, sigma88) are symmetric with square 1; Cm1644.gamma1644[A].Cm1644 = {-}transpose(gamma1644[A]), Cp1644.gamma1644[A].Cp1644 = +transpose(gamma1644[A])}
\NormalTok{[PASS] I71{-}V   omega1671 = [[0, {-}1], [1, 0]], omega1671.omega1671 = {-}1; chi1671 = i omega1671 is Hermitian with square 1 and anticommutes with every gamma1671[A]; (2 PL1671).(2 PL1671) = 2 (2 PL1671), (2 PL1671).(2 PR1671) = 0, gamma1671[A].PL1671 = PR1671.gamma1671[A]}
\NormalTok{[PASS] I71{-}U   with s = 1 {-} i Gamma17[0]: s.s\^{}dagger = 2 and s.chi1671 = {-}omega1644.s, so Uc1671 = s/sqrt(2) is unitary and maps PL1671 to PR1644 and PR1671 to PL1644}
\NormalTok{[PASS] I71{-}S   Cp1671 = gamma1671[4] and Cm1671 = gamma1671[4].omega1671 = diag(1\_8, {-}1\_8) are symmetric; Cp1671.gamma1671[A] = +transpose(gamma1671[A]).Cp1671, Cm1671.gamma1671[A] = {-}transpose(gamma1671[A]).Cm1671}
\NormalTok{[PASS] I{-}D1    (gamma1644\^{}A k\_A)\^{}2 = eta1644\^{}(AB) k\_A k\_B and (gamma1671\^{}A k\_A)\^{}2 = eta1671\^{}(AB) k\_A k\_B for 25 random integer vectors k: both Dirac operators square to the wave operator of their spacetime}
\NormalTok{[PASS] I44{-}D2  gamma1644[A] = [[0, taubar1644[A]], [tau1644[A], 0]] and tau1644[A].taubar1644[B] + tau1644[B].taubar1644[A] = 2 eta1644\^{}(AB) 1\_8: the equation (gamma1644\^{}A d\_A {-} m1644) dirac1644 = 0 is the pair taubar1644\^{}A d\_A psiR1644 = m1644 psiL1644, tau1644\^{}A d\_A psiL1644 = m1644 psiR1644}
\NormalTok{[PASS] I44{-}D3  4+4: K1644.K1644 = {-}(m1644\^{}2 + |k|\^{}2 {-} |kappa|\^{}2) 1 and K1644 + K1644\^{}dagger = 2 i kappa\_t gamma1644\^{}4 gamma1644\^{}t for 25 random integer (m1644, k, kappa): the x4{-}evolution is unitary exactly when the field is independent of x5, x6, x7}
\NormalTok{[PASS] I44{-}D4  W1644 = i Cm1644.gamma1644[4] is Hermitian, W1644.W1644 = 1, trace W1644 = 0: the conserved density dirac1644\^{}dagger W1644 dirac1644 of the Spin(4,4){-}invariant current is indefinite}
\NormalTok{[PASS] I71{-}D2  tau1671[4] = 1\_8; for the seven space indices tau1671[j] is antisymmetric with tau1671[i].tau1671[j] + tau1671[j].tau1671[i] = {-}2 delta\_ij; gamma1671[4] = [[0, 1], [1, 0]], gamma1671[j] = diag({-}tau1671[j], tau1671[j]); alpha1671\^{}j = gamma1671[4].gamma1671[j] = [[0, tau1671[j]], [{-}tau1671[j], 0]] is symmetric}
\NormalTok{[PASS] I71{-}D3  7+1: h1671(p) = p\^{}j alpha1671\^{}j + m1671 beta1671 satisfies h1671.h1671 = (|p|\^{}2 + m1671\^{}2) 1, and in the chiral basis EL1671\^{}dagger.h1671.EL1671 = {-}i p\^{}j tau1671\^{}j, ER1671\^{}dagger.h1671.ER1671 = +i p\^{}j tau1671\^{}j, EL1671\^{}dagger.h1671.ER1671 = ER1671\^{}dagger.h1671.EL1671 = m1671 (10 random integer cases): (i d\_4 + tau1671\^{}j d\_j) phiL1671 = m1671 phiR1671, (i d\_4 {-} tau1671\^{}j d\_j) phiR1671 = m1671 phiL1671}
\NormalTok{[PASS] I{-}P1    symmetry (+1 symmetric, {-}1 antisymmetric) of C.gamma[A1]...gamma[Ar] by rank r = 0..8: [1, {-}1, {-}1, 1, 1, {-}1, {-}1, 1, 1] for Cm1644 and for Cm1671, [1, 1, {-}1, {-}1, 1, 1, {-}1, {-}1, 1] for Cp1644 and for Cp1671; in particular the rank{-}0 and rank{-}8 (Lorentz{-}scalar) matrices are symmetric, so the scalar pair bilinears of one anticommuting field vanish}
\NormalTok{[PASS] I71{-}P2  7+1, two identical particles at rest in an s wave: the antisymmetric pair matrices invariant under all 21 rotation generators of SO(7) form a space of dimension 0 (no rotation{-}singlet pair); those invariant under the ordinary rotations (x1 x2 x3) form a space of dimension 10 (ten pairs of ordinary spin 0)}
\NormalTok{[PASS] I44{-}P2  4+4 (fields independent of x5 x6 x7), two identical particles at rest in an s wave: pairs invariant under the rotations of the four space dimensions: 6; under the ordinary rotations: 10; under the ordinary rotations AND the rotations of the extra times: 1; under the rotations of the four space dimensions AND of the extra times: 0}
\NormalTok{[PASS] I44{-}P3  4+4: the two{-}particle content of the bilinear DeltaV1644\^{}0 = dirac1644\^{}T Cm1644 gamma1644[0] dirac1644 is a non{-}zero antisymmetric pair matrix that is invariant under the ordinary rotations and under the rotations of the extra times: it is the unique pair of I44{-}P2 with ordinary spin 0 and extra{-}time spin 0}
\NormalTok{[PASS] I71{-}P3  7+1: the two{-}particle content of DeltaV1671\^{}j = dirac1671\^{}T Cm1671 gamma1671[j] dirac1671 (particles at rest) is {-}tau1671[j] for the seven space indices, and 0 for the time index 4; the hidden{-}axis component DeltaV1671\^{}0 has the pair matrix {-}tau1671[0], one of the ten pairs of ordinary spin 0}
\NormalTok{[PASS] I71{-}P4  7+1: the antisymmetric pair matrices invariant under the ordinary rotations (x1 x2 x3) AND under the rotations of x5 x6 x7 form a space of dimension 1, spanned by tau1671[0]: exactly one s{-}wave pair is a singlet of both, the hidden{-}axis pair, as in 4+4 (I44{-}P2, I44{-}P3)}
\NormalTok{[PASS] I{-}P5    in both signatures the bilinear DeltaV\^{}0 = dirac\^{}T Cm gamma[0] dirac is invariant under the 21 generators Sigma\^{}(AB) with A, B in 1..7 (the Lorentz group of the seven coordinates other than x0, which contains the Lorentz group of x1 x2 x3 x4), and is changed by each of the 7 generators Sigma\^{}(0B)}
\end{Highlighting}
\end{Shaded}

\hypertarget{appendix-b.-the-88-blocks}{%
\section{Appendix B. The 8×8 blocks}\label{appendix-b.-the-88-blocks}}

\textbf{\(\sigma_{88}\) and the matrices \(\tau_{1644}^A\),
\(\bar\tau_{1644}^A\) of 4+4} (4×4 blocks are separated by rules;
\(\tau_{1644}^0=\bar\tau_{1644}^0=1_8\)):

\[
\sigma_{88} = \left(\begin{array}{rrrr|rrrr}
\phantom{-}\cdot & \phantom{-}\cdot & \phantom{-}\cdot & \phantom{-}\cdot & \phantom{-}1 & \phantom{-}\cdot & \phantom{-}\cdot & \phantom{-}\cdot \\[-1.8pt]
\phantom{-}\cdot & \phantom{-}\cdot & \phantom{-}\cdot & \phantom{-}\cdot & \phantom{-}\cdot & \phantom{-}1 & \phantom{-}\cdot & \phantom{-}\cdot \\[-1.8pt]
\phantom{-}\cdot & \phantom{-}\cdot & \phantom{-}\cdot & \phantom{-}\cdot & \phantom{-}\cdot & \phantom{-}\cdot & \phantom{-}1 & \phantom{-}\cdot \\[-1.8pt]
\phantom{-}\cdot & \phantom{-}\cdot & \phantom{-}\cdot & \phantom{-}\cdot & \phantom{-}\cdot & \phantom{-}\cdot & \phantom{-}\cdot & \phantom{-}1 \\[-1.8pt]
\hline
\phantom{-}1 & \phantom{-}\cdot & \phantom{-}\cdot & \phantom{-}\cdot & \phantom{-}\cdot & \phantom{-}\cdot & \phantom{-}\cdot & \phantom{-}\cdot \\[-1.8pt]
\phantom{-}\cdot & \phantom{-}1 & \phantom{-}\cdot & \phantom{-}\cdot & \phantom{-}\cdot & \phantom{-}\cdot & \phantom{-}\cdot & \phantom{-}\cdot \\[-1.8pt]
\phantom{-}\cdot & \phantom{-}\cdot & \phantom{-}1 & \phantom{-}\cdot & \phantom{-}\cdot & \phantom{-}\cdot & \phantom{-}\cdot & \phantom{-}\cdot \\[-1.8pt]
\phantom{-}\cdot & \phantom{-}\cdot & \phantom{-}\cdot & \phantom{-}1 & \phantom{-}\cdot & \phantom{-}\cdot & \phantom{-}\cdot & \phantom{-}\cdot
\end{array}\right)
\]

\[
\tau_{1644}^{1} = \left(\begin{array}{rrrr|rrrr}
\phantom{-}\cdot & \phantom{-}\cdot & \phantom{-}\cdot & \phantom{-}\cdot & \phantom{-}\cdot & \phantom{-}\cdot & \phantom{-}\cdot & \phantom{-}1 \\[-1.8pt]
\phantom{-}\cdot & \phantom{-}\cdot & \phantom{-}\cdot & \phantom{-}\cdot & \phantom{-}\cdot & \phantom{-}\cdot & \phantom{-}1 & \phantom{-}\cdot \\[-1.8pt]
\phantom{-}\cdot & \phantom{-}\cdot & \phantom{-}\cdot & \phantom{-}\cdot & \phantom{-}\cdot & -1 & \phantom{-}\cdot & \phantom{-}\cdot \\[-1.8pt]
\phantom{-}\cdot & \phantom{-}\cdot & \phantom{-}\cdot & \phantom{-}\cdot & -1 & \phantom{-}\cdot & \phantom{-}\cdot & \phantom{-}\cdot \\[-1.8pt]
\hline
\phantom{-}\cdot & \phantom{-}\cdot & \phantom{-}\cdot & \phantom{-}1 & \phantom{-}\cdot & \phantom{-}\cdot & \phantom{-}\cdot & \phantom{-}\cdot \\[-1.8pt]
\phantom{-}\cdot & \phantom{-}\cdot & \phantom{-}1 & \phantom{-}\cdot & \phantom{-}\cdot & \phantom{-}\cdot & \phantom{-}\cdot & \phantom{-}\cdot \\[-1.8pt]
\phantom{-}\cdot & -1 & \phantom{-}\cdot & \phantom{-}\cdot & \phantom{-}\cdot & \phantom{-}\cdot & \phantom{-}\cdot & \phantom{-}\cdot \\[-1.8pt]
-1 & \phantom{-}\cdot & \phantom{-}\cdot & \phantom{-}\cdot & \phantom{-}\cdot & \phantom{-}\cdot & \phantom{-}\cdot & \phantom{-}\cdot
\end{array}\right)
\]

\[
\tau_{1644}^{2} = \left(\begin{array}{rrrr|rrrr}
\phantom{-}\cdot & \phantom{-}\cdot & \phantom{-}\cdot & \phantom{-}\cdot & \phantom{-}\cdot & \phantom{-}\cdot & -1 & \phantom{-}\cdot \\[-1.8pt]
\phantom{-}\cdot & \phantom{-}\cdot & \phantom{-}\cdot & \phantom{-}\cdot & \phantom{-}\cdot & \phantom{-}\cdot & \phantom{-}\cdot & \phantom{-}1 \\[-1.8pt]
\phantom{-}\cdot & \phantom{-}\cdot & \phantom{-}\cdot & \phantom{-}\cdot & \phantom{-}1 & \phantom{-}\cdot & \phantom{-}\cdot & \phantom{-}\cdot \\[-1.8pt]
\phantom{-}\cdot & \phantom{-}\cdot & \phantom{-}\cdot & \phantom{-}\cdot & \phantom{-}\cdot & -1 & \phantom{-}\cdot & \phantom{-}\cdot \\[-1.8pt]
\hline
\phantom{-}\cdot & \phantom{-}\cdot & -1 & \phantom{-}\cdot & \phantom{-}\cdot & \phantom{-}\cdot & \phantom{-}\cdot & \phantom{-}\cdot \\[-1.8pt]
\phantom{-}\cdot & \phantom{-}\cdot & \phantom{-}\cdot & \phantom{-}1 & \phantom{-}\cdot & \phantom{-}\cdot & \phantom{-}\cdot & \phantom{-}\cdot \\[-1.8pt]
\phantom{-}1 & \phantom{-}\cdot & \phantom{-}\cdot & \phantom{-}\cdot & \phantom{-}\cdot & \phantom{-}\cdot & \phantom{-}\cdot & \phantom{-}\cdot \\[-1.8pt]
\phantom{-}\cdot & -1 & \phantom{-}\cdot & \phantom{-}\cdot & \phantom{-}\cdot & \phantom{-}\cdot & \phantom{-}\cdot & \phantom{-}\cdot
\end{array}\right)
\]

\[
\tau_{1644}^{3} = \left(\begin{array}{rrrr|rrrr}
\phantom{-}\cdot & \phantom{-}\cdot & \phantom{-}\cdot & \phantom{-}\cdot & \phantom{-}\cdot & \phantom{-}1 & \phantom{-}\cdot & \phantom{-}\cdot \\[-1.8pt]
\phantom{-}\cdot & \phantom{-}\cdot & \phantom{-}\cdot & \phantom{-}\cdot & -1 & \phantom{-}\cdot & \phantom{-}\cdot & \phantom{-}\cdot \\[-1.8pt]
\phantom{-}\cdot & \phantom{-}\cdot & \phantom{-}\cdot & \phantom{-}\cdot & \phantom{-}\cdot & \phantom{-}\cdot & \phantom{-}\cdot & \phantom{-}1 \\[-1.8pt]
\phantom{-}\cdot & \phantom{-}\cdot & \phantom{-}\cdot & \phantom{-}\cdot & \phantom{-}\cdot & \phantom{-}\cdot & -1 & \phantom{-}\cdot \\[-1.8pt]
\hline
\phantom{-}\cdot & \phantom{-}1 & \phantom{-}\cdot & \phantom{-}\cdot & \phantom{-}\cdot & \phantom{-}\cdot & \phantom{-}\cdot & \phantom{-}\cdot \\[-1.8pt]
-1 & \phantom{-}\cdot & \phantom{-}\cdot & \phantom{-}\cdot & \phantom{-}\cdot & \phantom{-}\cdot & \phantom{-}\cdot & \phantom{-}\cdot \\[-1.8pt]
\phantom{-}\cdot & \phantom{-}\cdot & \phantom{-}\cdot & \phantom{-}1 & \phantom{-}\cdot & \phantom{-}\cdot & \phantom{-}\cdot & \phantom{-}\cdot \\[-1.8pt]
\phantom{-}\cdot & \phantom{-}\cdot & -1 & \phantom{-}\cdot & \phantom{-}\cdot & \phantom{-}\cdot & \phantom{-}\cdot & \phantom{-}\cdot
\end{array}\right)
\]

\[
\tau_{1644}^{4} = \left(\begin{array}{rrrr|rrrr}
\phantom{-}\cdot & \phantom{-}\cdot & \phantom{-}\cdot & \phantom{-}\cdot & \phantom{-}\cdot & \phantom{-}1 & \phantom{-}\cdot & \phantom{-}\cdot \\[-1.8pt]
\phantom{-}\cdot & \phantom{-}\cdot & \phantom{-}\cdot & \phantom{-}\cdot & -1 & \phantom{-}\cdot & \phantom{-}\cdot & \phantom{-}\cdot \\[-1.8pt]
\phantom{-}\cdot & \phantom{-}\cdot & \phantom{-}\cdot & \phantom{-}\cdot & \phantom{-}\cdot & \phantom{-}\cdot & \phantom{-}\cdot & -1 \\[-1.8pt]
\phantom{-}\cdot & \phantom{-}\cdot & \phantom{-}\cdot & \phantom{-}\cdot & \phantom{-}\cdot & \phantom{-}\cdot & \phantom{-}1 & \phantom{-}\cdot \\[-1.8pt]
\hline
\phantom{-}\cdot & -1 & \phantom{-}\cdot & \phantom{-}\cdot & \phantom{-}\cdot & \phantom{-}\cdot & \phantom{-}\cdot & \phantom{-}\cdot \\[-1.8pt]
\phantom{-}1 & \phantom{-}\cdot & \phantom{-}\cdot & \phantom{-}\cdot & \phantom{-}\cdot & \phantom{-}\cdot & \phantom{-}\cdot & \phantom{-}\cdot \\[-1.8pt]
\phantom{-}\cdot & \phantom{-}\cdot & \phantom{-}\cdot & \phantom{-}1 & \phantom{-}\cdot & \phantom{-}\cdot & \phantom{-}\cdot & \phantom{-}\cdot \\[-1.8pt]
\phantom{-}\cdot & \phantom{-}\cdot & -1 & \phantom{-}\cdot & \phantom{-}\cdot & \phantom{-}\cdot & \phantom{-}\cdot & \phantom{-}\cdot
\end{array}\right)
\]

\[
\tau_{1644}^{5} = \left(\begin{array}{rrrr|rrrr}
\phantom{-}\cdot & \phantom{-}\cdot & \phantom{-}\cdot & \phantom{-}\cdot & \phantom{-}\cdot & \phantom{-}\cdot & -1 & \phantom{-}\cdot \\[-1.8pt]
\phantom{-}\cdot & \phantom{-}\cdot & \phantom{-}\cdot & \phantom{-}\cdot & \phantom{-}\cdot & \phantom{-}\cdot & \phantom{-}\cdot & -1 \\[-1.8pt]
\phantom{-}\cdot & \phantom{-}\cdot & \phantom{-}\cdot & \phantom{-}\cdot & \phantom{-}1 & \phantom{-}\cdot & \phantom{-}\cdot & \phantom{-}\cdot \\[-1.8pt]
\phantom{-}\cdot & \phantom{-}\cdot & \phantom{-}\cdot & \phantom{-}\cdot & \phantom{-}\cdot & \phantom{-}1 & \phantom{-}\cdot & \phantom{-}\cdot \\[-1.8pt]
\hline
\phantom{-}\cdot & \phantom{-}\cdot & \phantom{-}1 & \phantom{-}\cdot & \phantom{-}\cdot & \phantom{-}\cdot & \phantom{-}\cdot & \phantom{-}\cdot \\[-1.8pt]
\phantom{-}\cdot & \phantom{-}\cdot & \phantom{-}\cdot & \phantom{-}1 & \phantom{-}\cdot & \phantom{-}\cdot & \phantom{-}\cdot & \phantom{-}\cdot \\[-1.8pt]
-1 & \phantom{-}\cdot & \phantom{-}\cdot & \phantom{-}\cdot & \phantom{-}\cdot & \phantom{-}\cdot & \phantom{-}\cdot & \phantom{-}\cdot \\[-1.8pt]
\phantom{-}\cdot & -1 & \phantom{-}\cdot & \phantom{-}\cdot & \phantom{-}\cdot & \phantom{-}\cdot & \phantom{-}\cdot & \phantom{-}\cdot
\end{array}\right)
\]

\[
\tau_{1644}^{6} = \left(\begin{array}{rrrr|rrrr}
\phantom{-}\cdot & \phantom{-}\cdot & \phantom{-}\cdot & \phantom{-}\cdot & \phantom{-}\cdot & \phantom{-}\cdot & \phantom{-}\cdot & -1 \\[-1.8pt]
\phantom{-}\cdot & \phantom{-}\cdot & \phantom{-}\cdot & \phantom{-}\cdot & \phantom{-}\cdot & \phantom{-}\cdot & \phantom{-}1 & \phantom{-}\cdot \\[-1.8pt]
\phantom{-}\cdot & \phantom{-}\cdot & \phantom{-}\cdot & \phantom{-}\cdot & \phantom{-}\cdot & -1 & \phantom{-}\cdot & \phantom{-}\cdot \\[-1.8pt]
\phantom{-}\cdot & \phantom{-}\cdot & \phantom{-}\cdot & \phantom{-}\cdot & \phantom{-}1 & \phantom{-}\cdot & \phantom{-}\cdot & \phantom{-}\cdot \\[-1.8pt]
\hline
\phantom{-}\cdot & \phantom{-}\cdot & \phantom{-}\cdot & \phantom{-}1 & \phantom{-}\cdot & \phantom{-}\cdot & \phantom{-}\cdot & \phantom{-}\cdot \\[-1.8pt]
\phantom{-}\cdot & \phantom{-}\cdot & -1 & \phantom{-}\cdot & \phantom{-}\cdot & \phantom{-}\cdot & \phantom{-}\cdot & \phantom{-}\cdot \\[-1.8pt]
\phantom{-}\cdot & \phantom{-}1 & \phantom{-}\cdot & \phantom{-}\cdot & \phantom{-}\cdot & \phantom{-}\cdot & \phantom{-}\cdot & \phantom{-}\cdot \\[-1.8pt]
-1 & \phantom{-}\cdot & \phantom{-}\cdot & \phantom{-}\cdot & \phantom{-}\cdot & \phantom{-}\cdot & \phantom{-}\cdot & \phantom{-}\cdot
\end{array}\right)
\]

\[
\tau_{1644}^{7} = \left(\begin{array}{rrrr|rrrr}
-1 & \phantom{-}\cdot & \phantom{-}\cdot & \phantom{-}\cdot & \phantom{-}\cdot & \phantom{-}\cdot & \phantom{-}\cdot & \phantom{-}\cdot \\[-1.8pt]
\phantom{-}\cdot & -1 & \phantom{-}\cdot & \phantom{-}\cdot & \phantom{-}\cdot & \phantom{-}\cdot & \phantom{-}\cdot & \phantom{-}\cdot \\[-1.8pt]
\phantom{-}\cdot & \phantom{-}\cdot & -1 & \phantom{-}\cdot & \phantom{-}\cdot & \phantom{-}\cdot & \phantom{-}\cdot & \phantom{-}\cdot \\[-1.8pt]
\phantom{-}\cdot & \phantom{-}\cdot & \phantom{-}\cdot & -1 & \phantom{-}\cdot & \phantom{-}\cdot & \phantom{-}\cdot & \phantom{-}\cdot \\[-1.8pt]
\hline
\phantom{-}\cdot & \phantom{-}\cdot & \phantom{-}\cdot & \phantom{-}\cdot & \phantom{-}1 & \phantom{-}\cdot & \phantom{-}\cdot & \phantom{-}\cdot \\[-1.8pt]
\phantom{-}\cdot & \phantom{-}\cdot & \phantom{-}\cdot & \phantom{-}\cdot & \phantom{-}\cdot & \phantom{-}1 & \phantom{-}\cdot & \phantom{-}\cdot \\[-1.8pt]
\phantom{-}\cdot & \phantom{-}\cdot & \phantom{-}\cdot & \phantom{-}\cdot & \phantom{-}\cdot & \phantom{-}\cdot & \phantom{-}1 & \phantom{-}\cdot \\[-1.8pt]
\phantom{-}\cdot & \phantom{-}\cdot & \phantom{-}\cdot & \phantom{-}\cdot & \phantom{-}\cdot & \phantom{-}\cdot & \phantom{-}\cdot & \phantom{-}1
\end{array}\right)
\]

\[
\bar\tau_{1644}^{1} = \left(\begin{array}{rrrr|rrrr}
\phantom{-}\cdot & \phantom{-}\cdot & \phantom{-}\cdot & \phantom{-}\cdot & \phantom{-}\cdot & \phantom{-}\cdot & \phantom{-}\cdot & -1 \\[-1.8pt]
\phantom{-}\cdot & \phantom{-}\cdot & \phantom{-}\cdot & \phantom{-}\cdot & \phantom{-}\cdot & \phantom{-}\cdot & -1 & \phantom{-}\cdot \\[-1.8pt]
\phantom{-}\cdot & \phantom{-}\cdot & \phantom{-}\cdot & \phantom{-}\cdot & \phantom{-}\cdot & \phantom{-}1 & \phantom{-}\cdot & \phantom{-}\cdot \\[-1.8pt]
\phantom{-}\cdot & \phantom{-}\cdot & \phantom{-}\cdot & \phantom{-}\cdot & \phantom{-}1 & \phantom{-}\cdot & \phantom{-}\cdot & \phantom{-}\cdot \\[-1.8pt]
\hline
\phantom{-}\cdot & \phantom{-}\cdot & \phantom{-}\cdot & -1 & \phantom{-}\cdot & \phantom{-}\cdot & \phantom{-}\cdot & \phantom{-}\cdot \\[-1.8pt]
\phantom{-}\cdot & \phantom{-}\cdot & -1 & \phantom{-}\cdot & \phantom{-}\cdot & \phantom{-}\cdot & \phantom{-}\cdot & \phantom{-}\cdot \\[-1.8pt]
\phantom{-}\cdot & \phantom{-}1 & \phantom{-}\cdot & \phantom{-}\cdot & \phantom{-}\cdot & \phantom{-}\cdot & \phantom{-}\cdot & \phantom{-}\cdot \\[-1.8pt]
\phantom{-}1 & \phantom{-}\cdot & \phantom{-}\cdot & \phantom{-}\cdot & \phantom{-}\cdot & \phantom{-}\cdot & \phantom{-}\cdot & \phantom{-}\cdot
\end{array}\right)
\]

\[
\bar\tau_{1644}^{2} = \left(\begin{array}{rrrr|rrrr}
\phantom{-}\cdot & \phantom{-}\cdot & \phantom{-}\cdot & \phantom{-}\cdot & \phantom{-}\cdot & \phantom{-}\cdot & \phantom{-}1 & \phantom{-}\cdot \\[-1.8pt]
\phantom{-}\cdot & \phantom{-}\cdot & \phantom{-}\cdot & \phantom{-}\cdot & \phantom{-}\cdot & \phantom{-}\cdot & \phantom{-}\cdot & -1 \\[-1.8pt]
\phantom{-}\cdot & \phantom{-}\cdot & \phantom{-}\cdot & \phantom{-}\cdot & -1 & \phantom{-}\cdot & \phantom{-}\cdot & \phantom{-}\cdot \\[-1.8pt]
\phantom{-}\cdot & \phantom{-}\cdot & \phantom{-}\cdot & \phantom{-}\cdot & \phantom{-}\cdot & \phantom{-}1 & \phantom{-}\cdot & \phantom{-}\cdot \\[-1.8pt]
\hline
\phantom{-}\cdot & \phantom{-}\cdot & \phantom{-}1 & \phantom{-}\cdot & \phantom{-}\cdot & \phantom{-}\cdot & \phantom{-}\cdot & \phantom{-}\cdot \\[-1.8pt]
\phantom{-}\cdot & \phantom{-}\cdot & \phantom{-}\cdot & -1 & \phantom{-}\cdot & \phantom{-}\cdot & \phantom{-}\cdot & \phantom{-}\cdot \\[-1.8pt]
-1 & \phantom{-}\cdot & \phantom{-}\cdot & \phantom{-}\cdot & \phantom{-}\cdot & \phantom{-}\cdot & \phantom{-}\cdot & \phantom{-}\cdot \\[-1.8pt]
\phantom{-}\cdot & \phantom{-}1 & \phantom{-}\cdot & \phantom{-}\cdot & \phantom{-}\cdot & \phantom{-}\cdot & \phantom{-}\cdot & \phantom{-}\cdot
\end{array}\right)
\]

\[
\bar\tau_{1644}^{3} = \left(\begin{array}{rrrr|rrrr}
\phantom{-}\cdot & \phantom{-}\cdot & \phantom{-}\cdot & \phantom{-}\cdot & \phantom{-}\cdot & -1 & \phantom{-}\cdot & \phantom{-}\cdot \\[-1.8pt]
\phantom{-}\cdot & \phantom{-}\cdot & \phantom{-}\cdot & \phantom{-}\cdot & \phantom{-}1 & \phantom{-}\cdot & \phantom{-}\cdot & \phantom{-}\cdot \\[-1.8pt]
\phantom{-}\cdot & \phantom{-}\cdot & \phantom{-}\cdot & \phantom{-}\cdot & \phantom{-}\cdot & \phantom{-}\cdot & \phantom{-}\cdot & -1 \\[-1.8pt]
\phantom{-}\cdot & \phantom{-}\cdot & \phantom{-}\cdot & \phantom{-}\cdot & \phantom{-}\cdot & \phantom{-}\cdot & \phantom{-}1 & \phantom{-}\cdot \\[-1.8pt]
\hline
\phantom{-}\cdot & -1 & \phantom{-}\cdot & \phantom{-}\cdot & \phantom{-}\cdot & \phantom{-}\cdot & \phantom{-}\cdot & \phantom{-}\cdot \\[-1.8pt]
\phantom{-}1 & \phantom{-}\cdot & \phantom{-}\cdot & \phantom{-}\cdot & \phantom{-}\cdot & \phantom{-}\cdot & \phantom{-}\cdot & \phantom{-}\cdot \\[-1.8pt]
\phantom{-}\cdot & \phantom{-}\cdot & \phantom{-}\cdot & -1 & \phantom{-}\cdot & \phantom{-}\cdot & \phantom{-}\cdot & \phantom{-}\cdot \\[-1.8pt]
\phantom{-}\cdot & \phantom{-}\cdot & \phantom{-}1 & \phantom{-}\cdot & \phantom{-}\cdot & \phantom{-}\cdot & \phantom{-}\cdot & \phantom{-}\cdot
\end{array}\right)
\]

\[
\bar\tau_{1644}^{4} = \left(\begin{array}{rrrr|rrrr}
\phantom{-}\cdot & \phantom{-}\cdot & \phantom{-}\cdot & \phantom{-}\cdot & \phantom{-}\cdot & -1 & \phantom{-}\cdot & \phantom{-}\cdot \\[-1.8pt]
\phantom{-}\cdot & \phantom{-}\cdot & \phantom{-}\cdot & \phantom{-}\cdot & \phantom{-}1 & \phantom{-}\cdot & \phantom{-}\cdot & \phantom{-}\cdot \\[-1.8pt]
\phantom{-}\cdot & \phantom{-}\cdot & \phantom{-}\cdot & \phantom{-}\cdot & \phantom{-}\cdot & \phantom{-}\cdot & \phantom{-}\cdot & \phantom{-}1 \\[-1.8pt]
\phantom{-}\cdot & \phantom{-}\cdot & \phantom{-}\cdot & \phantom{-}\cdot & \phantom{-}\cdot & \phantom{-}\cdot & -1 & \phantom{-}\cdot \\[-1.8pt]
\hline
\phantom{-}\cdot & \phantom{-}1 & \phantom{-}\cdot & \phantom{-}\cdot & \phantom{-}\cdot & \phantom{-}\cdot & \phantom{-}\cdot & \phantom{-}\cdot \\[-1.8pt]
-1 & \phantom{-}\cdot & \phantom{-}\cdot & \phantom{-}\cdot & \phantom{-}\cdot & \phantom{-}\cdot & \phantom{-}\cdot & \phantom{-}\cdot \\[-1.8pt]
\phantom{-}\cdot & \phantom{-}\cdot & \phantom{-}\cdot & -1 & \phantom{-}\cdot & \phantom{-}\cdot & \phantom{-}\cdot & \phantom{-}\cdot \\[-1.8pt]
\phantom{-}\cdot & \phantom{-}\cdot & \phantom{-}1 & \phantom{-}\cdot & \phantom{-}\cdot & \phantom{-}\cdot & \phantom{-}\cdot & \phantom{-}\cdot
\end{array}\right)
\]

\[
\bar\tau_{1644}^{5} = \left(\begin{array}{rrrr|rrrr}
\phantom{-}\cdot & \phantom{-}\cdot & \phantom{-}\cdot & \phantom{-}\cdot & \phantom{-}\cdot & \phantom{-}\cdot & \phantom{-}1 & \phantom{-}\cdot \\[-1.8pt]
\phantom{-}\cdot & \phantom{-}\cdot & \phantom{-}\cdot & \phantom{-}\cdot & \phantom{-}\cdot & \phantom{-}\cdot & \phantom{-}\cdot & \phantom{-}1 \\[-1.8pt]
\phantom{-}\cdot & \phantom{-}\cdot & \phantom{-}\cdot & \phantom{-}\cdot & -1 & \phantom{-}\cdot & \phantom{-}\cdot & \phantom{-}\cdot \\[-1.8pt]
\phantom{-}\cdot & \phantom{-}\cdot & \phantom{-}\cdot & \phantom{-}\cdot & \phantom{-}\cdot & -1 & \phantom{-}\cdot & \phantom{-}\cdot \\[-1.8pt]
\hline
\phantom{-}\cdot & \phantom{-}\cdot & -1 & \phantom{-}\cdot & \phantom{-}\cdot & \phantom{-}\cdot & \phantom{-}\cdot & \phantom{-}\cdot \\[-1.8pt]
\phantom{-}\cdot & \phantom{-}\cdot & \phantom{-}\cdot & -1 & \phantom{-}\cdot & \phantom{-}\cdot & \phantom{-}\cdot & \phantom{-}\cdot \\[-1.8pt]
\phantom{-}1 & \phantom{-}\cdot & \phantom{-}\cdot & \phantom{-}\cdot & \phantom{-}\cdot & \phantom{-}\cdot & \phantom{-}\cdot & \phantom{-}\cdot \\[-1.8pt]
\phantom{-}\cdot & \phantom{-}1 & \phantom{-}\cdot & \phantom{-}\cdot & \phantom{-}\cdot & \phantom{-}\cdot & \phantom{-}\cdot & \phantom{-}\cdot
\end{array}\right)
\]

\[
\bar\tau_{1644}^{6} = \left(\begin{array}{rrrr|rrrr}
\phantom{-}\cdot & \phantom{-}\cdot & \phantom{-}\cdot & \phantom{-}\cdot & \phantom{-}\cdot & \phantom{-}\cdot & \phantom{-}\cdot & \phantom{-}1 \\[-1.8pt]
\phantom{-}\cdot & \phantom{-}\cdot & \phantom{-}\cdot & \phantom{-}\cdot & \phantom{-}\cdot & \phantom{-}\cdot & -1 & \phantom{-}\cdot \\[-1.8pt]
\phantom{-}\cdot & \phantom{-}\cdot & \phantom{-}\cdot & \phantom{-}\cdot & \phantom{-}\cdot & \phantom{-}1 & \phantom{-}\cdot & \phantom{-}\cdot \\[-1.8pt]
\phantom{-}\cdot & \phantom{-}\cdot & \phantom{-}\cdot & \phantom{-}\cdot & -1 & \phantom{-}\cdot & \phantom{-}\cdot & \phantom{-}\cdot \\[-1.8pt]
\hline
\phantom{-}\cdot & \phantom{-}\cdot & \phantom{-}\cdot & -1 & \phantom{-}\cdot & \phantom{-}\cdot & \phantom{-}\cdot & \phantom{-}\cdot \\[-1.8pt]
\phantom{-}\cdot & \phantom{-}\cdot & \phantom{-}1 & \phantom{-}\cdot & \phantom{-}\cdot & \phantom{-}\cdot & \phantom{-}\cdot & \phantom{-}\cdot \\[-1.8pt]
\phantom{-}\cdot & -1 & \phantom{-}\cdot & \phantom{-}\cdot & \phantom{-}\cdot & \phantom{-}\cdot & \phantom{-}\cdot & \phantom{-}\cdot \\[-1.8pt]
\phantom{-}1 & \phantom{-}\cdot & \phantom{-}\cdot & \phantom{-}\cdot & \phantom{-}\cdot & \phantom{-}\cdot & \phantom{-}\cdot & \phantom{-}\cdot
\end{array}\right)
\]

\[
\bar\tau_{1644}^{7} = \left(\begin{array}{rrrr|rrrr}
\phantom{-}1 & \phantom{-}\cdot & \phantom{-}\cdot & \phantom{-}\cdot & \phantom{-}\cdot & \phantom{-}\cdot & \phantom{-}\cdot & \phantom{-}\cdot \\[-1.8pt]
\phantom{-}\cdot & \phantom{-}1 & \phantom{-}\cdot & \phantom{-}\cdot & \phantom{-}\cdot & \phantom{-}\cdot & \phantom{-}\cdot & \phantom{-}\cdot \\[-1.8pt]
\phantom{-}\cdot & \phantom{-}\cdot & \phantom{-}1 & \phantom{-}\cdot & \phantom{-}\cdot & \phantom{-}\cdot & \phantom{-}\cdot & \phantom{-}\cdot \\[-1.8pt]
\phantom{-}\cdot & \phantom{-}\cdot & \phantom{-}\cdot & \phantom{-}1 & \phantom{-}\cdot & \phantom{-}\cdot & \phantom{-}\cdot & \phantom{-}\cdot \\[-1.8pt]
\hline
\phantom{-}\cdot & \phantom{-}\cdot & \phantom{-}\cdot & \phantom{-}\cdot & -1 & \phantom{-}\cdot & \phantom{-}\cdot & \phantom{-}\cdot \\[-1.8pt]
\phantom{-}\cdot & \phantom{-}\cdot & \phantom{-}\cdot & \phantom{-}\cdot & \phantom{-}\cdot & -1 & \phantom{-}\cdot & \phantom{-}\cdot \\[-1.8pt]
\phantom{-}\cdot & \phantom{-}\cdot & \phantom{-}\cdot & \phantom{-}\cdot & \phantom{-}\cdot & \phantom{-}\cdot & -1 & \phantom{-}\cdot \\[-1.8pt]
\phantom{-}\cdot & \phantom{-}\cdot & \phantom{-}\cdot & \phantom{-}\cdot & \phantom{-}\cdot & \phantom{-}\cdot & \phantom{-}\cdot & -1
\end{array}\right)
\]

\textbf{The matrices \(\tau_{1671}^A\) of 7+1} (\(\tau_{1671}^4=1_8\)):

\[
\tau_{1671}^{0} = \left(\begin{array}{rrrr|rrrr}
\phantom{-}\cdot & \phantom{-}1 & \phantom{-}\cdot & \phantom{-}\cdot & \phantom{-}\cdot & \phantom{-}\cdot & \phantom{-}\cdot & \phantom{-}\cdot \\[-1.8pt]
-1 & \phantom{-}\cdot & \phantom{-}\cdot & \phantom{-}\cdot & \phantom{-}\cdot & \phantom{-}\cdot & \phantom{-}\cdot & \phantom{-}\cdot \\[-1.8pt]
\phantom{-}\cdot & \phantom{-}\cdot & \phantom{-}\cdot & -1 & \phantom{-}\cdot & \phantom{-}\cdot & \phantom{-}\cdot & \phantom{-}\cdot \\[-1.8pt]
\phantom{-}\cdot & \phantom{-}\cdot & \phantom{-}1 & \phantom{-}\cdot & \phantom{-}\cdot & \phantom{-}\cdot & \phantom{-}\cdot & \phantom{-}\cdot \\[-1.8pt]
\hline
\phantom{-}\cdot & \phantom{-}\cdot & \phantom{-}\cdot & \phantom{-}\cdot & \phantom{-}\cdot & -1 & \phantom{-}\cdot & \phantom{-}\cdot \\[-1.8pt]
\phantom{-}\cdot & \phantom{-}\cdot & \phantom{-}\cdot & \phantom{-}\cdot & \phantom{-}1 & \phantom{-}\cdot & \phantom{-}\cdot & \phantom{-}\cdot \\[-1.8pt]
\phantom{-}\cdot & \phantom{-}\cdot & \phantom{-}\cdot & \phantom{-}\cdot & \phantom{-}\cdot & \phantom{-}\cdot & \phantom{-}\cdot & \phantom{-}1 \\[-1.8pt]
\phantom{-}\cdot & \phantom{-}\cdot & \phantom{-}\cdot & \phantom{-}\cdot & \phantom{-}\cdot & \phantom{-}\cdot & -1 & \phantom{-}\cdot
\end{array}\right)
\]

\[
\tau_{1671}^{1} = \left(\begin{array}{rrrr|rrrr}
\phantom{-}\cdot & \phantom{-}\cdot & \phantom{-}\cdot & \phantom{-}\cdot & \phantom{-}\cdot & \phantom{-}\cdot & \phantom{-}\cdot & \phantom{-}1 \\[-1.8pt]
\phantom{-}\cdot & \phantom{-}\cdot & \phantom{-}\cdot & \phantom{-}\cdot & \phantom{-}\cdot & \phantom{-}\cdot & \phantom{-}1 & \phantom{-}\cdot \\[-1.8pt]
\phantom{-}\cdot & \phantom{-}\cdot & \phantom{-}\cdot & \phantom{-}\cdot & \phantom{-}\cdot & -1 & \phantom{-}\cdot & \phantom{-}\cdot \\[-1.8pt]
\phantom{-}\cdot & \phantom{-}\cdot & \phantom{-}\cdot & \phantom{-}\cdot & -1 & \phantom{-}\cdot & \phantom{-}\cdot & \phantom{-}\cdot \\[-1.8pt]
\hline
\phantom{-}\cdot & \phantom{-}\cdot & \phantom{-}\cdot & \phantom{-}1 & \phantom{-}\cdot & \phantom{-}\cdot & \phantom{-}\cdot & \phantom{-}\cdot \\[-1.8pt]
\phantom{-}\cdot & \phantom{-}\cdot & \phantom{-}1 & \phantom{-}\cdot & \phantom{-}\cdot & \phantom{-}\cdot & \phantom{-}\cdot & \phantom{-}\cdot \\[-1.8pt]
\phantom{-}\cdot & -1 & \phantom{-}\cdot & \phantom{-}\cdot & \phantom{-}\cdot & \phantom{-}\cdot & \phantom{-}\cdot & \phantom{-}\cdot \\[-1.8pt]
-1 & \phantom{-}\cdot & \phantom{-}\cdot & \phantom{-}\cdot & \phantom{-}\cdot & \phantom{-}\cdot & \phantom{-}\cdot & \phantom{-}\cdot
\end{array}\right)
\]

\[
\tau_{1671}^{2} = \left(\begin{array}{rrrr|rrrr}
\phantom{-}\cdot & \phantom{-}\cdot & \phantom{-}\cdot & \phantom{-}\cdot & \phantom{-}\cdot & \phantom{-}\cdot & -1 & \phantom{-}\cdot \\[-1.8pt]
\phantom{-}\cdot & \phantom{-}\cdot & \phantom{-}\cdot & \phantom{-}\cdot & \phantom{-}\cdot & \phantom{-}\cdot & \phantom{-}\cdot & \phantom{-}1 \\[-1.8pt]
\phantom{-}\cdot & \phantom{-}\cdot & \phantom{-}\cdot & \phantom{-}\cdot & \phantom{-}1 & \phantom{-}\cdot & \phantom{-}\cdot & \phantom{-}\cdot \\[-1.8pt]
\phantom{-}\cdot & \phantom{-}\cdot & \phantom{-}\cdot & \phantom{-}\cdot & \phantom{-}\cdot & -1 & \phantom{-}\cdot & \phantom{-}\cdot \\[-1.8pt]
\hline
\phantom{-}\cdot & \phantom{-}\cdot & -1 & \phantom{-}\cdot & \phantom{-}\cdot & \phantom{-}\cdot & \phantom{-}\cdot & \phantom{-}\cdot \\[-1.8pt]
\phantom{-}\cdot & \phantom{-}\cdot & \phantom{-}\cdot & \phantom{-}1 & \phantom{-}\cdot & \phantom{-}\cdot & \phantom{-}\cdot & \phantom{-}\cdot \\[-1.8pt]
\phantom{-}1 & \phantom{-}\cdot & \phantom{-}\cdot & \phantom{-}\cdot & \phantom{-}\cdot & \phantom{-}\cdot & \phantom{-}\cdot & \phantom{-}\cdot \\[-1.8pt]
\phantom{-}\cdot & -1 & \phantom{-}\cdot & \phantom{-}\cdot & \phantom{-}\cdot & \phantom{-}\cdot & \phantom{-}\cdot & \phantom{-}\cdot
\end{array}\right)
\]

\[
\tau_{1671}^{3} = \left(\begin{array}{rrrr|rrrr}
\phantom{-}\cdot & \phantom{-}\cdot & \phantom{-}\cdot & \phantom{-}\cdot & \phantom{-}\cdot & \phantom{-}1 & \phantom{-}\cdot & \phantom{-}\cdot \\[-1.8pt]
\phantom{-}\cdot & \phantom{-}\cdot & \phantom{-}\cdot & \phantom{-}\cdot & -1 & \phantom{-}\cdot & \phantom{-}\cdot & \phantom{-}\cdot \\[-1.8pt]
\phantom{-}\cdot & \phantom{-}\cdot & \phantom{-}\cdot & \phantom{-}\cdot & \phantom{-}\cdot & \phantom{-}\cdot & \phantom{-}\cdot & \phantom{-}1 \\[-1.8pt]
\phantom{-}\cdot & \phantom{-}\cdot & \phantom{-}\cdot & \phantom{-}\cdot & \phantom{-}\cdot & \phantom{-}\cdot & -1 & \phantom{-}\cdot \\[-1.8pt]
\hline
\phantom{-}\cdot & \phantom{-}1 & \phantom{-}\cdot & \phantom{-}\cdot & \phantom{-}\cdot & \phantom{-}\cdot & \phantom{-}\cdot & \phantom{-}\cdot \\[-1.8pt]
-1 & \phantom{-}\cdot & \phantom{-}\cdot & \phantom{-}\cdot & \phantom{-}\cdot & \phantom{-}\cdot & \phantom{-}\cdot & \phantom{-}\cdot \\[-1.8pt]
\phantom{-}\cdot & \phantom{-}\cdot & \phantom{-}\cdot & \phantom{-}1 & \phantom{-}\cdot & \phantom{-}\cdot & \phantom{-}\cdot & \phantom{-}\cdot \\[-1.8pt]
\phantom{-}\cdot & \phantom{-}\cdot & -1 & \phantom{-}\cdot & \phantom{-}\cdot & \phantom{-}\cdot & \phantom{-}\cdot & \phantom{-}\cdot
\end{array}\right)
\]

\[
\tau_{1671}^{5} = \left(\begin{array}{rrrr|rrrr}
\phantom{-}\cdot & \phantom{-}\cdot & -1 & \phantom{-}\cdot & \phantom{-}\cdot & \phantom{-}\cdot & \phantom{-}\cdot & \phantom{-}\cdot \\[-1.8pt]
\phantom{-}\cdot & \phantom{-}\cdot & \phantom{-}\cdot & -1 & \phantom{-}\cdot & \phantom{-}\cdot & \phantom{-}\cdot & \phantom{-}\cdot \\[-1.8pt]
\phantom{-}1 & \phantom{-}\cdot & \phantom{-}\cdot & \phantom{-}\cdot & \phantom{-}\cdot & \phantom{-}\cdot & \phantom{-}\cdot & \phantom{-}\cdot \\[-1.8pt]
\phantom{-}\cdot & \phantom{-}1 & \phantom{-}\cdot & \phantom{-}\cdot & \phantom{-}\cdot & \phantom{-}\cdot & \phantom{-}\cdot & \phantom{-}\cdot \\[-1.8pt]
\hline
\phantom{-}\cdot & \phantom{-}\cdot & \phantom{-}\cdot & \phantom{-}\cdot & \phantom{-}\cdot & \phantom{-}\cdot & \phantom{-}1 & \phantom{-}\cdot \\[-1.8pt]
\phantom{-}\cdot & \phantom{-}\cdot & \phantom{-}\cdot & \phantom{-}\cdot & \phantom{-}\cdot & \phantom{-}\cdot & \phantom{-}\cdot & \phantom{-}1 \\[-1.8pt]
\phantom{-}\cdot & \phantom{-}\cdot & \phantom{-}\cdot & \phantom{-}\cdot & -1 & \phantom{-}\cdot & \phantom{-}\cdot & \phantom{-}\cdot \\[-1.8pt]
\phantom{-}\cdot & \phantom{-}\cdot & \phantom{-}\cdot & \phantom{-}\cdot & \phantom{-}\cdot & -1 & \phantom{-}\cdot & \phantom{-}\cdot
\end{array}\right)
\]

\[
\tau_{1671}^{6} = \left(\begin{array}{rrrr|rrrr}
\phantom{-}\cdot & \phantom{-}\cdot & \phantom{-}\cdot & -1 & \phantom{-}\cdot & \phantom{-}\cdot & \phantom{-}\cdot & \phantom{-}\cdot \\[-1.8pt]
\phantom{-}\cdot & \phantom{-}\cdot & \phantom{-}1 & \phantom{-}\cdot & \phantom{-}\cdot & \phantom{-}\cdot & \phantom{-}\cdot & \phantom{-}\cdot \\[-1.8pt]
\phantom{-}\cdot & -1 & \phantom{-}\cdot & \phantom{-}\cdot & \phantom{-}\cdot & \phantom{-}\cdot & \phantom{-}\cdot & \phantom{-}\cdot \\[-1.8pt]
\phantom{-}1 & \phantom{-}\cdot & \phantom{-}\cdot & \phantom{-}\cdot & \phantom{-}\cdot & \phantom{-}\cdot & \phantom{-}\cdot & \phantom{-}\cdot \\[-1.8pt]
\hline
\phantom{-}\cdot & \phantom{-}\cdot & \phantom{-}\cdot & \phantom{-}\cdot & \phantom{-}\cdot & \phantom{-}\cdot & \phantom{-}\cdot & \phantom{-}1 \\[-1.8pt]
\phantom{-}\cdot & \phantom{-}\cdot & \phantom{-}\cdot & \phantom{-}\cdot & \phantom{-}\cdot & \phantom{-}\cdot & -1 & \phantom{-}\cdot \\[-1.8pt]
\phantom{-}\cdot & \phantom{-}\cdot & \phantom{-}\cdot & \phantom{-}\cdot & \phantom{-}\cdot & \phantom{-}1 & \phantom{-}\cdot & \phantom{-}\cdot \\[-1.8pt]
\phantom{-}\cdot & \phantom{-}\cdot & \phantom{-}\cdot & \phantom{-}\cdot & -1 & \phantom{-}\cdot & \phantom{-}\cdot & \phantom{-}\cdot
\end{array}\right)
\]

\[
\tau_{1671}^{7} = \left(\begin{array}{rrrr|rrrr}
\phantom{-}\cdot & \phantom{-}\cdot & \phantom{-}\cdot & \phantom{-}\cdot & -1 & \phantom{-}\cdot & \phantom{-}\cdot & \phantom{-}\cdot \\[-1.8pt]
\phantom{-}\cdot & \phantom{-}\cdot & \phantom{-}\cdot & \phantom{-}\cdot & \phantom{-}\cdot & -1 & \phantom{-}\cdot & \phantom{-}\cdot \\[-1.8pt]
\phantom{-}\cdot & \phantom{-}\cdot & \phantom{-}\cdot & \phantom{-}\cdot & \phantom{-}\cdot & \phantom{-}\cdot & -1 & \phantom{-}\cdot \\[-1.8pt]
\phantom{-}\cdot & \phantom{-}\cdot & \phantom{-}\cdot & \phantom{-}\cdot & \phantom{-}\cdot & \phantom{-}\cdot & \phantom{-}\cdot & -1 \\[-1.8pt]
\hline
\phantom{-}1 & \phantom{-}\cdot & \phantom{-}\cdot & \phantom{-}\cdot & \phantom{-}\cdot & \phantom{-}\cdot & \phantom{-}\cdot & \phantom{-}\cdot \\[-1.8pt]
\phantom{-}\cdot & \phantom{-}1 & \phantom{-}\cdot & \phantom{-}\cdot & \phantom{-}\cdot & \phantom{-}\cdot & \phantom{-}\cdot & \phantom{-}\cdot \\[-1.8pt]
\phantom{-}\cdot & \phantom{-}\cdot & \phantom{-}1 & \phantom{-}\cdot & \phantom{-}\cdot & \phantom{-}\cdot & \phantom{-}\cdot & \phantom{-}\cdot \\[-1.8pt]
\phantom{-}\cdot & \phantom{-}\cdot & \phantom{-}\cdot & \phantom{-}1 & \phantom{-}\cdot & \phantom{-}\cdot & \phantom{-}\cdot & \phantom{-}\cdot
\end{array}\right)
\]

\hypertarget{appendix-c.-files}{%
\section{Appendix C. Files}\label{appendix-c.-files}}

\begin{longtable}[]{@{}
  >{\raggedright\arraybackslash}p{(\columnwidth - 2\tabcolsep) * \real{0.4000}}
  >{\raggedright\arraybackslash}p{(\columnwidth - 2\tabcolsep) * \real{0.6000}}@{}}
\toprule\noalign{}
\begin{minipage}[b]{\linewidth}\raggedright
File
\end{minipage} & \begin{minipage}[b]{\linewidth}\raggedright
Content
\end{minipage} \\
\midrule\noalign{}
\endhead
\bottomrule\noalign{}
\endlastfoot
\texttt{tools/nb2txt.py} & converts a Mathematica notebook to readable
text, cell by cell (no Mathematica needed) \\
\texttt{tools/extract\_nb\_matrices.py} & extracts the matrices printed
in notebook output cells \\
\texttt{data/notebook\_matrices\_4+4.txt},
\texttt{data/notebook\_matrices\_7+1.txt} & the matrices as printed by
Mathematica \\
\texttt{code/spacetimes-rs/src/gamma.rs} & construction of all matrices
from the notebooks' definitions \\
\texttt{code/spacetimes-rs/src/verify.rs} & checks R, C, L, G, N, B, V,
P, S, K \\
\texttt{code/spacetimes-rs/src/dirac.rs} & checks D (the two Dirac
equations) \\
\texttt{code/spacetimes-rs/src/pairs.rs} & checks Q (fermion pairs) \\
\texttt{code/spacetimes-rs/src/imat.rs}, \texttt{fmat.rs} & exact
integer / Gaussian-integer matrices; double-precision matrices \\
\texttt{code/scripts/verify\_matrices\_independent.py} & independent
re-computation in Python (checks I) \\
\texttt{data/gamma1644.txt}, \texttt{data/gamma1671.txt} & the project's
matrices as plain integers \\
\texttt{data/matrix\_checks.txt}, \texttt{data/matrix\_checks.csv} & the
result of every check of the Rust program \\
\texttt{data/matrix\_checks\_independent.txt} & the result of every
check of the Python script \\
\texttt{data/pair\_classification.csv},
\texttt{data/pair\_classification.txt} & the pair states by multiplet \\
\texttt{provenance/generated/matrices.md} & the matrix displays inserted
into this document \\
\texttt{docs\_src/dirac\ matrices.src.md} & the source of this
document \\
\end{longtable}

\end{document}
